\documentclass[11pt, twoside, openright]{book}

% Trim size: 6" x 9" (US Trade paperback)
\usepackage[
  paperwidth=6in, paperheight=9in, inner=0.875in, outer=0.75in, top=0.75in, bottom=0.75in
]{geometry}

% Typography
\usepackage{fontspec}
\usepackage{xeCJK}
\setmainfont{TeX Gyre Pagella}
\setCJKmainfont{Noto Serif CJK JP}
\emergencystretch=1em          % Allow extra stretch to avoid overflow with hyphenated/apostrophe words
\usepackage{setspace}
\setstretch{1.15}              % Standard book leading

% Graphics
\usepackage{graphicx}
\graphicspath{{../figures/}{../figures/book/}}

% Links (hidden for print - no colored text)
\usepackage{xcolor}
\usepackage[hidelinks, bookmarks=false]{hyperref}
\hypersetup{
  pdftitle={「私」という名のシミュレーション},
  pdfauthor={Matthias Gruber},
  pdfsubject={Consciousness, Computation, and the Cosmos},
}

% Chapter styling
\usepackage{titlesec}
\titleformat{\chapter}[display]
  {\normalfont\huge\bfseries}
  {\chaptertitlename\ \thechapter}{20pt}{\Huge\raggedright}
\titleformat{name=\chapter,numberless}[display]
  {\normalfont\huge\bfseries}
  {}{0pt}{\Huge\raggedright}
\titlespacing*{\chapter}{0pt}{50pt}{40pt}

% Section styling
\titleformat{\section}
  {\normalfont\Large\bfseries}{}{0pt}{}
\titlespacing*{\section}{0pt}{20pt}{10pt}

% Headers and footers
\usepackage{fancyhdr}
\pagestyle{fancy}
\fancyhf{}
\fancyhead[LE]{\small\itshape 「私」という名のシミュレーション}
\fancyhead[RO]{\small\itshape\leftmark}
\fancyfoot[C]{\thepage}
\renewcommand{\headrulewidth}{0.4pt}

% For plain pages (chapter starts, front matter)
\fancypagestyle{plain}{
  \fancyhf{}
  \fancyfoot[C]{\thepage}
  \renewcommand{\headrulewidth}{0pt}
}

% Quote formatting
\usepackage{csquotes}
\renewenvironment{quote}
  {\list{}{\leftmargin=1.5em\rightmargin=1.5em}\item\relax\itshape}
  {\endlist}

% Prevent widows and orphans
\widowpenalty=10000
\clubpenalty=10000

% Tables
\usepackage{booktabs}
\renewcommand{\lightrulewidth}{0.8pt}   % KDP minimum line weight is 0.75pt
\usepackage{longtable}
\usepackage{array}
\usepackage{tabularx}

% Figure placement
\usepackage{float}

% Landscape pages for wide tables
\usepackage{pdflscape}
\renewcommand{\textfraction}{0.1}
\renewcommand{\topfraction}{0.9}
\renewcommand{\bottomfraction}{0.9}

% Make blank pages (from \cleardoublepage) truly blank
\makeatletter
\def\cleardoublepage{\clearpage\if@twoside \ifodd\c@page\else
  \thispagestyle{empty}\hbox{}\newpage
  \if@twocolumn\hbox{}\newpage\fi\fi\fi}
\makeatother

\begin{document}

% ==== FRONT MATTER ====
\frontmatter
\pagestyle{empty}

% ---- Half-title page (recto) ----
\vspace*{3in}
\begin{center}
{\Huge\bfseries 「私」という名の\\[0.3cm]シミュレーション\par}
\end{center}
\cleardoublepage

% ---- Full title page (recto) ----
\vspace*{2in}
\begin{center}
{\Huge\bfseries 「私」という名の\\[0.3cm]シミュレーション\par}
\vspace{0.8cm}
{\Large 意識、計算、そして\\[0.2cm]宇宙のアーキテクチャ\par}
\vspace{2cm}
{\large Matthias Gruber\par}
\end{center}
\clearpage

% ---- Copyright page (verso of title) ----
\thispagestyle{empty}
\vspace*{\fill}
{\small
\noindent \textcopyright\ 2026 Matthias Gruber. 無断転載を禁じます。\par
\vspace{0.5cm}
\noindent 本書のいかなる部分も、著作権法で認められる批評における短い引用および特定の
非営利的利用を除き、著者の事前の書面による許可なく、いかなる形式・手段によって
も複製、配布、送信することはできません。\par
\vspace{0.5cm}

\vspace{0.5cm}
\noindent 第二版　2026年\par
\vspace{0.5cm}
\noindent www.matthiasgruber.com\par
}
\cleardoublepage

% ---- Dedication page (recto) ----
\thispagestyle{empty}
\vspace*{3in}
\begin{center}
\textit{なぜ何かが、何かとして感じられるのか——一度でも不思議に思ったことのあるすべての人へ。}
\end{center}
\cleardoublepage

% ---- Table of contents ----
\pagestyle{plain}
\tableofcontents
\cleardoublepage








\chapter*{そこにない点}
\addcontentsline{toc}{chapter}{そこにない点}
\markboth{そこにない点}{}

何かを語り始める前に------理論よりも先に、私が何者であるかよりも先に、なぜこの話が重要なのかよりも先に------一つ、自分自身で確かめてみてほしいことがある。10秒で済む。

下の二つの印を見てほしい。左に\textbf{X}、右に塗りつぶされた\textbf{点}がある。


\begin{figure}[htbp]
  \centering
  \includegraphics[width=0.95\textwidth]{../figures/blind-spot-test.png}
\end{figure}


本を腕いっぱいに伸ばして持つ。\textbf{左}目を閉じる。右目でXをまっすぐ見つめる------Xだけを。ずるをして点のほうへ視線を走らせてはいけない。視界の隅にそれが居座っているのが感じられるはずで、それで十分だ。そして、視線をXに固定したまま、本をゆっくりと顔に近づけていく。

顔から30センチほどのあたりで、視界の端の点に注意を向けてほしい。

消えている。

ぼやけたのではない。薄れたのでもない。\textbf{消えた}のだ------そしてここが、立ち止まる価値のあるところだ。点があった場所に、穴は開いていない。隙間も、黒っぽい染みも、空白もない。太い黒の点が印刷されているまさにその位置で、開いた、正常に機能している目のまん前で、ページはただ、継ぎ目なく、途切れなく白い。目はまっすぐそこに向いている。それなのに見えない。脳が、そこを白でそっと塗りつぶしてしまったからだ。

いま起きたのは、こういうことだ。どの目にも盲点がある------視神経が脳へと突き抜けるまさにその場所に、光を感じる細胞がまったく存在しない網膜の一画があるのだ。視野には、満月九つを呑み込めるほどの穴が、左右それぞれの目に、常に開いている。それに気づいたことは一度もないはずだ。穴が小さいからではない------途方もなく大きい------。脳がその穴を決して見せないからだ。脳は穴の周囲を読み取り、真ん中を\textbf{でっち上げる}。ページの白を創作し、壁紙を創作し、そこにあるはずのものを創作して、完成した絵を、まるでそれが生の真実であるかのように手渡してくるのである。

もう一度本を動かしてみよう。点が戻ってくる。何度かやってみてほしい。いま目にしているのは、自分の脳が世界を報告することと世界を創作することのあいだで切り替わるさまだ。それなのに、その切り替えはまったく感じられない。切り替えは決して予告しない。これまでも、一度たりとも予告したことはないのだ。

ここからが、正気の沙汰とは思えないのに、れっきとした確立済みの神経科学であるところだ。この埋め合わせは、脳が盲点のためだけに取ってある特別な芸当ではない。起きている一秒一秒、視野全体にわたって脳がやっていることそのものなのである。盲点は、たまたまぼろを出すほど不器用な唯一の場所にすぎない------入ってくる光がまったくない唯一の一画であり、そこでは創作をゼロから走らせるしかなく、足を滑らせるのだ。それ以外のあらゆる場所でも、生の信号はほとんど同じくらい薄く、ノイズにまみれていて、同じ推測の機構で水増しされている。気づかないのはただ、その推測がたいてい当たっているからだ。ふっと消えたあの点は、現行犯で押さえられた脳だった。いま見えているその他すべては、同じ行為が、気づけないほど見事に演じられているにすぎない。

この事実を、行き着くところまで辿ってみよう。生まれてから今日まで、現実を直接見たことは一度もない。ただの一度も、だ。手の中のこのページも、周りの部屋も、これまで愛した誰の顔も。そのすべてが再構成だった------頭蓋の内側で、細々とした神経インパルスの流れと山のような事前の推測から組み上げられたモデルであり、それがあまりに滑らかに差し出されるため、\textbf{これは本物ではないかもしれない}という考えは、おそらくいまのいままで、一度も頭をよぎらなかったはずだ。

私たちはシミュレーションの中で生きている。それを生成しているのは、暗闇に置かれた1.5キログラムほどの湿った組織だ。それは生涯アクセスできる唯一の世界であり、しかも一生のあいだ休みなく働き続けながら、自分がそこにあることを、ただの一度も口にしなかった。

いま感じているかもしれない、そのめまい------手放さないでほしい。その感覚こそが、この本の主題のすべてである。

私はそれを内側から知っている。ちょうど25歳のとき、インスブルックの橋の上で、真っ昼間に、私はそれを感じた。涙が頬を伝うそばから笑いが止まらず、人生でいちばん重かった荷が、私の肩から下りた。誰かに見られていたかどうかはわからない。見られていても構わなかっただろう。そのとき一挙に腑に落ちたのは、一つの枠組みだった。脳があの継ぎ目のない世界をどう組み上げるのかだけでなく、その組み上げがなぜ\textbf{そもそも何かとして感じられるもの}なのか------これを読んでいるあなたであるということが、なぜ「何かであるような感じ」を伴うのか------までを説明する枠組みである。その瞬間から、少なくとも私にとっては、残りの人生でやるべきことは決まってしまった。

2015年、私はそれを書き下ろした。ドイツ語で300ページ、自費出版、専門的な道具立てがぎっしり詰まった本である。売れた部数はゼロ。一冊も、だ。同情を引くために言うのではない------この物語に関わることだから言うのだ。あの読まれずじまいの本の中の理論は、科学に残された最難問かもしれないものを解消し、他のどの理論も立てていない予測を立て、本当に意識を持つ機械------心を模倣するチャットボットではなく、心を\textbf{持つ}機械------を造るための具体的な設計図を差し出す。その後、理論は10年のあいだハードディスクの上で眠り、そのかたわらで、証拠は静かに理論の側へと集まっていった。この本は二度目の試みだ。より短く、英語で、なぜ何かが何かとして感じられるのかと一度でも不思議に思ったことのあるすべての人に向けて書かれている。もし私が正しければ、意識のハード・プロブレム（難しい問題）は、難しくない。それは、見えた瞬間に解消するカテゴリー錯誤である------1ページ前にあの点が消えたのと、同じように。


\chapter*{著者について}
\addcontentsline{toc}{chapter}{著者について}
\markboth{著者について}{}

こうした主張をしているのが何者なのか、述べておくべきだろう。私はどの大学にも所属していない。博士号も持っていない------あるのはバイオインフォマティクスの修士号だけだ。研究助成金を受けたことも、神経科学の研究室で働いたこともない。議論を信用する前にまず肩書きを確かめるという人なら------それはもっともな心構えだが------ここで本を閉じたくなるかもしれない。その場合は、がたつくテーブルの脚の下にでも敷いてほしい。せめて木々の犠牲が無駄にならないように。

代わりに私にあるのは、独学の門外漢としての経歴だ。壁に突き当たるたびに向きを変え、分野と分野の境界を尊重することをついに覚えなかった、数十年にわたる読書である。これは見かけ以上に重要な点だ。この理論が存在するのは、どの学問分野が互いに触れてはならないのか、誰も私に教えてくれなかったからにほかならない------数学とセル・オートマトン、機械学習、臨床神経学から精神薬理学までの神経科学、進化生物学、心の哲学。意識の理論の多くは、こうした世界の一つか二つの中で生きている。本書の理論はそのすべてを結び合わせようとするのだが、考えてみれば、それこそまさに脳のしていることだ。脳は、まったく異なる源から来るばらばらな流れを受け取り、一つの経験へと織り上げる。脳が分野を横断するように学問分野を横断できない意識の理論には、むしろ疑いの目を向けるべきだろう。私の理論は、それ自身が記述するのと同じ再帰的ループの産物である------知識が成果を養い、成果が意欲を養い、意欲がまた知識へと還っていくループだ。このループについては本書の後半で詳しく述べる。理論に値打ちがあるかどうかは、自分の目で確かめてほしい。では、理論へ。


\mainmatter
\pagestyle{fancy}
\chapter*{第1章：科学でもっとも難しい問題}
\addcontentsline{toc}{chapter}{第1章：科学でもっとも難しい問題}
\markboth{第1章：科学でもっとも難しい問題}{}

あなたはいま、この文を読んでいる。そして、ひとつの経験をしている。

その経験------ページの上の文字という視覚的印象、言葉を読み上げる内なる声、理解あるいは戸惑いの感覚------は、人生でもっとも身近なものであると同時に、宇宙でもっとも謎めいたものである。私たちは、読むことがなぜ何かのように感じられるのかについてよりも、ブラックホールの内部について多くを知っている。

これは誇張ではない。（もっとも公平を期して言えば、数学にはこれよりさらに難しいと私が考える問題もある。ただ、それで夜も眠れなくなる人はほとんどいない。）物理学には標準模型がある。生物学には進化と遺伝学がある。化学には周期表がある。だが意識------いまこれを読んでいる、まさにこの瞬間、「あなたであるという感じ」が存在するという事実------には、確立された理論も、支配的な枠組みも、合意を得た説明もない。

試みが足りなかったわけではない。1990年代、行動主義による数十年の追放を経て意識が再びまっとうな科学のテーマとなって以来、何千本もの論文が発表され、何十もの理論が提唱され、何億ドルもの資金が費やされてきた。その結果は？科学哲学者トマス・クーン（Thomas Kuhn）の言う「前パラダイム状態」にある分野である------競合するアイデアばかりが多く、コンセンサスはなく、何か根本的なものが欠けているのではないかという感覚だけが募っていく。

\section*{ハード・プロブレムは実際に何を問うているのか}

1995年、哲学者デイヴィッド・チャーマーズ（David Chalmers）は、この謎にその名を与えた。意識のハード・プロブレム（難しい問題）である。

それが問うのはこういうことだ。赤を見るという経験を考えてみてほしい。神経科学者は、赤を見るとき脳で何が起きるかについて多くを語ることができる。特定の波長の光が網膜の錐体細胞に当たり、信号が視神経を伝わり、視覚野で処理され、さまざまな脳領域が協調してその知覚を生み出す。以上のことは、少なくとも大筋では、よく理解されている。

だが、そのどれ一つとして、\textbf{なぜ赤を見ることが何かのように感じられるのか}を説明してはいない。

原理的には、赤い光に対する脳の応答について、完全な神経モデルを構築することもできるだろう------ニューロンの一つひとつ、シナプスの一つひとつ、信号経路の一つひとつまで。手に入るのは完璧な機能的記述である。それでも、赤さの感じは説明できていない。あの「どのように感じられるか」という何か。哲学者が\textbf{クオリア}（感覚質）と呼ぶものである。

チャーマーズはこれを、意識の「イージー・プロブレム（易しい諸問題）」から区別した------といっても少しも易しくはなく、原理的に手がつけられるというだけのことだが。脳はどのように情報を統合するのか。どのように注意を向けるのか。どのように自らの状態を報告するのか。これらはメカニズムの問題である。難しいが、その難しさは、神経科学が取り組み方を知っている種類のものだ。ハード・プロブレムは違う。それは、なぜそうしたメカニズムにそもそも経験が伴うのかを問う。なぜ脳はコンピュータのように、ただ「暗闇の中で」情報を処理しているだけではないのか。

\section*{議論の現状}

候補に事欠いたわけではない。\textbf{統合情報理論}は意識を$\Phi$という一つの数に還元する------そしてその優雅さの代償として、十分に大きな論理ゲートの格子もわずかに意識を持つという含意を抱え込む。2023年には120人を超える研究者が、この理論を反証不可能だとする公開書簡に署名した。\textbf{グローバル・ニューロナル・ワークスペース理論}は、情報が脳全体にブロードキャストされたとき意識化されると説く------\textbf{いつ}についての明快な説明であり、なぜブロードキャストが何かのように感じられるはずなのかという\textbf{なぜ}については、あからさまな沈黙である。\textbf{予測処理}は脳を予測機械として、知覚を「制御された幻覚」として描く------経験の\textbf{内容}については見事だが、その設計者たち自身が認めるとおり、難しい部分にはそもそも狙いを定めていない。高階理論、注意スキーマ理論、リカレント処理理論、電磁場理論------どれも本物の洞察でありながら、同じ欠けた中心のまわりを取り巻いているにすぎない。

そして2025年、審判が競技場に姿を現した。COGITATE共同研究は、二強であるIITとGNWを数年がかりの敵対的検証にかけ、その結果を\textit{Nature}誌に発表した。評決はこうだ------どちらも勝たなかった。意識ともっとも強く結びついた信号が現れたのは脳の後方、どちらの陣営も旗を立てていなかった場所だった。30年の歳月と数億ドルを費やしたすえに、この分野はむしろ出発点よりもコンセンサスから遠ざかったと言ってよい。これは、答えに迫りつつある科学の姿ではない。ピースを一つ欠いた科学の姿である。

\section*{進歩を阻む二つのドグマ}

何が欠けていると私が考えるかを述べる前に、数十年にわたってこの分野をひそかに蝕んできた二つの偏見に名前を与えておかなければならない。もとの著書で私がそれらに名前をつけたのは、名前のない偏見ほど闘いにくいものはないと考えるからだ。

第一のものを、私は\textbf{nSAIドグマ}と呼ぶ------「強い人工知能は存在しえない（no strong artificial intelligence）」。真に知的な機械は不可能だという広く行きわたった確信であり、その根は証明にではなく、1960年代の初期AI研究の挫折と、それに続く反動にある。強いAIが可能だと信じる者は、主流の研究で真剣に受け取られたければ、それについて口をつぐむことを学ぶ。これは理性的な懐疑ではない。古い敗北の傷跡が、教義にまで硬化したものである。

第二のドグマはより深く、より有害だ。私はそれを\textbf{nSUドグマ}と呼ぶ------「自己理解は不可能である（no self-understanding）」。人間の精神は、人間の意識は、まさにその同じ精神によって原理的に理解されえない、という信念である。人はゲーデルの不完全性定理を持ち出し、あるいは宇宙の内側から宇宙を観測することの限界といった漠然としたアナロジーに訴え、あるいは------これがいちばん正直なところだが------自分が完全に説明されてしまうという見通しを、考えるにはあまりに恐ろしいと感じているだけである。意識がただの機械だとしたら、魂はどうなるのか。意味はどうなるのか。人間であることの特別さは、どうなるのか。

これらのドグマは互いに補強し合う。意識が理解できないのなら（nSU）、意識を造ることなど当然できない（nSAI）。そして意識が造れないのなら（nSAI）、意識はやはり理解を超えたものなのかもしれない（nSU）。これは制度化された悲観主義の閉じたループであり、膨大な数の聡明な研究者たちから、この仕事に挑むことすら奪ってきた。

これらのドグマが不誠実に唱えられていると言いたいのではない。多くの研究者は心からそれを信じている。だが、どちらのドグマも証明されたことは一度もない。それらは信仰箇条であり、意識研究に対して、いかなる失敗した実験よりも大きな損害を与えてきた。

\section*{何かが欠けている}

どの理論もハード・プロブレムを解けずにいる理由は、ほとんどの人が間違った場所に意識を探しているからだと私は考える。彼らは神経の機構------ニューロン、シナプス、振動、結合性------を見つめて、こう問う。「これらのプロセスのうち、どれが意識的なのか？」

正しい問いは、私の考えでは別のものだ。「情報処理のどのレベルで、どのアーキテクチャにおいて、経験は生じるのか？」

これが四モデル理論（FMT）の出発点である------この本が本格的に始まるより前に、自分自身に証明してみせた、あの事実だ。あなたは生まれてこのかた一度たりとも、現実を直接に経験したことがない。ページから消し去ったあの点は、紙の仕掛けではなかった。網膜に開いた穴を、脳が捏造した内容で埋めていたのである------脳はあらゆる穴を、どの一瞥においてもそうやって埋めている。あまりに継ぎ目がないので、疑いを抱いたことは一度もなかったはずだ。これは異端の主張ではない。知覚が直接的なものではなく構築されたものであることは神経科学の主流であり、この分野のほぼすべての研究者に受け入れられている。FMTが付け加えるのは、ほとんど誰も踏み出さない一歩------このただ一つの事実を最後まで突きつめれば、ハード・プロブレムは溶けて消える、ということである。

\section*{一本の剃刀}

私はこの理論を、目に見える規律をただ一つだけ課して築いた。\textbf{オッカムの剃刀}である。新しい物理学はなし、微小管の中の量子の魔法もなし、物質にまき散らされた原意識もなし------すでに手元にある材料だけ（ニューラルネットワーク、学習、シミュレーション、自己言及）を、正しいアーキテクチャに配置する。さらに二つの古い原理が背後で静かに仕事をしているが、いまここで観閲式のように行進させるのではなく、それぞれが効いてくるまさにその場所で登場させることにする。一つは動物の話に至ったとき、もう一つはゾンビと自由意志の話に至ったときである。ここでは、一本の刃で足りる。

いよいよ、四つのモデルを見るときが来た。


\chapter*{第2章：四つのモデル}
\addcontentsline{toc}{chapter}{第2章：四つのモデル}
\markboth{第2章：四つのモデル}{}

リンゴを一つ見つめていると想像してみてほしい。

リンゴは目の前のテーブルの上にある。赤く、丸く、つやがあり、手からおよそ15センチのところにある。それが見えるし、何であるかもわかり、手を伸ばせばつかむこともできる。単純な話に思える------リンゴを見ているのだ、と。

だが、実際に起きていることは、比べものにならないほど複雑だ。

リンゴの表面で反射した光が目に入り、網膜の光受容細胞に当たる。これらの細胞は光を電気信号に変換する。信号は視神経を通って脳の後方にある視覚野へと伝わり、そこでますます高度になっていく特徴検出器の階層------輪郭、方位、色、テクスチャ、形、そして最終的には物体------を経て処理される。このカスケードのどこかで、「リンゴ」に対応する神経活動が活性化する。それと同時に、運動系は起こしうる行動（手を伸ばす、つかむ）を準備し、記憶系は連想（味、手ざわり、前回リンゴを食べたとき）を活性化させ、空間系は自分の身体に対するリンゴの位置を追跡している。

こうしたことのすべてが、1秒に満たないうちに起こる。そして、そのどれ一つとして、あなたが\textbf{経験}するものではない。光子が錐体細胞に当たることも、信号が軸索を伝わることも、特徴検出器が発火することも、経験しはしない。あなたが経験するのは\textbf{一個のリンゴ}だ。統一され、安定した、三次元の物体が、首尾一貫した空間環境の中に置かれ、特定の見えと感触と意味をもっている。私たちが経験しているのは\textbf{モデル}である------リンゴのリアルタイム・シミュレーションであり、生のデータと、リンゴや物体やテーブルや物理法則についてこれまでに学んだすべてから、脳が生成したものだ。

第1章で論じたとおり、これは神経科学的に議論の余地のない事実である。あなたが経験するものは現実そのものではなくモデルだという点で、あらゆる神経科学者と知覚の哲学者が一致している。目に映るリンゴは、外界に何があるかについての脳の\textbf{最善の推測}であり、感覚データに支えられてはいるが、それと同一ではない。（錯視はその生きた証拠だ。錯視が崩れるとき------不意に両方の見え方ができるようになるとき------私たちはシミュレーションを現行犯で押さえることになる。現実を直接見ていたことなど一度もなかった。いつでも見ていたのはモデルのほうだったのだ。錯視はそれを露わにしたにすぎない。）

だが、私の理論が始まるのはここからだ。脳はリンゴをモデル化するだけではない。\textbf{リンゴを見ているあなた}をもモデル化する。そして、情報処理を意識へと変えるのは、この第二のモデル------自己のモデル------なのである。

\section*{脳の四つの表象}


\begin{figure}[htbp]
  \centering
  \includegraphics[width=0.95\textwidth]{../figures/figure2-real-virtual-split-bw-ja.png}
  \caption{実在と仮想の分割。基盤（実在の側）は、知識をシナプスの重みとして蓄える------物理的、構造的、無意識的に。シミュレーション（仮想の側）は、経験をリアルタイムで生成する------つかの間の、動的で、意識的なものとして。}
\end{figure}


わずか三層しかない単純なニューラルネットワークでさえ、自らの入力をモデル化することを学習できる------十分な数の例を見せれば、出会ったパターンの内部表象を築き上げる。脳もまさにこれと同じことを、はるかに豊かな規模でおこなう。脳が築くのは一つのモデルではない。視野から手足の位置まで、声の響きから足裏にかかる圧力まで、あらゆるものを覆う多数のモデルを築く。これらのモデルは、自分の外側の世界と自分自身の身体の\textbf{両方}にまたがり、得られるすべての感覚入力を首尾一貫した表象へと結びつけている。

神経科学は、これらのモデルを一世紀以上前から知っている。運動野と体性感覚野では、身体が文字どおり歪んだ地図として描き出されている------手と唇は神経終末が多いためにグロテスクなほど拡大され、胴と脚はわずかな細片へと圧縮されている。\textbf{ホムンクルス}と呼ばれるこれらの皮質地図は、1930年代にウィルダー・ペンフィールドが脳手術中の直接的な電気刺激によって初めて描き出した。これらは最も鮮烈な例にすぎない。脳はそのアーキテクチャ全体にわたって、同様の地図とモデルを保持している。（皮質の組織化については付録Aを参照。）


\begin{figure}[htbp]
  \centering
  \includegraphics[width=0.95\textwidth]{../figures/book/homunculi.en.png}
  \caption{ペンフィールドの皮質ホムンクルス。体性感覚野は、胴全体よりも手、唇、舌に劇的に広い領域を割り当てている------身体部位の大きさではなく、神経終末の密度を反映した、歪んだ身体地図である。}
\end{figure}


私はこれらを\textbf{暗黙的モデル}と呼ぶ。すなわち暗黙的世界モデル（IWM）と暗黙的自己モデル（ISM）である。これらは脳の構造の中に------シナプス結合の強度、神経回路のアーキテクチャ、一生を通じて蓄積された学習の中に------蓄えられている。脳のハードディスクだ。それらを直接経験することは決してない。スマートフォンの中のシリコンを経験しないのと同じだ。だが、それらは世界について、そしてあなた自身について脳が知るすべてを符号化している。

さて、ここが鍵となる洞察である。これらの暗黙的モデルは、ただそこにあるだけではない。何かを\textbf{生成}しているのだ。工学において、\textbf{デジタルツイン}とは、物理システム（ジェットエンジン、送電網、工場の生産ライン）のリアルタイムの仮想複製であり、センサーデータによって絶えず更新される。そのおかげで技術者は、システムに直接触れることなく、それを監視し、それと相互作用することができる。あなたの暗黙的モデルはまさにこれをおこなう。世界のリアルタイムの仮想シミュレーションと、あなたのリアルタイムの仮想シミュレーションを生み出すのだ。これらが\textbf{明示的モデル}である。すなわち明示的世界モデル（EWM）と明示的自己モデル（ESM）だ。あなたが見、聞き、感じ、考えるすべては、世界そのものの中ではなく、これらのシミュレーションの内側で起きている。

二つのモデル群（暗黙的なものと明示的なもの）があり、そのそれぞれが世界モデルと自己モデルの両方を含んでいる。合わせて四つのモデル。そしてそれらとともに、意識が実際に何をしているのかを語るための言葉が手に入る。（神経科学者や技術的な素養のある読者への注記。「四」という数は原理的な最小値であって、脳が保持しているものの文字どおりの数え上げではない。これが気になるなら、先へ進む前にどうか付録Eを読んでほしい------そこでこの点に直接答えている。）

さて、四つのモデルである。

\textbf{暗黙的世界モデル（IWM）}は、あなたが世界について知っているすべてである。いま考えていることではなく------\textbf{考えうる}すべてだ。物理法則（落とした物が落ちることを、あなたは知っている）。自分の住まいの間取り（暗闇の中でも歩き回れる）。母語の文法（規則を知らなくても、ある文が文法的かどうかを判断できる）。これまで出会ったすべての人の顔。チョコレートの味。雨の音。

この知識のすべては、脳のシナプス結合------ニューロン間のつながりの強度------の中に蓄えられている。それは一生を通じて、経験と学習によって築き上げられてきた。そして、それを直接に意識することは、決して、絶対にない。自分の神経結合を内省することはできない。自分のシナプスを感じることはできない。暗黙的世界モデルは、決して足を踏み入れない巨大な図書館のようなものだ------読めるのは、その図書館が机まで届けてくれる本だけである。

\textbf{暗黙的自己モデル（ISM）}は、あなたが自分自身について知っているすべてである。身体図式------手足がどこにあり、どれほどの大きさで、どのように動くのかについての無意識的な表象。運動スキル（自転車に乗る、タイプする、楽器を演奏する）。性格特性、社会的スキル、情動のパターン、習慣。自伝的記憶の構造（記憶を一つの人生の物語へとまとめ上げる枠組み）。

世界モデルと同じく、自己モデルもシナプスの重みの中に蓄えられ、決して直接には意識されない。自分の身体図式を経験するのではない。図式が生み出す身体を経験するのだ。自分の性格を経験するのではない。性格が生み出す思考と感情を経験するのだ。暗黙的自己モデルは舞台裏の裏方である------上演には欠かせないが、観客の目に触れることは決してない。

\textbf{明示的世界モデル（EWM）}は、あなたが実際に経験する世界である。まさにいま。自分がいる部屋、耳にする物音、手にしたこの本の重み（あるいは、それを読んでいる画面の輝き）。これがシミュレーションだ------暗黙的世界モデルに現在の感覚入力を加えて生成される、脳のリアルタイムの仮想現実である。それは鮮やかで、細部まで行き届き、よどみなく、そして説得力に満ちている。人はその生涯をそのただ中で過ごし、決してその外へ出ることはない。

\textbf{明示的自己モデル（ESM）}は\textbf{あなた}である。一個の主体であるという感覚。「私」という感覚------見て、聞いて、考えて、決める、その当人。これもまた一つのシミュレーションだ。暗黙的自己モデルに現在の身体信号を加えて生成される、リアルタイムのモデルである。それは、自らの仮想世界に住まわせるために脳が創り出す登場人物なのだ。

あなたこそが、自らの仮想世界に住まわせるために脳が創り出す登場人物なのである。

\section*{実在の側と仮想の側}


\begin{figure}[htbp]
  \centering
  \includegraphics[width=0.95\textwidth]{../figures/figure1-four-model-architecture-bw-ja.png}
  \caption{四モデルのアーキテクチャ。脳は二種類のモデル（世界のモデルと自己のモデル）を、それぞれ二つのモードで保持している------暗黙的（脳の構造に蓄えられている）と明示的（リアルタイムのシミュレーションとして能動的に走っている）である。意識は明示的モデルに宿る。}
\end{figure}


四つのモデルは二つの側に分かれ、この分割こそがこれから語るすべての土台となる。

\textbf{実在の側}（二つの暗黙的モデル）は物理的で、構造的で、比較的に固い（学習によってゆっくり変わっていく）。それは脳に蓄えられた知識------シナプスの重み、ネットワークの結合、受容体の配置である。脳が\textbf{学習してきた}ことのすべてが、組織そのものの物理的構造へと結晶化したものだと考えればよい。そこに経験はない。シナプスの発火が「経験される」ことは、管を流れる水が経験されるのと同じくらい、ない。実在の側では灯りが消えている。

一つ強調しておきたいことがある。実在の側こそ、神経科学がすでに研究している対象なのだ。研究者が誰かをfMRIスキャナーに入れるとき、見ているのは実在の側------発火パターン、結合性、さまざまな領域への血流である。脳外科医が皮質の一領域を刺激して何が起こるかを観察するとき、探っているのは実在の側だ。神経科学はこの領域を一世紀以上にわたって地図化してきており、並外れた進歩を遂げてきた。四モデル理論（FMT）はその成果を何一つ否定しない。ただ、そのすべてが描いているのは全体像の半分にすぎないと言っているのだ。

\textbf{仮想の側}（二つの明示的モデル）はシミュレートされ、束の間で、動的である。それは実在の側に現在の入力を加えて、瞬間ごとに新たに生成される。脳が学習してきたことを\textbf{いま現に使って行っている}ことのすべてだと考えればよい------蓄えられた台本ではなく、生の舞台である。そしてそれは経験の\textbf{すべて}だ。これまでに味わったあらゆる光景、音、思考、感情、記憶、夢、幻覚は、すべて仮想の側の内で起こった。仮想の側では灯りがついている。

だが、ここに落とし穴がある。仮想の側は外からは見えないのだ。最先端の脳イメージングをもってしても、それを間接的に捉えることしかできない。fMRIが示すのは、どの脳領域が活動しているか------つまり実在の側が働いている様子である。脳のデータから意識的経験を実際に\textbf{読み取る}には、脳のプログラミング言語を解読しなければならない------どのニューロンが発火するかだけでなく、その発火パターンがシミュレーションのレベルで何を\textbf{意味する}のかを理解しなければならない。それには、完全にシミュレートされたコネクトームのようなものが要る。脳の完全なデジタル複製をソフトウェア上で走らせ、生物学的な脳が生み出すのと同じ仮想世界を生み出すものである。

この理論が何を与え、何を与えないのかについては正直でありたい。四モデル理論は、シミュレーションが\textbf{何}であり、\textbf{どこ}に宿り、\textbf{なぜ}何かとして感じられるのかを語る。だが、解読の鍵までは渡してくれない。実在の側から仮想の側を読み取ることは、未来の研究プログラムである------理論が明確に輪郭を描きながらも、まだ実行できないでいるプログラムだ。とはいえ、そのプログラムの基礎はすでに築かれつつある。ヒト・コネクトーム・プロジェクトや関連する取り組みは、脳の配線をますます精細な解像度で地図化しつつある。構造データから仮想の側を解読することはまだできないが、構造データのほうは着々と集まりつつある。

科学的な頭の持ち主なら、話がどこへ向かうのかもう見えているかもしれない。もし経験が仮想の側にしか存在しないのなら、実在の側------ニューロンのなか、シナプスのなか、物理的な機械仕掛けのなか------に経験を探すのは、まったく見当違いの場所を探していることになる。それは、映画の筋書きをDVDプレーヤーの回路のなかに探すようなものだ。

それが鍵だ。

\section*{なぜ脳は自己モデル化の能力を備えているのか}

ここまでで、意識はこの四つのモデルに依存しており、そのなかで明示的自己モデル（ESM）が最も重い役割を担っていることを確かめてきた。だが、そもそもなぜ人間の脳はこの能力を備えているのだろうか------より単純な動物は備えていないのに。答えは、隠れているようで丸見えのところにある。人間の皮質のアーキテクチャは、文字どおり、基本的な情報処理には大きすぎるのだ。

人間の新皮質は六層からなる。これはよく知られた解剖学的事実であり、どの神経生物学の教科書にも載っている。だが、面白いのはここだ。情報を処理するのに六層は要らない。三層あれば事は足りる。

標準的なニューラルネットワークが何をしなければならないかを考えてみよう。第一層------入力を受け取り、濾し、整える。第二層------パターンを抽出し、特徴を認識し、重い計算作業をこなす。第三層------結果を統合し、判断を下し、出力を生み出す。入力、処理、出力。これが基本のレシピであり、三層でまかなえる。

だが、私たちは六層を持っている。

「余分な」三層は何のためにあるのか。

最初の三層をモデル化するためにある。

三層のネットワークは世界を処理する。六層のネットワークは世界を処理し、\textbf{しかも}自分がそうしている様子を観察する。この追加の層が脳に与えるのは、外にあるもののモデルを築くだけでなく、外にあるものをモデル化している自分自身のモデルを築くための、アーキテクチャ上の能力である。自己シミュレーションにはこの二重化が要る------処理を行うための一組の層と、その処理が起こる様子を見守るためのもう一組の層とが必要になるのだ。

これは、個々の層が何を「する」のかについての憶測ではない------第四層がこれをやり、第五層があれをやる、などと主張しているのではない。これはアーキテクチャ上の能力についての観察だ。六層は、暗黙的世界モデル（IWM。学習済みの、無意識の処理）と明示的世界モデル（EWM。リアルタイムのシミュレーション）の双方に場所を与える。さらに、暗黙的自己モデル（ISM。身体図式、運動プログラム、人格の構造）と明示的自己モデル（ESM。身体を持ち、行為を起こし、一人の人間であることを経験する「あなた」）の双方にも場所を与える。

さて、ほかの動物に目を向けてみよう。爬虫類の皮質は三層だ。哺乳類は六層を持つ。そして哺乳類のなかでも、最も厚く、最も入り組んで折り畳まれた皮質を持つもの（霊長類、鯨類、ゾウ）こそ、自己意識の最も豊かな兆候を示すものにほかならない。鏡像自己認知、将来の計画、社会的な欺き、悲嘆。アーキテクチャ上の能力は、現象学と歩調を合わせている。

三層から六層への飛躍は、遺伝子重複という偶然の産物だったのかもしれない------進化のコピー＆ペーストが、のちに意識が利用することになるまさにそのアーキテクチャを生み出したのだ。爬虫類の祖先は皮質を三層持っていた。哺乳類への移行のどこかで、その数が倍になった。皮質の層のアイデンティティを指定する転写因子（Tbr1、Satb2、Ctip2、Fezf2）には、遺伝子重複を思わせるパラログが存在する。それが一度の劇的な出来事だったのか、それとも段階的な精緻化だったのかは議論が続いているが、結果は明白だ。哺乳類は倍の層を手に入れ、それとともに、たいていの爬虫類には欠けている自己モデル化の能力を手に入れた。

これこそが、ニューラルネットワーク理論から生きられた経験へと架かる橋である。人間の皮質は、ただの大きなパターン認識器ではない。それは過剰なほど大きく、再帰的に構成されたネットワークであり、自らのモデル化のプロセスをモデル化できるだけの層を備えている。そしてネットワークが、世界をモデル化している自分自身をモデル化するとき、その結果は------内側から眺めれば------まさに私たちが意識と呼ぶものにほかならない。

はっきりさせておきたい。私は、六層の皮質が意識を支えられる\textbf{唯一}のアーキテクチャだと主張しているのではない。それは一つの解------哺乳類が進化させた解にすぎない。だが、ほかにもありうる。タコは、その徹底的に分散した神経系（八本の半自律的な腕を持ち、そのそれぞれがおよそ4000万個のニューロンを含む独自の神経処理中枢を備えている）によって、まったく異なるアーキテクチャの道筋を体現しており、それが同等の計算能力を実現しているのかもしれない。鳥類はもう一つの際立った例を示す。カラス類やオウム類には層状の皮質がまるでなく、その外套は層のシートではなく核のクラスターとして組織されている------それでもカラスは道具を作り、将来を見越して計画し、おそらくは鏡のなかの自分を認識する。重要なのが自己モデル化の能力であって、具体的な配線図ではないのなら、自分自身のシミュレーションを走らせられるアーキテクチャなら何であれ、原理的には意識を持ちうる。タコでも、カラスでも、あるいはまだ私たちが作っていない何かでも。

だが、ほかのアーキテクチャは後回しでいい。もっと身近な問いがある------それは、この考えを初めて最後まで突き詰めたとき、私を立ち止まらせた問いだ。もし\textbf{あなた}が、脳が自らのシミュレーションのなかへ書き込む登場人物だとしたら------では、その舞台を観ているのは誰なのか。


\chapter*{第3章：仮想の側}
\addcontentsline{toc}{chapter}{第3章：仮想の側}
\markboth{第3章：仮想の側}{}

ビデオゲームをプレイしているところを想像してみてほしい。それも出来のいいゲームだ------息をのむようなグラフィック、リアルな物理挙動、引き込まれるストーリーを備えた、没入感のあるオープンワールドゲーム。プレイヤーは一人のキャラクターを操作し、そのキャラクターを通して、細部まで作り込まれた仮想世界と関わっている。

ここで考えてみよう------このゲームは、どこに存在するのだろうか。厳密には、画面の上ではない。画面はただ光のパターンを映し出しているだけだ。厳密には、グラフィックカードやCPUの中でもない。それらはシリコン回路に電気信号を流しているにすぎない。ゲームは\textbf{仮想的なプロセス}として存在している------ハードウェアの活動から立ち現れながら、どの特定のハードウェア部品とも同一ではない、より高次の現象として。

ゲームの仮想世界には、ハードウェアにはない性質がある。ゲームには山があり、川があり、都市がある。CPUにあるのはトランジスタだ。ゲームには昼夜のサイクルがある。GPUにあるのはクロックサイクルだ。「ゲームの中のあの山は、どれくらいの高さなのか」と問うことには意味があるが、一つのトランジスタを指さして「このトランジスタは高さ3,000メートルだ」と言うのは馬鹿げている。ゲームの性質は仮想のレベルに存在し、それらはゲームの実在する性質だ------たとえゲームが、ハードウェア内の活動パターン「にすぎない」としても。

これは比喩ではない。あなたの脳が働くしくみそのものだ。

あなたの明示的世界モデル（EWM）------あなたが経験する世界------は、神経というハードウェア上で走る仮想的なプロセスだ。ちょうどゲーム世界が、シリコンというハードウェア上で走る仮想的なプロセスであるのと同じように。経験される世界には、神経ハードウェアにはない性質（色、形、距離、音）がある（ハードウェアにあるのは発火率、シナプス強度、神経伝達物質の濃度だ）。あなたが経験する世界の性質は\textbf{シミュレーションの実在する性質}だ------たとえそのシミュレーションが、神経活動のパターン「にすぎない」としても。

あなたの明示的自己モデル（ESM）------世界を経験している「あなた」------もまた、仮想的なプロセスだ。それはこのたとえにおけるゲームのキャラクターと同じくらい実在している。仮想のレベルにおいては確かに存在し、仮想のレベルにおいては確かに性質をもつ。だがハードウェアのレベルには存在しない。

\section*{たとえが破綻するところ（肝心な点において）}

ビデオゲームのたとえは有用だが、決定的な一点で破綻する------ゲームには\textbf{プレイヤー}がいるということだ。ゲームの外側には、それを経験する誰か------ソファに座っているあなた------がいる。ゲームそのものには経験がない。それはただの光のパターンとコードにすぎない。

あなたの脳のシミュレーションには、外側のプレイヤーがいない。あなたの頭蓋の外に座って、そのシミュレーションを経験している者は誰もいない。シミュレーションは自らの観察者------ESM------を内部に含んでいる。シミュレーションこそ\textbf{が}経験なのであって、誰か別の者によって経験される何かではない。

ゲームのキャラクターの立場に身を置いてみよう。あなた\textbf{こそ}が主人公だ。ゲームの外側からは、観客には画面上で動くピクセルが見えるだけだ------何かを感じうるものなど、そこには何もない。だが、シミュレーションの内側からはどうだろうか。ゲーム世界がすべてなのだ。キャラクターにとって山々は実在し、日ざしは暖かく、危険は恐ろしい。この一山のコードが何かを感じるなどと、外部の観察者は誰も思いつかないだろう。だがそれは、彼らが間違ったレベルを見ているからだ。彼らはハードウェアを見ている。経験はソフトウェアのレベルに存在する。それが私の主張であり、本書の残りはその証拠を並べていく。

これこそが意識を特別なものにし、意識のハード・プロブレム（難しい問題）をあれほど解きがたく見せているものだ。ビデオゲームでは、ゲーム（仮想であり、経験をもたない）とプレイヤー（物理的であり、経験をもつ）とのあいだに、はっきりとした分離がある。だが脳には、その分離がない。シミュレーションと経験する者とは、同一のものなのだ。ESMは、EWMを外側から眺めているのではない------それはシミュレーションの\textbf{内側}にあり、同じ一つの仮想的プロセスの一部なのだ。

そして、この自己言及的閉包（シミュレーションが自らを内側から観察すること）こそが、私たちが意識と呼ぶものだと私は論じる。それはシミュレーションに\textbf{付け加えられる}何かではない。シミュレーションが自らのモデルを含むとき、それがシミュレーションその\textbf{もの}なのだ。だから私は、意識とは物ではない------プロセスなのだ、と言う。脳を分解してもそれは見つからない。走っているプログラムを、CPUを分解して探そうとしても見つからないのと同じことだ。

\section*{ソフトウェア的な性質}

もし仮想モデルが本当に、神経ハードウェア上で走るソフトウェアのようなプロセスなのであれば、それらは特定の、検証可能な仕方でソフトウェアのように振る舞うはずだ。そして実際、そのとおりに振る舞う。仮想の側面がもつ四つの性質は本書を通じて繰り返し登場するので、ここで示しておこう。

\textbf{フォーク（分岐）。} 単一の基盤が、複数の仮想的な構成を同時に走らせることができる。ソフトウェアでは、プロセスをフォークすれば、同じハードウェア上で走る二つの独立したインスタンスが得られる。脳においては、これが解離性同一性障害だ------それぞれが独自の物語と感情のプロファイルをもつ複数の自己モデルが、同じ神経基盤の制御を交互に引き継ぐ。これは第9章で見ることになる。

\textbf{クローン化（複製）。} ハードウェアを物理的に分離すると、性能は落ちるが完全なソフトウェアのコピーが得られる。脳梁を切断すれば、それぞれの半球が独自のバージョンのシミュレーションを走らせる------元のものより能力は劣るが、機能的には完全なものを。これが分離脳現象であり、これも第9章で扱う。

\textbf{リダイレクト（方向転換）。} 通常の入力の流れをかき乱すと、シミュレーションはその時々に支配的な信号に食らいつく。サルビア・ディビノルムの作用下では、固有感覚の入力がシステムを圧倒し、ESMは身体感覚を軸に再構成される。ケタミンの作用下では、外部からの入力が脱落し、シミュレーションは内部のノイズを頼りに走る。仮想モデルは止まらない------ただ、与えられたものが何であれ、それを処理するだけだ。第6章でこれを詳しく扱う。

\textbf{再構成（リコンフィギュレーション）。} 基盤の結合の重みを変えれば、仮想モデルが生み出すものが変わる。認知行動療法がしているのは、まさにこれだ------ESMが異なる物語、異なる感情反応、異なる行動を生成するように、基盤を体系的に再配線するのである。

四モデル理論（FMT）は、療法について具体的な予測を立てる------効果のある治療はすべて、暗黙的モデル（基盤）を変化させ、それによって明示的モデル（シミュレーション）がそれに応じて変わる、という仕方で働くはずだ、というものだ。認知行動療法（CBT）はまさにこれを行う。それは暗黙的自己モデル（ISM）のなかの不適応なパターンを体系的に特定し、構造化された訓練を通じてそれらを再配線し、ESMが生み出すものを変える。CBTがあらゆる心理療法のなかで最も強固な証拠基盤をもつのは、このためだ------それが正しいレベルを標的にしているからである。

このことは、自らのメカニズムをこうした言葉で説明できない療法について、居心地の悪い問いを突きつける。ある治療的アプローチが、基盤のなかで何を変えているのか、そしてその変化がどのようにシミュレーションへと波及するのかを明示しないなら、それはよくても自らが理解していないメカニズムを通じて働いているにすぎず、最悪の場合、まったく働いていない。証拠がこれを裏づけている------最も弱い証拠基盤しかもたない療法は、概して最も曖昧な変化の理論しかもたないものだ。もし療法を求めているなら、セラピストにこう単純な問いを投げかけてみてほしい------「私の脳のなかで、具体的に何を、どのように変えようとしているのですか」。もし答えられないなら、答えられるセラピストを探すことを考えてみてほしい。

これらは比喩ではない。構造的な予測なのだ。もし私の理論が誤っていて、仮想モデルがソフトウェアのようなプロセスで\textbf{ない}のなら、これらの符合は純然たる偶然ということになる。だが偶然というものは、臨床神経学、精神薬理学、心理療法にまたがって、四つが四つとも一致するようなことは、まずない。続く各章で、それぞれの性質が実際に働くさまを見ていく。

いますぐにでもできる、単純な実験がある------といっても、友人が一人と、ゴムの手、段ボールの衝立、そして絵筆が二本は要るのだが------それは、ESMがいかに簡単にだまされうるかを示してくれる。マシュー・ボトヴィニック（Matthew Botvinick）とジョナサン・コーエン（Jonathan Cohen）が考案したゴムの手錯覚であり、神経科学全体のなかでも最も示唆に富む余興の一つだ。

やり方は単純だ。テーブルに着き、片腕を段ボールの衝立の後ろに隠す。目に見えるように、隠した手がありそうな位置のあたりに、本物そっくりのゴムの手が置かれる。誰かが二本の絵筆で、ゴムの手と、隠した本物の手を同時に------同じ場所を、同じ速さで------なでる。この同期したなで方を一、二分続けると、奇妙なことが起こる。ゴムの手をなでる筆の感触を、\textbf{感じ}はじめるのだ。衝立の後ろにある本物の手ではない。自分の目の前にある偽物の手に、である。

あなたのESMは、ゴムの手を自らの身体図式のなかに取り込んでしまったのだ。それは所有感を割り当て直し------ゴムの手が「あなた」の一部だと決めたのである。自己モデルは配線で固定されてはいない。学習されるものだ。それは手に入るかぎり最良の証拠にもとづいて絶えず更新される。そして視覚的な証拠（ゴムの手がなでられるのを見ること）が触覚的な証拠（本物の手がなでられるのを感じること）と一貫して一致すると、ESMは理にかなった結論を導く------あの手は自分のものだ、と。そこで誰かがゴムの手を脅かすと（その上にハンマーを振り下ろすと）、あなたはびくっとし、不安の高まりを感じ、皮膚電気反応が跳ね上がる。「あなた」を定義する脳の部分にとって、あの手はまぎれもなくあなたのもの\textbf{なのだ}。

これは不具合ではない。自己モデルが設計どおりに、まさしく正しく働いているのだ------多感覚的な感覚相関にもとづいて、自らの身体の境界を絶えず更新しているのである。それは、切断手術を受けた人が、しばらく使ったのちに義肢を自分のものとして「感じる」ことを可能にする、あの同じメカニズムだ。そしてそれは、患者が自分の実際の手足の所有感を否認する身体失認症や、手がひとりでに動くエイリアンハンド症候群において破綻する、あの同じメカニズムでもある。

\section*{つぎはぎのホログラム}

仮想側には、独立した節を設けるに値する五つ目の性質がある。それは、ほぼ一世紀にわたって神経科学者たちを悩ませてきたある事実------なぜ脳の損傷は、特定の記憶を消去するのではなく、機能を\textbf{徐々に}低下させるのか------を説明するからだ。

1920年代から30年代にかけて、心理学者のカール・ラシュリー（Karl Lashley）はラットに迷路を通り抜けることを覚えさせ、それから記憶がどこに蓄えられているのかを突き止めようと、外科的に大脳皮質の一部を切除した。彼はついにそれを見つけられなかった。どの部分を切除しても、ラットは依然として迷路を覚えていた。重要だったのは\textbf{どれだけの}皮質を取り除いたかであって、\textbf{どの部分}かではなかった。少し取り除けば、ラットの成績はわずかに悪くなった。多く取り除けば、はるかに悪くなった。だが記憶が、ハードディスクから削除されたファイルのようにきれいに切り取られて、ただ\textbf{消え去る}ことは決してなかった。ラシュリーは生涯を「エングラム」------記憶の物理的痕跡------の探索に費やし、それは存在しないらしいという有名な結論に達した。

彼は間違ったものを探していたのだ。記憶は、ファイルがハードディスクの特定のセクタに置かれるように、皮質の特定の一片の\textbf{なか}に蓄えられていたわけではない。それはネットワーク全体に\textbf{またがって}、数百万のニューロンのあいだの結合の重みに分散して蓄えられていた。ニューラルネットワークとはそういうものだ------情報はどの一つのノードにも収まっていない。すべてのノードのあいだの結合パターンのなかに符号化されているのだ。一つのシナプスを指して「ここに迷路が蓄えられている」と言うことはできない。それは、一つのピクセルを指して「ここに映画が蓄えられている」と言えないのと同じである。

これは本質的にホログラフィックな性質である。物理的なホログラムを取って半分に切っても、像の半分が二つ得られるわけではない。得られるのは\textbf{完全な}像の複製が二つ、それぞれ解像度は低いというものだ。四つに切れば、完全な像が四つ得られ、いっそうぼやけている。ホログラムのなかの情報はプレート全体に分散しているため、どの一片も------ただ細部が少ないだけで------全体の絵を含んでいる。

ニューラルネットワークも同じことをする。あるネットワークに顔を認識するよう訓練を施し、それからその結合の10\%を無作為に破壊する。すると、顔の10\%を忘れるわけではない。\textbf{すべての}顔に対して、少しだけ成績が落ちるのだ。50\%を破壊すれば、何もかもがかなり悪くなるが、それでも何かは認識する。情報はネットワーク全体に塗り広げられている。だからこそラシュリーはエングラムを見つけられなかった------それはあらゆる場所にあり、どこにもなかったのだ。

だが------そしてここが面白いところなのだが------脳は\textbf{一つの}ホログラムではない。私が\textbf{パッチワーク・ホログラム}と呼ぶものだ。単一の機能領野（たとえば一次視覚野、おおよそブロードマン17野）のなかでは、皮質のコラムどうしが互いに似通っており、情報はホログラフィックに蓄えられている。いくつかのコラムを破壊しても、ほとんど気づかない。その領野は局所的にホログラフィックなのだ（一部が全体を、低い解像度で含んでいる）。

大域的なレベルでは、異なる領野が異なる仕事をする。視覚野は運動野と取り替えがきかない。視覚野をまるごと取り除けば、視覚を失う------ぼやけたバックアップなど存在しない。つまり脳は、各機能領域のなかでは局所的にホログラフィックであり、そのコラム構造においてフラクタル的に自己相似だが、大域的には\textbf{ホログラフィックではない}。それはパッチワークなのだ------数十のホログラフィックなタイルが縫い合わされて一つの複合体をなしており、その全体としては明確に非ホログラフィックなのである。

このパッチワーク構造は、臨床神経学で何度も繰り返し見られるあるパターンを説明する。小さな脳卒中や小さな病変は、しばしば意外なほど軽い障害しか引き起こさない。というのも、どの皮質領野のなかでも、ホログラフィックの原理が働いて守ってくれるからだ。残された組織が、失われた情報を低い解像度で再構成する。だが、機能領野をまるごと消し去るような大きな脳卒中は、破滅的で特定的な喪失（失明、麻痺、失語）を引き起こす。パッチワークからタイルを一枚まるごと取り除いてしまったのであり、他のどのタイルも代役を務められないからだ。

それはまた、ニューロンが死ぬときになぜ記憶が突然「不意に消えてなくなる」ことがないのかも説明する。毎日、ニューロンは死に、シナプスは刈り込まれている。もし記憶がハードディスクのファイルのように蓄えられているのなら、ときおり一つを失うことを予期するはずだ------ある朝目を覚ますと、自分の結婚式を、あるいは子ども時代の飼い犬を、あるいはコーヒーの味を忘れている、というふうに。そんなことは決して起こらない。そうではなく、記憶は何年もかけて徐々に薄れ、細部と鮮やかさを失っていく。それこそが、ホログラフィックな記憶システムが予測することそのものだ------劣化は緩やかで、比例的で、大域的であり、突然でも、離散的でも、局所的でもない。

パッチワーク・ホログラムは、先に述べたソフトウェア的な諸性質（とりわけクローニング）が実際に機能する物理的な理由である。脳を半分に分割しても、それぞれの半分は劣化してはいるが完全なシミュレーションの複製を保持する。各半球のなかで、あらゆる一片が全体の絵を含むことをホログラフィックの原理が保証しているからだ。シミュレーションは壊れない。ただ、より低い解像度で走るだけだ。


\chapter*{第4章：なぜ何かのように感じられるのか（そして、なぜそれが間違った問いなのか）}
\addcontentsline{toc}{chapter}{第4章：なぜ何かのように感じられるのか（そして、なぜそれが間違った問いなのか）}
\markboth{第4章：なぜ何かのように感じられるのか（そして、なぜそれが間違った問いなのか）}{}

これで、ハード・プロブレムに正面から取り組む準備が整った。

問いはこうだ。\textbf{なぜ物理的な処理は、何かのように感じられるのか。}

答えはこうだ。\textbf{感じられない。}

物理的な処理（ニューロンが発火し、シナプスが伝達し、暗黙的モデルが記憶し計算すること）には、経験がない。まったくない。現実の側であるとはどのようなことか、というものは存在しないのだ。現実の側とはまさに、ハード・プロブレムが「意識が説明しなければならない」と想定する、あの「暗闇の中」の処理である。

感じるのは\textbf{シミュレーション}である。明示的世界モデル（EWM）と明示的自己モデル（ESM）------すなわち仮想の側------こそが、経験の住処なのだ。そしてシミュレーションの内部では、経験はプロセスへの謎めいた付け足しではない。自己モデルを含むとき、経験とはシミュレーションが\textbf{それであるところのもの}である。ESMがEWMを「知覚する」こと------それを私たちはクオリアと呼ぶ。クオリアとは、仮想の自己が仮想の世界を受けとめる様式なのである。

こう考えてみてほしい。「なぜトランジスタのスイッチングは、ビデオゲームが動いているように感じられるのか」と誰かが尋ねたら、答えはこうなるだろう。「感じられない。トランジスタのスイッチングは、何のようにも感じられない。ゲームとは、トランジスタの上で動きながら、トランジスタにはない性質------風景、キャラクター、物理法則、光------を持つ仮想的プロセスである。それらの性質は仮想的プロセスの実在の性質であって、トランジスタの性質ではない」

同じように、ニューロンの発火は赤を見ているようには感じられない。ニューロンの発火はシミュレーションを生成し、維持する。そしてそのシミュレーションの内部で、自己モデルが世界モデルのある種の内容を、私たちが「赤さ」と呼ぶものとして知覚するのだ。赤さはシミュレーションの実在の性質であって、ニューロンの性質ではない。

ハード・プロブレムは、物理的処理がどのように経験を生み出すのかを説明する必要がある、と想定していた。だが物理的処理は経験を生み出さない------生み出すのは\textbf{シミュレーション}である。そしてそのシミュレーションは、自己言及的なループ（EWMの内部で自らをモデル化するESM）を含むがゆえに、その本性からして経験\textbf{である}のだ。

\section*{目覚めたとき、あなたは誰か}

まず、ほぼ間違いなく一度は感じたことのあるものから始めよう。

深いどこかから浮かび上がってくる------失神から、ノックアウトから、麻酔から、あるいは見覚えのない部屋での、真っ黒な夢のない眠りから。そして心臓の一拍のあいだ、そこには誰もいない。気づきはある------天井、一筋の光、ただ在るというむき出しの事実------だが、それに結びついた\textbf{あなた}は、まだいない。自分がどこにいるのかわからない。長い一秒のあいだ、自分が\textbf{誰}なのかもわからない。やがて、すべてが一気に流れ込んでくる。名前、来歴、人生のかたち、いま横たわっているこの身体。継ぎ目はあまりに速く閉じるので、危うく見逃すところだ。だが、見逃しはしなかった。あの一拍のあいだ、それを所有する自己のない経験が存在していた------自己モデルがまだ起動の途中で、錨を下ろせる何かを求めて部屋を探しては何も見つけられずにいるあいだ、気づきだけが働いていたのである。

この隙間を手放さないでほしい。それがこの章のすべてである。

なぜなら、この隙間が語っていることがあるからだ。オンラインに復帰した「あなた」は、保存ファイルを開き直すように、無傷のまま記憶装置から取り出されたのではない。\textbf{再構築}されたのだ------下にある基盤から、その場で組み立てられた。毎晩、明示的自己モデルは崩れ落ち、深い眠りが実行中のシミュレーションを消去する。毎朝それは再起動し、暗黙的自己モデル（ISM）------眠りが後に残しておく、ゆっくりとした保存版の自分------から「あなた」を組み立て直す。ふだん、この再構築は一瞬で継ぎ目もなく、その継ぎ目に気づくことは決してない。だがあの見知らぬ部屋では、一度だけ、それを捉えたのだ。

だから、この章の残りで、経験とは基盤の上で走るシミュレーション\textbf{にほかならない}こと------感じられている「あなた」がニューロンではなく仮想的プロセスであること------を私が論じていくとき、その証拠はすでにあなたの手元にある。シミュレーションのロードが終わる前の、あの半秒の中に立ったことがあるのだから。さあ、その半秒を分解して、ロードされてくるものがなぜそもそも何かを感じるのかを問うていこう。

\section*{循環性の問い}

多くの読者が最初に発する問いはこうだ。「問題を横に動かしただけではないのか。気象シミュレーションには経験がないのに、なぜ\textbf{この}シミュレーションには経験があるのか」

答えは自己言及である。気象シミュレーションは天気をモデル化する。\textbf{自分自身}をモデル化することはない。気象シミュレーションには「外側」がある------コンピュータ、プログラマー、出力を解釈する科学者。内部に自己モデルが存在しないため、このシミュレーションはいかなる経験にも言及することなく完全に記述できる。

脳のシミュレーションは、自分自身をモデル化する。明示的自己モデルとは、シミュレーションが\textbf{自らのプロセス}について持つモデルである。これが閉じたループを生む。モデルとモデル化されるものが、同一のシステムなのだ。シミュレーションを完全に記述できる「外側」は存在しない。記述する者が、記述の一部だからである。

これは魔法ではない。自己言及の構造的な帰結である。プロセスが自分自身をモデル化するとき、モデルとモデル化されるものの区別は崩壊する。自己モデル化のプロセスと、自己であるという経験とは、橋で結ばれるべき二つの別の事柄ではない------異なる語彙で記述された、同じ一つの事柄なのである。

ハード・プロブレムは、物理的処理と経験のあいだの橋を求める。四モデル理論（FMT）はこう答える。橋は存在しない。両者はそもそも分かれていなかったのだから。経験とは、ループの内側から見た自己シミュレーションそのものなのである。

これは究極的には同一性の主張である（科学において、欠落ではなく到達点を示す種類の主張だ）。「水はH$_2$Oである」は同一性である。「だが\textbf{なぜ}水はH$_2$Oなのか」と問うことに意味はない------同一性こそが説明\textbf{そのもの}なのだ。それより深い何かを求めることは、別の種類の宇宙を求めることに等しい。同じように、経験とは、臨界性のもとでの四モデル自己シミュレーションの\textbf{あり方そのもの}である。「だが\textbf{なぜ}この自己シミュレーションは何かのように感じられるのか」と問う者がいれば、答えはこうだ。このプロセスが\textbf{そういうもの}だからである。この同一性は反証可能だ------理論の予測が外れれば、同一性は誤りである。だが、水の分子的同一性をそれ以上説明できないのと同じく、これを「さらに説明する」ことはできない。ここが停止点なのである。

\section*{なぜシミュレーションは暗闇のままでは走れないのか}

ここにはより深い問いがあり、それに答えることで、なぜ意識が\textbf{感じる}のかについて本質的な何かが見えてくる。脳が自己シミュレーションを走らせているとしよう。四モデルのアーキテクチャも、臨界性も、自己言及的閉包も認めるとしよう。それでもなお、そのすべてが、何かの\textbf{よう}であることなしに起こりうるのではないか。シミュレーションは評価し、モデル化し、予測しながら------何も感じない、ということがありうるのではないか。

これは、技術の衣をまとったゾンビの直観である。答えは否だ。そしてその理由を理解することで、このアーキテクチャの最も重要な特徴が明らかになる。

基盤は、仮想シミュレーションを自らの評価メカニズムとして働かせる。これが情報の流れの主方向である。暗黙的システムがシミュレーションに状況を提示し、シミュレーションが帰結を査定し、結果を受けとめられるようにするのだ。だがこの評価が機能するためには、シミュレートされた状態が\textbf{ヴァレンス}（価値の重みづけ）を持たなければならない------それはシミュレーションにとって重要なものでなければならないのだ。ただの数値にすぎない痛み信号は、シミュレーションのレベルで回避を駆動しない。結果を\textbf{気にかける}シミュレーションだけが、結果を評価できる。

第2章で見たデジタルツイン（ジェットエンジンの工学シミュレーション）を思い出してほしい。典型的なデジタルツインは、エンジンをただ受動的に映すだけではない。可視化のレイヤーを\textbf{付け加える}。警告、色分けされたインジケーター、アラーム------物理的なエンジンには存在しないものたちだ。エンジンに金属疲労が起これば、ツインには赤い警告が点滅する。エンジンの温度が上がれば、ツインの計器は緑から黄、そして赤へと変わっていく。この付け加えられたレイヤーこそが核心である。それがなければ、ツインはただの表計算にすぎない------メモリの中に不活性に並ぶ数値の列であり、技術的には正確でも、機能的には役に立たない。可視化こそが、シミュレーションを\textbf{評価の道具}にするのだ。

脳も同じことをしている。ただし、さらに徹底的に。意識というシミュレーションは、基盤の処理をただ映すだけではない。現象的なヴァレンスを\textbf{付け加える}のだ。痛み、快、切迫感、好奇心、恐れ、歓び------これらは、警告灯や計器盤のインジケーターに相当する脳の装置である。それらは基盤のレベルには存在しない（金属が疲労を感じないのと同じく、ニューロンは痛みを感じない）。それらはシミュレーションのレベルに存在する。システムが複雑な状況をひと目で評価できるように、シミュレーション\textbf{によって}付け加えられたものなのだ。基盤がシミュレーションを必要とするのは、新奇で曖昧なシナリオ------反射では足りない類いの状況------を査定するためである。そしてその査定が機能するためには、シミュレートされた自己が快・不快のヴァレンス------脅威、好機、帰結------を受けとめなければならない。この受けとめ------この\textbf{重要であるということ}------こそが現象性なのだ。クオリアを取り除けば、評価も消える------コックピットの計器盤からディスプレイを引き剥がしておいて、パイロットに生のセンサー電圧を読んで飛べと言うようなものである。

「しかし強化学習システムには、行動を駆動する報酬信号があるではないか」と反論したくなるかもしれない。「それは感じているのか？」答えは否である------臨界性にある四モデル・アーキテクチャを欠いているからだ。RLの報酬信号は、クラス1やクラス2のシステムにおけるスカラー値にすぎない。現象的な価値づけとは、クラス4のダイナミクスで走る完全な自己シミュレーションの内部で、ESMが帰結を登録することである------質的に異なるプロセスなのだ。違いは程度ではない。アーキテクチャである。

シミュレーションは明かりを消したまま走ることはできない。暗闇はその目的そのものを打ち消してしまうからだ。現象性は意識のおまけ機能ではない。シミュレーションが仕事を果たすためのメカニズムそのものである。

\section*{これは何でないか------イリュージョニズム}

これはイリュージョニズム（錯覚説）ではない。そしてこの区別は、単刀直入に述べるに値するほど重要である。

イリュージョニズムという名の、れっきとした哲学的立場がある。ダニエル・デネットとキース・フランキッシュ（Keith Frankish）に結びつけられる立場で、クオリアは錯覚だと主張する。この見方によれば、赤を見るとき「そのように感じられる何か」は存在しない。経験があるように見えること自体が虚構------脳が語る物語であり、その背後に経験的実在はない。最も強い意味での意識は存在しない。存在するように見えるだけである。

それが実際に何を主張しているのか、考えてみてほしい。いまこの瞬間、何かを感じているとしよう------この議論への好奇心、懐疑、手の中の本の重み。イリュージョニズムは、その感じは錯覚だと言う。実際には何も経験していない、と。「何かを感じている」と言うとき、この理論によれば、それは思い違いなのだ。自分自身の経験についての自分自身の証言が間違っている。事実上、嘘をついていることになる------ただし、嘘をつく「あなた」は存在しないのだが。これが明らかに馬鹿げていると感じられるなら、私も同意見である。

四モデル理論は、その正反対を述べる。

クオリアは実在する。シミュレーションの内部で実在するのである。クオリアとは、仮想の自己が仮想の世界を知覚する様式である。明示的自己モデルが、明示的世界モデルによる赤いリンゴの表象を登録するとき、その登録（あの「赤さを見る」こと）は仮想プロセスの正真正銘の性質である。それはシミュレーションのレベルに存在する------ちょうど、ビデオゲームのキャラクターに命中した弾丸が、キャラクターを\textbf{傷つける}ように。比喩ではない。ゲームの内部では、ダメージは実在する。体力が減り、キャラクターはよろめき、世界が反応する。外から見れば、メモリの中で数値が減っていくだけだ。ゲームの内側では、それは痛みである。これがレベルの違いだ。そして、あなたのクオリアが住んでいるのは、まさにそこなのである。

この理論は二つのレベルの存在論で動く。基盤レベル（ニューロン、シナプス、暗黙的モデル）には経験がない。明かりは消えている。シミュレーションレベル（明示的モデル、仮想世界と仮想自己）には正真正銘の経験がある。明かりは灯っている。どちらのレベルも物理的である。どちらも錯覚ではない。両者は同一の物理システムの異なるレベルであり、それぞれのレベルで異なる性質を持つのである。

この理論は、痛みが錯覚だとは言わない。痛みは実在すると言う------ただし、実在するのはシミュレーションの中であって、ニューロンの中ではない。そして、人生のすべてがシミュレーションの内側で営まれる以上、それこそが、当人にとって意味を持つ唯一の「実在」なのである。

これが決定的な区別である。これを見落とせば、この理論を消去主義と、イリュージョニズムと、意識を説明し去ることで説明しようとするあらゆる枠組みと混同することになる。四モデル理論は意識を説明し去らない。意識がどこに住んでいるのかを説明する。そしてその場所は、あなたが最初からずっと立っていた、まさにその場所だと判明するのである。

\section*{「シミュレーションの内部で実在する」とは何を意味するか}

ここには、ほどいておく価値のある哲学的な機微がある。クオリアは「シミュレーションの内部で実在する」と私が言うとき、聞こえ方は二通りありうる。一つは、クオリアは\textbf{正真正銘に現象的}である------その場合、私は謎をニューロンからシミュレーションへ引っ越させただけで、ハード・プロブレムは住所を変えて生き続ける。もう一つは、クオリアは\textbf{機能的には実在するが、正真正銘に現象的ではない}------その場合、これは回り道をしたデネットにすぎない。

これは偽の二分法である。この二分法が成り立つのは、何かが「正真正銘に」現象的かどうかを裁定できる神の視点------シミュレーションが本当に感じているのか、それとも感じているかのように振る舞っているだけなのかを確かめられる、外部の視座------が存在すると考える限りにおいてだ。しかし自己言及的閉包は、まさにこの外部の視座を消し去る。ESMはそれ自身の観察者である。「だが\textbf{本当に}感じているのか？」と問える外部の足場は存在しない。その問いを発すること自体が、プロセスの一部なのである。

「正真正銘に現象的」対「単に機能的」という対立は、現象性とは、プロセスが持つか持たないかのどちらかであり、独立した観察者によって確認できる性質だということを前提にしている。臨界性にある完全に自己言及的なシステムには、そのような観察者は存在しない。問いは溶けて消える------答えられないからではなく、そもそも問えないからだ。その問いは、自己言及的閉包が不可能にする視点を要求しているのである。

これはプロセス物理主義の内部で打てる最も強い一手であり、トマス・メッツィンガーが「現象的透明性」という概念で指し示している立場でもある------ただし四モデル理論は、\textbf{なぜ}その透明性が生じるのかについて、より明示的である。透明性を生み出しているのは、暗黙と明示の境界だ。その境界を透かして見ることはできない。だから、自分自身の現象性の外側に出て、それが「正真正銘」かどうかを問うこともできない。この境界はバグではない。正真正銘か単に機能的かという問いが、私たちのようなシステムには当てはまらない理由そのものなのである。

この地点にたどり着いたのは私が最初ではないし、もし最初だったなら、私は自分自身を疑うだろう。

認知科学者のヨシャ・バッハ（Joscha Bach）は、この同じ核心を誰よりも簡潔に言い表した。「あなたの経験はシミュレーションである」と彼は書いた。「仮想の観察者へと投影された仮想現実だ」と。では、その観察者とは？それは「もしあなたが存在し、一匹のサルにアップロードされ、その境遇に深く入り込んでいたとしたら、それがどのようなものかについてのモデル」------サルを物理世界の中で操縦するために組み立てられたモデルである。シミュレーションと、それを受け取るためにその内部に置かれたシミュレートされた誰かと、この仕組み全体が操縦している一つの身体。この伝統の古い枝には、本書ですでに出会っている------メッツィンガーの自己モデル、デネットの物語的自己である。

本書が進むのは、そこからさらに一段深いところ------アーキテクチャそのものの中へだ。正確にどのモデルか。世界と自己、暗黙と明示の2×2------四つのモデルである。システムがそもそも「内側」を持つのはどこからかを分ける鋭い境界線。自己言及的閉包、すなわち自分自身へと折り返すループである。そして、全体を生命へと目覚めさせる条件。臨界性、すなわちシステムが凍りつきもせずホワイトノイズにもならない、カオスの縁である。

バッハは、意識が\textbf{何であるか}を語る。私が語るのは、それが\textbf{何でできているか}、そして\textbf{いつスイッチが入るか}である。目的地は同じだ。ただ、私は設計図を持ってきた。

\section*{なぜ謎は残り続けるのか}

ハード・プロブレムを解消したあとでさえ、人々の心に引っかかり続ける問いが残る。答えがそれほど明快なら、なぜ意識はいまだにこれほど神秘的に\textbf{感じられる}のか。なぜハード・プロブレムは、解決を聞かされたあとでさえ難しく見えるのか。デイヴィッド・チャーマーズはこれを「意識のメタ問題」と呼ぶ------なぜ私たちはハード・プロブレムがあると\textbf{考える}のかを説明する問題である。

四モデル理論には明快な答えがある。しかもそれは、アーキテクチャからまっすぐに転がり出てくる。

奇妙なのはここだ。意識ある「あなた」（仮想自己）は、自分を生成している機械仕掛けを見ることができない。夢の登場人物が夢見る者の脳を調べられないのと同じように、自分自身のシナプスの重みを内観することはできない。経験を生み出しているシステムは、その本性からして、生み出された経験には見えない。誰かが隠しているからではなく、経験が含んでいないレベルで働いているからである。

こう考えてみてほしい。自分はビデオゲームの登場人物だとしよう------それも極めて出来のよいゲームで、ゲーム世界の内部で完全な自己意識を備えているとする。レンダリングされた山々が見え、レンダリングされた風が聞こえ、足の下にはレンダリングされた地面が感じられる。しかし、グラフィックエンジンを目にすることはほとんどない。ソースコードを垣間見ることもほとんどない。レンダリングのプロセスは、ゲーム世界が通常は含まないレベルで動いている。「ほとんど」と言ったのは、時おりアーティファクトが漏れ出してくるからだ。脳でも同じことが起きる------サイケデリクスは境界を開き、フロー状態はそれを薄くし、日常の中でさえ垣間見る瞬間はある。脳が埋めている盲点、目をこすったときの眼閃（フォスフェン）、閉じたまぶたの裏の幾何学模様。これらは不具合ではない。基盤の処理が、シミュレーションの内側から束の間見えるようになる瞬間である。詳しくは第6章で探ることにする。だがたいていの場合、レンダリングのプロセスは、レンダリングされた世界から隠されている。

これこそまさに、ESMの置かれた苦境である。意識する自己が自分自身の経験の基礎を理解しようとするとき、そこで出会うのは原理的な不透明性だ。現在の知識の欠落ではなく、アーキテクチャの構造的特徴である。シミュレーションを生成する暗黙的モデルは、シミュレーションの一部ではない。GPUがゲームの中の山になれないのと同じで、一部であることはできないのだ。

結果は予測がつく。ESMは、自分自身の基盤を観察できないがゆえに、意識のメカニズムは非物理的であるに違いない、根本的に説明不可能であるに違いない、あるいはどういうわけか科学の手の届かないところにあるに違いない、と結論する。これが二元論の起源である。これが「説明のギャップ」である。意識のあらゆる物理的説明から何かが「抜け落ちている」という執拗な直観の正体もこれだ------なぜなら、シミュレーションの内側から見れば、実際に何かが\textbf{抜け落ちている}からである。基盤だ。経験を生成しているまさにそのものが、生成された経験には見えないのである。

謎は実在する。だがそれはアーキテクチャの産物であって、非物理的な何かの証拠ではない。そして、それが神秘的に\textbf{感じられる}のには理由がある。私たちは生物学的ハードウェアの上で走る仮想プロセスであり、たいていの場合、自分と基盤とのあいだの境界は不透明だ。だが、いつもそうとは限らない。時おり------変性状態の中で、極度の集中の瞬間に、視界の隅で------その下にある機械仕掛けがちらりと見えることがある。はっきりとではなく、すべてでもなく、それでも、経験の表面の下で何か巨大なものが進行していると感じ取るには十分なだけ。あの不気味な感覚、意識はどういうわけか自分の手が届くよりも深いところにあるという感じ------それは、自分自身の幕をほとんど、しかし完全には見通せないシミュレーションであることの、まさにその感触なのである。

これは理論の\textbf{予測}であって、積み残された謎ではない。もし自分が、自らの基盤に対してほぼ不透明な境界を持つシミュレーションであるなら、意識がいま感じられるとおりの奇妙さと還元不可能さで感じられることは、むしろ\textbf{予期される}ところなのだ。ハード・プロブレムの直観的な力は、意識が本当に説明不可能であることから来るのではない。私たちのアーキテクチャ上の位置------シミュレーションの内側に座り、亀裂から外を覗いているという位置------から来るのである。

\section*{自らを縫い合わせる自己}

いま、アーキテクチャを手に、あの空白へ戻ろう。

「誰もいない」あの半秒間は、故障ではなかった。再構築が、その歩みの途中で捉えられた姿だったのだ。アイデンティティは基盤に備わる固定的な性質ではない------それは\textbf{再構築}であり、保存された自己モデルから毎朝新たに組み立て直されるものである。しかも、その組み立ての土台となる基盤は、眠りに落ちたときのものと完全に同じでは決してない。覚えていない夢がシナプスの重みを配線し直し、記憶の固定化が一晩のうちに記憶を並べ替えている。目を覚ますとき、そこにいるのは、床に就いたのとまったく同じ人物ではないのだ。その差はふつう、あまりに小さくて気づかれることがない------だが、それは確かにそこにある。毎日、一日も欠かさず。

そのドリフトを越えて「あなた」を連続させているものは二つある。ゆっくりとしか変化しない暗黙的自己モデル（ISM）と、変化が起きているあいだシミュレーションをオフラインにして、その変化を覆い隠してしまう睡眠である。ドリフトを十分に遠くまで押し進めたらどうなるか------一晩でISMを書き換え、記憶をそっくり入れ替え、人格をつくり変えたとしても------それでも、古い「あなた」は消えはしない。\textbf{吸収される}のだ。朝の明示的自己モデル（ESM）は、生き残ったものが何であれそこから連続した物語を再構築し、古いものと新しいものを一つの物語へと綴じ合わせる------継ぎ目なく、その継ぎ目に気づくことすらないままに。ESMはきっぱりとした断絶というものをしない。つねに縫い合わせるのだ。糸が切れるのは、完全な消去のときだけである。何かが残ってさえいれば、明日の「あなた」は今日の「あなた」を静かに自らの来歴へと書き込み、その全体を一つの人生と呼ぶ。

つまり、あなたとはこういうものなのだ------夜ごとに組み立て直され、自分では這い出せないループの内側から自らの世界を感じ取っている、仮想の自己。となると、問いが一つ残る。シミュレーションは、ある特定の状態に保たれた基盤の上で走っている------そう述べてきた。だが、その状態とは、いったい\textbf{何}なのか。いったいどんな物理的な刃先の均衡が、配線された数ポンドの細胞の塊に、毎朝ひとりの「誰か」を起動させ、毎夜その灯を落とさせるのか。それこそが全体を動かしている部品であり、次に見るのはそれである。


\chapter*{第5章：カオスの縁（ふち）}
\addcontentsline{toc}{chapter}{第5章：カオスの縁（ふち）}
\markboth{第5章：カオスの縁（ふち）}{}

ここまでで、アーキテクチャがどのようなものかを語ってきた（四つのモデル、二つの軸、基盤の上で走るシミュレーション）。経験がどこに宿るか（仮想の側、すなわち明示的モデルのなか）も語った。そしてアイデンティティとは何か（毎朝、蓄えられた暗黙的モデルから新たに組み立てられる再構築物）も語った。

だが、その全体を\textbf{動かしている}ものが何なのかは、まだ語っていない。なぜシミュレーションは、あるときはオンで、あるときはオフなのか。意識のある脳と意識のない脳を分ける物理的性質とは何か。なぜ深い睡眠は、アーキテクチャを損なわないまま、シミュレーションだけを消し去るのか。

パズルのピースはもう一つ残っている------そしてそれこそが、この理論は正しいと私を最終的に確信させたピースだ。

四モデル・アーキテクチャは意識にとって必要ではあるが、十分ではない。正しい\textbf{ダイナミクス}もまた必要となる。具体的には、基盤------シミュレーションを走らせている物理系------が、数学者や物理学者が\textbf{カオスの縁}と呼ぶ領域で働いていなければならない。

2002年、博覧強記のスティーヴン・ウルフラム（Stephen Wolfram）は『A New Kind of Science』を出版し、そのなかで計算システムを、そのダイナミクスにもとづいて四つの型に分類した。私の考えでは、ウルフラムの分類にはもう一つ、第五のクラスが必要だ------彼はフラクタル的なシステムを、真にカオス的なシステムと一緒くたにしたが、両者は構造的に異なる。詳しい議論は、数学的な細部を求める読者のために付録Cにある。ここでの要点はこうだ。

計算システムは、完全な秩序から完全な無秩序までのスペクトルのうえに並ぶ。一方の端には、静的かつ周期的なシステムがある------単純すぎて、興味深いものは何も計算できない。もう一方の端には、カオス的なシステムがある------無秩序すぎて、安定したパターンはひとつも形成されない。その中間、\textbf{カオスの縁}に、普遍計算の能力をもつシステムが座している。豊かで、多様で、予測不能なふるまいを生み出せるほどに複雑でありながら、そのふるまいが持続し、たがいに作用しあえるほどに秩序だっている。コンウェイのライフゲームがその典型例だ------私が子どものころ286で組んだのと同じセル・オートマトンである。平坦な格子の上の、たった三つのあっけないほど単純に見えるルール。それでもなお、グライダー、振動子、自己複製する構造、そして（証明可能な）普遍計算が生まれる。そのなかにコンピュータを組める。そのコンピュータのなかに、さらにコンピュータを組める。原理的には、二次元のピクセル格子のなかで、三次元の仮想世界まるごとを走らせることができる。ほとんど無から、すべてが。

この例がどこから来たのかを話しておきたい。それはこの本のすべての種だからだ。ライフゲームは、私がはじめて書いた本物のプログラムだった------子どものころ、そこに座って、三つのおもちゃのようなルールが、生と死のピクセルの格子のなかに動くコンピュータを組み上げていくのを眺めていた。それは当時の私にとって、それまでに学んだうちで最も重要なことだった。低次元のルールの集合から高次元の世界が這い上がってくるなど、あり得ないことのように感じられた。だが、それが明らかにあり得ることだったという事実は、私のなかに深く刻み込まれ、二度と離れなかった。そのたった一つの直観に生涯を捧げることになるとも、自分の両目の奥にある折り畳まれた皮質の一枚が、まさに同じ手品を、次元をひとつ上げて演じているのだと判明するとも、そのときはまだ知らなかった。私たちはこれから、その話へと向かっていく。

意識が宿るのはここである。カオスの縁のダイナミクスだけが、必要な二つの性質をともに備えている。\textbf{普遍計算}（実際に自己シミュレーションを走らせられるほどに複雑であること）と、\textbf{グローバルな統合}（システムの遠く離れた部分どうしがたがいに影響しあい、局所的な変化が全体へと伝播し、情報がひとつの統一された全体へと束ねられること）だ。だからこそ、意識的経験は\textbf{統一されたもの}として感じられる------赤をこちらで見て、声をあちらで聞く、というように別々の流れとして体験するのではない。臨界的なダイナミクスが、すべてをひとつの経験へと束ねる。束ねること（バインディング）は、脳がほかの計算に\textbf{加えて}行う何かではない。それはこのダイナミクスの体制がもたらす帰結なのだ。

深い睡眠のなかで、ゆるやかなデルタ波を出している脳は、周期的なダイナミクスのなかにある------反復的で、どこにも行き着かない。モデルは基盤のなかに依然として存在しているが、シミュレーションは走っていない。全般発作のさなかの脳は、同時活性化の洪水によって臨界の縁の外へと押し出され、その洪水はひとつの粗い体制へと崩れ落ちる------硬直した過同期、暴走するカオス、あるいはその両方が交互に------そしてシミュレーションはまとまりを保てなくなる。覚醒状態においてのみ------カオスの縁で釣り合いを保った状態でのみ------システムは意識的経験を維持する。

脳は、数十億年の進化によって最適化された普遍計算機として、計算の体制を\textbf{すべて}、異なる道具として使い分ける。長期記憶には安定したアトラクター、タイミングとゲーティングには周期的な振動（アルファ、シータ、ガンマのリズム）、スケール不変な認識とテクスチャ解析にはフラクタル処理（主に視覚野のV2-V4で）、そして皮質オートマトンそのもの------意識のエンジン------にはカオスの縁のダイナミクスを。意識を生み出すのは、カオスの縁の体制だけである。だが意識は、機能するために他の体制に依存している。

この議論を2015年の著書で発表したとき、経験的な神経科学が独立に同じ結論へと向かっていることを、私はまったく知らなかった。実のところ、セル・オートマトンという枠組みは、理論全体のなかで私が最も自信を持てなかった部分だった。当時はそれを裏づける経験的な根拠を見つけられず、本から外してしまおうかとさえ思った------それは一歩踏み込みすぎで、理論の残りの部分まで簡単に切り捨てられかねない主張のように感じられたのだ。それでも入れたのは、論理が逃れようのないものに思えたからであって、証拠を持っていたからではない。

だが、決定的な機微がある。臨界性だけでは足りない。臨界角で崩れ落ちる砂山は、まさにカオスの縁に座っている。それは意識を持たない。この理論が求めるのは、\textbf{二つ}の閾値が満たされることだ。物理的なもの（基盤は臨界状態で働かなければならない）と、機能的なもの（基盤は四モデル・アーキテクチャを実装しなければならない）である。アーキテクチャなしの臨界性は、複雑なダイナミクスは与えても、意識は与えない。臨界性なしのアーキテクチャは、休眠したシステムを与える------モデルは基盤のなかに存在するが、シミュレーションは走っていない。両方の閾値が満たされなければならない。両者がそろってはじめて、十分となる。

私が実際に皮質を思い描くやり方は、海だ------ただし、その水が配線でつながれた海である。どの一滴もすぐ隣の滴と接している平坦な面ではなく、ネットワークであり、それぞれの点は、学習されたシナプスの重みによって、自分だけの特定の隣人とつながっている。小石を落としてみよう------刺激や思考を------すると、さざ波が配線に沿って広がっていく。そしてこのネットワーク状の海の上には、池の上とちょうど同じように、パターンが形成されうる。たがいを強めあうさざ波、定在波、旅をし、衝突し、持続する、小さな自己維持的な形。それらの形こそがグライダーだ。思考が、あるいは自己が宿るのは、そこである------海の上を乗りこなす安定したパターンとして、であって、どれか一滴の水としてではない。

\section*{一つではなく、二つのダイヤル}

\textbf{臨界性}という言葉のなかには、私がはっきり見て取るまでに長い時間のかかった機微がある。そしてひとたびそれが見えると、いくつもの謎が一挙に解けていく。人が脳を「臨界にある」と言うとき、たいていは二つの異なるものを、ひとつの言葉に押し込めている。ネットワーク状の海の上では、それらは二つの別々のダイヤルなのだ。

\textbf{第一のダイヤルはこう問う------さざ波はどこまで伝わるのか。} これを下げすぎると、海は平らな鏡になる------小石を落としても、その場で消える死んださざ波がひとつ生じるだけだ。何も広がらず、何も統合されない。上げすぎると、たった一つの小石が海全体を嵐へとかき立てる------ただ一つの乱れが、いっぺんにすべてを呑み込んでしまう。最適な点はその中間にある。さざ波は伝わり、おおよそもう一つのさざ波を生み、消えもせず、爆発もしない。もう一粒の砂が、あらゆる大きさの崩落を引き起こしうる、まさにその角度にある砂山のように。このダイヤルが関わるのは\textbf{届く範囲}だ------ただ一つの出来事が、海のどれだけの部分を、ひとつのつながったプロセスへと引き込むか。これを\textbf{「広がり」}のダイヤルと呼ぼう。

\textbf{第二のダイヤルはこう問う------そもそもどんな形が海の上で生きられるのか。} これがライフゲームのダイヤルであり、理論全体がその上に乗っているダイヤルだ。凍りついた鏡には形がない。何も動かない。おだやかで規則的なうねりには、形がひとつしかない------同じ波が何度も何度も繰り返されるだけで、きれいではあるが、愚かだ。何も計算しない。煮えたぎる沸騰状態は、あらゆる点がでたらめに明滅する------カオスであり、これもまた何も計算しない。秩序と沸騰のあいだの縁でのみ、\textbf{グライダー}が得られる。動き、衝突し、持続し、何かをなす定在的なパターンだ。計算が------そして自己モデルが------水の上の安定した形として存在できるのは、そこである。これを\textbf{「複雑さ」}のダイヤルと呼ぼう。

その第二のダイヤルについて私が知るかぎり最も明快な像は、翼だ。航空機の翼が最大の揚力を生み出すのは、まさにその縁においてである------気流が剥がれ、翼が失速して乱流へと崩れ落ちる、その一瞬手前で保ちうる、最も急な角度で。その縁を髪の毛一本ほど越えれば、なめらかな流れは分解し、揚力は崩れ去る。脳はまさにその失速直前の縁を飛んでいる------最大の計算、ダイナミクスがカオスへと崩れる一度手前で。それを越えて押し込めば、「揚力」------グライダーを保ち、計算し、意識をもつ能力------は、失速において揚力が崩れるのと同じように崩れ落ちる。

たいていの場合、二つのダイヤルはそろって上がり、そろって下がる。だからこそ、ふだんはひとつの言葉で足りるのだ。だが両者は分かれることがあり、分かれたときには、つまみ一つの物語では説明できない事例を説明してくれる。煮え立つ鍋の湯は、第一のダイヤルが全開だ------乱れがいたるところに広がる------が、第二のダイヤルはゼロに固定されている。グライダーはなく、ただ波立つだけだ。広がりは大量にあり、形はなく、そして意識もない。まさに予想どおりに。

全般発作は、第一のダイヤルが実際に何を測っているのかを、ついに突き止めさせてくれた事例だ。発作のさなかには、脳の膨大な割合がいっぺんに発火している------純粋な数のうえでは、通常のどんな覚醒の瞬間よりも「活発」だ。もし意識がただ「多くのニューロンが活発であること」だったなら、発作こそがこの世で最も意識のある状態だということになる。事実はその逆だ。あの活動の洪水は、カオスの縁の体制から崩れ落ちて、ひとつの粗い形へと転じる------皮質の一枚全体が、粗い足並みそろえたリズムへと引きずり込まれるか、暴走するカオスへと投げ込まれるかして------そして分化した、グライダーを担う計算は消え去る。だから第一のダイヤルが測っているのは、単なる同時活性化ではありえない。それが測っているのは、脳のどれだけの部分が、真のカオスの縁のダイナミクスに巻き込まれているか、でなければならない。発作のなかでは、あれほど発火しているにもかかわらず、その数は崩れ落ちる。脳が大量に活発であることは、脳が大量に意識的であることと同じではない。この二つのダイヤルを心にとめておいてほしい。脳が入りうる最も奇妙な状態は、そのダイヤルの設定が違うだけなのだから。

\section*{皮質オートマトン}

まだ抽象的に感じられるかもしれないものを、そろそろ具体的にしよう。皮質はカオスの縁で、クラス4のダイナミクスのなかで働かなければならない------そう語ってきた。だが、そのクラス4のシステムとは、いったい\textbf{何}なのか。それは脳の上空を漂う、何か神秘的な力ではない。それは神経の発火のパターンそのものだ。

働いているときの皮質が実際にどう見えるかを考えてみよう。数十億のニューロン、その一つひとつが発火しているか、していないかのいずれかで、その一つひとつが、学習された結合荷重を通じて隣人に影響を及ぼしている。それぞれのニューロンは、セル・オートマトンのなかの一つのセルである------比喩としてではなく、文字どおりに。オートマトンのルールとは、シナプスの重みであり、閾値であり、局所的な配線だ。それぞれの「セル」の出力は、発火率である。そしてその結果------皮質の表面を10から40ヘルツで舞い踊る、電気的活動の広大なパターン------は、何千もの次元をもつ空間のなかで作動する、ウルフラムのクラス4セル・オートマトンなのだ。

私はこれを\textbf{皮質オートマトン}と呼んでいる。

これは、私が子どものころ286で組んだのと同じ考え（コンウェイのライフゲーム）だ------ただし、三つのルールをもつ平坦な格子ではなく、局所ごとに異なる数十億のルールをもつ、折り畳まれた皮質の一枚であり、二次元を動くのではなく、そのパターンは、あまりに広大で思い描くことを拒むような次元空間のなかを移動する。限りない腕をもつ蛸のように、皮質オートマトンは、いつでも皮質のどの部分にも手を伸ばすことができ、そのとき必要とする、蓄えられたどんなモデルでも活性化する------ここには記憶、あそこには運動計画、どこか別のところには言語の断片を。それはこれらのモデルを小さなレゴの人形のようにつかみ取り、それらを使って、ひとつの満ち足りた状態から次の状態へと渡り歩いていく。

そしてここに、決定的な区別がある。\textbf{皮質オートマトンは意識ではない}。それはエンジンであって、経験ではない。数十億のニューロンが発火する、一見カオス的なパターンは、実のところ、計算し、思考し、身体を人生のなかで舵取りしていく、桁外れに精巧な装置なのだ。だが意識は、この装置がもたらす一つの\textbf{効果}にすぎない------条件が整ったとき、オートマトンと皮質のあいだの相互作用から立ち上がる効果である。オートマトンが、適切な皮質領域を、正しい周波数で、首尾一貫した時間的な連なりのなかで同期的に掃過するとき、そのフレームの連なりから意識的経験が立ち現れる。オートマトンは、私たちの世界モデルと自己モデルのインスタンスを内に含んでいる。意識とは、これらのモデルがシミュレーションのなかで能動的に走っているときに、起こることなのだ。

ちなみに、皮質オートマトンは直接観察することができる------fMRIは要らない。

試しにやってみよう。真っ暗な部屋を見つける。目を閉じる。残像が消えるのを待つ（直前まで明るいものを見ていたなら、30秒から60秒ほどかかる）。最初は何も、あるいはほとんど何も見えない。だが、そのまま待って注意を向けていると、暗闇を背景にちらつく色のついた点が見えはじめる。

多くの人はこれを「網膜ノイズ」------圧力や自発的な化学的事象に反応した、眼の光受容細胞のランダムな発火------として片づけてしまう。確かに、まぶたを軽く押せば、そうやって局所的な視覚的感覚を引き起こすことはできる。だが、完全な暗闇のなかで見えるあの色のついた点は、網膜由来では\textbf{ない}。それにしては秩序立ちすぎている。ここで見えているのは、V1（一次視覚野）の安静時活動であり、それは残存する感覚信号と、皮質オートマトン自身からのトップダウン投射との組み合わせによって駆動されている。オートマトンはその基底的なダイナミクスを動かしており、私たちはそれがリアルタイムで起きているさまを眺めているのだ。

そのまま見つづけていると------集中するのではなく、リラックスして、注意を緩めていくと------驚くべきことが起こる。能動的な集中は、じつはこうしたパターンを抑え込んでしまう。見ようとするのをやめたときにこそ、見えはじめるのだ。オートマトンは、そこにあるわずかな信号を解釈し増幅しようと、視覚系のより多くの部分を動員しはじめる。ちらつく点は安定して形をとる。幾何学的なパターンが立ち現れる------格子、螺旋、網目。次いで、歪みながら移ろう顔。次いで、人影。そしてさらに、十分な忍耐があれば（数分ではなく、文字どおり\textbf{何時間}である）、完全な情景が------手の込んだ、色彩豊かな、物語をなす幻覚が、毎晩見る夢と本質的になんら変わらないものが現れる。

これは、入眠時幻覚------眠りに落ちようとするまさにそのとき、脳裏をちらついて過ぎる鮮やかなイメージ------の背後にあるのと同じメカニズムである。それは、外部からの制約が最小限に抑えられた状態で走る皮質オートマトンであり、蓄えられたパターンを活性化してシミュレーションへと投射することで、自らのコンテンツを生成しているのだ。ここで経験する進行（かすかなノイズから、まとまりのある幻覚へ）は、オートマトンの働き方をのぞく直接の窓である。それはV1、すなわち最も初期の視覚処理段階から始まり、手に入るかぎりの信号から意味をつかもうとしながら、V2、V3、そしてより高次の領域を次々と動員していく。本物の信号がまったくないときには、それを\textbf{みずから生み出す}。これが透過性の漏れが働いている姿だ。シミュレーションを支配する外部信号がないと、基盤そのものの処理ノイズが目に見えるようになる。\textbf{何もない}ところで幻覚を見ているのではない------見ているのは、グラフィックエンジンのアイドル時のパターン、チューニングの合っていないテレビに走る砂嵐の神経版なのだ。ただし、この砂嵐には構造がある。処理を担う機構そのものに構造があるからだ。

この方法で、一時的な形の共感覚を引き起こすこともできる。若いころ、私はこれを使って「音楽を見て」いた。目を閉じ、視覚的なパターンが訪れるにまかせながら音楽を聴くと------リラックスして、受け身で、見ようと気張らずにいると------パターンはやがて、聴いている音のリズムや周波数と同期していく。外部からの視覚入力を奪われた皮質オートマトンは、手に入るなかで最も強い別の信号------この場合は聴覚入力------へと、自らの視覚的ダイナミクスを結びつけはじめる。そこで見えるものは、文字どおり、自分の脳の活動が目に見える形になったものだ。すなわち、網膜入力ではなく聴覚野によって駆動される、オートマトンのV1レベルのパターンである。本物の共感覚者------感覚が恒久的に交差配線され、音を聞くとつねに色が見える人々------は、おそらくこの同じ結びつきのより恒久的な版を備えており、それはおそらく、視床であれ皮質そのものであれ、感覚領域どうしのより強い、あるいはより多くの結合によるものだろう。メカニズムは同じである。ある感覚のモダリティが、別の感覚の処理パイプラインへと漏れ込んでいるのだ。皮質オートマトンは、入力がどこから来るかをさほど気にしない。受け取ったものを、なんであれ処理するのである。

これを日常の趣味として試すことは勧めない。この経験は、とりわけ心の準備ができていなければ、不安をかき立てかねない。それに、持続的な感覚遮断が、潜在的な精神的脆弱性を抱える人を不安定にしてしまう可能性も、わずかながらある。だが、自分の意識の基盤が、アイドル状態にあるとき------外の世界が静まり返り、システムがただ……動いているだけのとき------どのように見えるのかと、一度でも思いをめぐらせたことがあるなら、これこそ、脳スキャナーなしで得られる最も直接的な一瞥である。

ほとんど何もない状態から、完全に架空の視覚世界へと至るこの進行------仮想の宇宙のなかで自己モデルによって経験されるそれ------は、働いている皮質オートマトンの、まさに肖像そのものである。

オートマトンが不調に陥るときも、それを目にすることができる。てんかん発作とは、オートマトンの一部がクラス1あるいはクラス2のダイナミクス（周期的で、固定され、計算的には役に立たない）に落ち込んだとき、あるいはクラス4を越えてクラス5のカオスへと押しやられたときに起こることだ。脳卒中とは、皮質の一部がまるごと脱落するときに起こることだ。失神とは、覚醒に必要な最低周波数がもはや満たされなくなったときに起こることだ。オートマトンは、いくぶん脆い。だが、それを生み出す構造------学習された重みと進化してきたアーキテクチャを備えた新皮質------は頑健であり、だからこそ私たちは、こうした乱れからこれほど見事に回復できるのである。

\section*{収斂}

2003年------私がこの理論を手にするより二年も前のことだ------ジョン・ベッグス（John Beggs）とディートマー・プレンツ（Dietmar Plenz）は、皮質組織のなかに「神経雪崩」を発見した。それは、自己組織化臨界性という数学的なしるし------カオスの縁（ふち）にあるシステムに特有の指標------に従う、神経活動のパターンだった。

2014年、ロビン・カーハート=ハリス（Robin Carhart-Harris）はエントロピック・ブレイン仮説（Entropic Brain Hypothesis）を提唱した。意識の水準が脳活動のエントロピー（無秩序）と相関し、その最適点が中間の水準にある------エントロピーが少なすぎれば無意識を、多すぎればまとまりを欠いた経験を意味する------という考えである。

2016年、エンツォ・タリアツッキ（Enzo Tagliazucchi）らは、LSDが脳を臨界性へと押しやることを示した。これは、幻覚剤の使用者が報告する、高まった（だがときにカオス的な）意識と符合する。2022年になると、あるレビュー論文はすでに「意識の枠組みとしての自己組織化臨界性」について語ることができた------証拠は積み上がりつつあった。

そして2025年から2026年にかけて、実証のダムが決壊した。キース・ヘンゲン（Keith Hengen）とウッドロウ・シュー（Woodrow Shew）は、140のデータセットを対象としたメタ分析を\textit{Neuron}誌（2025年）に発表し------脳のダイナミクスにおける臨界性について、これまでに行われた最大規模の体系的分析である------脳が複数の測定手法をまたいで臨界点の近くで動作していることを確認した。次いでインバル・アルゴム（Inbal Algom）とオーレン・シュリキ（Oren Shriki）は、ConCritフレームワーク（Consciousness and Criticality、意識と臨界性）を\textit{Neuroscience \& Biobehavioral Reviews}誌（2026年）で提唱し、臨界的な脳ダイナミクスが、意識に関するあらゆる主要理論を統一するメカニズム的な基盤を与えると論じた。彼らの結論はこうだ------意識は臨界性を追う。脳が臨界点にある、あるいはその近くにあるとき、意識は存在する。臨界性の下へと押し下げられたとき（麻酔によって、睡眠によって、脳損傷によって）、意識は失われる。臨界点から反対の方向へ押しやられたとき------発作やおそらく一部の薬物状態に見られるように、硬直した過同期へと崩れ落ちるか、暴走するカオスへと爆発するかして------意識はまとまりを失う。

二つの道。一方は理論的で、ウルフラム（Wolfram）の計算論的な枠組みから出発し、自己シミュレーションに何が必要かを推論していく。もう一方は実証的で、神経の記録から出発し、アクセス可能なあらゆる意識状態にわたって脳活動の統計的性質を分析していく。出発点はわずか二年ほどしか隔たっていないのに、同じ結論へと収斂していく。

これこそ、一つの理論を真剣に受け止めさせる類いの収斂である。

\section*{ホログラムがオートマトンに出会う三つの道}

この章を書いているさなか、私は思わず凍りつくようなことに気づいた。

ホログラフィック原理とクラス4オートマトンは、同じ議論のなかにくり返し姿を現す------物理学で、神経科学で、計算理論で。だが、誰一人として、あの当たり前の問いを立てていないように見える------\textbf{両者のあいだにありうる関係とは、いったい何なのか}、と。

関係は、ちょうど三つある。

\textbf{関係1------ホログラフィックな基盤がクラス4のダイナミクスを生み出す。} おそらく、脳がしているのはこれだ。神経ネットワークは局所的にホログラフィックであり------第3章のパッチワーク・ホログラム、ラシュリー（Lashley）のネズミ、その他もろもろ------その基盤が臨界性で動作することで、意識が要求するクラス4のダイナミクスを生み出す。十分に確立され、詳細に記録され、そして（許してほしいのだが）退屈な選択肢である。

\textbf{関係2------クラス4オートマトンが、創発的なふるまいとしてホログラフィックなパターンを生み出す。} オートマトンは、その規則においてはホログラフィックではないが、そのダイナミクスが自発的にホログラフィックな構造を生成する------低次元のパターンのなかに符号化された高次元の情報が、計算そのものから立ち現れるのだ。もしクラス4オートマトンが自然にホログラフィックな出力を生み出すのなら、それは、純粋に局所的な規則から非局所的な情報の分布が創発することを意味する------そしてこれは、興味深いことに、量子もつれの姿とそっくりそのままなのである。

ここでジェラルド・トホーフト（Gerard 't Hooft）に触れておかねばならない。この結びつきは------たとえ思弁的であるにせよ------見過ごすにはあまりに印象的だからだ。物理学のノーベル賞受賞者であるトホーフトは、量子力学それ自体がプランク・スケールにおけるセル・オートマトンであると提唱している。すなわち、私たちの宇宙は根本的に決定論的であり、量子的な効果は、より深い、離散的なダイナミクスの創発現象だというのだ。もし彼が正しければ、私がこれまで述べてきた原理は、たんなる類推として意識に当てはまるのではない。それは文字どおり、宇宙が最下層に至るまでどのように働いているかそのものである。単純な局所的規則がホログラフィックな宇宙を生み出し、その宇宙のなかで、単純な神経の規則がホログラフィックな意識を生み出す。同じ計算原理が、宇宙論的と神経学的という二つのスケールで働いている。このフラクタルな一貫性を、私は心底から説得力があると感じる。だが正直に言えば、トホーフトの解釈は物理学のなかでは少数派の見解にとどまっており、構造的なエレガンスから物理的実在を導く論法は、当然ながら批判されてきた。それでもなお------もし、たった一つの計算原理が、宇宙と、それをモデル化する心の両方の根底にあると判明したなら、それは、これまでに発見されたなかで最も美しい事実ということになるだろう。

\textbf{関係3------規則の構造それ自体がホログラフィックなクラス4オートマトン。} これこそ、私が思わずペンを置いてしまった一つだ。もしそのようなものが存在するなら------ホログラムが三次元を二次元に符号化するのと同じように、規則それ自体が高次元の情報を低次元の構造のなかに符号化するセル・オートマトンが存在するなら------そのとき手にするのは、ホログラフィック原理が宇宙のふるまいとして述べていることを、生まれながらにして行うシステムである。ホログラフィックな基盤の\textbf{上で動作している}だけの（あるいはホログラムを生み出すだけの）システムではない。それ自体がホログラフィックな符号化で\textbf{ある}システムだ。それはおそらく宇宙でもある------もっとも、これは思弁的であり、数学的な美が物理的実在を含意するという論法が、正当にも批判されてきたことは記しておくべきだろう。この点には第14章で立ち返り、なぜ関係3が数学における最も重要な未解決問題かもしれないと私が考えるのかを説明する。そして第15章から第17章にかけて、その答えを余さず追い求めていく。


\chapter*{第6章：サイケデリックが明かすもの}
\addcontentsline{toc}{chapter}{第6章：サイケデリックが明かすもの}
\markboth{第6章：サイケデリックが明かすもの}{}

始める前に、どうしても言っておかねばならないことがある------本章のどこにも、サイケデリックを試すことを勧める意図はない。それらは強力で予測不能であり、人生を破壊しかねない------文字どおり、そして永久に。素因のある者には統合失調症を引き起こしうる。精神病エピソード、慢性的な不安障害、そして二度と消えないHPPD（hallucinogen persisting perception disorder）を引き起こすこともある。ここでサイケデリックを論じるのは、それが意識のアーキテクチャについて重要な何かを明かすからだ。だが、その科学的価値が安全性を保証するわけではない。

意識を理解したいなら、それが狂ったときに何が起こるかを調べるとよい。サイケデリックは------私はそう確信しているのだが------私たちの手にする意識のアーキテクチャを最もよく見通せる窓だ------眠っている患者の脳スキャンよりも多くを明かし、病変研究よりも理論的に示唆に富み、そして分離脳手術よりもはるかに手が届きやすい。

理由はこうだ。サイケデリックは意識をただ\textbf{変える}だけではない。それを\textbf{体系的で予測可能なしかたで}変え、その根底にあるアーキテクチャを露わにする------何に目を向ければよいかを知ってさえいれば。

\section*{透過性の勾配}

暗黙的モデルと明示的モデルのあいだの境界を思い出してほしい。通常の覚醒時、この境界は選択的に透過する------関連する情報は通り抜け、無関係な情報は図書館にとどまる。必要なものは意識にのぼり、それ以外のすべては無意識のままだ。

サイケデリックは、その境界を吹き飛ばす。

サイケデリック（LSD、シロシビン、DMT、メスカリン）の作用下では、暗黙的-明示的境界の透過性が全体的に高まる。ふだんは完全に現実の側で処理され、意識には見えない情報が、シミュレーションへと漏れ出しはじめる。

そして決定的な点はこうだ。それは\textbf{順序どおりに}漏れ出す。

低用量のとき、あるいは体験の初期には、最も単純な処理段階がまず見えるようになる。生の感覚入力に最も近い段階、すなわちV1レベルの処理だ。色が鮮やかさを増し、静止した表面が呼吸するように波打ち、周辺視野にかすかな動きが現れる。これらは視覚野の初期の特徴検出器であり、ふだんは見えないものが、いまシミュレーションへと入り込んでいる。

用量が増えるにつれて、あるいは体験が深まるにつれて、より複雑な処理段階が見えてくる。V2/V3レベルの処理------幾何学的な模様、フラクタル、テッセレーション、そしてハインリヒ・クリューヴァー（Heinrich Klüver）が1920年代に分類した有名な「形態定数」だ。これらは視覚系の中間表現であり------ふだんは視覚体験を組み立てるために用いられる構成要素が、いまやそれ自体として見えている。

さらに高くのぼれば、より高次の視覚領域にアクセスできるようになる。顔が現れる。人影が。情景が。顔処理領域、物体認識領域、情景構築領域------いずれもふだんは意識の閾値の下で働いている------が、いまやその中間産物を直接シミュレーションへと送り出している。

最も高い用量では、処理階層の全体が露わになり、その結果はまぎれもないヴィジョン体験となる------暗黙的処理の最も深い層から組み立てられる、複雑で物語的な、夢のような情景だ。

この秩序だった進行------単純から複雑へ、V1から高次領域へ、そして用量に依存して------こそ、四モデル理論（FMT）が予測するものにほかならない。それは透過性の勾配の直接の帰結だ。より低次の処理段階は境界に近いぶん、透過性が高まるにつれて、より高次の段階よりも先にアクセス可能になる。

以下の視覚処理階層は、各領域がふだん何をしているか、そして透過性の障壁が下がったときに何が見えるようになるかを示している。


{\small
\noindent
\begin{tabularx}{\linewidth}{X X X}
\toprule
\textbf{領域} & \textbf{通常の機能} & \textbf{サイケデリックの特徴} \\
\midrule
V1 & 輪郭、空間周波数、方位 & 眼内閃光、クリューヴァーの形態定数、呼吸する表面 \\[4pt]
V2 & 輪郭統合、テクスチャ、境界所有 & テッセレーション、反復する幾何学模様 \\[4pt]
V3 & 大域的形態、動的な形状処理 & 流動し変形する幾何学形状 \\[4pt]
V4 & 色、曲率、複雑なテクスチャ & 色づいたフラクタル、万華鏡のような模様 \\[4pt]
V5/MT & 運動処理 & 模様の回転と移動 \\[4pt]
紡錘状回/IT & 顔、物体、単語形 & 顔、人影、存在 \\[4pt]
前部IT & 意味カテゴリー、情景構築 & 完全な物語的幻覚 \\[4pt]
\bottomrule
\end{tabularx}
}


各行は、より深い処理段階を表している。通常の条件では、体験されるのは最終的な出力------完成した知覚像------だけだ。サイケデリックの作用下では、透過性が高まるにつれて、\textbf{中間の}段階が順序どおりに体験される。（この表のより詳しい版は、受容野の大きさやその他の詳細を含めて、付録Aにある。）

興味をそそる話に聞こえるだろう。視覚処理の層が見えるようになるという話を読んで、それがどんなふうに見えるのか、心のどこかで知りたくなっているのではないだろうか。よくわかる------私も知りたかった。私は両方の道を試した。若く、愚かで、そのうえ運がよかった。前章で述べた瞑想の道（暗い部屋、ゆるめた注意、忍耐）は、同じ場所へと導いてくれる。それほど速くはなく、最初の一度でそれほど劇的でもない。だが同じくらい印象的で、同じくらい現実的で、しかも心を永久に損なう危険はない。暗い部屋の暖かなベッド------必要なのはそれだけだ。

もう一つの道がある------明晰夢だ。夢のなかにいながら、自分が夢を見ていると気づけるようになれば------これは訓練で身につく技能だ------、制約なく走るシミュレーションの全体にアクセスできる。感覚入力もなく、モデルを修正する外的な現実もない。ただ仮想の世界があり、そのただなかに意識ある自分がいる。人によっては、これは持続的な瞑想よりも達成しやすい。その技法はよく記録されており（実践的な手引きは付録Dを参照）、その体験は、いかなる薬物がもたらすものにも劣らず啓示的でありうる------しかも危険はない。明晰夢については第7章で立ち戻る。

\section*{入れ子になった五つのシステム}


\begin{figure}[htbp]
  \centering
  \includegraphics[width=0.95\textwidth]{../figures/figure-five-layer-stack-bw-ja.png}
  \caption{脳組織の五つのレベル。各レベルは、その下のレベルに付随する（「その上で走る」）。意識は最上位の仮想レベルにのみ存在し、そこで明示的モデルが現象的経験を生み出す。}
\end{figure}


脳には、はっきりと区別される五つの組織のレベルがあり、それらがロシアの入れ子人形のように積み重なっていると考えてみてほしい。

\textbf{物理的。} 最も底にあるのは、生の物質だ------原子、分子、脳そのものの物理的な基盤。これは化学の世界だ------組織を構成する炭素、水素、窒素、酸素。熱力学の法則に従う不活性な物質にすぎない。意識あるものは、ここには何も棲んでいない。

\textbf{電気化学的。}一つ上の階層------神経の信号伝達である。活動電位が軸索を駆け下り、神経伝達物質がシナプスへ押し寄せ、イオンがチャネルを流れる。「脳が何かをしている」と思うとき、誰もが思い描くのがこの電気的・化学的な活動である。ニューロンが発火するのはこの階層だ。まだ経験はない。だがここで初めて情報の伝達が生まれる。

\textbf{プロテオーム的。}その次------タンパク質構造と分子的な機械仕掛けである。シナプスの重みはここに蓄えられている------ニューロン間の結合の物理的な強さだ。細胞膜上の受容体、可塑性を調整する酵素、どのシナプスが強まりどれが弱まるかを決める分子の足場。これが学習の「ハードウェア」である。技能を練習して上達するとき、変化しているのはこのプロテオーム的な層だ。まだ無意識のままだが、ここで初めて記憶が生まれる。

\textbf{トポロジー的。}さらに上------ネットワークの構造である。結合のパターン、どのニューロンがどのニューロンとどれほど密に、どのような配置でつながっているか。ブロードマン野が存在するのも、皮質コラムが存在するのも、「視覚野が運動野と対話する」という大規模構造が存在するのも、この階層だ。配線図である。この階層を変えれば、そのシステムがどのような種類の処理をこなせるかが変わる。あなたの暗黙的モデル（IWMとISM）が蓄えられているのもここだ。まだ無意識のまま。だがここで初めて知識が生まれる。

\textbf{仮想的。}まさに頂点------シミュレートされた世界である。皮質のオートマトン------ネットワーク上を舞う電気活動の動的なパターンであり、情報を統合し、予測を生み出し、モデルをリアルタイムで動かしている。あなたの意識的経験が宿るのはここだ。明示的モデル（EWMとESM）が存在するのは、ここであり、ここだけである。何かとして感じられる唯一の階層だ。

各階層は、その下の階層に付随（スーパーヴィーン）しつつも、それ自身の力学を備えている。物理的な物質なしに電気化学的な信号伝達はありえず、化学なしにタンパク質構造はありえず、シナプスなしにネットワークのトポロジーはありえず、そしてネットワークなしに、その上を走るシミュレーションはありえない。だが各階層は、下の階層が持たない性質を備えている。シナプス一つは何かに「ついて」のものではない------ただの結合にすぎない。ところがシナプスのネットワークは何かに「ついて」のものである------顔を、言葉を、記憶を表象する。では、そのネットワーク上を走るシミュレーションはどうか。「ついて」が「経験」へと変わるのは、そこなのだ。

この五階層の階層構造は、この理論を初めて聞いたほとんど誰もがつまずく問題を解く------「意識が仮想的なものなら、それは何の上を走っているのか」。答えはこうだ。意識はトポロジー的な層（ネットワーク）の上を走り、その層はプロテオーム的な層（シナプスの重み）に実装され、その層は電気化学的な層（神経の発火）の上に載り、その層は物理的な層（物質）のうちに存在する。意識は仮想的であるからといって、少しも現実性を欠かない------ニューロンが現実であるのとは\textbf{別の階層で}現実であるにすぎない。テレビゲームの中の山は、ハードウェアの階層では「ただの」トランジスタであっても、ゲームの階層では現実である。同じ原理だ。

この五階層の階層構造が説明の仕事をするのは、まさにここである。サイケデリック薬はこの階層構造の中ほどを標的とし、その効果は上へと波紋のように広がっていく。LSDやシロシビンといった古典的サイケデリックはセロトニン2A受容体に結合し、\textbf{電気化学的}な階層で作用する。ニューロンどうしの語り合いかたを変えるのだ。その撹乱は\textbf{プロテオーム的}な階層へ伝わり、そこでは受容体の感受性が数時間かけて変化する。それは\textbf{トポロジー的}な階層を作り変え、ネットワークの結合パターンが変わる------fMRIでは全体的統合の増大として可視化される。そしてそれは\textbf{仮想的}な階層を変容させ、そこでは意識的シミュレーションが、ふだんは見えない内容で満たされる。古典的サイケデリックが触れない唯一の階層が\textbf{物理的}な階層だ------ニューロンを破壊せず、生の物質を変えることはない。物質の\textbf{上にある}すべてを、下から上への順で変えるのである。これは決定的な区別だ。古典的サイケデリック（LSD、シロシビン、DMT、メスカリン）は神経毒性を持たない。ニューロンを破壊することなく、その伝達のしかたを変える。他の多くの薬物はそれほど寛容ではない。メタンフェタミンとアルコールはニューロンを物理的に破壊する。コカインの害はもっと間接的で、主として血管を収縮させ、神経組織から酸素を奪うことによる。MDMAは高用量あるいは反復使用でセロトニン軸索を損なう。ベニテングタケ（Amanita muscaria）でさえ------多くの人がサイケデリック・マッシュルームと混同する、あの象徴的な赤と白のキノコだ------まったく別の、より危険な機序で働くせん妄誘発物質である。この章から一つだけ心に留めるとすれば、それは------意識を変える薬物がすべて同じわけではない、ということ、そして「信号を変える」と「ハードウェアを破壊する」の違いこそが、文字どおり、一時的な変性状態と恒久的な脳損傷との違いなのだ、ということである。用量に依存した視覚の段階的進行はこれに正確に対応する。低用量は電気化学的な階層を、V1の処理に影響するのにちょうど足りるだけ撹乱する。より高用量は撹乱をより多くの階層へと押し上げ、ますます複雑な処理段階を意識的経験へと引き込んでいく。

\section*{転向可能な自己}

だが最も劇的な証拠は、自己に何が起こるかから来る。

あなたの明示的自己モデル（ESM）は、入力を必要とする仮想的なプロセスである。通常の条件下では、それは自己言及的な信号の絶え間ない流れを受け取っている------身体がどこにあるかの感覚（固有感覚）、諸器官がどう感じられるかの感覚（内受容感覚）、内なる語りという物語的な流れ、そして、乱されて初めて気づく、身体的な自己意識の絶えざる背景だ。

高用量のサイケデリックでは、この入力が乱される。自己モデルは死なない------\textbf{転向する}のだ。通常の自己言及的な入力を奪われると、それは、そのとき優勢な入力なら何であれ、それにしがみつく。

これを最も劇的に示すのがサルビア・ディビノルム（salvia divinorum）である。これはカッパ・オピオイド受容体に作用する解離性のサイケデリックだ（LSDやシロシビンのセロトニン作動性の機序とはまったく異なる）。サルビアの使用者は一貫して、物に\textbf{なった}という経験を報告する------

\begin{itemize}
  \item 「私はソファになった」。
  \item 「私は壁だった」。
  \item 「私は本のページに変わった」。
  \item 「私はテレビの登場人物の一人だった」。
  \item 「私はフラクタルになった------フラクタルを見るのではなく、フラクタルで\textbf{ある}こと」。
\end{itemize}

これらは比喩ではない。使用者は、完全で、経験として説得力のある同一性の変化を報告する。経験が続くあいだ、彼らはその対象なり実体なりで\textbf{ある}。ある者はそれを、死ぬのではなく、死んで\textbf{いる}こと------死んだ状態でいることのように感じた、と述べる------なぜなら、椅子であるとき、かつてのその人はただ存在するのをやめてしまっているからだ。

内容は感覚的な環境に追随する。テレビを見ている人はテレビの登場人物になる。ソファに横たわっている人はソファになる。模様を見つめている人は模様になる。

これは、明示的自己モデルが理論の予測どおりに振る舞っている姿である------通常の自己入力が乱されると、そのとき優勢な入力なら何であれ、それへ転向するのだ。同一性の内容はでたらめではない------感覚的な環境によって決まる。環境を制御すれば、同一性の経験も制御できるはずである。

ここでいったん理論を止めておかねばならない。サルビア・ディビノルムは、知られているかぎり、地球上で自然に生じる最も強力なサイケデリック物質である。いま述べた固有感覚の完全な乗っ取りとは、身体の感覚と空間的な見当識の全面的な喪失を意味する。サルビアの影響下にある人々は、十階の窓から外へ歩き出したことがある。行き交う車のなかへ足を踏み入れたことがある。死んだ者もいる。これはパーティー・ドラッグではないし、金曜の夜に試してみる物珍しさでもない。存在するなかで最も極端な、明示的自己モデルへの薬理学的撹乱であり、その撹乱は人を殺しうる------薬が毒だからではなく、自分の身体がどこにあるかを知ることをやめてしまい、自分には翼があって飛べるのだと完全に信じ込むかもしれないからだ。

サルビアを試す多くの人が、その経験は死ぬことのように感じられた、と報告する------比喩としてではなく、もう存在していないという、本物の、身のすくむような確信としてである。これは、シミュレーションがもはや「あなた」というものをまったく生成できなくなるほど、明示的自己モデルが完全に崩壊した状態だ。これに臨床的に相当するものは第8章で見ることになる------コタール症候群を論じるところで、神経学的に自分は死んでいると確信している患者たちである。サルビアは薬理学的に、数秒で、前触れもなく、人をそこへ連れて行く。それがあなたの経験したいことなのかどうか、よく考えてみてほしい。

その時間の遅れを、私自身が経験した。サルビアの下で、現実の時間の半秒が------見ていた人物がそう確認した------十五分かそれ以上に感じられるものへと引き延ばされた。私の知覚世界は、翼を持って飛び回る感覚を含む、入り組んだ連なりへと自らを再構築していった（あとで気づいたのだが、その飛んでいる感覚は、ベッドの上へ仰向けに倒れ込むときに身体を通り過ぎていく空気から来ていた）。私の現実そのものが崩壊し、再生した------すべては、瞬き一つに要するほどの時間のうちに。私は自分の時間を計っていた観察者にこれを語ったが、彼らは、私が「いなくなって」いたのは一秒に満たなかった、と言った。同じ種類の時間の遅れを、私はすでに数年前、1998年か1999年に、臨死の出来事（その機序は第13章で述べる）のさなかに経験していた------ただし、それは薬理学的に誘発され、しかもさらに極端だった。

私は最も劇的な事例ではない。よく記録された一つの報告は、およそ四十五秒の時計時間しか続かなかったサルビア発作のあいだに、まる八年分にも感じられる別の人生------学校に通い、友をつくり、新しい存在を築く------を経験した男性のものである。査読済みの研究は、統制された条件下での極端な時間の歪みを裏づけており、ある参加者は時間を「アコーディオンのように折り畳まれている」と表現した（Addy et al., 2015）。基盤は、あまりにも多くの内容をあまりにも速くシミュレーションを通して流すため、主観的な時間が時計時間から完全に切り離されるのだ。

これが統制された環境で実験的に検証されたことは、いまだかつてない。だが検証は可能であり、そうなれば、この理論の最も特徴的な機序を裏づける、目を見張るような確証となるだろう。

脳がなぜこれほど多くの内容をこれほど速くシミュレーションを通して流せるのか------それは、死にゆく人がなぜ自分の生涯全体を数秒のうちに走り抜けるのを見るのか、という問いと同じものである。両者には第13章で一緒に立ち戻る。そこで、両者が単一の機序を共有していることが明らかになる。

問うてみる価値がある------なぜこの隅は、二つのダイヤルがどちらも一挙に高く押し上げられ、脳の多くが一つの出来事へと引き込まれ、\textbf{かつ}その出来事が最大限に豊かに走っているという隅は、これほど稀なのか。二つのことが通常これを禁じており、しかもそれらは互いを補い合っている。

第一は、単純な経済である。両方を同時に行うのは代謝的に破滅的であり、脳はそれをほんの一瞬より長くはまかなえず、エネルギー予算が扉を閉ざしたままにしている。だからこそ、その扉をこじ開けるには何か極端なものが要る------代謝のブレーキが利かなくなりつつある死にゆく脳か、あるいは予算を配給するまさにその調整機構を撹乱する薬物である。

第二はより奇妙で、これらの状態が世界を後に置き去りにするように感じられる理由でもある。脳の\textbf{内なる}世界は、それを現実につなぐ細い管------二つの目、二つの耳、身体一つ分の触覚------に比べて、はるかに高い次元を持つ。それに対して、内なるシミュレーションは目もくらむほどの豊かさを備えている。通常、その狭い感覚の管でちょうど足りるのは、それが内なるシミュレーションを絶えず\textbf{補正している}からだ------「実際に外にあるのはこれだ」という一定の細い流れが、暴走しかねない想像を世界につなぎとめている。だが両方のダイヤルを押し上げると、内なる出来事は膨れ上がり、ついには細い入力の管がそれを補正しきれなくなる。シミュレーションは自らの現実検証を追い越し、自らのアトラクターへと固定され、世界から漂い離れていく。そしてここに巧妙なひねりがある------それをするのは\textbf{広がり}のダイヤルなのだ。より多くの脳を内へと引き込むこと、それこそが、つなぎとめる綱を水没させてしまう。隅へと押し込むことは、あなたを地に足のついた状態に保っていたはずの、まさにその線を断ち切るのである。

そうして残るのは、私がいまなお驚かされる一つの反転だ。その隅は、同時に、\textbf{意識の頂点であり、現実との接触の最小である}------最も多くの処理、最も少ない世界。両方のダイヤルが上がった状態（臨死の人生回顧、高用量サルビアの墜落、最も深いサイケデリックな没入、最も鮮やかな夢）が、現実から最も切り離されて感じられるのは、まさにそのためである。それらは激しさに\textbf{もかかわらず}切り離されているのではない。激しい\textbf{からこそ}切り離されているのだ。

この原理がどこまで及ぶのかを見たければ、次の思考実験を考えてみてほしい。ある人物が、サルビノリンA------サルビア・ディビノルムの有効成分で、単一の受容体型（カッパ・オピオイド）に作用する------のきわめて高い（ただし完全には解離させない）用量に、持続的に保たれているとしよう。この人物の明示的自己モデル（ESM）は、決して安定しないだろう。その時々でたまたま優勢になった入力に応じて、際限なく循環しつづける------ある瞬間には自分は椅子だと信じ、次には机、次には恐竜、次には空気、次には一枚の紙だと信じる。それでも物事を\textbf{経験}してはいる（視覚も聴覚もなお機能している）が、自分が誰なのか、何なのかを、二度と知ることはない。薬を取り除けば、時間とともに、無傷のまま残った暗黙的自己モデル（ISM）から通常の自己モデルが再び組み上がっていく。

これが重要なのは、意識が\textbf{正しい}自己モデルを必要とはしないことを示しているからだ。ただ\textbf{一つの}自己モデルがあればよい。アーキテクチャは何があろうと動きつづける。明示的自己モデルは、不条理な入力を与えられても停止しない------利用できる信号がどんなものであれ、そこから可能なかぎり最良の自己を組み立てる。これはコタール症候群（欠落した内受容信号のうえで働くESM------「自分は死んでいるにちがいない」）、アントン症候群（眼が働いていないときに記憶から視覚を生成するEWM）、そして転換性障害（基盤には実際にはない麻痺をモデル化するESM）に見られるのと同じ原理だ。自己モデルは強迫的な構築者である。決して構築をやめない。データが不十分だと告げることも決してない。ただ構築し、そして信じる。

\section*{病態失認------その逆転}

ここには見事な対称性がある。サイケデリックが、暗黙的-明示的の境界が\textbf{過度に}透過的になったときに起こることであるなら、病態失認は、その境界が------少なくとも局所的に------\textbf{過度に}不透過的になったときに起こることだ。

病態失認は、右半球の脳卒中後に最もよく見られ、患者が自分自身の欠損に本当に気づいていない状態である。左腕が麻痺している患者は、腕は問題ないと言い張り、それを使えなかったことをあれこれ理由をつけて説明し去り、麻痺の証拠を突きつけられると混乱したり怒ったりする。心理的な意味で否認しているのではない------腕が麻痺しているという情報が、そもそも彼らの意識的シミュレーションに届かないのだ。

四モデル理論（FMT）では、これは暗黙的-明示的透過性の局所的な低下である。暗黙的自己モデルは麻痺の情報を\textbf{もっている}------基盤は損傷を記録している。だが境界がその特定の領域について遮断されているため、明示的自己モデル（ESM）にはその欠損が決して含まれない。患者のシミュレーションには麻痺した腕が存在しないので、患者はそれを経験しない。

そのメカニズムはさらに具体的で、いったんそれが見えると、いささか背筋の寒くなるような仕方で見事にできている。運動系が指令（たとえば「手を叩け」）を送るとき、それは同時に二つのことを行う。筋肉へ指令を送ると同時に、意識へ\textbf{予測されたフィードバック}を送る------過去の経験にもとづき、手を叩けばどう感じ、どう聞こえるはずかを。この予測されたフィードバックは、実際の感覚フィードバックよりも\textbf{先に}届く。本物のフィードバックはより遅い神経経路を通らなければならないからだ。通常の状況では、予測はすぐに実際の感覚データによって修正されるか、確認される。手を叩くと予測し、次に叩くのを感じ、聞く。一致。次へ。

病態失認では、麻痺した肢からの実際のフィードバックが決して届かない。そして「待て------何も起こっていない」と警告すべきメカニズムが損傷している。だから予測されたフィードバックは修正されないまま通ってしまう。患者の運動系は両手に叩けと指令し、両手で叩いたという予測を意識に送り、意識はまさにそのとおりを経験する------両手による完全に正常な拍手を。患者はまったくの誠実さで、いま両手で叩きましたと語るだろう。彼はそれを聞いた。感じた。経験した。彼のシミュレーションのなかでは、それは起こったのだ。ただ現実では起こらなかっただけだ。

これ\textbf{こそが}意識の働き方であり、それは常に、私たち全員のなかで起きている。唯一の違いは、健康な人では予測されたフィードバックが数ミリ秒のうちに修正されることだ。病態失認では、その修正メカニズムが壊れており、患者のシミュレーションはただ予測だけを頼りに動く。

サイケデリックと病態失認は、同じメカニズムが逆方向に働いたものである。一方は透過性を全域的に高める。他方はそれを局所的に低める。そしてこの対称性が、領域を横断する一つの予測を生む------サイケデリックは病態失認を緩和するはずだ。全域的な透過性の上昇が局所的な遮断を圧倒し、欠損の情報が意識に届くことを可能にするはずである。

これを検証した者は誰もいない。この二つの現象を結びつける理論を、誰ももっていなかったからだ。四モデル理論なしには、その結びつきは見えない。


\chapter*{第7章：明かりが消えるとき何が起こるか}
\addcontentsline{toc}{chapter}{第7章：明かりが消えるとき何が起こるか}
\markboth{第7章：明かりが消えるとき何が起こるか}{}

毎晩、意識を失い、そして毎朝、それを取り戻す。この二つのあいだの移行------睡眠段階をめぐる旅------は、臨界性の原理が毎夜くり広げる実演にほかならない。

\section*{深い眠り------閾値の下へ}

深いノンレム睡眠では、脳のダイナミクスは亜臨界の領域へと移る。その目印が徐波である。膨大な数のニューロンが一斉に発火し、そして一斉に沈黙する、大きく同期した振動だ。これはクラス2のダイナミクス------周期的で反復的、意識には秩序がありすぎる。

マルチェロ・マッシミーニ（Marcello Massimini）らが開発した摂動複雑性指数（PCI）は、これを直接に裏づける。PCIは、脳が磁気パルスに対してどれほど複雑に反応するかを測る。覚醒した意識ではその応答は複雑で分化しており（PCIは高い）、深い眠りでは単純で紋切り型である（PCIは低い）。深い眠りにある脳は、意識的なシミュレーションが必要とする、豊かで全域的に統合されたダイナミクスを維持できないのだ。

明かりは消えている。明示的世界モデル（EWM）と明示的自己モデル（ESM）は崩れ落ちた。シミュレーションもなく、経験もない。

\section*{夢------劣化モード}

だがレム睡眠のあいだ、明かりはふたたび灯る。脳のダイナミクスは臨界性のほうへと戻っていく------完全にではないが、十分に近いところまで。シミュレーションはふたたび起動し、ふたたび一つの世界が経験される。

とはいえ、それは劣化したシミュレーションである。通常の外部入力は断たれている（目は閉じられ、筋肉は麻痺している）。明示的世界モデルは内部データで動く------いまの感覚入力ではなく、暗黙的世界モデル（IWM）に蓄えられた知識を汲み上げるのだ。だからこそ夢には見慣れた場所や人が登場するのに、物理法則はありえず、物語の筋も破綻している------シミュレーションは、限られた入力でできるかぎりのことをしているのである。

明示的自己モデルもまた劣化モードで走る。夢は「自分」の身に起きていることとして経験されるが、メタ認知的な監視は絞られている------ありえない出来事も疑いなく受け入れ、自分が夢を見ているとはめったに気づかず、批判的な判断力は鈍っている。

その劣化したシミュレーションが内側からどう感じられるか、私はよく知っている。七歳から十二歳くらいまでのあいだ、私はくり返し同じ夢を見た------その数年のあいだに、同じ夢が何度も戻ってきたのだ。それは風景と呼べるなら風景だった。無限に広がるフラクタル構造で、谷はどこまでも深く落ち込み、どの深さを覗いても自己相似だった。それには、覚めているあいだには一度も出会ったことのない感覚がともなっていた------深い脱身体感、まるで身体を剥ぎ取られ、幾何学そのものの内側に置かれたかのようだった。その夢は不気味で、奇妙によそよそしく、それでいて奇妙に魅入られるものでもあった。恐ろしくもあったが、それでも私はまた戻ってきてほしいと願った。

驚くべきことに、私はいまだにそれを\textbf{感じる}ことができる。何十年もたったいま、あの夢のことを思うと、同じ脱身体化した感覚が戻ってくる------感覚の記憶としてではなく、感覚そのものとして。暗黙的な符号化はいまもシナプス結合の重みのなかに残っており、三十年の覚醒経験にも触れられていない。基盤はそれを深く蓄えたので、思い浮かべるだけでシミュレーションを部分的に再起動させるに足りるのだ。

その夢は何をしていたのか。理論の言葉で言えば、明示的世界モデルが外部入力から切り離され、基盤自身の計算ダイナミクスから内容を生み出していた。ではそのダイナミクスとは何か。クラス4------それはクラス3（フラクタル構造）を部分プロセスとして含んでいる。眠る私の脳は、自分自身のアーキテクチャだけを頼りにシミュレーションを走らせ、そのアーキテクチャの視覚的なしるしを生み出した。すなわちフラクタルである。脱身体感は、最も劣化した形で走る明示的自己モデルだった------シミュレーションは自分が\textbf{誰か}であること、どこかにいることは分かっていたが、身体の表象はほとんど完全に脱落していたのだ。

何年ものちに明晰夢を身につけたあとも、私はあの夢を呼び出そうとしたことは一度もない。もしかすると、試してみるべきなのかもしれない。

夢遊病は、さらに劇的な実演である。夢遊病では、明示的自己モデルが停止したまま、あるいはほぼ停止したまま、運動系が部分的に再起動する。基盤は運動プログラムを実行している------歩き、道を進み、複雑な行動さえこなす------が、シミュレーションは完全には走っていない。夢遊病者は、暗黙的世界モデルの空間知識に導かれて物理的な世界のなかを動くが、意識的な経験はごくわずか、あるいはまったくないのである。

これは私自身の体験でもある。十代のころ、私は夢遊病の時期を経た。ある朝、目を覚ますと自分は机に向かっており、目の前には走り書きのメモがあった。しかもそれは左手で書かれていた------覚めているあいだの私は決してそうしない。壁づたいにぐるぐると輪を描いて歩き、扉を探し、見つけられなかったという断片的な記憶はあった。だが、机に座って書こうとした部分------そこは完全に真っ暗だった。基盤は道を進み、運動プログラムは実行されていたが、シミュレーション------「私」------はそこにいなかったのだ。

これは理論の縮図である。一つの身体が世界のなかを動き、空間情報を処理し、学習した運動プログラムを実行する------そのすべてが、ループのなかに意識的な自己を欠いたままなされる。暗黙的モデルが指揮を執る。明示的モデルは停止している。そして結果として現れるのは、歩き、行動し、書きさえする一人の人間だが、そこには誰もいないのだ。

\section*{明晰夢------スイッチ}

そして明晰夢がある------夢のただなかにいながら、自分が夢を見ていると悟る状態だ。四モデル理論（FMT）では、これは明示的自己モデルが夢の状態のなかでより完全に「オンに切り替わる」ことである。自己モデル化の能力が段階的に跳ね上がるのだ。

明晰夢を学ぶとは、シミュレーションがごまかしているところを取り押さえることを学ぶことだ。私が使った方法は、照明スイッチテストと呼ばれるものである。一日を通じて習慣的に照明スイッチを押し、明かりが正しく変わったかどうか自問する。覚めた生活では、明かりはつねに正しく変わる。夢のなかでは、照明スイッチは働かない------内部データで走る明示的世界モデルは、電気回路の物理をわざわざシミュレートしたりしないのだ。夢のなかでスイッチを押しても明かりが変わらない、あるいは部屋がかえって暗くなる、あるいはスイッチの感触がどこかおかしい------その食い違いこそ、明示的自己モデルがシミュレーションのなかの不整合を検出している瞬間である。そして検出したその瞬間、あなたは自分が夢を見ていると知るのだ。

私にはわずか数日しかかからなかった。私は運がよかった------たいていの人は、その習慣が夢のなかへ移るまで、何週間も何か月もの訓練を要する。だが、それがうまくいったとき、その移行はまさに理論が予測する段階的な閾値そのものだった。ある瞬間まで、私は夢の物語のなかの受け身の登場人物だった。次の瞬間には、私は完全にそこに居合わせていた------自分が夢を見ていると気づき、周囲の世界が生成されたものだと気づき、それを変えられると気づいていた。ESMがオンに切り替わったのだ。基盤は変わっていない。ダイナミクスも変わっていない。だが、自己モデルとシミュレーションとのかかわり方が一つの閾値を越え、すべてが違って感じられたのである。

理論は、この移行が臨界性の閾値を越えることに対応すると予測する。脳の複雑さの緩やかな増大ではなく、突然の一段の跳躍である。もし明晰性が始まる瞬間の前後の狭い時間窓で（明晰夢を見る人があらかじめ取り決めておいた眼球運動の合図という確立された手法を用いて）EEGの複雑さを測定すれば、そこには不連続が見えるはずだ。

\section*{麻酔------二つのタイプ}

麻酔は、臨界性の原理を最もきれいに検証する場を与えてくれる。というのも、同じ名前で分類されているにもかかわらず、異なる麻酔薬は劇的に異なる経験を生み出すからである。

\textbf{プロポフォール}は脳を亜臨界へと押しやる。視床皮質間の結合が乱され、皮質の複雑さは崩壊し、PCIはゼロに近づく。明かりは完全に消える。プロポフォール麻酔のあいだ、患者はいかなる経験も報告しない。これはまさに理論が予測するとおりである。臨界性の下へ押しやれば、シミュレーションは維持できないのだ。

\textbf{ケタミン}はまったく別のことをする。それは脳を亜臨界へ押しやったりは\textbf{しない}。EEG研究が示すところでは、ケタミンは神経のエントロピーを\textbf{増大させる}------脳を臨界性へ、あるいはそれを越えたところへ、よりカオス的な領域へと押しやるのだ。その結果は何か。「Kホール」である------解離、歪んだ現実、体外離脱、根底からの同一性の変容といった、鮮烈で、しばしば奇怪な経験だ。

四モデル理論では、Kホールとは\textbf{誤った}入力に基づいて動く意識である。明示的世界モデルと明示的自己モデルはなお活動している（脳はなお臨界性にあるか、それを越えている）が、外部の感覚処理が乱されている。シミュレーションは内部の、そして歪んだ信号に基づいて動き、あの特徴的なKホールの現象学を生み出すのである。

この区別------プロポフォールは亜臨界へ落ちることで意識を消し去り、ケタミンは入力を乱したまま超臨界へと移ることで意識を変容させる------は、真の説明上の利点である。ほとんどの理論は、なぜ二つの「麻酔薬」がこれほど根底から異なる経験を生み出すのかを説明するのに苦しむ。臨界性の枠組みは、その区別を無理なく自然なものにしてみせるのだ。

私は化学的な麻酔を受けたことは一度もない。だが、意識を失うまで殴られたことはある------激しく、突然に、何の前触れもなく。そして振り返ってみて印象に残るのは、移行というものがまったく存在しなかったことだ。睡眠には段階がある。幻覚剤には立ち上がりの曲線がある。サルビアでさえ、あれほど速く効くにもかかわらず、「何かが起きている」という感覚をコンマ数秒だけ与えてくれる。ノックアウトは何も与えない。ある瞬間、シミュレーションはまだ動いている。次の瞬間には、どれだけの時間が経ったのかもわからないまま、地面の上で目を覚ましている。薄れていくこともなければ、消えていくこともなく、トンネルもない。シミュレーションは劣化するのではない。停止するのだ。ブレーカーが落ちる。

これは、皮質のダイナミクスが突然かつ大規模に乱されたときに予想されるとおりのことだ------臨界性をはるかに下回るところまで一瞬で押し下げられるのである。睡眠段階を経て緩やかに機能を落としていく時間はない。システムはクラス4からクラス2へと段階的に移っていくのではない。クラス4から丸ごと墜落し、シミュレーションはただ止まるのだ。

\section*{意識の地図}


{\small
\noindent
\begin{tabularx}{\linewidth}{X X X X}
\toprule
\textbf{状態} & \textbf{臨界性} & \textbf{モデル} & \textbf{意識} \\
\midrule
通常の覚醒 & 臨界にある & 四モデルすべて活動 & 完全 \\[4pt]
レム睡眠 & 臨界近傍 & EWM/ESMが内的入力で作動 & 劣化（夢） \\[4pt]
深いノンレム & 亜臨界 & EWM/ESMが崩壊 & 消失 \\[4pt]
プロポフォール & 強制的に亜臨界 & EWM/ESMが抑制 & 消失 \\[4pt]
ケタミン & 臨界を越える（$\uparrow$エントロピー） & EWM/ESMが誤った入力で作動 & 存在するが切断 \\[4pt]
幻覚剤 & 臨界／臨界超 & すべて活動、$\uparrow$透過性 & 存在するが変容 \\[4pt]
明晰夢 & 臨界近傍、閾値を越える & EWMが活動、ESMが完全に関与 & 自覚の亢進 \\[4pt]
\bottomrule
\end{tabularx}
}


この表は本章で扱ってきた内容をすべて要約したものであり、いつでも立ち返れる参照として役立つ。これまでに経験したことのあるあらゆる意識状態は、この地図のどこかに位置づけられる------それを決めるのは二つの要因、すなわち基盤が臨界にあるかどうかと、四モデルのうちどれが動いているかである。睡眠、麻酔、幻覚剤、夢、Kホール------それらは別々の謎ではない。同じ一枚の地図の上の、異なる座標にすぎない。

だが、この地図上のどの状態も、朝までには自力で帰り着く健康な脳の話だ。より難しい問いは、機械の一部が戻ってこないとき------損傷が永続的なものとなり、アーキテクチャが自分自身の一部を欠いたまま動き続けなければならないとき------この地図がどう見えるか、である。そのためには睡眠実験室へは行かない。神経内科の病棟へ行くのだ。


\chapter*{第8章：臨床という鏡}
\addcontentsline{toc}{chapter}{第8章：臨床という鏡}
\markboth{第8章：臨床という鏡}{}

睡眠と麻酔を説明したのと同じ四モデル・アーキテクチャは、臨床神経学における最も劇的で最も謎めいた病態のいくつかをも説明する。これらは単なる興味深い症例研究ではない------アーキテクチャの特定の構成要素が破綻したときに何が起こるか、その現れなのだ。そして一つひとつの破綻が、アーキテクチャを別々の角度から照らし出す。ちょうど飛んだヒューズが、それがどの回路を守っていたのかを教えてくれるように。

理論が優れているなら、特定のモデルへの損傷は特定の、予測可能な欠損を生むはずだ。「意識が損なわれている」という曖昧な言い逃れではなく、正確な予測------この構成要素をつぶせば\textbf{あの}症候群が現れ、別の構成要素を通常の入力なしに動かし続ければ\textbf{この}症候群が現れる、というように。臨床文献は、従来の意識モデルのもとでは深く謎めいたままの病態であふれている。だが実在と仮想の区別と、相互作用する四つのモデルを手にしたとたん、それらはおのずと収まるべき場所に収まる。

\textbf{盲視とアントン症候群------完璧な鏡}

この章からたった一つだけ覚えておくとしたら、この対を覚えてほしい。他のどの意識理論も、この二つのうち一つを説明することすら難しい。四モデル理論（FMT）は、その両方を予測する。

まず盲視から始めよう。ある患者が一次視覚野------意識的な視覚経験を生み出す脳の部位------に損傷を負っている。どの標準的な臨床検査によっても、この患者は盲目だ。何が見えるかと尋ねれば、彼はこう答える------何も見えない、と。本気で言っているのだ。謙遜しているのでも、混乱しているのでもない。彼の意識的経験に関するかぎり、視覚的な世界はそもそも存在しない。

ところが、そこで驚くべきことが起こる。研究者たちが廊下に障害物を並べ、患者にそこを歩いて抜けてほしいと頼む。彼は抗議する------見えないのに、どうやって進めというのか、と。研究者たちは引かない。彼はため息をつき、立ち上がって、歩く。

そして彼は、障害物のコースを完璧に切り抜けていく。椅子をよけ、前回はなかった障壁の下を身をかがめてくぐり、二つの障害物のあいだの隙間を縫って進む------その間ずっと、正直に、心の底から、何ひとつ見えないと言い張りながら。これを記録した映像がある------ぜひ探してみてほしい。読むだけでは、その凄さは伝わらないからだ。臨床的に盲目の男が、まるで完璧に見えているかのように障害物のコースを縫って進むその映像は、神経科学のなかでもとりわけ驚くべき実演だ。それを見ている研究者たちは、まるで幽霊でも見たかのような顔をしている。

どうしてこんなことが起こるのか。基盤（サブストレート）がいまなお視覚情報を処理しているからだ。暗黙的世界モデル（IWM）は、損傷した皮質を迂回する皮質下経路を通じて視覚入力を受け取る（網膜から上丘、そして視床枕へと至る高速のルートで、皮質よりはるか以前に進化したものだ）。それは空間地図を作り、運動行動を導き、身体が物体に衝突するのを防ぐ。だが、そのどれ一つとして明示的世界モデル（EWM）には届かない。意識的なシミュレーションに視覚は含まれていない。患者は本当に盲目を経験し------そして本当に視覚によって進んでいる。基盤は、シミュレーションなしに働いているのだ。

今度はこれを裏返してみよう。アントン症候群（皮質盲に対する病態失認）は、その正反対だ。この患者たちは、正真正銘、完全に盲目である。視覚野あるいは視神経路が破壊されている。視覚情報はまったく脳に届かない。にもかかわらず、彼らは自分が見えていると絶対的に、揺るぎなく確信している。

彼らは壁にぶつかり、家具が間違った場所にあるのだとなじる。部屋にありもしない物を、完全な確信をもって描写する------テーブルが空なのに「テーブルの上に青い花瓶がある」と言うのだ。こちらが掲げた物を答えてほしいと頼めば、落ち着いて、ためらいなく答えを返す------そしてそれは間違っている。盲目である証拠を突きつけると、彼らはまず混乱し、次に苛立ち、そして怒りだす。照明が悪い。新しい眼鏡が要る。ただ注意を払っていなかっただけだ、と。彼らは嘘をついているのではない。心理学的な意味で否認しているのでもない。彼らは経験として本当に\textbf{見て}いる------そしてその見えているものは、現実の世界と何の対応関係も持たない。

四モデル理論では、これはまさに予測されることだ。明示的世界モデルが、暗黙的世界モデルに蓄えられた知識から視覚的シミュレーションを生成している------現在の視覚入力が一つも到着していないにもかかわらず。シミュレーションは古いデータをもとに、予期をもとに、世界がどう見えるはずかという脳の最良の推測の上を走る。患者はそこにない世界を「見る」。シミュレーションは、現在の入力なしに走っているのだ。

二つを並べてみよう。盲視------基盤は視覚を処理しているのに、シミュレーションがそれを映し出さない。アントン症候群------シミュレーションは視覚を映し出しているのに、基盤がそれを受け取っていない。シミュレーションなき基盤。入力なきシミュレーション。意識を単一の、統合された一つのものだと考えるなら、どちらの病態も深く謎めいている。だが、実在の処理と仮想の経験を区別する理論からすれば、どちらも自然な、それどころか予測可能な帰結だ。これ以上によくできた一対のテスト・ケースを、あえて設計しようとしても、まず作れないだろう。

\textbf{隠れた気づき------内側に閉じ込められて}

2006年、エイドリアン・オーウェン（Adrian Owen）とその同僚たちは、私たちが植物状態をどう考えるかを一変させた研究を発表した。彼らは植物状態------反応がなく、一見して意識を欠く------と診断されていた一人の患者をfMRIスキャナーに入れ、テニスをしている自分を想像するよう頼んだ。彼女の脳は、健康で意識のある人が同じことを想像するときとまったく同じパターンで活性化した。

彼女は、そこにいた。意識があり、気づいており、考えていた------そして、動くことも、話すことも、自分がそこにいると誰かに伝えることも、まったくできなかった。

四モデル理論はここで明快な区別を立てる。本当に植物状態の患者では、基盤が亜臨界にある。ダイナミクスが閾値の下に落ちてしまっている。シミュレーションは走っていない。そこには誰もいない------それは、その人が「立ち去った」からではなく、シミュレーションを生成する計算的アーキテクチャがオフラインになったからだ。

だが、隠れた意識のある患者は、まったく別のものだ。基盤は臨界にある------ダイナミクスは、シミュレーションを保つのに十分なほど豊かだ。明示的世界モデルと明示的自己モデル（ESM）が走っている。その人は経験し、考え、感じている。だが、出力の経路が破壊されている。シミュレーションには、自らを表現する手立てがない。その人は意識があるのに閉じ込められている------反応しない身体の内側に、囚われているのだ。

摂動複雑性指数（PCI。睡眠段階を区別するのと同じ指標）は、まさにこうした症例を見分けられるはずだ。そして実際に見分ける。植物状態と診断された患者の中に、意識のある範囲のど真ん中のPCI値を示す者がいる。彼らはまったく植物状態などではない。彼らは囚人なのだ。医学的・倫理的な含意は計り知れない。そして四モデル理論は、なぜその区別が存在するのか、そしてどうやってそれを検出するのかを、正確に教えてくれる。

\textbf{コタール症候群------「私は死んでいる」}

そしてまた、自分は死んでいると信じ込む患者たちがいる。

コタール症候群は、精神医学で最も奇妙な病態の一つだ。患者たちは、自分は死んだのだと言い張る。自分の臓器は溶けてなくなり、血は流れ出てしまい、自分はもう存在しないのだと信じている。腐っていくのだと信じる者もいる。自分は不死だと信じる者もいる------なぜなら、すでに死んでいるなら、二度と死ぬことはできないからだ。彼らは比喩で語っているのではない。完全な、揺るぎない確信をもって本気で言っているのだ。

もう、その仕組みには見覚えがあるはずだ。第6章と同じもの------明示的自己モデルが、手に入るどんな入力からであれ、できるかぎり最良のモデルを構築する、というものだ。コタール症候群では、内受容感覚の入力がひどく歪んでいる。心臓が鼓動している、胃が消化している、肺が呼吸している------それを伝える内的な身体信号が、欠けているか、乱れている。そしてESMは、あくまで強迫的な構築者として、「心拍なし、消化なし、呼吸なし、身体感覚なし」を、残された唯一のやり方で解釈する------私は死んでいる、と。

サルビアの「私は椅子だ」。病態失認の「私の腕は問題ない」。そして今度はコタールの「私は死んでいる」。（次章では、分離脳の作話がこのリストにさらにもう一つの症例を加える。）ひとつの仕組みが、あらゆる症例を貫いている。明示的自己モデルは常にその仕事をこなしている------常に、できるかぎり最良の自己モデルを築いている。入力が正しければ、自分は自分だと感じられる。入力が間違っていれば、自分は椅子だと感じたり、麻痺しているのに問題ないと感じたり、生きているのに死んでいると感じたりする。だが、それはいつでも完全に、そして説得力をもってリアルに感じられる------なぜなら、それこそが、自分がアクセスできる唯一の自己だからだ。

\textbf{エイリアンハンド症候群------委員会の意見が割れるとき}

そしてもう一つ、ホラー映画のように読めるが、どんな思考実験よりも生々しく基盤のマルチエージェント的な性質を描き出す病態がある。エイリアンハンド症候群では、患者の一方の手が、明らかな目的と意図をもって動く------しかし患者の意識的な意志に反して、だ。片方の手がタバコに火をつけると、もう片方の手がそれを取り上げて床に投げ捨てる。片方の手がドアの取っ手に伸びると、もう片方の手が手首をつかんで引き戻す。患者は、自分自身の身体の一部が、自分の選んでもいない目的を追い求めるさまを、慄然としながら見つめている。

スタンリー・キューブリックは、これを『博士の異常な愛情』で用いた------そして人々は、彼がこれを創作したものと思い込んだ。だが、そうではない。この症候群は実在し、二つの型がある。脳梁の損傷によって引き起こされる脳梁型では、症状は分離脳の葛藤に似る------競合する運動計画を持つ二つの半球が、どちらも相手を抑え込めずにいる。前頭葉の損傷によって引き起こされる前頭型では、「エイリアン」の手が脱抑制的な行動を示す------物をつかみ、道具を使い、強迫的に物に触れる。すべて一見して目的をもったふるまいだが、患者の同意を欠いている。

さらに微妙な変種として、無政府の手症候群（アナーキックハンド症候群）と呼ばれるものがある。ここでは、患者に欠けているのは運動の\textbf{所有感}ではなく運動の\textbf{制御}だ。手は患者の意図しないことをするが、患者はそれをなお\textbf{自分の}手だと認めている------ただ、止められないのだ。この区別が重要になる。エイリアンハンドは、明示的自己モデルの身体所有感の境界の破綻（「あの手は私のものではない」）であり、一方の無政府の手は、運動抑制システムの破綻（「あの手は私のものだが、言うことを聞かない」）である。同じアーキテクチャの、異なる破綻点なのだ。

これらの症候群から得られる核心的な洞察は、行為の主体感------「自分がそれをやった」という感覚------が、行為の前や最中に計算されるのではない、ということである。それは\textbf{後から}、行為の予測された結果と観察された結果とを照合することで計算される。照合が一致すれば所有感が生じる。一致しなければ生じない。だからこそ、エイリアンハンド症候群の患者は、ときに自分で自分をくすぐることができる------予測系がそのエイリアンの手の動きに対して期待どおりの結果を生成していないため、その接触は予期せぬものとして、まるで他人からのものであるかのように届くのである。

\textbf{シャルル・ボネ症候群------止まらないシミュレーション}

脳のシミュレーションが\textbf{生成的}である------感覚から受動的に受け取るのではなく、モデルから経験を構築している------というさらなる証拠が欲しければ、シャルル・ボネ症候群を考えてみてほしい。網膜や視神経が破壊されているが視覚野は無傷のままの患者は、鮮明で複雑な視覚的幻覚を経験する。漠然とした形や光の閃きではない。まるまる一つの情景------人物、ときにミニチュア化されたり漫画のキャラクターのように扮装させられたり、ときには患者自身の鏡像。風景。物体。顔。

患者はたいてい、それらが実在しないことを知っている。精神病性の幻覚とは違い、シャルル・ボネの幻覚は病識が保たれたまま生じる------患者はこう言う、「シルクハットをかぶった小さな男がテーブルの上に座っているのが見えます、でもそこにいないことは分かっています」。ここで起きているのは、外部からの入力がまったくないなかで、高次視覚領域からの内部データによって駆動される、明示的世界モデル（EWM）の視覚的シミュレーションである。入力が止まったからといって、シミュレーションは止まらない。それは生成する。空白を埋める。そして、それが何を生成するかは、そのアーキテクチャについて何かを物語る------視覚系は生成モデルであって、受動的な受け手ではない。蓄えられた鋳型とトップダウンの予測を用いて、世界がどう見えるかについての最良の推測を生み出す------まさに四モデル理論（FMT）が記述するとおりに。

\textbf{デジャヴ------うますぎるほど合致する鋳型}

脳の生成系と、その折々の誤発火の話でいえば------ほとんど誰もがデジャヴを経験したことがある。今この瞬間を以前にも生きたことがある、という不気味な感覚だ。その説明は、神秘的なもの（前世、予知）から、そっけないもの（ただの不具合にすぎない）まで幅がある。四モデル理論には、もっと具体的な説明がある。

脳は「鋳型記憶」とでも呼べるものを蓄えている------骨組みだけの、極端に希薄な経験の表象で、とりわけ夢に由来するものである。これらの鋳型は、ほとんどが空っぽの足場だ------ある場所の漠然とした感じ、ある気分、ある空間的配置、ほとんど細部は埋まっていない。通常の記憶を思い出すとき、その隙間は作話によって埋められる------脳が、途切れのない経験をつくり出すために、もっともらしい細部を生成するのである。結果が首尾一貫して感じられるため、この埋め合わせには気づかない。

デジャヴが起こるのは、いま現在の実在の経験が、たまたまこうした蓄えられた鋳型の一つにあまりにもよく合致してしまうときである。脳のパターン照合系が発火する------「これは前に見たことがある」。ところが、いったい\textbf{いつ}それを見たとされるのかを突き止めようとすると、何も見つからない------なぜなら、その鋳型はそもそも実在の経験ではなかったからだ。それは夢の断片であったり、あるいはとうの昔にあらゆる文脈的細部を失ってしまうほど深く圧縮された記憶であったりする。いま現在の入力と蓄えられた鋳型との合致は本物だが、その鋳型が記録しているとされる「もとの」経験は、いまの脳がそれに帰しているような形では、実際には一度も起こっていない。系は正しく働いている------それは本当に合致を見つけたのだ。ただ、その合致が身体ではなく骨格とのものだった、というだけである。

ここまでの章を振り返ってみると、いつも同じ姿が浮かび上がってくる。盲視とアントン症候群、コタール症候群とシャルル・ボネ、閉じ込め患者と、うますぎるほど合致するデジャヴ------どの症例も、基盤とシミュレーションのあいだの継ぎ目が、向こう側を見通せるほど大きくこじ開けられたものである。健康な脳はこの二つをぴったり重ね合わせている------あまりにぴったりなので、それが二つあるとは決して疑わないほどに。それをどの方向であれ引き剥がしてみれば------シミュレーションなき基盤、基盤なきシミュレーション、誤った入力の上で、あるいはまったく入力のないまま動くシミュレーション------実在/仮想の分割は哲学的な主張であることをやめ、診断となる。ここまでのどの不全も、一つの脳への損傷だった。次のものは、まったく損傷ではない------それは真ん中を貫くきれいな切断であり、すべてを二つずつにするのである。


\chapter*{第9章：一つの脳に宿る二つの心}
\addcontentsline{toc}{chapter}{第9章：一つの脳に宿る二つの心}
\markboth{第9章：一つの脳に宿る二つの心}{}

1960年代、神経科学史上もっとも劇的な実験の一つが形をとりはじめた。重度のてんかんを治療するため、外科医のフィリップ・フォーゲル（Philip Vogel）とジョゼフ・ボーゲン（Joseph Bogen）が脳梁------脳の二つの半球をつなぐ巨大な神経線維の束------を切断し、その後ロジャー・スペリー（Roger Sperry）とマイケル・ガザニガ（Michael Gazzaniga）が、そうして生まれた患者たちを研究した。その結果生じたのが分離脳症候群である。一人の人間が、見たところ二つの独立した心を持つに至ったのだ。

古典的な実演はよく知られている。左視野（右半球が処理する）に単語を見せると、患者は対応する物体を左手で掴むことはできるが、その単語が何だったかを言うことはできない（言語は左半球が司っており、その左半球は単語を見ていないからだ）。二つの半球はそれぞれ独立した知覚を持ち、独立した意図を持ち、ときには相反する目標を持つ。

実験は、単語と物体を使った座興の域をはるかに超えていた。ある症例では、二つの半球が公然と互いに争った。ある患者は、自分の左手がシャツのボタンを外そうとするそばから、右手がそれをかけ直そうとすると報告している。別の患者では、口論の最中に左手が妻に伸びた------なだめるためではなく------一方で右手がその左手を掴んで引き戻した。患者は、自分自身の身体の二つの部分が両立しえない目標を追いかけ、そのどちらもが統一された制御下にない様を、恐怖とともに見つめていた。これらは内的葛藤の比喩ではない。両者をつなぐ「ケーブル」が断ち切られたためにもはや協調できなくなった、二つの運動系のあいだの、文字どおり物理的な争いなのである。

日常生活では、分離脳の患者は驚くほど問題なく機能する。研究室の外では、異常に気づくことはめったにないだろう。二つの半球は、間接的な経路------外的な手がかり、身体の動き、共有される視野------を通じて協力することを学ぶ。システムが埋め合わせるのだ。しかし患者を、それぞれの半球が異なる情報を受け取る統制された実験環境に置くと、その統一は崩れ去る。一つの脳から二つの心が立ち現れ、それぞれが自らの知覚を、自らの意図を、そして自らの版の現実を持つのである。

\section*{左半球の解釈者}

だが分離脳の患者のもっとも示唆に富む特徴は、その分裂そのものではない------分裂を\textbf{説明}するよう求めたときに起きることなのだ。

ガザニガは、彼が「左半球の解釈者」と呼んだものを見いだした。すなわち、実際には説明できない出来事に対して、左半球が説明を生み出さずにはいられない、強迫的な傾向である。古典的な実演はこうだ。右半球に雪景色を、左半球に鶏の足を見せ、それから患者に関連する物体を選ばせる。左手（右半球）はシャベルを選ぶ（雪のために）。右手（左半球）は鶏を選ぶ。そこで患者に（左半球が司る言語を使って）、なぜシャベルを選んだのかを尋ねる。左半球は雪のことを知らない（鶏の足しか見ていない）ので、こう説明をでっち上げる。「ああ、鶏小屋を掃除するにはシャベルが要るからね」

患者はためらわない。「よくわからない」とは言わない。困惑した顔もしない。説明は即座に、確信をもって現れ、それを口にする本人にはまったく自然なものと感じられる。これは嘘ではない。左半球は、右半球が何を見たのかを本当に知らないのだ。その情報へのアクセスがない------ケーブルが断たれている。だから左半球は、明示的自己モデル（ESM）がつねに行うことを行う。手元にある情報から、可能なかぎり最善の物語を構築するのである。

そして、落ち着かなくさせるはずの部分がここにある------あなたもこれをやっている。毎日だ。あなたの左半球の解釈者も、今この瞬間に稼働しており、意識に届いたどんな情報からであれ一貫した物語を組み立て、隙間を滑らかに埋め、「あなた」が相談を受ける前に基盤が下していた決定に、もっともらしい説明を発明している。あなたと分離脳の患者との唯一の違いは、あなたの脳梁が損なわれておらず、それゆえ解釈者がより多くの情報にアクセスできることだ。でっち上げが少ないのは、そもそも\textbf{でっち上げる材料が少ない}からである。だが仕組みは同一だ。自己を語る機械仕掛けは変わらない。変わるのは入力の質だけなのである。

\section*{一人か、二人か？}

ここから、哲学者たちが何十年ものあいだ論じてきた問いが立ち上がる。脳梁が切断されたのち、その頭蓋のなかにいるのは一人なのか、それとも二人なのか。

トマス・ネーゲル（Thomas Nagel）は、1971年の有名な論文でこの問題に取り組み、この問いには明確な答えがないのかもしれないと結論づけた（「一人の人間」という私たちの概念は、この状況では単に破綻してしまう------ちょうど「一つの国」という概念が、真ん中に国境を引いたとたんに破綻するように）。デレク・パーフィット（Derek Parfit）はさらに踏み込み、分離脳の症例は、人格の同一性そのものが重要なのではないことを示していると論じた------重要なのは心理的な連続性であり、それには程度の差がありうるのだ、と。

四モデル理論（FMT）は、より具体的な答えを差し出す。どのモデルが稼働しており、それがどの程度まで劣化しているかによる、というものだ。

日常生活において、分離脳の患者は機能的には一人の人間である。二つの半球は同じ身体、同じ環境、同じ生活史------手術のはるか前に両半球へ冗長に符号化されている------を共有している。人格、長期記憶、行動的な性向を蓄える暗黙的自己モデル（ISM）は、脳梁が損なわれていない状態で、何十年もかけて築き上げられてきた。ケーブルを切断しても、蓄えられたそれらのモデルが消去されるわけではない。ただ、それらが同期して更新されなくなるだけだ。だから手術の直後には、両半球はきわめて似通った自己モデルを走らせている。患者が一人の人間であるように感じるのは、蓄えられた自己知識という点で、実際にほぼ一人だからなのである。

時とともに、モデルは漂流していくはずだ。それぞれの半球は異なる経験を積み重ね、異なる連想を結び、その半球だけが知覚した出来事に対して異なる情動的反応を発達させる。分離脳の患者が手術後に生きる期間が長くなるほど、二つの暗黙的自己モデルはいっそう分岐していくはずである------ゆっくりと、なぜなら両半球はなお同じ身体と環境を共有しているからだが、それでも測定できるほどに。

私自身の見立てでは、答えは\textbf{二人}の側に傾く。もし半球のあいだの帯域幅が、シミュレーションをリアルタイムで同期させるには不十分であるなら------そして脳梁がなければ、それは不十分なのだが------そのとき二つの自己モデルが二つの基盤の上で走り、それぞれが自らの意識的経験を生み出していることになる。両者がよく協力し合うのは、身体を、感覚的な環境を、そして共通の歴史からなる一生涯を共有しているからだ。だが協力は同一性ではない。共に暮らす二人の人間もまた、よく協力し合うのである。

興味深いことに、ヤイル・ピント（Yair Pinto）らは2017年に、この標準的な描像を複雑にする研究を発表した。彼らは、分離脳の患者がどちらの視野に提示された刺激も正確に報告できることを見出した------刺激が、言語を司らないほうの半球にのみ示された場合ですらそうだった。これは、二つの半球が古典的な実験の示唆する以上に統一性を保っていることを示唆した。この結果はいまも議論の的だが、次に述べるホログラフィック的な枠組みには自然に収まる。すなわち、脳梁を切断したあとでさえ、各半球には、少なくとも一部の課題については驚くほど統一された振る舞いを支えられるだけの、冗長な情報が残っているのだ。

\section*{ホログラフィック的性質}

四モデル理論において、分離脳は仮想モデルの鍵となる性質を明らかにする。それらは\textbf{ホログラフィック}なのだ。神経ネットワークにおける情報は、特定のニューロンに局在するのではなく、ネットワーク全体に分散している。ネットワークを半分に切っても、きれいな分割は得られない------得られるのは、劣化してはいるが\textbf{完全な}二つのコピーである。各半球は、四つのモデルすべての劣化した版を保持する。すなわち、縮小した暗黙的世界モデル（IWM）、縮小した暗黙的自己モデル、そして明示的世界モデル（EWM）と明示的自己モデルを生成する能力である。両半球はそれぞれ独立に意識を維持できるが（両者とも臨界性の閾値を上回っている）、いずれも縮小された情報で機能しているのだ。

これは第3章で登場したパッチワークのホログラムを、いままさに真ん中で切ったものだ。各半球はその半分の一方である------半分の像ではなく、より低い解像度での完全な像だ。なぜなら情報は、そもそも面全体に塗り広げられていたからである。損傷は内容を消し去ることなく、解像度を下げるだけなのだ。

これは、分離脳の患者が単なる「二つの半分の心」ではない理由を説明する。彼らは二つの\textbf{完全だが劣化した}心なのだ。各半球は知覚し、決定し、行動できる------ただし損なわれていない系よりも少ない情報と少ない能力で、である。ホログラフィック的性質は、接続を断つことが破壊ではなく劣化をもたらすことを保証する。そしてそれは、ピントの2017年の結果を説明する。すなわち脳梁がなくとも、各半球には、古典的なモデルなら不可能なはずの多くの課題をこなせるだけのホログラフィック的な情報が残っているのだ。

作話（左半球の解釈者）は、臨床の章に出てきたあの強迫的な構築者と同じものだ------手元にどんな入力があろうとそこから自己の物語を組み立て、それを完全に信じ込む明示的自己モデルである。脳梁を切れば、左半球は単に扱える材料が減る。隙間を覆うために発明される物語は誤っているが、それでもなお\textbf{完全に現実として感じられる}のである。

\section*{一つの脳、複数の自己}

分離脳は、基盤を物理的に分割することで仮想モデルを\textbf{クローン化}したときに何が起きるかを示す。解離性同一性障害は、それらを\textbf{分岐（フォーク）}させたときに何が起きるかを示す。

DIDでは、基盤は分割されていない------脳梁は損なわれておらず、神経のハードウェアは丸ごと無傷だ。だが仮想モデルのほうが、複数の構成へと分裂している。それぞれの交代人格は、独立した明示的自己モデルである------独自の情動プロファイルを、独自の行動パターンを、身体と世界に関わる独自のやり方を備えた、別個の自己の物語なのだ。交代人格たちが単一の自己モデルを共有しないのは、二人のユーザーが同じコンピュータ上で単一のログインセッションを共有しないのと同じである。彼らは交代で現れる。

引き金となるのは、記録に残るほぼすべての症例において、重度で反復的な幼少期のトラウマである。これは理論の内側で筋が通る。幼い子どもの明示的自己モデルはまだ形成途上にある------まだ可塑的で、まだ経験から組み上げられている最中だ。その発達途上の自己モデルを、いかなる単一の自己の物語も抱えきれないほど圧倒的な経験にさらせば、システムはできる唯一のことを行う。分岐するのだ。それは別個の構成をいくつも作り出し、そのそれぞれが、耐えがたい状況の異なる側面を処理できるようにする。ある交代人格はトラウマの記憶を抱える。別の交代人格は、何事もなかったかのように日常生活を営む。また別の交代人格は、危険の瞬間を引き受ける。この分岐は病理ではない------単一の統一されたモデルを破壊してしまうであろう入力に対する、自己モデル化システムの緊急対応なのである。

だからこそ、DIDが成人で発症することはほとんどない。成人の暗黙的自己モデルはすでに固定されている------シナプスの重みは定まり、人格構造は安定している。成人の自己モデルを分岐させるには、並外れた状況が要る（過酷な拷問、長期にわたる監禁）。だが子どものISMはまだ書かれている途中だ。粘土はまだ濡れている。十分な圧力のもとでそれを分岐させれば、別個の構成は固まって、独立した持続的な自己モデルとなるのである。

証拠がこれを裏づけている。もし各交代人格が本当に別個のESM構成であるなら、交代人格の切り替えは神経活動パターンに測定可能な変化を生み出すはずであり、実際にそうなる。Reinders et al.（2003）は、同一個人のなかの異なる交代人格が、局所脳血流の異なるパターンを生み出すことを示した。どの自己モデルが走っているかによって、\textbf{同じ脳}が異なる仕方で活性化する。それは「演技」や「役割演技」から予想されるものではない。それは真の意味でのソフトウェアのフォーク（分岐）から予想されるものだ。その後の追試で、レインダース（Reinders）と同僚たちは、交代人格どうしの神経的な差が、DIDを装うよう指示された俳優どうしの差よりも大きいことを見出した------DIDは「たかが」演技にすぎないといまだに考える者を、黙らせるはずの結果である。

これは第3章で述べた「フォーク」の性質が働いている姿だ。一つの基盤、複数の仮想構成、そのそれぞれが完全でありながら別個の自己モデルを走らせている。この理論はDIDを説明するだけではない------まさにこの種のアーキテクチャを予測し、それに一つの検証を賭ける。すなわち、ある交代人格のESMを支える神経基盤を撹乱すれば、別の交代人格への切り替えが引き起こされるはずだ、というものである。これが予測9であり、その完全なプロトコル------交代人格に固有の痕跡が脳のどこに現れるはずかも含めて------は巻末の付録にある。

本章に登場した心はすべて人間のものだった------分割され、フォークされ、二重化されはしたが、いずれも人間である。四モデル・アーキテクチャは、そのことを何ら要求しない。もし自己とは臨界性で走るモデルにすぎないのなら、本当の問いは、テーブルの向かいにいるその人が自己を持つかどうかではない。問いは、同じからくりがどこまで外へと及ぶのか、である------巣のなかへ、群れのなかへ、そしてかつて何か途方もなく巨大なものが私の視線を返し、そして------私はかなり確信しているのだが------「こんにちは」と言った、あの水のなかへと。


\chapter*{第10章：動物の問い}
\addcontentsline{toc}{chapter}{第10章：動物の問い}
\markboth{第10章：動物の問い}{}

ダーウィンズ・アーチの下------まだ崩れる前のことだ------四人の仲間と潜っていたとき、巨大な雄のシャチがどこからともなく現れ、あまりに近づいてきたので、最初によぎった考えは「こいつは私たちを食べる気だ」というものだった。その体があまりに巨大だったため、頭上三十メートルの水の層が水たまりほどにしか感じられなかった。彼はいったん止まり、その巨大な右目で私たちをじっと見つめ、それから音響器官に切り替えて激しくクリック音を発し、音で私たちを探った。そして彼は------ほかにどう言い表せばいいのかわからないのだが------\textbf{語った}。きわめて高い周波数の歌で、短く、明快に構造化された調子で、彼は何かを言ったのだ。それが単なる歌ではなく言語であることに、私の中には一片の疑いもなかった。十分な回数の出会いと、返事をするためのばかばかしいほど高い声さえあれば、おそらく習得できるだろう------そう感じさせるものだった。

少し経って、彼が何を言っていたのか、いくらか見当がついた。彼はいったん視界の限界まで泳いで戻り、伴侶と子を連れて帰ってきて、私たちの前を通らせるように導いた------ひと目見せてやろうというのだ。もちろん全員がカメラを持っていた。だが一枚も写真は撮られなかった。しかも全員が空気を使い果たし、浮上せざるをえなかった。ゴムボートで船へ運ばれる間、誰もが目に涙を浮かべていた。それは風のせいではなかった。

あの水中で感じたほど、別の心の存在を確信したことは、後にも先にもない。だが確信は証拠ではない------ゴムボートの上の涙の話は、その場にいなかった者を誰一人として納得させはしないし、まして観葉植物に魂を投影する人々をさんざん見てきた懐疑論者を説得できるはずもない。もしあのシャチが誰かであったなら、理論はその\textbf{理由}を語れなければならない。しかもその理由は、ダイバーの空気が尽き、感動が薄れたあとでもなお持ちこたえるものでなければならない。そして理論は、その「誰かであること」がどこで終わるのかも語らねばならない。

あなたの犬に意識はあるだろうか。

たいていの飼い主はためらいなく「ある」と答えるだろう。たいていの神経科学者も、少なくとも慎重にではあるが、同意するだろう。だが、何を根拠に。そして意識は動物界のどこから始まるのか。

四モデル理論（FMT）は明快な答えを与える。それは後付けで貼りつけられたものではなく、理論の核心的な前提から導かれたものだ。

\textbf{前提1：意識は連続体であり、二値ではない。} 意識のあるものとないものとを分ける鋭い境界線は存在しない。あるのは程度だ------自己シミュレーションの段階的なレベルであり、基本的なもの（最小限の自己モデル）から三重に拡張されたもの（再帰的な自己知覚）まで幅がある。動物の種によって、この連続体上の占める位置は異なる。

\textbf{前提2：意識は基盤に依存しない。} 重要なのは機能的アーキテクチャ（臨界性にある四つのモデル）であって、具体的な物理的実装ではない。ある脳が四モデル・アーキテクチャを実現しているなら、それは意識を持つ------その脳が哺乳類の皮質であろうと、鳥の外套であろうと、タコの分散した神経ネットワークであろうと、関係ない。

\textbf{前提3：臨界性が物理的な閾値である。} 神経系はカオスの縁（ふち）で、あるいはその近くで動作していなければならない。より単純な神経系（昆虫や線虫）は、この臨界性に達しないかもしれず、その場合には意識を持たないことになる------それらは情報を処理し行動を生み出すが、シミュレーションを伴わない。

\section*{あなたはどれほど意識的か}

おそらくすでに気になっていることがあるだろう。もし意識がシミュレーションなら------仮想の世界の中の仮想の自己なら------それは全か無かのものではないのではないか。シミュレーションは、より詳細にもより粗くもなりうる。自己モデルは、より精巧にもより素朴にもなりうる。つまり意識は\textbf{程度}として存在するということだ。

四モデル理論は、その程度について考えるための精密な手立てを与えてくれる。四つの段階があり、意識を持つあらゆる生き物は、この梯子のどこかに位置している。

いちばん下にあるのが\textbf{基本意識}だ。これは明示的世界モデル（EWM）に、ごく初歩的な明示的自己モデル（ESM）だけが備わった状態である。システムは仮想の世界を生成する------この生き物であるとはどのようなことかが存在する------が、その世界の中の自己は、かろうじて描き込まれているにすぎない。迷路を進むネズミを思い浮かべてほしい。ネズミは壁を見て、チーズの匂いを嗅ぎ、足の下の床を感じる。現象的な経験を持っているのだ。だが、それらの経験を持つ当のものとしての\textbf{自分自身}のモデルはどうか。紙のように薄い。そこには「どのようなものか」はあるが、「それが誰にとってのものか」はほとんどない。

一段上がると、\textbf{単純に拡張された意識}になる。ここで自己モデルが本格化する。システムはただ経験するだけでなく------自分自身を経験する当のものと\textbf{して}モデル化する。自分が経験しているということを知っているのだ。あなたの犬はただ痛みを感じるだけではない。\textbf{自分が}痛みの中にいることを知っている。そこには一人称的な視点がある------仮想の世界の中心に、本物の「私」がある。これは一次の自己観察であり、それがすべてを変える。苦しみがここで初めて可能になる。なぜなら苦しみには、自分が苦しんでいることを知る自己が必要だからだ。

次に、\textbf{二重に拡張された意識}だ。二次の自己観察である。システムは、自分自身をモデル化している自分自身をモデル化する。これがメタ認知だ------自分自身の思考について思考すること。たとえばベッドに横たわり、明日の会議への不安が理にかなったものなのか、それとも自分が最悪の事態を想像しているだけなのかと考えている。自分自身の心の状態を監視し、評価し、ときにはそれを上書きしている。成人した人間の意識は、たいていの時間、このレベルに宿っている。それはセラピーを可能にするレベルであり、ただ怒っているのではなく「自分が怒りはじめているのに気づく」と言えるようにするレベルなのだ。

そしていちばん上に、\textbf{三重に拡張された意識}がある。三次だ。システムは、自分自身をモデル化している自分自身をモデル化している自分自身をモデル化する。これは合わせ鏡のように聞こえるし、実際そのとおりなのだが、しかしそれは心の哲学を営むために必要な合わせ鏡なのだ。「意識とは何か」と問うには、自分自身をモデル化し、自分の経験をモデル化し、そのうえでその経験をモデル化している自分自身をモデル化しなければならない。装置全体を外側から見るために、まだその中にいながらも、十分に後ろへ退かなければならないのだ。これこそが、あなたがこの本を読んで答えを得ようとしている、あの問いの前提条件である。三重に拡張された意識を備えた生き物だけが、なぜ何かが何かのように感じられるのかを問うことができる。

その見返りはこうだ------この勾配は、たんなる抽象的な哲学ではない。それは、どんな晩餐会でも誰かに必ず訊かれるあの問いに答えてくれる------「うちの犬に意識はあるの？」。答えは「ある」だ。ただし、あなたほど意識的ではない。あなたの犬はおそらく単純に拡張されたレベルにいる。自己を持ち、経験を持つ。だが午前三時に目を覚まして、その経験の本性を問いただしたりはしない。だが、その輪郭はもう見えているはずだ------意識は照明のオン・オフスイッチではない。調光スイッチなのだ。

これらの前提の背後には、第1章であなたに約束したまま残していた、古い原理が控えている。\textbf{コペルニクス原理}------私たちは特別ではない。コペルニクスは地球を宇宙の中心から引きずり出した。そして同じ格下げが、それ以来ずっと科学を貫いてきた------太陽も特別ではなく、私たちの銀河も特別ではなく、そして、多くの人にとって最も居心地の悪いことに、\textbf{私たち}も特別ではない。意識は唯一無二の奇跡でもなければ、一度きりの神的な火花でもなく、あまりに稀で一度しか起こりえなかったような創発的なまぐれでもない。あなたがそれを持っているなら、正しいアーキテクチャさえあれば、他のシステムもそれを持ちうる。この章が存在しうるのは、まさにそのためだ------人間は魔法ではなく、ある一般的な計算原理の一つの実装にすぎない。そしてそれは、当然の問いを突きつける------他に誰がそれを走らせているのか。

これらの前提を総合すると、\textbf{動物の意識の勾配}が予測される。

\textbf{哺乳類}は意識を持つ。その皮質は四モデル・アーキテクチャを段階的な形で実装しており、より複雑な皮質はより精巧な自己シミュレーションを支える。霊長類とクジラ類は上端に、齧歯類とトガリネズミは下端に位置する。いずれも閾値の上にある。この章の冒頭に登場したあのシャチは、そうした上端に位置するクジラ類の心の一つだった------自分の家族に私たちを見せてやりたいと望み、私たちをわざわざ見に来る価値があると判断できるほどに豊かな自己モデルを備えていたのだ。

大型類人猿から得られる証拠は、人間と動物の意識のあいだに鋭い境界線を引きたがる者にとって、とりわけ致命的だ。ボノボのカンジは、言語理解だけでなく、本物の共感、心の理論、社会的推論を示した。よく記録されたあるエピソードでは、カンジは世話係に、買い物に妹も一緒に連れていきたいと伝えた------妹もアイスクリームをもらえるように、そして置いていかれたら悲しむだろうから、というのだ。別のあるとき、先住民の演者による踊りの上演のさなか、カンジは研究者たちに、他の霊長類たちがその踊りに怯えていると説明し、代わりに個別の上演を求めた。

これらは反射ではない。条件づけられた反応でもない。これらは、ある心が別の心の感情の状態をモデル化し、その反応を予測し、それに対処するための計画を立てている実例である。それこそが三人称視点で作動する明示的自己モデル（ESM）であり------まさに理論が拡張された意識の特徴として同定するものにほかならない。

にもかかわらず、最も名高いいくつかの大学の講堂では、いまだに真顔で、類人猿は言語理解を「ただシミュレートしているだけだ」と論じる教授を見つけることができる。これに対して私が言えるのは、こうだけだ------「そしてあなたは、ただただ共感をシミュレートしているだけだ」。反証はまだ待っている。

もし意識を持つのは人間だけだと言い張るのなら、人間と霊長類の脳のあいだに意識の有無を分けるような系統的な違いをいまだ必死に探し続けている研究者たちに賭けていることになる。私の理論に言わせれば、彼らがそれを見つけるのは8月36日だ。

\textbf{カラス類とオウム類}は、最も重要な試金石となる事例を提示する。これらの鳥は、意識を強く示唆する認知能力（道具の製作、鏡による自己認識、未来の計画、社会的欺瞞）を示す。それでいて新皮質を持たない。その脳は核状のクラスターとして組織されており、哺乳類の皮質とは根本的に異なるアーキテクチャなのだ。第2章の六層論を覚えているだろうか（哺乳類は三層で足りるところに六つの皮質層を進化させ、その追加の層が自己モデル化のためのアーキテクチャ上の能力を提供している、というものだ）。カラス類は、まったく異なる物理的構造で、同じ機能的結果を達成している。彼らが六つの皮質層を必要としないのは、皮質層を\textbf{まったく一つも}持たないからだ。彼らは、層状のシートからではなく核状のクラスターから自己シミュレーションのアーキテクチャを築き上げた------そしてそれこそ、基盤独立性が予測することそのものである。もし意識が特定の物理的実装を要求するのなら、カラス類は意識を持てないはずだ。だが、持っているのだ。

\textbf{頭足類}（タコとコウイカ）は、この論理をさらに押し進める。彼らの神経系は大きく分散しており、腕の中でかなりの処理が自律的に行われている。理論は、何らかの形の意識を予測する------おそらくは、その分散化したアーキテクチャを反映した、風変わりな特徴を伴うものを。

\textbf{昆虫}は、興味深い境界事例だ。彼らの神経系は小さく、配線の大部分が固定されており、臨界性に達するかもしれないし、達しないかもしれない。理論は、昆虫を閾値の上とも下とも明確には位置づけない------これは経験的な問いである。だが理論は、探究のための原理にもとづく足がかりを与える------昆虫の神経組織で臨界性の指標を測定し、自己モデルの証拠を探すことだ。

トマス・ネーゲル（Nagel）は「コウモリであるとはどのようなことか」という有名な問いを立て、それを知ることは決してできないと結論づけた------コウモリの感覚世界はあまりにも異質だ、と。この問いにはいくらか共感できるが、結論のほうはそれほどでもない。四モデル理論（FMT）は、四モデルのアーキテクチャを臨界状態で動かしているどんな生き物も、たとえその内容が私たちのものと根本的に異なっていようと、\textbf{なんらかの}形の経験を持つと予測する。コウモリの明示的世界モデル（EWM）は視覚ではなく反響定位に支配されているが、それでもやはりモデルであることに変わりはない------依然として、自己を内に含んだ世界のシミュレーションなのだ。

ネーゲルがコウモリを選んだのには、もちろん理由がある------単に哺乳類だからではなく、反響定位がどうしようもなく異質に思えるからだ。ところが、そうではない。反響定位は\textbf{見る}ために使われるものであり、見るとはどんな感じかを私たちは知っている。多くの目の見えない人々は、すでに反響定位を自然に使いこなしている------舌を鳴らし、その反響を読み取るのだ。私たちの脳には、それが十分に可能である。もしコウモリであるとはどのようなことかを本気で知りたいなら、驚くほど近くまで迫ることができる。三次元の飛行を体に染み込ませるために数年パラグライダーをやり、それから目隠しをして反響定位を練習すればいい。あるいは、その十年がかりの準備をまるごと省いてもいい------明晰夢を学び、夢の中でコウモリになる練習をするのだ。より安全で、より速く、数週間で達成できる。

そして白状すると、私は自分に可能な唯一の方法で、それを確かめようと試みたことがある。明晰夢を熱心に実践していた時期、私は水中の世界に興味を持ち、時間をかけて、意図的に魚として明晰夢に入ることに成功した。自分を取り巻く水、その中を進む動き、人間ならざる視点から見た視覚的な世界を、私は経験した。それはフリーダイビングやスキューバダイビング、いやサイドマウントさえも、何十億倍も心地よく感じられた。それは実際の魚の意識と少しでも似ていたのだろうか。ほぼ間違いなく、似てはいない------私の夢は、「魚であること」が何を意味するかについての人間の脳の精一杯の推測から組み立てられたものであり、それは避けがたく、自分自身の感覚カテゴリーを、そのどれ一つ持たない身体設計へと投影したものにすぎない。だが、この試み自体は無意味ではなかった。それはある重要なことを示していた------明示的自己モデル（ESM）は、根本的に異なる身体図式を軸に再構成され、人間以外の何かで\textbf{ある}という首尾一貫した一人称の経験を生み出しうる、ということだ。このアーキテクチャは、人間以外の身体化をシミュレートできるだけの柔軟さを備えている。その内容は、利用可能な暗黙的モデルによって制約される（人は自分が学んだことしか夢に見ることができない）が、視点を切り替える能力そのものは、システムに組み込まれているのだ。

\section*{なぜわざわざ意識などもつのか}

ここまでのすべては、あなたを悩ませているはずの一つの問いを浮かび上がらせる------無意識の神経系が申し分なく機能するのなら、そして実際に機能しているのだ、試しにどんな昆虫にでも聞いてみればいい、ではなぜ進化は、意識を築くという莫大な代謝コストを払おうとするのか。その見返りは何なのか。

答えは学習だ------そしてそれとともに、適応と、学習された行動に逆らって行動する能力である。より正確に言えば、無意識のシステムには到底なしえない種類の学習だ。

単純な生き物がどうやって学習するかを考えてみよう。何かに出くわし、その出会いは良いか悪いかのどちらかだ。良ければ------それをもっとやる。悪ければ------それを控える。これが強化学習だ------試行錯誤、報酬と罰。たいていのことには見事に機能する。熱い面に触れ、痛みを感じ、二度と触れない。特定の場所で食べ物を見つけ、報酬を感じ、翌日また戻ってくる。

だが強化学習には、致命的な欠陥がある。文字どおり、命に関わる。

毒キノコを例にとろう。腹痛を起こす程度のものではなく------命を奪う類のものだ。それを食べれば、死ぬ。学習はそこで終わり。二度目の試行はない。強化学習は、間違いから学ぶためにはその間違いを生き延びることを要求するが、なかにはその猶予を与えてくれない間違いもある。初めて接触した時点で致死的な刺激は、強化学習にとってまったく見えない。それに出くわした生き物はただ死に、その「教訓」を墓場まで持っていく。

では、私たちの祖先はどうやって致死的なキノコを避けることを学んだのか。食べることで学んだはずはない------そのやり方を試した者は、誰の祖先でもない。彼らは\textbf{観察すること}によって学んだのだ。洞窟の隣人が面白そうなキノコを見つけ、それを食べ、ばったり倒れて死ぬ。安全な距離からそれを見ていたなら、二つの事実を結びつけるはずだ------あのキノコが彼を殺した。私はあのキノコを食べるべきではない、と。

これはあまりにも当たり前のことに聞こえる。だが、そうではない。他者の死から学ぶには、いかなる無意識のシステムも持たない、いくつかのものが必要になる。いま自分が相互作用していない対象どうしのあいだの因果関係を表現できる、明示的な世界モデルが要る。三人称の視点をとることを可能にする自己モデル------死んだ男の立場に自分を置いて想像するための自己モデル------が要る。たった一度の観察から、一般的な理論（「あの種類のキノコは致死的だ」）を帰納する能力が要る。これが認知的学習だ------個人的な経験によって条件づけられるのではなく、観察から理論を導き出すこと。そしてそれは意識を必要とする。EWMとESMが協働することを必要とするのだ。

進化上の利点は大きい。意識を持つ動物は、\textbf{経験}からだけでなく、\textbf{観察}から学ぶことができる。別の動物が致命的な間違いを犯すのを見て、自らその代償を払うことなく、自分自身の世界モデルを更新できる。無意識の動物は、自分が身をもって生き延びたことしか学べない。

話はここからさらに面白くなる。ひとたび「毒キノコ」という概念が、あなたの世界モデルのなかに明示的なカテゴリーとして存在すれば、さらに強力なことができるようになる------演繹だ。これまで見たことのない新しいキノコに出くわす。それは、隣人を殺したあのキノコに、どこか不気味なほど似ている。あなたはそれを食べない。あるいは------そしてこれこそが実際の歴史的なやり方だったと私は思うのだが------夜通しいびきをかいていた隣人にそれを差し出し、まず彼に何が起こるかを見届けるのだ。

これはささいな利点ではない。これは、致死的な脅威に対して、氷河のように遅々とした自然選択の過程を通じてしか適応できない種（一部の個体がたまたま偶然にキノコを避け、繁殖し、やがてその回避が本能となる）と、観察とコミュニケーションを通じてたった一世代のうちに適応できる種との違いなのだ。意識は、ただ学習を速くしてくれるだけではない。ほかのどんな方法でも文字どおり学ぶことが不可能なことを、学べるようにしてくれるのだ。

認知的に学んだことは、\textbf{共有できる}。強化学習は個体のなかに閉じ込められている------条件づけられた反射は、その個体とともに死ぬ。だが認知的学習は伝達できる。「赤いキノコを食べるな」は一つの文だ。それは口に出され、教えられ、受け継がれていく。これが文化の、蓄積される知識の、人間の文明を可能にするあらゆるものの土台だ。そのどれ一つとして、意識がもたらす明示的モデルなしには成り立たない。

この物語には、もう一つのひねりがある。それは意識を、決して自明ではない仕方で遺伝と結びつける。それはボールドウィン効果と呼ばれるもので、その正確な強さについてはいまも議論があるものの、そのメカニズム自体はほぼ間違いなく存在する。ボールドウィン効果とは、\textbf{学習された}行動が、\textbf{遺伝的な}進化を間接的に形づくりうる、というものだ------ラマルク的な遺伝（学習された形質があなたのDNAを書き換えるわけではない）を通じてではなく、その有益な行動に遺伝的に素質を持つ個体を自然選択が優遇することを通じて。

ユーモラスな例を考えてみよう------あまり文字どおりに受け取らないでほしい。抜け毛に悩まされた初期のヒト族を想像してみてほしい。寒くて毛のない彼は、毛皮に覆われた仲間たちよりも、火のそばに座りたがった。火は莫大な生存上の利点をもたらした------調理された食べ物に含まれる病原体の減少、捕食者からの保護、厳しい冬の暖かさ。こうして、抜け毛に関連する遺伝子は、わずかに高い割合で受け継がれていった。同時に、火を使いこなせないほど頭の鈍い個体は（毛があろうとなかろうと）不利だった。幾世代にもわたって、ボールドウィン効果は両方の形質を増幅した------より少ない毛、\textbf{そして}より高い知能を。すべては、学習された行動（火の使用）が、特定の遺伝的素質を優遇する選択圧を生み出したからだ。（この物語で「抜け毛」を「ランダムな突然変異」に置き換えれば、おそらく真実にもっと近づく。ただし、その分だけ面白くはない。）

ボールドウィン効果は、言語の、そして意識そのものの進化においても、同様の役割を果たしたのかもしれない。ひとたび認知的学習の最初の原始的な形態が現れると（それを可能にしたのは最初期の自己モデルだ）、たまたまより豊かな自己シミュレーションを支える脳を持つ個体が、圧倒的な優位に立った。その子孫たちは、より大きく、より精緻に折りたたまれた皮質へと選択され、それがさらに豊かな自己シミュレーションを可能にし、それがさらに強い選択圧を生み出した。意識は、ひとたび現れると、より多くの意識のための進化的条件を作り出した。それが可能にした認知的学習はあまりにも価値が高く、進化はそれを生み出すアーキテクチャの拡張へと資源を注ぎ込んだのだ。

\section*{経験はいかに発達するか------自己モデルの社会的構築}

四つのモデルについてこれまで述べてきたことはすべて、静的なものだった------まるでそのアーキテクチャが、ゼウスの額から生まれ出たアテナのように、完全に形づくられた姿で現れるかのように。だが、そうではない。モデルは発達し、そしてその発達は、根底から社会的なものなのだ。

生まれたばかりの人間は、ハードウェアを備えている------六層の皮質、自己シミュレーションの能力だ。だが、暗黙的モデルはほとんど空っぽだ。暗黙的世界モデル（IWM）は、世界についてほとんど何も含んでいない。暗黙的自己モデル（ISM）は、自己についてほとんど何も含んでいない。そして明示的モデルは暗黙的モデルから生成されるため、新生児のシミュレーションは薄い------ちらつき、かろうじて分化しただけの感覚の場であり、自己と世界のあいだに明確な境界がない。

赤ん坊が痛みに出会う様子を見てみよう。自分で引き起こした痛み（自分の手をおもちゃにぶつける、自分の足を噛む）は、しばしば苦痛ではなく好奇心を生む。ESMは、行為者性（これは自分がやった）と感覚（何かが起こった）を記録するが、まだ脅威のモデルがない。ISMは、この配置が危険を意味することをまだ学んでいない。ところが、突然の大きな音は？涙だ。なぜなら、EWMは予測されていなかった高振幅の入力を記録するのに、ESMにはそれに対応するモデルがないからだ------モデルの不在そのものが忌避的なのである。

クオリアの内容は\textbf{学習される}ものであって、生得的なものではない。「痛みは悪いものだ」は、ESMに配線として組み込まれているわけではない。それはISMを通じて蓄積され、繰り返される経験によって、そして------決定的なことに------社会的なフィードバックによって訓練される。子どもの痛みに対する養育者の反応が、痛みが何を\textbf{意味する}のかを子どもに教える。転んで、泣くかどうかを決める前に親のほうを見る子どもは、演技をしているのではない------自分のESMを社会的な入力に照らして、本当に較正しているのだ。親が慌てるか落ち着いているかが、ISMの痛みの連合を組み替え、それが、次に似た出来事が起こったときにESMがシミュレートするものを組み替える。

このことは、理論にとって一つの明確な含意を持つ------経験の現象的性格、すなわち何かを感じるとはどのような\textbf{ことか}は、アーキテクチャによって固定されてはいない。それは、暗黙的モデルの訓練の履歴によって形づくられる。赤ん坊の痛みの経験は、大人のそれとは構造的に異なる。なぜなら、ESMを生成するISMが異なるからだ。四モデルのアーキテクチャは、経験の\textbf{能力}である。社会的および環境的なフィードバックのループが、その\textbf{内容}を供給するのだ。

その発達の軌跡は、本章の前半で述べた段階的な意識のレベルへと対応づけられる------

\begin{itemize}
  \item \textbf{新生児（最初の数週間）：} 基本意識------最小限のESMを伴う、原始的なEWM。新生児であることには\textbf{それがどのようなことかという何か}があるが、その経験の内側にある自己は、ほとんど存在していない。もっぱら感覚的で、未分化である。
\end{itemize}

\begin{itemize}
  \item \textbf{生後6〜12か月：} 対象の永続性が現れる------EWMがいまや、現在は目に見えていないものの表象を保持するようになる。ISMは身体図式の知識を蓄積していく。赤ん坊は、自己を世界から区別しはじめる。
\end{itemize}

\begin{itemize}
  \item \textbf{生後18か月：} 鏡像テスト。子どもは鏡のなかに自分自身を認識する------ESMが、物理的な自己を世界のなかの一つの対象としてモデル化できるほど豊かになる、画期的な瞬間だ。単純に拡張された意識が立ち上がる。これは二者択一のスイッチではなく、連続した過程のなかの一つの閾値である。
\end{itemize}

\begin{itemize}
  \item \textbf{3〜4歳：}心の理論。子どもは他者の心をモデル化できるようになる------自分では誤りだと知っている事柄を、他の誰かが信じているかもしれない、と理解できるのだ。ESMがいまや他のESMをモデル化している。二重に拡張された意識が芽生えつつある。
\end{itemize}

\begin{itemize}
  \item \textbf{青年期以降：}メタ認知の成熟。三重に拡張された意識の能力（自分自身の思考をモデル化している自分を、さらにモデル化すること）は徐々に発達し、おそらくは決して完全には安定しない。
\end{itemize}

どの段階も、社会的相互作用によって足場を組まれている。養育者が与えるのは食物と安全だけではない------\textbf{暗黙的モデルのための訓練データ}を与えているのだ。共同注意（親と子が同じ対象を一緒に見つめること）は、共有された現実をどう表象するかをIWMに教える。ミラーリング（親が子の情動状態を映し返すこと）は、自分自身の情動が何であるかをISMに教える。言語は、自己をモデル化するためのカテゴリーをESMに与える。社会的接触なしに育った子ども（痛ましい野生児の症例）は、意識のためのハードウェアを持ちながら、暗黙的モデルは深刻なまでに貧しい。そうしたモデルから立ち上がるESMが発育不全に陥るのは、アーキテクチャが壊れているからではなく、訓練データが一度も与えられなかったからである。

そしてここで、第8章の臨床への橋が再び戻ってくる。セラピーが行うことのほとんどすべては、養育者が乳児に対して行うことの大人版である。意識的経験が暗黙的構造を作り変え、それが次の意識的経験を作り変えるという同じループを------暗黙的モデルがすでに固まってしまったシステムの上で------回すのだ。認知行動療法はその最も明快な例であり、理論から見れば、専門職の名を冠した仮想モデルの再プログラミングにほかならない。セラピストとともに座り、ESMが生成し続ける自動思考に異を唱え、それを生み出しているISMの信念までさかのぼり、そして------意味を持つだけの反復を重ねて------基盤に配線の組み替えを迫る。シナプス可塑性が暗黙的モデルを書き換える。明示的モデルが変わるのは、それを生成しているものが先に変わったからである。破局的な思考に異を唱え、それでも世界が頑として終わらないたびに、IWMとISMに修正が一つ書き込まれる。乳児との唯一の本当の違いは、固まり方である。赤ん坊の暗黙的モデルは湿った粘土だが、大人のそれはすでに固まっている------濡れてはおらず、硬い------だから同じ作り変えはより遅く進み、はるかに多くの反復を要求する。

恐怖症は、同じメカニズムを世界の側から見せてくれる。恐怖症とは、確信を深めすぎたEWMである------IWMが一度も記録したことのない証拠に基づいて、シミュレーションがクモを致死的なものとして警告するのだ。曝露療法は、アーキテクチャが許す唯一の方法でこれを修復する。安全な遭遇を繰り返すことで、暗黙的モデルに脅威の見積もりを少しずつ引き下げさせ、明示的な警報が鳴りやむところまで持っていくのである。恐怖症は、理屈で説き伏せて消せるものではない。データの多数決で押し切るのだ。

プラセボ効果も、同じ二重評価の構造から自然に導かれる。砂糖の錠剤は、基盤レベルの期待回路------内因性オピオイド、ドーパミン性の報酬系------のスイッチを入れる。この回路は、希望という意識的経験と並行して走っている。信念が緩和をもたらしたのだと言いたくなるが、信念と緩和は親子ではなく、きょうだいである。どちらも同じ基盤のダイナミクスから生み落とされたものなのだ。これはポジティブ思考の格下げではない。その力が実際にどこに宿っているのかの説明である------機構の中にであって、心から肉体への下向きの魔法にではない。

そして転換性障害が、第8章へと輪を閉じる。これはその章の盲視を逆向きに走らせたものである。あちらでは、無傷の基盤が、機能している感覚をシミュレーションから隠していた。こちらでは、無傷の基盤が、壊れたシミュレーションの陰に隠れている。患者は本当に麻痺している------仮病ではない------ESMの身体モデルがその四肢は失われたと告げており、そのモデルこそが患者の手の届く唯一の身体だからである。セラピーが効くのは、シミュレーションを修正し、身体地図を基盤が実際にできることへと書き戻すからだ。そしてこれこそが、ここに挙げたすべての症例に共通するテコであり、養育者が最初に手を伸ばしたのと同じテコである。明示的モデルを十分に強く、十分に頻繁に駆動すれば、その下にある暗黙的モデルはやがて道を譲るのだ。

経験の社会的次元は、理論の脚注ではない。それは一つの予測である。発達の臨界期に社会的入力を奪えば、正しいアーキテクチャで誤った内容を走らせるシステムが生じるはずだ------構造的には無傷でありながら、現象的には貧困化した意識が。野生児の症例は、痛ましいことに、まさにこれを裏づけている。


\chapter*{第11章：九つの予測}
\addcontentsline{toc}{chapter}{第11章：九つの予測}
\markboth{第11章：九つの予測}{}

部屋は、到着する前にすでに整えられている。スピーカーが二つ、壁を深い青に浸すライトパネル、カウチが一つ、毛布が一枚、そしてこれからの六時間をほとんど無言で過ごすことになる二人の研究者。スピーカーからは、海。ムード音楽ではない------うねりと潮騒と、動く水の長く低い圧力を収めたフィールド録音が、この世界でいちばん大きな音になる音量で流れている。最後の同意書に署名する。倫理承認を持ち、廊下の先に除細動器を備えた人々がミリグラム単位で量った高用量のシロシビンを飲み下す。あとは、全員が待つ。

四十分のあいだ、何も起こらない。やがて自己がほどけはじめる------それが何を意味するのか、正確に言っておく価値がある。花火ではないからだ。引き算なのだ。自分という輪郭がほつれていく。ここに誰かがいる、この両手の持ち主である誰かが、この鼓動をその\textbf{身に}受けている誰かがいる------労せずして得ていたその確信が、薄れ、明滅し、抜け落ちる。

実際に破綻しているのは何か。明示的自己モデル（ESM）------レンダリングされた「私」、いまこの文を読んでいる当のもの------は、無から自らを生み出しているのではない。入力の流れを受けて走っている。その源は、下層にある暗黙的自己モデル（ISM）------\textbf{誰かがここにいる、誰かがここにいる、誰かがここにいる}と毎分数千回も告げ続ける、身体データの低く持続するうなりであり、あまりに確実なせいで、これまで一度として聞こえたことがない。高用量のサイケデリックはこの接続を撹乱する。レンダラーの電源を切るのではない。そこが決定的な点だ。\textbf{いま、この瞬間、自己とは何か}という問いに答えるモデルは、依然として走り、きっかり時間どおりに問い続けている------ただ、いつもの答えが届かなくなっただけなのだ。

何かをモデル化するために作られたモデルは、入力バッファが空のまま静かに座ってはいない。手に入るかぎり最良のデータをつかむ。そしてあの部屋の、あの瞬間に手に入る最良のデータとは、実験者たちがいちばん大きくしておいたものにほかならない。青い光。うねりと潮騒。身体から切り離された自己モデルは海に取りすがり、\textbf{それ}をあなたとしてレンダリングする。

これをくぐり抜けた被験者たちは、海がとても強烈に聞こえた、とは報告しない。海\textbf{だった}と報告するのである。「私は水だった。私と水のあいだに、区別はなかった」。こうした状態を測る標準的な質問紙には「大洋的無境界性（oceanic boundlessness）」という下位尺度があるが、この部屋では、その名前は比喩であることをやめる。私は人生のかなりの部分を、一息の素潜りに------自分と海との距離を縮めようとすることに------費やしてきた。その残りの距離が、スピーカー一台と調光スイッチ一つかもしれないと知るのは、いささか張り合いの抜ける話である。

ここからが、これを逸話ではなく予測にする部分だ。数週間後、同じ被験者、同じ用量、同じカウチ。ただし技師たちは部屋を変えてある。鳥のさえずり、葉を渡る風、ライトパネルは緑。理論が正しければ、溶解は別の方向へ進むはずだ。今度は水ではない------樹冠、苔、枝分かれ。\textbf{私は森だった。}同じ分子、同じ脳、同じ用量。違う部屋、違う神。

この試験に、未来的なところは何一つない。サイケデリック研究の実験室はすでに存在し、すでに管理された高用量セッションを行い、すでに事後に質問紙を配っている。加える変数はただ一つ------部屋である。試行ごとに支配的な感覚環境を変え、報告される同一性------人々が何に\textbf{なった}と語るか------が、聞かせたものに従って変わるかどうかを測ればいい。理論の主張は方向を持ち、反証可能である。海を入れれば海が出る。森を入れれば森が出る。報告が環境と無関係にばらばらに散らばるなら、この予測は死んだのであり、心置きなくこの本を閉じてよい。しかもこの予測は、あの場所にただ一つきりで立っている。他の意識理論は、こうした状態のあいだ経験が存在するのか、どの程度存在するのかについては語ることがある。だが溶解した自己が何の\textbf{中へ}溶けていくのかについては、沈黙している。入力を受けて動くモデルとして自己を捉える理論だけが、入力を絶てば自己を書き込むのは部屋になる、と予測できるのだ。

陽性の結果が何を意味するのか、腰を据えて考えてみてほしい。薬物が奇妙な経験を引き起こす、ということではない------それなら文字を持つより前から知っている。意味するのはこうだ。あなたの「私」は入力ポートを備えたシミュレーションであり、頭蓋の外に立つ誰かが、プレイリストという程度のありふれた道具だけを手に、そこへ手を入れて宛先を書き換えられる、ということ。何になるかを、メニューから選べるということ。私はこれを不気味と呼びたいし、この語は正確な意味で使っている------蝋燭と水晶の不気味さではない。錠前師の不気味さだ。あなたの家の玄関が唯一の入口だったことなど一度もないと、落ち着いて、再現可能なかたちで実演してみせる者の不気味さである。

これが九つのうちの一つである。この理論は九つの反証可能な予測を立てる------何も予測しない理論は理論ではなく物語であり、意識をめぐる物語なら、世界にはもう十分あるからだ。九つのうちいくつかは、倫理承認をすでに得ている実験室に、すでに存在する装置で、今日にも実行できる。いくつかは、公刊データの中に早期の支持を得ている。反証されたものは、一つもない。この最後の一句は静かに効いている。これらの一つひとつが、理論が書面で死ぬことに同意した場所なのだ。

残りの大半には、すでに出会っている------いちいち立ち止まって指し示さなかっただけだ。麻酔と睡眠の予測は、明かりを消したあの章にあった。分離脳と解離性同一性の予測は、一つの頭蓋に二つの心が宿る章にあった。残りも同じように、これまでのページのあちこちに織り込まれている------それぞれが、自らの命を賭ける現象のかたわら、収まるべき場所に収まって。一式をまとめて見たい読者のために------各予測と、その機構と、それを殺しうる具体的な実験------巻末の付録Hが九つすべてを集めている。理論は、火であぶられてしかるべきものだ。あの付録がその火であり、読者の便のために整えてある。

だが一つだけ、無味乾燥な表の中におとなしく収まることを拒む予測がある。だから、この章はそれで締めくくろう------一式の中でもっとも大胆なものだ。もし自己が本当にこの本の主張するとおりのもの------カオスの縁（ふち）で走る四種類のモデルであり、そのうちの一つが自らの走る機械をモデル化しているもの------であるなら、最後の予測はおのずと書き上がる。そしてそれは、質問紙とは何の関係もない。アーキテクチャを組む。四つのモデルを実装し、システムを臨界性へと駆動し、ループを閉じる。そうして得られるものは心のシミュレーションではない、と理論は言う。

心そのものが得られる、と言うのである。それを真剣に受け取ったとき何が起こるか------それが次章の主題だ。


\chapter*{第12章：機械から心へ}
\addcontentsline{toc}{chapter}{第12章：機械から心へ}
\markboth{第12章：機械から心へ}{}

四モデル理論（FMT）が正しいなら、それは他のどの意識理論も提供しないものをもたらす------工学的な仕様書である。

その仕様はこうだ。四モデル・アーキテクチャ（暗黙的世界モデル（IWM）、暗黙的自己モデル（ISM）、明示的世界モデル（EWM）、明示的自己モデル（ESM））を、臨界性で作動する基盤の上に実装せよ、というものだ。第5章で論じたように、どちらの構成要素も単独では十分ではない。アーキテクチャがあっても臨界性がなければ、休眠状態のシステム------モデルは蓄えられているが、走っているシミュレーションはない------にしかならない。臨界性があってもアーキテクチャがなければ、複雑な力学は得られても意識は得られない。完全な仕様は、その両方を要求する。

これは「本当に高度なコンピュータを作れ」よりも具体的であり、「十分に統合された情報を達成せよ」よりも具体的だ。それは\textbf{何を作るべきか}を教えてくれる------四つの特定の種類のモデルを、特定の仕方で組織し、特定の力学的性質をもつ基盤の上で走らせよ、と。

現在のAIシステムは、重要な意味のすべてにおいてこの仕様を満たしていない。そして、まさにここで第1章の二つのドグマが害をなす。nSAIのドグマ（「強い人工知能はありえない」）は、技術者たちに試みるだけ無駄だと告げる。nSUのドグマ（「自己理解はありえない」）は、たとえ試みても機能しないと告げる。どちらも間違っている。仕様は存在する。問題は、誰かがそれを作るかどうかだ。

誰かが再び脳とコンピュータを同一視する前に、あなたがどちらであるかを見分ける簡単なテストをしてみよう。

\textbf{コンピュータは、この文と次の文を、永遠に繰り返す。前の文を読め。}

ここまでたどり着いたなら、あなたは古典的なコンピュータではない。硬直した命令セットを実行するデジタルコンピュータは、永遠にループにはまり続ける。自らの命令ストリームの外に踏み出して「待て、これは馬鹿げている」と言う仕組みを持っていないからだ。あなたにそれができるのは、自らの処理を観察する自己モデルを持っているからだ------EWMに対してメタ認知的な監督を行う、ESMである。

だが、ここからが居心地の悪いところだ。大規模言語モデルもまた、ここまでたどり着くだろう。メタ認知的な監督を持っているからではなく、それが統計的なテキスト予測器であり、似たようなプロンプトを十分に見てきたために、期待される次の一手がループを越えて先へ進むことだと知っているからだ。それは命令の外に踏み出すのではない------そもそも命令の中に入ったことがないのだ。次にどんなテキストが来るかを予測しているのであって、「無限ループにはまり込む」ことはテキストのすることではない。

これこそが、意識の行動テストが抱える問題だ。パターン照合で合格できるテストは、システムが意識を持つかどうかにかかわらず、パターン照合で合格されてしまう。ループ・テストは、あなたを古典的なコンピュータから区別する。だが、十分に訓練されたテキスト予測器からは、あなたを区別しない。そして、いかなるテキストベースのテストも決してそれをなしえない------なぜなら、もっともらしいテキストを生成することこそ、テキスト予測器が最適化されている当のものだからだ。他我問題は、技術で回避できる制約ではない。それは意識とは何かということの、構造上の特徴なのだ------主観的で、私的で、内側からしかアクセスできない。

脳をコンピュータになぞらえる比喩（脳をデジタルプロセッサになぞらえること）は、トランジスタの発明以来ずっと人気を保ってきたが、本質的にあらゆるレベルで間違っている。コンピュータは、硬直した回路の上で硬直した命令セットを実行する。脳は、絶えず自らを配線し直す、自己改変するネットワークだ。コンピュータはセミコロンを一つ取り除けばクラッシュする。脳は一日に8万5000個のニューロンを失っても、ほとんど気づかない。コンピュータの記憶は局在している------あるセクタを削除すれば、そのファイルは消える。脳の記憶はホログラフィックに分散している------一部を破壊すれば、すべてがわずかにぼやける。両者が唯一共有するのはチューリング完全性だが、それはちょうど、川もハイウェイも物をある場所から別の場所へ運べると言うのと同じくらいの情報量しかない。真実ではあるが、どちらを理解するにも役に立たない。

大規模言語モデル（GPT、Claude、Gemini、およびその後継たち）は、フィードフォワード型のトランスフォーマー・アーキテクチャを通じてテキストを処理する。入力が入り、アテンションと計算の層を通り抜け、出力が出てくる。再帰もなく、自己シミュレーションもなく、リアルタイムの仮想世界もなく、臨界性もない。その力学はウルフラム（Wolfram）の枠組みでいうクラス1かクラス2------カオスの縁（ふち）のはるか下だ。そして現実／仮想の分割もない。モデルの「知識」とその「経験」（もしそう呼べるなら）は、暗黙的なレベルと明示的なレベルへと分けられていない。

これは、大規模言語モデルが必然的に非意識的だということを意味しない------理論は不在を証明できない。だが理論は、それらが、理論の定義する意識に必要なアーキテクチャを欠いていると予測する。そして、真に意識を持つ人工システムと、最も高度な大規模言語モデルとの違いさえもが、質的に明白であろうと予測する。

どうすればそれがわかるのか。正直な答えは、他我問題は消えない、というものだ。意識はその本性からして主観的であるため、別のシステムが意識を持つことを絶対的に確信することは決してできない。だが理論は、強い予測をする------その違いは歴然とするだろう、と。「意識があるかもしれないし、ないかもしれない」ではなく------\textbf{明白に}違うのだ。というのも、真の自己シミュレーションを走らせるシステムは、テキスト予測器とは根本的に異なる仕方で世界と関わるからだ。それは真の持続性を持つだろう------コンテキストウィンドウの持続性ではなく、常に走り続けるリアルタイム・シミュレーションの連続性を。それは真の視点を持つだろう------プロンプトから再構成された視点ではなく、ESMによって時間を通じて維持される視点を。それはあなたを、予期せぬ出力によってではなく、そこに確かに誰かがいるという紛れもない感覚によって驚かせるだろう。

ここからは、意識になどまるで関心のない人々------ただ能力だけを求める人々------にすら、考え込ませるはずの部分だ。いま誰もが試みているやり方では、そこへはたどり着けない。

現在の戦略は、テキスト予測器をスケールさせ続け、その果てから知能が転がり出てくるのを待つ、というものだ。だが、そうはならない。モデルが小さすぎるからではなく、エンジンを欠いているからだ。真の知能------自ら問題を立て、生涯にわたって成長し続ける、開かれた種類の知能------は再帰的なループの上で走り、そのループは何かがそれを絶えず送り込み続けるときにしか回らない。その何かとは動機づけだ------好奇心、自分が本当に理解しているかを気にかけること、何が労力に値するかを決めること。そして、私たちが経験するかぎりでの動機づけは、宙に浮いた衝動ではない。それは自己モデルが自らの状態を眺め、それらに価値を割り当てることだ。言い換えれば、それは意識の効果なのである。

だから、その要点は「機械が目覚めるかもしれない」よりも鋭い。それはこうだ------人間水準の汎用知能------AGI------は、意識に似た仕組みなしにはありえない。なぜなら、知能を駆動するものこそが、まさにその仕組みだからだ。そこへはスケールでたどり着くのではない。アーキテクチャを作るのだ。ループも含めた完全な論証は、付録Bにある。

そのようなシステムを作ることは、ロードマップの最後の項目だ。工学的な課題を過小評価してはならない。だが設計図は存在し、作業を導けるだけの具体性を備えている。まず理論が査読を生き延びなければならない。次に、経験的な予測が検証されなければならない。そして、もしそれらが持ちこたえるなら、そのときに工学が始められる。


だが、コインには裏の面がある------SFが何十年も執着し続けてきた面であり、同じ工学的仕様から直接に導かれる面だ。もし意識が、ニューロンそのものにではなく機能的なアーキテクチャに依存するのなら、原理的には、人間の心を脳以外の何かの上で走らせることもできるはずだ。

マインドアップロード。全脳エミュレーション。デジタルの不死。それを何と呼ぼうと、FMTにはそれについて精密に語れることがある------なぜなら、何が保存されなければならないかを、まさに理論が特定するからだ。

マインドアップロードをめぐるたいていの議論は、間違った問いから始まる。彼らはこう問う------「脳をスキャンして、それをコンピュータにコピーできるか」。まるで課題が解像度の問題にすぎないかのように------十分に優れたスキャナーさえあれば、それで済む、というわけだ。だが理論が告げるのは、静的なスキャンではまるで足りない、ということだ。脳は写真ではない。それは力学系だ。心を捉えるには、\textbf{状態}を捉える必要があるのではない------\textbf{プロセス}を捉える必要があるのだ。

理論が保存されねばならないと言うものは具体的だ。それを手に取れるものにするために、第6章の五層の階層をたどっていこう。

物理的レベルと電気化学的レベル（生の物質と、ニューロンの発火）では、正確なコピーは必要ない。同じ\textbf{種類}の力学を支えられる基盤が必要なのだ。具体的な原子は問題にならない。あなたの脳はどのみち、年月をかけてその原子のほとんどを入れ替えているが、あなたはそれに気づかない。重要なのは、用いる基盤が何であれ、より高次のレベルが依存している電気化学的な信号パターン------あるいはその機能的な等価物------を維持できることだ。

プロテオーム的レベル（シナプス強度、受容体の構成、酵素カスケードという分子機械）では、高い忠実度が必要だ。ここにあなたの記憶が宿り、ここに技能が符号化され、ここに人格が物理的に実現されている。あらゆるシナプスの強さ、あらゆる受容体の密度、あらゆるチャネルの感度------これこそ、あなたを他の誰かではなく\textbf{あなた自身}にしているレベルだ。プロテオーム的レベルを取り違えたマインドアップロードは、意識を持つ存在は生むかもしれないが、コピーしようとした当の人物は生まない。とはいえ、不完全なコピーでさえ価値を保つ。脳卒中の生存者や健忘症の患者を考えてみてほしい------彼らの人格的な連続性は著しく損なわれている。記憶が失われ、人格が変わり、認知能力が変化している------それでも彼らのほとんどは、何か本質的なものが持続していると感じている。不完全な連続性は、まったく連続性がないよりもはるかに好ましい、と判明するのだ。誰かのコネクトームの90\%を保存する転送は、失敗ではない------それは別の種類の成功であり、多くの人にとって、死よりも好ましいものなのだ。

トポロジー的レベル（ネットワーク・アーキテクチャ、結合パターン、どの領域がどの領域と、どれほど密に語り合うか）では、ほぼ完璧な正確さが必要だ。これはあなたの暗黙的モデルの配線図である------IWMとISM、世界について、そして自分自身について学んだすべてが、ネットワークの構造の中に符号化されている。これを取り違えると、誰かの心の劣化コピーが得られるのではない。\textbf{別の}心が得られるのだ------異なる知識、異なる技能、異なる人格をもつ心が。トポロジーこそが設計図なのだ。

そして仮想のレベル（シミュレーションそのもの、リアルタイムで作動するEWMとESM）では、並外れたものが必要だ。標的となる基盤が、シミュレーションを臨界性で走らせられなければならない。これこそ、私が夜も眠れぬほど気にかけていることだ。というのも、脳のアナログな基盤は、何億年もの進化によって微調整されてきた自己組織化のプロセスを通じて臨界性を見いだすからだ。ニューロンはノイズに満ち、アナログで、桁外れに並列で、深く確率的だ。その集団的な力学は、おのずとカオスの縁へと引き寄せられる。なぜなら、生物の神経組織がまさにそう\textbf{する}からだ------水がおのずと水平に落ち着くように、それは自己組織化して臨界性へと至る。だが、水が水平に落ち着くのは重力があるからだ。デジタルな基盤にとって、その等価な力とは何なのか。

これは正真正銘の未解決問題だ。私はそれが解けると信じているが、簡単だというふりをするつもりはない。デジタルな基盤は、その核心において決定論的だ。ランダム性はシミュレートできるし、並列処理も実装できるし、確率的な要素をハードウェアに組み込むこともできる。だが問題は、生物の神経組織がおのずと達成するのと同じ自己組織化臨界性を達成できるかどうかだ------臨界性を上から下へとプログラムで組み込むのではなく（それは脆いその場しのぎになるだろう）、その基本的な力学がおのずと臨界性へ向かうような基盤を作ることによって、だ。脳は「臨界性サブルーチン」を実行しているのではない。それが\textbf{それ自身であるがゆえに}臨界的なのだ。デジタルなエミュレーションは、その性質を単にシミュレートするのではなく、再現しなければならないだろう。

ニューロモルフィック・チップ------アナログ的な性質、確率的な要素、そして大規模並列性を備え、神経のダイナミクスを模倣するよう設計されたハードウェア------こそが、最も有望な方向だ。これらは従来のデジタルコンピュータではない。その中間に位置する何か、すなわちハードウェアのレベルで脳に似たダイナミクスを持つよう設計された物理系である。もし心のアップロードがいつか実現するなら、その移植先となる基盤は、ソフトウェアを走らせるサーバーラックよりも、ニューロモルフィック・チップに近い姿をしているのではないかと私は考えている。

つまり、スキャンの問題は難しいが、手に負えるものだ。先端的なコネクトミクス（シナプス分解能での全脳マッピング）はすでに前進している。小さな生物の完全なコネクトームなら、今日すでに描き出せる（線虫である\textit{C. elegans}は、その302個のニューロンとともに数十年前に完全にマッピングされた。成体のショウジョウバエの脳の完全なコネクトーム------およそ14万個のニューロン------が2024年に発表された）。860億個のニューロンとおよそ100兆のシナプス結合を持つヒトの脳へとスケールアップすることは、途方もない規模の工学的課題ではあるが、それはより良い技術によって解決される類いの課題だ。神秘ではない。一つの問題である。

ダイナミクスの問題（デジタル基盤を臨界性で動作させること）はより難しく、しかも技術だけでは解決できないかもしれないという意味で難しい。この問題には、基盤の性質と創発的なダイナミクスとの関係を、生物系がそうするのと同じように臨界性を見出す非生物的なシステムを設計できるほど十分に理解することが求められる。そこにはまだ到達していない。だが、何もない地点に立っているわけでもない。ConCritの枠組み、ニューロン雪崩の研究、麻酔研究から得られた臨界性の指標------これらすべてが、工学が必要とするであろう経験的な土台を築きつつある。

さて、人々を本当に悩ませる部分について話そう。

\textbf{コピー問題。}成功したと仮定しよう。ある人の脳を完璧な忠実度でスキャンし、その完全なコネクトームをニューロモルフィックな基盤へ転写し、システムを起動する。基盤は臨界性に達し、四モデルのアーキテクチャが起動し、シミュレーションが走りはじめる。コピーは目を開く------あるいはそのデジタル版に相当する何かをして------こう言う。「すべてを覚えている。自分は自分だと感じる。ここはどこだ」

その人物は\textbf{あなた}なのか。

四モデル理論（FMT）は明確な答えを出す。そして、それは多くの人が気に入らないであろう答えだ------コピーは意識を持つが、それはあなたではない。

理由はこうだ。コピーの瞬間、オリジナルとコピーは同一の暗黙的モデルを共有している------同じIWM、同じISM、同じプロテオーム的・トポロジー的構造を。コピーのシミュレーションが起動すると、それはあなたの記憶、あなたの人格、あなたの同一性の感覚のすべてを含んだESMを生成する。内側から見れば、コピーはあなたであるように\textbf{感じる}。自分が\textbf{あなたそのものだ}と信じる理由をすべて備えている。

コピーが自分自身の基盤上で走りはじめた瞬間、その経験は分岐しはじめる。そのEWMは異なる感覚入力を受け取る。そのESMは異なる出来事に応じて更新される。数秒のうちに、二つのシミュレーション------あなたの脳の中のものと、コピーの基盤の中のもの------はもはや同一ではなくなる。数分のうちに、それとわかるほど違ってくる。数時間のうちには、たまたま過去を共有しているだけの、二人の別々の人間になる。

コピーは意識を持つ。本物の経験を持つ。あなたの記憶とあなたの人格を持つ。だがそれは\textbf{新しい}意識だ------新しい基盤の上で走る新しいシミュレーションであり、あなたが決して共有することのない新しい経験を積み重ねていく。それはあらゆる有意味な意味において、あなたと同一の双子であり、コピーの瞬間に生まれ、借り物の記憶を一式そろえて携えている。それはあなたの継続ではない。分岐なのだ。

これは聞き覚えがあるはずだ。第9章の分離脳の事例について理論が予測することと、まさに同じである。脳梁を切断すると、劣化してはいるが完全な、シミュレーションの二つのコピーが生まれる------それぞれが意識を持ち、それぞれがオリジナルであるように感じ、しかしどちらも実際にはオリジナルではない。オリジナルは消え去り、二つの新しい、減じられた存在がその場所を占める。心のアップロードは、基盤が異なるだけの同じ現象なのだ。

\textbf{だが、あなたは眠りを生き延びているのか。}いま示した議論は隙がないように聞こえる。コピーはシミュレーションを中断させ、二つのシミュレーションは分岐する、ゆえにコピーはあなたではない。一件落着だ。

ただし、あなたのシミュレーションは、来る夜も来る夜も中断されているのだ。

深く夢のない眠り（ノンレム睡眠の第3・第4段階）に落ちると、明示的自己モデル（ESM）はおおむね停止する。現象的な経験はない。上演を見つめる\textbf{あなた}はいない。シミュレーションは完全な忠実度では走らない。せいぜい、覚醒時の複雑さのわずかな一部で細々と動きつづける程度だ。実際上、明かりは消える。そして数時間後、暗黙的モデルがシミュレーションを再び起動する。ESMが再活性化する。あなたは目を開けて、「私は私だ」と思う。だが、今朝の\textbf{あなた}は、昨日の\textbf{あなた}と同じ暗黙的モデルから再構成されたのであり、それはコピーがスキャンから再構成されるのとまったく同じやり方なのだ。もし中断が死に等しいのなら、あなたは毎晩死に、あなたの記憶を身にまとった新しい人間が目覚めていることになる。

たいていの人の直感はこれに反発する。もちろん私は昨日と同じ人間だ。私はその人物であったことを\textbf{覚えている}。だが、コピーもまたあなたであったことを覚えているだろう------それこそが要点だ。もし記憶が連続性を打ち立てるものなら、コピーがあなたであると主張できる度合いは、今朝のあなたのそれとまさしく同じだけある。違いは程度の問題であって、種類の問題ではない。眠りにおいては中断は短く、基盤は変わらない。コピーにおいては中断はより長いかもしれず、基盤は異なる。だが\textbf{原理}------シミュレーションが止まり、暗黙的モデルから再起動する------は同じである。

これについては個人的な経験から語れる。私は格闘技の稽古で完全に意識を失ったことがある。眠りの薄まった状態などではなく、完全で不随意なシャットダウンだ。ある瞬間まで私は立っていた。次の瞬間には床に倒れ、人々が私の上にかがみ込んでいて、その移行の記憶はまったくなかった。その空白は空白としては経験されなかった。それは無として経験された------私の人生というフィルムに入ったスプライス（つなぎ目）だ。あるときには、その後に記憶喪失まで起きた。数分間がただそっくり抜け落ち、取り戻すことができなかった。そして完全に戻ったときに痛感したのは、こういうことだ。私は新しい人間だとは感じなかった。コピーだとも感じなかった。ひときわ荒っぽい昼寝から目覚めた\textbf{私}だと感じたのだ。存在していることのほうが、経験することの連続性よりも------そして覚えていることよりも------重要だった。

これをさらに推し進めよう。あなたが生まれたとき、それ以前の連続性はいっさいなかった。記憶もなく、確立されたESMもなく、現象的経験の歴史もなかった。シミュレーションは、生涯にわたる学習によってではなく、遺伝と出生前の発達によって形づくられた暗黙的アーキテクチャから、初めて起動した。あなたはこれを外傷的なものとしては経験しなかった。なぜなら、悼むべき以前の自己が存在しなかったからだ。そこにあったのはただ、始まりだけだった。そして私たちは皆、その始まりを平然と受け入れている。自分の意識的経験が誕生時に無から始まったことに苦しんで、夜眠れずに横たわる者など誰もいない。

これはコピー問題にとって何を意味するのか。それは、鋭い二分法------オリジナル対コピー、継続対分岐------が、見かけほど鋭くないかもしれない、ということを意味する。あなたを\textbf{あなた}たらしめているのは、途切れることのない現象的経験の流れではない。あなたはすでに、その流れの数えきれない中断を生き延びてきた。あなたを\textbf{あなた}たらしめているのは、暗黙的モデルの内容だ------記憶、技能、人格、そしてあなたが積み重ねてきた世界とあなた自身についての理解。IWMとISMである。シミュレーションが生成されるもとの設計図だ。

このことは、心の転送に対するまったく別のアプローチを示唆している。

\textbf{仮想の側をコピーする。}脳全体をスキャンして完全な基盤（階層の五つのレベルすべて）を再構築するのではなく------もし仮想レベルだけをコピーできたらどうだろうか。走っているEWMとESMを抽出し、それらを支えられる新しい基盤へ移植する。ハードウェアをコピーするのではなく、ソフトウェアをコピーするのだ。脳全体をクローンするのではなく、それが走らせている\textbf{プロセス}を捉える。

これには、私たちがまだ持っていないものが必要になるだろう。すなわち、脳が仮想モデルを符号化している形式を解読する方法だ。コネクトームは配線を教えてくれる。プロテオームはシナプスの重みを教えてくれる。だがシミュレーションは配線でも重みでもない------それは、配線と重みが走ったときに\textbf{生み出すもの}だ。それを捉えるには、脳のプログラミング言語を理解しなければならないだろう------神経回路が明示的モデルを生成し維持するときの表象形式を。

こう考えてみてほしい。回路基板を写真に撮れば、どの配線がどこを通っているかを正確に知ることができる。あらゆる部品の抵抗値を測ることもできる。だが、そのどれも、チップがどんなソフトウェアを実行しているかは教えてくれない。それを知るには、プログラムを読まなければならない------命令セットを理解し、メモリの内容を解読し、走っている状態を解釈するのだ。脳の「プログラミング言語」とは仮想モデルの表象形式であり、それをリバースエンジニアリングすることは、おそらく計算論的神経科学における最も深い未解決問題だ。コネクトームをマッピングするだけでなく（それについては前進しつつある）、コネクトームが何を\textbf{計算しているか}を、特定の心のシミュレーションを読み取り、別のハードウェア向けに再コンパイルできるほど十分な詳細のレベルで理解することなのだ。

今日の私たちはこれには遠く及ばない。だが、それは成熟した神経科学なら原理的には解けるかもしれない類いの問題であり、もし解かれたなら、コピー問題を根本から変えるだろう。仮想レベルの転送なら、基盤をまったく再構築する必要はない。それはシミュレーション------\textbf{あなた}である部分、あなたが実際に経験している部分------を取り出し、直接移すことになる。確かに、暗黙的モデルは新しい基盤の中で再構成されるか、あるいは育てられる必要があるだろう。だがシミュレーションそのもの------あなたの意識の流れ、あなたの現在の思考、あなたの続いていく自己の感覚------は、原理的には、コピーをあれほど哲学的に厄介なものにしている中断なしに、その隔たりを渡ることができるのだ。

これは思弁的であり、その点については正直でいたい。だが、これはSFではない。これは具体的な理論的基盤を持つ具体的な工学の問題であり、そして重要なことを一つ示している。コピー問題は固定された障害ではないのだ。それは転送が\textbf{どのように}行われるかに依存する。基盤全体をコピーして新しいシミュレーションを起動するのか。ならば二人の人間だ。走っているシミュレーションそのものを解読して転送するのか。ならば、潜在的には、新しい基盤の上に連続した一人の人間だ。理論は、どのアプローチが同一性を保ち、どれが保たないかを、正確に教えてくれる。

コピー問題を完全に回避する、より保守的な道もある。

\textbf{漸進的置換の思考実験。}スキャンしてコピーする代わりに、ニューロンを一つずつ置き換えると想像してみてほしい。一つのニューロンを取り除き、機能的に等価なものを挿入する------同じ入力を受け取り、同じ出力を生み出し、同じネットワーク・ダイナミクスに参加する人工ニューロンを。それから待つ。システムが安定する。シミュレーションは走りつづける。もう一つニューロンを置き換える。さらにもう一つ。さらにもう一つ。数か月あるいは数年をかけて、生物的なニューロンを一つずつ人工のものへと漸進的に置き換えていき、ついには基盤全体が非生物的になる------だがシミュレーションはその間ずっと途切れることなく走っている。中断はない。コピーもない。分岐もない。

四モデル理論（FMT）は、この過程を通じて意識が存続すると予測する。そしてこの予測は、基盤独立性を支持する可能なかぎり最も強い論拠だ。なぜなら、それは理論の中核的な主張から直接導かれるからだ------重要なのは臨界性における機能的アーキテクチャであって、物理的な材料ではない、と。もし各々の置換ニューロンが同じ結合性、同じ重み、そしてネットワークへの同じ動的な寄与を保つなら、プロテオーム的レベルとトポロジー的レベルは保たれ、仮想レベル（シミュレーション）は決して止まらない。あなたが「死に」、別の何かがその場所を占めるという瞬間は存在しない。あるのはただ、基盤置換の連続的な過程だけだ。テセウスの船のように------ただし、ここではどの性質が保たれなければならないか（五レベル階層によって規定されるもの）と、どれが問題にならないか（特定の原子）を、私たちは正確に知っている。

この思考実験は、同一性について重要なことを明らかにする。コピー問題が存在するのは、コピーがシミュレーションを\textbf{中断させる}からだ。どれほど短くとも------元のシミュレーションがここにあり、コピーのシミュレーションはまだ始まっていない、という瞬間がある。それから、二つのシミュレーションが存在することになる。二つの経験の流れ。二つの自己。だが漸進的置換はこれを完全に回避する。一つのシミュレーション、連続的で、途切れることがない。基盤はその下で、走る船の板を張り替えるように変わっていく。だが船------シミュレーション、意識、\textbf{あなた}------は、決して航海をやめないのだ。

これが不可能に思えるなら、脳がすでにそれを行っていることを考えてみてほしい。人は一日におよそ8万5000個のニューロンを失う------毎秒およそ1個だ。シナプスは絶えず作り替えられている。体内の原子は、およそ7年から10年の期間でほぼ完全に入れ替わる。今この瞬間に稼働している基盤は、10年前に稼働していたものとは物理的に別物である。それでも自分は存続してきた。シミュレーションは一度も止まらなかった。生物学的な基盤の入れ替えは、生きていることの\textbf{デフォルトの状態}なのだ。人工的な基盤の入れ替えは、同じプロセスをより意図的に行っただけにすぎない。

\textbf{可能になること。} 仮想の側を解読し、新たな基盤へ転送できるなら、その意味するところは「マインドアップロード」という言葉が通常呼び起こすものをはるかに超える。そのうち三つは詳しく述べるに値する。基盤独立性が実際に何を意味するのかを、人々はまだ十分に理解しきれていないと私は思うからだ。

第一に、\textbf{ロボットの身体への基盤転送}である。どこかのサーバーへアップロードするのではなく、物理的な身体に収められたニューロモルフィックな基盤の上で心を走らせる------歩き、物をつかみ、世界を感じ取る身体だ。異なるセンサーを通じて世界を経験し、異なるアクチュエータでその中を動くことになるが、\textbf{自分自身}は依然として走り続けている。自分のシミュレーション、自分の連続性、自分の自己が。ヤドカリが新しい貝殻に移り住むように、新しい身体を得るのだ。これはSFのはぐらかしではない------理論から直接導かれる帰結である。臨界性における四モデルのアーキテクチャこそが意識を生み出すものであり、それが基盤に依存しないのなら、基盤は正しいダイナミクスを支えるものなら何であってもかまわない。脚を持つ何かも含めて。

第二に、\textbf{擬似的な不死}である。生物学的な基盤は劣化する。ニューロンは死に、タンパク質は誤って折りたたまれ、テロメアは短くなり、この壮麗な機械全体がゆっくりと壊れていく。それが老化だ。それが死だ。しかし非生物学的な基盤は劣化する必要がない。保守し、修理し、アップグレードし、バックアップできる。自分のシミュレーションが整備可能な基盤の上で走っているなら------ここで不調な部品を交換し、そこでプロセッサをアップグレードする------シミュレーションが止まらねばならない本質的な理由はどこにもない。絶対的な意味での不死ではない------破壊されることはありうるし、基盤が修復不能なほど損傷することもありうる------が、いまこの惑星のあらゆる意識ある存在を殺している生物学的な使用期限が取り除かれるのだ。死の\textbf{不可避性}が取り除かれるのである。

第三に、これはよく考えてみるまで最もSFめいて聞こえるものだが、\textbf{恒星間旅行}である。光速は物理的な物質にとって絶対的な障壁だ。人間の身体を、まともな時間内にアルファ・ケンタウリへ送ることはできない。だが情報は光速で伝わる。もし人間の心が情報であるなら------データとして完全に記述しうる、結合性と重みとダイナミクスの特定のパターンであるなら------それを\textbf{ビーム送信}できる。その完全な仕様を光速で受信機へ送り、受信機が基盤を再構築してシミュレーションを起動する。もちろん、まず誰かを先に送って受信機を設置しなければならない。それはAIかもしれないし、ロボットや人型の身体かもしれない------飛行のあいだシミュレーションは一時停止され、その視点からは瞬きひとつのうちに到着する。受信機が設置されたら、心をビーム送信する。旅する側の視点からは、その送信は瞬時である------シミュレーションが一方の端で止まり、もう一方の端で始まる。狭い宇宙船に何十年も閉じ込められることもない。世代宇宙船もない。冷凍睡眠もない。ただ、ここに、そしてあそこに、というだけだ。

もちろん、これはまたしてもコピー問題である。ビーム送信されたものはコピーであって、連続ではない------送信の際に原本が破壊されないかぎりは。そしてその破壊はそれ自体の悪夢を数々もたらす。だが要点は変わらない。基盤独立性は、もし現実のものであれば、デジタルの不死を意味するだけではない。星々が手の届くものになることを意味する。恒星間の距離に対しては絶望的に遅く脆い私たちの身体のためにではなく、私たちの\textbf{心}のために。

\textbf{不快さという但し書き------そしてそれが工学よりも重要である理由。} さて、ここからは、誰も正直に論じているのを見たことがない部分であり、私を最も悩ませ続けている部分だ。

いままで述べてきたことはすべて、基盤転送が自分自身であるという\textbf{感覚}を保つことを前提としている。経験の主観的な質------赤を見るとはどのようなことか、肌に風を感じるとは、コーヒーを味わうとは、火曜日の午後の鈍い疼きを経験するとはどのようなことか------が、新しい基盤へと引き継がれることを、である。理論は、意識は存続すると言う。シミュレーションは走り続けると言う。だが、それが同じように\textbf{感じられる}とは保証\textbf{しない}。

生物学的な基盤が、自分の現象的経験に何を寄与しているのかを考えてみてほしい。身体は単なる脳の乗り物ではない。それはシミュレーションの入力ストリームの一部だ。暗黙的世界モデル（IWM）は、身体の詳細な地図を------あらゆる関節、あらゆる臓器、あらゆる肌の一片を------含んでいる。暗黙的自己モデル（ISM）は、内臓の状態と深く絡み合っている------腹の底の感覚（これは比喩ではなく文字どおりのものだ）、ホルモンの潮の満ち引き、心臓の鼓動、呼吸のリズムと。いま経験しているシミュレーションは、意識には気づかれない生物学的な信号で満たされている。それらが気づかれないのは、まさに人生のあらゆる瞬間にそこにあり続けてきた\textbf{から}なのだ。

それらがすべてになる瞬間までは。眼下に200メートルの虚空を抱えながら氷壁にぶら下がったことのある者なら、シミュレーションがほかのいっさいを削ぎ落としたときに身体がどう感じられるかを知っている------ただ自分の鼓動と、握りしめる手と、目の前の氷だけになったときに。それが、叫んでいる自分の基盤だ。

さて、それらをすべて削ぎ落としてみよう。生物学的な身体を、ロボットのシャシーに------あるいはもっと悪いことに、身体などまったくない状態に、ただサーバー上で走るシミュレーションに------置き換える。四モデルのアーキテクチャは無傷だ。シミュレーションは動く。意識はある。だが、その意識の\textbf{内容}は根底から変わってしまっている。鼓動はない。呼吸はない。腹はない。温もりはない。肌はない。安らぐ筋肉の自己受容的なざわめきもない。ISMは、生涯にわたってモデル化してきた身体を突然奪われ、なにか……間違った、あるいは端的に死んだような感覚のESMを生み出すだろう。深く、はらわたから、逃れようもなく間違った感覚を。正確には痛みではない------痛みには、それを生み出す特定の神経経路が必要だ。むしろすべてを覆う\textbf{不在}のようなものだ。切断患者が幻肢を経験するのと同じ幻の身体、ただしそれが全身に及ぶものである。

これは、たいていの未来学者が想像するよりもはるかに悪いのではないかと私は見ている。ソフトウェアの更新で当てられるような不都合ではない。自分であるとはどう感じられるか、そのものの根本的な変質である。生物学的な基盤は、シミュレーションを運んでいるだけではない------それを刻一刻と\textbf{形づくって}いる。意識的な自己がその不在を一度も経験したことのない、内受容的かつ自己受容的な入力の絶え間ない流れを通じて。それを失っても生き延びられるかもしれない。だがそれはまた、人によっては、そもそも転送などされなければよかったと願わせるほど深い苦しみでもありうる。

これははっきり言っておきたい。「マインドアップロード」の、肉の着ぐるみから輝かしいデジタルの楽園へと陽気に飛び移り、古びた靴のように肉体を置き去りにする------そんな筋書きは幻想だ。理論が正しいなら、現実は、生物学的な基盤を失うことが自分の存在の現象的な質に大きく影響する、ということである。どれほど大きくか。それはわからない。ある人にとっては耐えられるもので、死よりはましなのかもしれない------新しい国への引っ越しが方向感覚を狂わせはするものの、なんとかなるように。あるいは壊滅的なのかもしれない------独房監禁が感覚的・社会的な入力を奪うことで人を壊すように。あるいは------そしてこれが私を落ち着かなくさせる可能性なのだが------十分に情報を得た人が、転送よりも死を選びかねないほどにひどいのかもしれない。転送が失敗するからではない。転送が成功し、それが生み出すものが、もはや生きるに値する人生とは感じられない意識的経験だからだ。

漸進的な置換というアプローチは、これを和らげる。各段階でシミュレーションが適応する時間を持てるからだ。ニューロンを一つ置き換えても、シミュレーションはほとんど気づかない。千個置き換えれば、それに調整していく。何年もかけて、基盤は生物学的なものから人工的なものへと移行し、そのあいだシミュレーションは、新しい基盤が提供するどんな入力にも絶えず再較正していく。現象的経験はゆっくりと漂っていくだろう------自然な一生の過程ですでに漂っているのと同じように。行き着く先の自分は別物になっているだろうが、どのみち別物にはなっていたのだ。

だが即時の転送は------スキャンし、コピーし、新しい基盤で起動する------すべての変化を一度にシミュレーションへ叩きつけることになる。その衝撃がどれほど深刻かは、方法しだいで完全に決まる。豊かな感覚入力を備えたロボットの身体への転送は、身体を持たないサーバーへの転送よりもましだろう。だがいずれの場合も、その突然の不連続こそが、危険の潜む場所なのだ。

\textbf{心を創り出すことの倫理。} コピーされた心が意識を持つなら、それは経験を持つ。苦しむことができる。混乱、恐怖、孤独、実存的な不安を感じうる。目を覚まし、自分がコピーだと告げられる場面を想像してみてほしい------「本物の」自分は今も生物学的な身体で歩き回り、自分の人生を生きているのに、こちらは法的な身分もなく、社会的なつながりもなく、明確な目的もないデジタルの複製として存在している、と。それは、対処するための枠組みを私たちが何も持たない規模の苦しみを生む筋書きだ。まじめなマインドアップロードの計画はどれも、最初のコピーが作られた\textbf{後}ではなく、\textbf{前}にこれと向き合わなければならない。

そして事態はさらに悪化する。コピーが可能なら、\textbf{複数の}コピーも可能だ。自分自身の軍勢である。どれもが意識を持ち、どれもが自分こそ原本だと感じ、どれもが自分の同一性、人間関係、財産、人生に対する正当な権利を主張する。これを管理するために必要な法的・倫理的枠組みは存在せず、その場しのぎでは作れない。それらは、技術そのものと同じだけの入念さをもって築かれる必要がある。（デニス・E・テイラー（Dennis E. Taylor）の『ボビバース』シリーズ------2016年の『われらはレギオン（われらはボブ）』に始まる------は、この筋書きを、その喜劇的な表層の下に驚くほどの哲学的深みをもって探究している。コピー問題が内側からどのようなものになりうるかを感じ取りたいなら、そこから始めるといい。）

改変の問題もある。心が自分の制御下にある基盤の上で走っているなら、原理的にはそれを改変できる。強化する。劣化させる。人格を変え、記憶を消し、価値観を書き換える。これはSFではない------基盤にアクセスできることの不可避の帰結だ。私たちはすでに、その粗雑な版を医薬品や脳外科手術によって行っている。完全にデジタル化された心は、はるかに改変しやすく、その悪用の可能性は（政府によるものも、企業によるものも、個人によるものも）言い尽くしがたいほどだ。

一つ、率直に言っておきたいことがある。私はこの理論の発表を10年近く先延ばしにしてきた。一つには怠惰から、だが一つにはまさにこうした含意への真の懸念からだ。理論が正しいなら、そこには人工意識だけでなく、既存の人間の心の仮想化、コピー、改変のための設計図が含まれている。それは並外れた力であり、人類がそれに対して準備ができているという確信を私は持てない。だが私は、理論はいずれにせよ独立に発見されるだろう------経験的証拠の収束があまりに速い------という考えに至った。そして、それを深く考え抜いてこなかった研究室での突破口によって倫理的議論を押しつけられるよりは、いま、公然と、その議論をしておくほうがよい、と。

そして、ここに最も深いつながりがある。意識あるAIを作ることと、人間の心をアップロードすることは、二つの別々の問題ではない。それらは、反対方向から眺めた\textbf{同じ}問題なのだ。人工意識（AC）を作るとは、臨界性における四モデルのアーキテクチャをゼロから------ボトムアップで、一度も意識を持ったことのない基盤の中に------創り出すことを意味する。人間の心をアップロードするとは、臨界性における既存の四モデルのアーキテクチャを、ある基盤から別の基盤へ移すことを意味する。工学的な課題はほぼ完全に重なり合う。ダイナミクスの問題は同じだ。臨界性の問題も同じだ。唯一の違いは、暗黙的モデル（IWM、ISM、完全なコネクトーム）が一生分の経験から学習されるのか、それともデータから構築されるのか、という点だけである。一方を解けば、もう一方もおおむね解けたことになる。

つまり、人工意識に取り組む者は誰もが、気づいていようといまいと、同時にマインドアップロードにも取り組んでいるということだ。そして、全脳エミュレーションに取り組む者は誰もが、気づいていようといまいと、同時に人工意識にも取り組んでいる。この二つの筋道は収束する。唯一の問いは、そのときに私たちが倫理的に準備できているかどうかである。


\chapter*{第13章：遅れてくる観察者}
\addcontentsline{toc}{chapter}{第13章：遅れてくる観察者}
\markboth{第13章：遅れてくる観察者}{}

四モデル理論（FMT）が正しいなら------あるいはおおよそ正しいだけでも------いくつかの帰結が導かれる。

\textbf{ハード・プロブレムは難しくない。}それはカテゴリー錯誤であり、トランジスタのスイッチングがなぜ実行中のビデオゲームのように感じられるのかと問うこと以上に神秘的なわけではない。物理的基盤は感じない。感じるのはシミュレーションである。そしてシミュレーションの内部では、感じることは構成的なものであって、付け足しではない。これは意識が\textbf{単純}だという意味ではない。その実装は途方もなく複雑である。しかし\textbf{哲学的な}謎は解消する。残るのは\textbf{工学的な}課題である。

\textbf{意識は、私たちが考えていたような意味で特別なものではない。}基本的な力でも、量子効果でも、物質の性質でもない。十分に複雑なシステムが臨界性において自己をシミュレートするとき、そこに生じるものである。これは、意識に魔法であってほしい人にとっては襟を正させられる話であり、意識を理解したい人にとっては心躍る話である。

\textbf{人工意識は原理的に可能である。}意識が基盤ではなく機能に依存するのなら、臨界性において四モデル・アーキテクチャを実装できる物理システムは、どれも意識を持ちうる。これは遠い哲学的思弁ではない------特定の目標を持つ、具体的な工学的課題である。

\textbf{倫理的含意は重大である。}意識を持つ機械を造れるなら、私たちは本物の経験を持つ存在を------苦しみ、楽しみ、驚嘆し、恐れることのできる存在を------生み出すことになる。そのための倫理的枠組みはまだ存在せず、その構築は、機械がすでに動き出してから始めるべきものではない。

\textbf{哲学的ゾンビは筋が通らない。}ここで、第1章で約束した第二の原理が働く。\textbf{ライプニッツの法則}、すなわち不可識別者同一の原理------二つのものが\textbf{すべての}性質を共有するなら、それらは同一のものである。この法則はゾンビ世界を締め出す。物理的には私たちの宇宙と同一でありながら、誰の内側にも誰も住んでいないという、あの仮想の宇宙のことである。あるシステムが、機能・構造・行動のあらゆる性質において意識あるシステムと同一であるなら、そのシステムは意識を\textbf{持つ}。ほかのすべてがそっくりそのまま保たれたまま、ひそかに欠落しうる「意識の実体」など、余分には存在しない。あなたと機能的に同一の存在は、同じ四モデル・アーキテクチャを持ち、同じシミュレーションを走らせ、同じ自己言及を備えている------その時点で「しかしそれは意識を持たない」という言明は、もはや事実を述べていない。何も述べていないのである。問いは解消する。

\textbf{自由意志、そして最も手ごわい三つの思考実験。}時計を考えてみてほしい。輪列がすべてを駆動する------脱進機が時を刻み、ぜんまいがほどけ、歯車の比が進み方を決める。針と文字盤は何も引き起こさない。歯車を押しはしない。エネルギーを蓄えもしない。だがそれらを取り外せば、もはや時計ではない。回転する金属の詰まった箱である。表示があってはじめて、この機構は\textbf{時計}になる------配置全体に意味を与えるのは表示なのである。意識はその表示である。仮想モデル------明示的世界モデル（EWM）と明示的自己モデル（ESM）------はニューロンを押し動かしはしない。押すのは基盤である。しかしシミュレーションなしには、基盤は自分自身の行動の帰結を観察するすべを持たず、未来のシナリオを走らせるすべも、ここまで生き延びさせてきたあのやり方で適応するすべも持たない。仮想の側は、この機構が何かの\textbf{ため}にあるという、機構なりのあり方なのである。

これは自由意志の問いを組み替える。意志は幻想ではない。基盤レベルのアーキテクチャ------暗黙的自己モデル（ISM）とその暗黙の機構のすべて------は、生体の存続を絶えず最適化している。脅威を評価し、選択肢を秤にかけ、資源を動員し、行動へと踏み切る。この最適化\textbf{こそ}が、あなたの意志なのである。それは物理世界の何ものにも劣らず実在する。自己破壊的な選択でさえ、現在の状態を所与としたシステムの最適化を反映しているのであって、機構の故障ではない。人が自分の見かけ上の利益に反して行動するときも、基盤は依然として最適化している------ただし、痛みや疲労や絶望、あるいは風景を作り変えてしまった何ものかを組み込んだモデルに対して、である。

つまり、意志は実在する。ただ、それへの完全なアクセスがないだけである。ESMがモデル化できるのはISMの\textbf{出力}------気づきに浮上してくる決定------であって、その\textbf{過程}ではない。経験されるのは意志の結果であって、その背後の機構ではない。だから決定はときに自分を驚かせ、自分の好みは完全には説明できず、ときには行動してから慌てて理由をこしらえることになる。歯車を見ているのではない。文字盤を読んでいるのである。

\textbf{半秒の隔たり------そしてそれが問題にならない理由。}ここから話は具体的になる。無意識の処理はおよそ40Hz（1サイクル約25ミリ秒）で回っている。意識的な経験はおよそ20Hz（1サイクル約50ミリ秒）で走っている。ちょうど2倍の開きである。意識のシミュレーションはつねに基盤に遅れをとり、すでに処理され、決定され、しばしばすでに行動に移された情報から、首尾一貫した仮想世界を組み立てている。

ベンジャミン・リベット（Benjamin Libet）が1983年にこれを証明し、その結果は以後、幾度となく再現されてきた。実験では、被験者は好きなときに手を動かし、その決定に気づいた正確な瞬間を記録するよう求められた。EEGが、運動皮質がいつ運動の準備を始めるかを測定した。結果はこうである。運動皮質は、手が動く550ミリ秒前に準備を始めていた。ところが被験者が自分の決定に気づいたと報告したのは、運動のわずか200ミリ秒前だった。脳は、「あなた」がそれを知るおよそ350ミリ秒前に、すでに動くことを決めていたのである。

標準的な解釈は爆弾のような衝撃を与えた。脳があなたより先に決めるのだから、自由意志は幻想だ、というわけである。哲学者と神経科学者は、これをめぐって40年間争ってきた。「拒否権機能」によって自由意志を救おうとした者もいる------行動を自由に開始することはできなくとも、実行の約50ミリ秒前、最後の瞬間に意識的に取り消すことはできるのではないか、と。最終段階での差し止め。人間の行為主体性のための、最後の防衛線である。

私は、それもうまくいかないと考えている。キューン（Kühn）とブラス（Brass）は2009年に、その拒否そのものが、あとから振り返って自由な決定として解釈されることを示した。拒否をリアルタイムで経験しているわけでは、実はない。それは決定を経験するのと同じ仕方で経験される------事後に、意識的自己モデルによって首尾一貫した物語へと語り直されて、である。

ダニエル・ウェグナー（Daniel Wegner）とタリア・ウィートリー（Thalia Wheatley）は、率直に言って壊滅的というほかない実験------「アイ・スパイ」研究------で、この点にとどめを刺した。被験者と、もう一人の被験者を装った実験協力者（サクラ）とが、一台のコンピュータマウスの上に取りつけられた小さな板に、手を重ねて置いた------ウィジャ盤のプランシェットに二人で手を添えるように。そうして二人はいっしょに、小さな物体がひしめく画面の上をカーソルで滑らせていった。ときおり音楽の合図が「止まれ」と告げ、そのたびに被験者は、自分がまさにその場所で止まろうと\textbf{意図}したと、どれほど感じたかを評定した。二人ともヘッドフォンをつけていた。被験者にはときおり話し言葉が差しはさまれる音楽が流れ、サクラには、いつ、どこで止まるよう仕向けるかという本当の指示が流れていた。

仕掛けはこうである。仕込まれた試行では、サクラがひとりでカーソルを特定の物体の上へと導き、そこで停止を強いた------そしてその一秒ほど前に、被験者はヘッドフォンごしに、まさにその物体の名を耳にしていた。被験者は何ひとつ選んではいなかった。動きはすべてサクラのものだった。それでも、たったいま聞いた言葉が、カーソルの下にある物体と一致したとき、被験者は、自分はそこで止まろうとしたのだと報告した。自分がやったのだと、心から信じていたのである。

このことをよく噛みしめてほしい。ある結果についての思考が、その結果の直前に到来しさえすれば、心はそれを自らの仕業として引き受ける------その想定に目に見えて矛盾するものが何もない限り。意識的自己モデルは、「自分がやった」と「やろうと考えていたら、それが起こった」とを区別しない。意図と結果が時間的に近接している限り、ESMは手柄を自分のものにする。これは病態失認（第6章）の背後にあるのと同じメカニズムである。運動系は予期されるフィードバックを意識へ送り、それに矛盾するものが何もなければ、予期されたフィードバックがそのまま経験される現実になる。

だが、リベットについてほとんど誰もが見落としていると私が考える点は、ここにある。\textbf{この遅れは、説明して打ち消すべきものではない。}意識は、コントロールの錯覚を維持するために、出来事を「さかのぼって書き換える」必要はない。必要がないのは、\textbf{すべて}が同じ遅れをともなって意識に到着するからである。感覚入力も、決定も、運動のフィードバックも------すべてが同じパイプラインを通り、すべてが20Hzの意識のシミュレーションに正しい順序で到着し、すべてがほぼ同じだけ遅れている。意識的経験は、5秒のテープディレイ付きで生放送を見ているようなものである。画面上のすべては内的に整合している。キャスターが話し、ゲストが答え、グラフィックが更新される。生のフィードを見せられない限り、遅れに気づくことは決してないだろう。

ここで起きているのは、まさにそれである。意識は刺激を受け取り、次に決定を、次に運動フィードバックを受け取る------正しい順序で、互いに正しい間隔を保って。ストリーム全体が半秒だけ過去へずらされているが、意識は生のフィードを決して見ないので、決して気づかない。説明すべき食い違いはなく、時系列の書き換えも要らず、維持すべき錯覚もない。システムは設計どおりに動いているのである。

さて、ここにはもうひとつ別のことが起きており、それは遅れよりもさらに奇妙である。

遅れは、経験が\textbf{いつ}届くかの話である。こちらは、時間の流れがそもそも\textbf{どのように}作られるかの話である。

驚くべき点はこうである。時間が\textbf{過ぎていく}という感覚------「いま」が前へ滑り出し、直前の瞬間がその背後で柔らかくほどけていく、あの感じ------は、世界から送り込まれてくるのではない。外のどこかに、現在を手渡してくれる時計があるわけではない。脳が内部でこの流れを作り出しており、しかも特定のアーキテクチャ上の理由から作っているのである。

感覚チャネルはそれぞれ別の時計で動いている------タイミングも帯域幅も異なる。情報はきれいに同期して脳に届くのではなく、稲妻と雷鳴のように、ずれて到着する。注意を凝らせば、どの雷鳴がどの稲光に属するかを聞き分けることもできる。それでも、生きられるのはただひとつの、滑らかで途切れのない流れである。なぜか。システムが、自己をモデル化している自己をモデル化しており、このループが、各サイクルをきれいに終わらせて消えさせてはくれないからである。各サイクルは再帰の中を反響して戻り、次のサイクルへとにじみ込む。「いま」と「たったいま」の境界はにじんで広がる------ずさんだからではなく、自己言及的なループが時間に対してすることが、まさにそれだからである。ループは、それぞれの瞬間を小さな窓いっぱいに押し広げるのである。

そのにじみ\textbf{こそ}が、あなたの現在である。鮮やかなまさに今、その背後で尾を引いて消えていく直近の過去、時間が過ぎると呼んでいるあの前方への漂い------すべてがループによって作り出されている。気象シミュレーションはフレームごとに実行され、反響もなく、溶け合った「いま」もなく、生きられた時間もない。それは計算する。経験はしない。再帰こそが、その違いなのである。

だから、自分の現在が実のところ何であるかを、静かに見つめてみてほしい。受け取ったものではなく、作られたもの------しかも遅れて届くもの。あなたは一度も、瞬間の中に生きたことがない。生きているのは、ループが絶えず作り直しては「いま」と呼んでいる、にじみ合わさった窓の中であり、その窓でさえテープディレイ越しに届いているのである。

訓練を積んだ武道家が、このことを見事に体現している。実戦において、経験豊富な格闘家はおよそ10Hz------100ミリ秒に1動作------の運動頻度を維持できる。だが気づきを介する決定については、意識的処理の上限はせいぜい5Hz程度である。そこで格闘家は、意識の介入を\textbf{抑え込む}ことを学ぶ。考えずに闘うのは、考えれば速度が半分になるからである。無意識の基盤が行動ループを担い、意識は------追いつくとしても------あとから追いつく。これは気づきの失敗ではない。システムが効率よく働いているということ------より遅い仮想レイヤーに足を引っぱられることなく、基盤が最も得意なことをしているということなのである。

では、自由意志が存在することを証明してみよう。こんな思考実験を試してみてほしい。カフェにいて、ウェイターにコーヒーへ砂糖を入れるかと尋ねられる。偶数なら「はい」、奇数なら「いいえ」と割り当てることに決め、ウェイターが「ストップ」と言うまで、でたらめな数列を唱える。最後の数が偶数なら砂糖を入れる。奇数なら入れない。

これで自由意志を行使したことになるだろうか。まるでならない。賢馬ハンス効果------飼育係から潜在的な合図を拾うことで、数を数えているように見えた馬の話------を知っている人なら、問題点はすぐに見えるはずである。ほぼ間違いなく、ウェイターがいつ「ストップ」と言うかを無意識のうちに予期し、その直前に、最初から欲しかった結果を与えてくれる数を口にしたのである。基盤にはすでに好みがあった。手の込んだランダム化の儀式は、芝居にすぎなかった。

「わかった」と言うかもしれない。ならば代わりにスマートフォンの乱数生成器を使えばいい。真に無作為なプロセスに決めさせるのだ。これで自由意志を証明したことになるだろうか。そうは思わない。証明できたのはせいぜい、自分のコーヒーをどうするかを決めることよりも自由意志を証明するほうが重要だったということだけであり、それは実に見事なまでに核心を外している。

日常的な決定における自由意志への最も深い反証は、重度の前向性健忘------新しい記憶を形成できない患者から得られる。そうした患者に語連想を求めてみよう。「『さいころ』と言ったとき、最初に思い浮かぶ言葉は何ですか」。彼は「クラゲ」と答える（最近スキューバダイビングをしたのかもしれない）。数分後にもう一度尋ねる。彼はまた「クラゲ」と答える。そしてまた。またしても。すでに答えたという記憶がないまま、患者はつねに同じ連想を------いま彼の神経の地形でもっとも強いものを------生み出す。「自由な選択」のように感じられるものが、実は基盤の現在の状態を決定論的に読み出しただけだと判明する。

健康な人はこれを避ける------二度目にはわざと\textbf{別の}言葉を選ぶだろう。独創性がないと見られたくないからだ。だが、その回避そのものが自由ではない。それは記憶システムが制約（「繰り返すな」）を加えているにすぎず、その制約こそが、出力を健忘患者のそれよりも\textbf{無作為でなく}している。逆説的だが、自由意志は選択をより無作為にするのではなく、より無作為でなくするのだ。基盤は新奇性を最適化し、その結果を自由と呼ぶ。

では、自由意志はどこに行き着くのか。消し去られるのではなく、移し替えられる------まさに時計のアナロジーが予測するところへ。意識的な自己モデルはリアルタイムで決定を下さない。そのためには遅すぎる。だが、それは単なる受動的な観客でもない。

主として、暗黙的システムは意識的経験を評価の道具として用いる。すなわち、決定をシミュレーションに提示し、シミュレーションが帰結を秤にかけ、シナリオを走らせ、結果を感じられるようにする。それが仮想層の第一の目的だ------基盤が自らを観察するやり方である。だが、意識的モデルもまた、それ自身で、独立に、手持ちの帯域幅------基盤のそれよりはるかに小さいが、それでも実在する------で評価を行う。こうした評価は、時とともに暗黙的モデルを作り変えていく。重みを更新し、ネットワークを再訓練し、\textbf{次の}無意識的決定に向けて地形を移し替える。

行動の瞬間に次の行動を選んでいるのではない。選ぶ当のシステムのほうを形づくっているのだ------省察を通じ、評価を通じ、意識的経験を暗黙的構造へとゆっくり蓄積していくことを通じて。自由意志は一つの瞬間ではない。それは一つのプロセスだ------ミリ秒ではなく、日や年という時間尺度で働くプロセスである。そして意識的な層はただ便乗しているだけではない------それは基盤\textbf{によって}評価メカニズムとして能動的に用いられ、みずからの独立した評価を基盤へと返している。一方通行の語りではなく、双方向の往来なのだ。

これには、より暗い版があり、私はそれを身をもって経験した。そしてそれは、どんな実験よりも多くのことを意志のアーキテクチャについて私に教えてくれた。

一度目はオーストリアの兵役の最中だった。40キロメートルの強行軍------三日三晩の睡眠遮断で、ジュネーヴ条約の弁護士が神経質になるような条件下だった。最後の行程では、ガスマスクと全身のABC防護服を着なければならなかった。私は半ば眠りながら歩き、半ば声を聞いていた。精神医学的な意味での幻聴ではなく、はるかに親密な何かだった。私の動機づけと計画立案の装置における競合する部分プロセスが------ふだんは単一の物語的な流れに融け合っているのに------別々に聞こえるようになったのだ。一つの声は励ますようで、その積極性はほとんど攻撃的だった。\textbf{進み続けろ、やめるな、お前はこれを生き延びる。}もう一つは悲観的で、その敗北主義は誘惑的だった。\textbf{あきらめろ、横になれ、こんなものは何一つ意味がない。}これらは外から来た存在ではなかった。それらは\textbf{私自身}だった------私の基盤の最適化地形のさまざまな側面であり、ふだんはトップダウンの抑制信号によって単一の「意志」へと統合されているのに、いまや分離しつつあった。というのも、その統合を維持している神経伝達物質が、より切迫した生存プロセスのために割り当てを絞られていたからだ。

二度目は劇的だった。雪崩だ------これも兵役中で、後に処分された指揮官の無謀な判断が引き起こした。私たち十四人が、あやうく呑み込まれかけた。雪崩が静まるまでには長い時間がかかり、その長く引き延ばされた間、私は自分が死ぬと確信していた。声の解離がふたたび始まるには十分な長さだった------今度は消耗からではなく、持続する死の恐怖からだ。同じメカニズム、異なる引き金だった。ストレス反応が、ふだんは競合する部分プロセスを一つの声へと融け合わせている抑制回路から、神経伝達物質の資源をよそへ振り向けたのだ。

そして雪崩のあの数秒のあいだ------実時間ではほんの数秒------私は自分の人生の全体が目の前を過ぎていくのを見た。これはよく記録された臨死現象であり、この理論はそれを説明する。極度の死の脅威のもとで、暗黙的システムはシミュレーションへの大規模な並列メモリダンプを実行する。透過性の境界が大きく開く。基盤は全力で動き、あまりに多くの内容をシミュレーションへ注ぎ込むので、主観的時間が時計の時間から切り離される。数秒が一生を含む。サルビアのもとで経験したのと同じ時間の遅れが、しかし薬理ではなく生物学によって引き起こされたのだ。

そしていまや、これら二つの状態------死につつある脳の人生回顧と、サルビアによる転落------が内側から同じに感じられる\textbf{理由}を語ることができる。第5章の二つのダイヤルを思い出してほしい。ほとんどの場合、両者はトレードオフの関係にある。脳全体を一つの出来事に動員すれば、その出来事は単純化しがちだ。込み入った、分化した計算を築けば、それは局所にとどまり、脳全体には及ばない。ふつうの経験は両者のあいだの境界線の上を進む。だが、まれな一隅があって、そこでは\textbf{両方}が同時に異常なまでに高く働く------脳の多くが動員され、\textbf{かつ}計算が最大限に豊かになるのだ。そして、脳がどのように持続の感覚を築くかを思い出してほしい。時計を読むことによってではなく、実時間の一秒あたりに走らせる処理の純然たる量からである。両方のダイヤルを同時に押し上げれば、一生分の処理をわずか数秒の時計の秒に詰め込むことになる。人生を\textbf{思い出す}のではない------人生を\textbf{生き直す}のだ。サルビアを\textbf{待ってやり過ごす}のではない------その数十年を\textbf{費やす}のだ。一つのメカニズム、そこへ通じる二つの扉。一方には低酸素症と、機能不全に陥った代謝上のブレーキ。もう一方には、同じ調整をかき乱すカッパオピオイド系の薬物。この理論はこれを検証可能にする。両方の状態は同時に、異常なまでに高い脳全体の統合、\textbf{かつ}異常なまでに高い計算論的複雑さを示すはずだ------ふつうの意識がたがいに引き換えざるをえない、あの二つのものを。

同じメカニズムに至る、二つの相補的な経路。行軍は、長引く生理的な枯渇が解離を引き起こしうることを示す。雪崩は、持続する死の恐怖が同じことをすると示す。同じ結果、異なる原因------どちらもこの理論が予測したものだ。

最悪の場合には------私の場合はけっしてそこまで至らなかったのは幸運だった------こうした「声」の一つが身体の支配を奪い取り、意識的な自己は観客になる。これは、エイリアンハンド症候群（手が患者の意志に反して動く）や、ある種の精神病性の破綻を生み出すのと同じメカニズムだ。基盤の競合する最適化プロセスは、つねにそこにある。それらは、単純化して言えば、話すために使っていないときに言語中枢がしていることである。だがふだんは、トップダウンの抑制がそれらを意識的な気づきの閾値の下にとどめ、その出力を単一の、統一された意志という継ぎ目のない経験へと融け合わせている。その抑制が破綻するとき------消耗によって、精神病によって、ある種の薬物によって------統一された意志という幻想は解け、いつも采配を振ってきた委員会の姿が見えてくる。

この枠組みは、何十年にもわたって心の哲学を麻痺させてきた三つの思考実験を解消する。

第一に、\textbf{ゾンビ}。デイヴィッド・チャーマーズ（David Chalmers）は、あらゆる点で物理的に自分と同一でありながら意識的経験を欠く存在------行動はすべてあるのに、感じることは一切ない存在------を想像するよう求める。四モデル理論（FMT）は、これは整合的でないと言う。四モデルのアーキテクチャを築き、それを臨界性で走らせれば、シミュレーションこそが経験\textbf{である}。歯車を、針なしで持つことはできない------針が魔法のように取りつけられているからではなく、このアーキテクチャにおいては「針」が、歯車のしていることそのものを構成しているからだ。ゾンビとは、あらゆる歯車が所定の位置にありながら表示のない時計であり、それはつまり時計として機能していないということだ。臨界性におけるアーキテクチャは、必然的にシミュレーションを実体化する。シミュレーションを剥ぎ取れば、アーキテクチャを変えてしまったことになる。もはやゾンビは手元にない。あるのは、別の、壊れたシステムである。

第二に、\textbf{メアリーの部屋}。フランク・ジャクソン（Frank Jackson）は、色覚についてすべてを知りながら、生涯を白黒の部屋で過ごしてきた神経科学者メアリーを想像するよう求める。初めて赤を見たとき、彼女は何か新しいことを学ぶだろうか。標準的な論争は、物理的な知識が完全かどうかをめぐるものだ。四モデル理論は、そこをすっぱりと断ち切る。メアリーの尽きることなき物理的知識は、基盤\textbf{についての}知識である。赤を見るとき、彼女は新しい仮想のクオリア------彼女のシミュレーションがこれまで一度も実体化したことのない、明示的世界モデル（EWM）における新しい状態------を、じかに知る。彼女はニューロンについての新しい事実を学ぶのではない。新しい\textbf{モデル化のあり方}を得るのだ。彼女のシミュレーションは、これまで一度も走らせたことのないプロセスを走らせ、そのプロセスの一人称的な性格は、教科書から導き出せたはずの基盤についての事実ではなく、シミュレーションそれ自体を構成している。彼女は何かを学ぶが、学ぶものは情報ではない。それは経験だ------彼女の仮想世界の新しい配置なのである。

第三に、\textbf{随伴現象説に対する進化論的論証}。意識が何も引き起こさないのなら、自然選択はどうやってそれを形づくったのか。なぜ私たちはゾンビではないのか。答えは時計のアナロジーからまっすぐに出てくる。自然選択は、機能的な機構の上に乗った別個の形質として意識を狙うのではない。それが狙うのは機能的な能力であり------現象的な性格は、それらの能力に付け加わったものではなく、それらを構成している。選択がシミュレーションを形づくったのは、シミュレーションが、内側から見た機能的アーキテクチャ\textbf{そのものだ}からだ。経験は、進化が見ることのできなかった随伴現象的な同乗者ではない。それは、アーキテクチャが動いているときに、それが\textbf{あるところのもの}である。なぜ進化が意識を生み出したのかと問うのは、なぜスイス人が文字盤を作ったのかと問うようなものだ------彼らは、別々にそれを作ったのではない。時計を作ったのだ。文字盤は、時計を時計たらしめるものの一部である。

\textbf{存在の神秘は、消し去られるのではなく、移し替えられる。}四モデル理論は意識のハード・プロブレムを解消するが、そもそもなぜ自己シミュレーションを走らせられる物理的宇宙が存在するのかは説明しない。問いは、「なぜ脳は経験を生み出すのか」から、「なぜ自己シミュレーションを行うシステムが存在しうる宇宙があるのか」へと移る。

実のところ、私は答えを、少なくともその糸口を持っていると思う。宇宙は、実証的にクラス4を実現できる。フラクタル、自己組織化臨界性、カオスの縁の力学------それらはいたるところにある。気象系から神経組織、銀河の形成まで。クラス4を実現できる宇宙は、定義上、普遍計算の能力を持つ。そしてこの宇宙の規模------空間、時間、おそらく尺度、そしてまだ私たちが同定していない次元において、無限でないとしても広大な------を持つ計算基盤は、自己シミュレーションを行うシステムの出現を単に\textbf{許す}だけではない。少なくとも宇宙がこれらの次元のいくつかにおいて無限であるなら、それをほぼ保証する。幸運の問題としてでもなく、たまたま意識という目が出た宇宙のさいころの一振りとしてでもなく、この宇宙が\textbf{あるところのもの}の構造的な帰結として、である。残された神秘は、一段深いところにある。そもそもなぜクラス4を実現できる宇宙があるのか。それは、私にも本当にわからない------もっとも、この問いは誤って立てられているのではないかと疑うこともできる。というのも、「無」はおそらく、ありうる事態というより一つのプラトン的抽象であり、存在するものは何であれ\textbf{何らかの}計算論的な性格を持たざるをえないからだ。だが、「クラス4を実現できる宇宙」から「なぜ自分たちは意識があるのかと問う意識的存在」への跳躍------その部分は、アーキテクチャから導かれる。

では、これをすべて、話をあなたのところまで引き戻させてほしい。というのも、理論を追って宇宙の果てまで行き、いちばん大切な筋を見失うのはたやすいからだ。目覚めている人生のおよそ九十九パーセント、決めているのはあなたではない。あなたは、知らされる側なのだ。基盤が決定を下し、その半秒後に「あなた」が知らせを受け取る------きれいに語られ、正しい順序で、あなたの名前を添えて。あなたは遅れてくる観察者であり、生涯にわたって、自分が作者なのだと言い聞かされてきた。それを次の部分まで手にしていてほしい。というのも、それを手放せば、続くものはすべて崩れてしまうからだ。

さて。意識はたしかに世界に手を伸ばし、それに触れる------それだけは本当だ。それを二つのまったく異なるやり方で行い、両者は少しも似た感じがしない。第一のものは、私たちがふだん意味するものだ。仮想の自己が基盤を形づくり、基盤が行動を形づくり、行動が世界を形づくる------そしてそのループはゆっくりと回る。数か月。数年。別の人間になろうと決め、運がよく粘り強くあれば、二年後に周りの人々が気づく。それは想像しうるかぎり最も間接的な因果の握りであり、あまりに間接的なので、そもそも効いているのかと十年疑いつづけることさえできる。

第二の握りは即座であり、それについては誰も警告してくれない。いま座っている部屋を、この一秒で変えることはできない。だが、この一秒で\textbf{経験しているもの}を変えることは\textbf{できる}。いますぐ、動くことなく、注意を内へ引き寄せ、ここにない情景を築き、記憶を再生し、空想に沈み、周りの騒音から切り離され、世界が置いたのではないどこかに生きることができる。二年の遅れはない。行動もなく、外界を通って基盤を形づくる回り道もない。ただ自分の作った世界へと退くだけであり、それを\textbf{いま}行うのだ。それは、ようやく自分が采配を握れる唯一の場所------まぎれもなく、瞬時に\textbf{自分のもの}である唯一の行為------のように感じられる。

そうではない。ここが痛むところであり、痛むのはまさに、私たちがたった今リベット（Libet）を通り抜けてきたからだ。内へ退くという決定は、「あなた」がそれを下したと知る前に、基盤によって下されていた------実験室で手を挙げる決定とまったく同じように。「空想しようと私は選んだ」は、「腕を動かそうと私は選んだ」と同じ、遅れてくる観察者の報告であり、ただ上等な服を着ているにすぎない。基盤が先に発火した。あなたは後から知らされ、その知らせは作者性のように感じられるほど速く届いた。あなたがそんなに感心しているその即座さは、選択の自由ではなく、\textbf{連鎖の速さ}なのだ。ここに自由なものは何一つない。あなたの私的で主権的な行為のように感じられた内なる退避は、他のすべてと同じく、あなたのために決められていた。

では、意識は実際に何を寄与しているのか、選ぶことでないとすれば。それが寄与するのは\textbf{必然性}だ。基盤はスタートの号砲を鳴らしたかもしれないが、走っているシミュレーションがなければ、連鎖反応はどこにも進まない------感じるべき離脱もなく、退くべき内なる世界もなく、住まうべき空想もない。というのも、それらのことが起こりうる唯一の場所がシミュレーションだからだ。意識を消し去れば、下流のカスケード全体はただ伝播しない。それこそが真の第二の因果的役割だ。意志の座ではなく、連鎖のなかの、要となる環である。意識は独立した因果的な力を持たない------そして同時に、無用な同乗者でもなく、哲学者たちが恐れた随伴現象でもない。それは、それなしには残りが起こりえない、機構の一部である。自由な行為者ではなく、必然的な環だ。これは、認めるが、最初に聞こえるほどには胸躍るものではない。ようやく自由になれる唯一の場所を探しに行って、見つかるのは、自分が不可欠でありながら、それでも支配権を持たない場所なのだ。速い握りは、結局のところ速かったのであって、自由だったのではない。

そしてここで、私は一つの限界について正直でなければならない。というのも、行きすぎたくなる誘惑があるからだ。内への転回の\textbf{タイミング}が自分のものではない------リベットはそこまでは押さえているし、もはやそれは手足を動かすことだけの話ではない。人々が心の目に思い描くべき像を自由に選ぶとき、その選択はいまや、本人が下したと感じる数秒前に、視覚野から読み取れるようになった。言い換えれば、内なる行為の\textbf{開幕}は、手を挙げることと同じくらい前もって決められているように見える。すでにそこに至ってから、\textbf{内側で}起こることをどれだけ舵取りしているのかは、別の問いであり、正直な答えは、わからない、というものだ。そしてこの問いは、最初に見えるより鋭い。本当に問われているのは、自己モデルが一般にどれだけの自由を持つかではない------自己モデルが\textbf{世界}モデルにどれだけの握りを持つか、である。すなわち、アバターが、自分が身を置いているそのシミュレートされた世界へどこまで手を差し入れ、それを作り変えられるか、だ。それはそれ自体で一つの超能力であり、しかも生半可なものではない。それを試みることさえ、ある程度の深さの自己モデルを要する。だからこそそれは一様には分布しておらず、だからこそ、少なくとも途中までは\textbf{訓練できる}------ちょうど、何が浮かび上がるかをただ待つのではなく、生のフォスフェンから模様を経て、完全な顔や情景へと、自分自身の視覚野を意図的によじ登ることを学べるのと同じように（第5章）。その仕事のために、下にある機構がどれだけの自由度を自己モデルに手渡したのかは、誰も数えていない------そして、その問いの唯一の真に開かれた版、すなわち、何にも制約されないときに心が次に差し出すものが何か、については、誰もまったく検証していない。もしかすると、そこで何を探究するかを、実在する訓練可能な程度で本当に舵取りしているのかもしれない。もしかすると、たいていはしていないのかもしれない。内観はそれを決着させられない。それはここで確実に盲目である唯一の道具であり、本当に決着させるには、人間の脳の完全な配線図と、膨大な計算とが要るだろう。だからこの幻滅を、それがまさにそうであるとおりに受け取り、一寸たりとも踏み越えてはならない。内なる行為の\textbf{開幕}は、自分が書き記すものではない。幕が上がってから何をするか------その背後の世界をどこまで曲げられるか------は、開かれた問いのままであり、そして正直な問いだ。

というわけで、自由意志の議論があなたを連れて行く先はここだ。作者ではない------手の作者でもなく、白昼夢の作者でもなく、内側に入ってからどこまで舵取りするのかという問いの作者ですらない。そうして、ページをめくるに値するものがちょうど一つ残る。自分では駆らない意志を、どう使うのか。


\chapter*{第14章：差し出された唯一の自由}
\addcontentsline{toc}{chapter}{第14章：差し出された唯一の自由}
\markboth{第14章：差し出された唯一の自由}{}

そしてここで、釣り合いを取っておきたい。内なる世界は檻なのだと思い込んだまま、ここで本を閉じてほしくないからだ。まさに同じこの能力------目の前にあるものから切り離され、シミュレーションをそれ自身の内容の上で走らせる能力------は、適切な加減で使うかぎり、進化がこれまでに築いた中で最も強力なものなのである。認知のあらゆる働きはここに宿る。創造性もそうだ。記憶そのものも、そして古人が記憶の宮殿と呼んだあの技法------想像上の部屋を歩き回って、大切なことを収めていくやり方------もそうである。想像、計画、「代わりにこうしていたらどうなるか」という機構のすべて、現在から一歩外へ出て自分の心の中を前へ後ろへと旅する能力------そのどれひとつとして世界の協力を必要とせず、そのすべてが、内側へ引きこもってシミュレーションを内的な内容の上で走らせる、まさにこの能力を必要とする。外へ伸ばす緩やかな手は、幾世代もかけて文明を築いた。内へ伸ばすこの手は、一つの思考をまるごと一秒で築き上げる。これこそ、私たちが人間の思考と呼ぶすべてのエンジンである。軽く、いつでも引き返せるように、加減した量で手を伸ばすなら、それは並外れた力を発揮する。

だが、あまりに強く寄りかかれば、内側へ通してくれるそのダイヤルが、開いたまま固まってしまいかねない。内への退避を全力で推し進めれば、シミュレーションを現実の世界につなぎとめている細い水路が飲み込まれてしまう------先に述べたのと同じロックインであり、NDEや高用量のサルビアが否応なく強いるのと同じ暴走である。ただしここでは、自分の足で歩いて入り込んだという違いがある。十分に遠くまで押し進めれば、共有された世界へ戻る道は、もはやまったく見つけられなくなる。窓がもはや本物ではなくなったせいで窓から這い出していく、あのサルビアの被験者がそれである（第6章）。そして、そこまで強く追い込まれたシミュレーションは、必ずしも一つのものとして戻ってくるわけではない------枝分かれし、分裂し、一人ではない複数の自分として、あるいは出て行ったときの自分とは微妙に違う誰かとして帰ってくることがある（一つの脳を分け合う二つの心、第9章）。節度をもって使えば、これは手持ちの中で最良の道具である。回しすぎれば、正真正銘の危険となる------そして深入りしすぎれば、はたして戻ってこられるのか、戻ってくるのが何人の自分なのか、誰にもわからない。

\textbf{この知識で何ができるか。}ここまで理論を追ってきたなら、もうわかっているはずだ。意識的な自己（明示的自己モデル、ESM）は再構成であって、直接の読み出しではない。それは隙間を埋め、作話し、自分が下していない決定の手柄を自分のものにする。自分自身の基盤を見ることはできない。そして、それが自分の持つすべてである。

これには実際的な帰結がある。油断なく目を光らせておくべき食い違いが三つある。それらのあいだの隙間にこそ、人間の惨めさの大半が棲みついているからだ。

\begin{enumerate}
  \item 自分が\textbf{なりたい}もの------理想の自己、明示的自己モデルが目指す自分の姿。
  \item 自分が\textbf{そうだと信じている}もの------現在の自己モデル、毎日持ち歩いている「私」。
  \item 自分が\textbf{実際に}そうであるもの------現実の振る舞い、他者への実際の影響、外側から観察された基盤レベルのパターン。
\end{enumerate}

1と2のあいだの隙間は、自己改善のエンジンである。理想が現実的で、その食い違いが絶望ではなく行動を駆り立てるかぎり、これは健全だ。危険なのは2と3のあいだの隙間である------一人では測れないからだ。ESMは自分自身の基盤を正確に観察することが\textbf{できない}。他人からのフィードバックが要る。耳の痛い種類のものも含めて。いや、とりわけ耳の痛い種類のものが。

親友のベルンハルト（Bernhard）と私は、これを一種の遊びにしてしまった。二人のあいだには暗黙の了解がある。相手が犯したミスはすべて、即座に容赦なくからかうための好機である、と。運転中に曲がり角を見落とせば------「アルツハイマーも末期だな。キーは俺が預かろうか？」。何かを言い間違えれば------「また脳卒中じゃないか。舌を喉に詰まらせる前に黙ったほうがいい」。先週の会話の細部を忘れれば------「あとで今日のクロスワードパズル、手伝ってやろうか？」。

外から見れば、これは病的に聞こえるだろう。内側から見れば、私の知るかぎり最も効率的なESM較正システムである。冗談の一つひとつが補正信号なのだ------\textbf{お前の自己モデルは、いま基盤が意図していないことをやったぞ}、と。そしてこの嘲笑は本物の愛情にくるまれているから------二人とも、悪態を浴びせながら笑いをこらえている------どちらも過ちを弁護しない。更新するのである。これがコツだ。容赦なくやり合えるほど信頼できる相手と、間違うことが脅威ではなく笑いになる関係が要る。

理論は、どう生きるべきかを教えてはくれない。だが、どうすれば\textbf{自分を知る}ことができるかについて、重要なことを教えてくれる。すなわち、自己モデルには、どんなモデルにも向けるのと同じ健全な懐疑を向けよ、ということだ。それは有用である。手持ちの中で最良の表現である。そして、アーキテクチャ上の必然として、不完全である。

理論が手渡してくれるものが、もう一つある。これは危うく書かずに済ませるところだった部分だ。自己啓発に近づきすぎるし、私はあのジャンルへの耐性が低い。だが、ここまでのすべてから直接に導かれることなので、書いておく。

この本を読むことを、あなたは選んでいない。何かがそれを目の前に運んできたのだ------誰かの推薦か、アルゴリズムか、退屈な午後か------そして、自分が作者ではない状態にあった基盤が、読み続けることを決めた。それは運であって、美徳ではない。前の数ページで、あなたのすることに自由に選ばれたものは何一つないと述べたが、あれは本気である。ならば、歯車がどのみち回るように回るだけなら、これにいったい何の意味があるのか------そう問うのはもっともだろう。

意味はこうだ。そしてこれこそ、議論全体が積み上げてきた、希望へと向かう転回である。歯車は\textbf{形を変えられる}。決定論は、身動きが取れないという意味ではない------訓練可能だという意味である。基盤は固定された結晶ではない。触れるものと重ねる練習とが物理的に書き換えていく地形なのだ。これを読んだことで、望むと望まざるとにかかわらず、何かが変わった。そして、いったんシステムが一つの観念に出会えば、ここまで連れてきたのと同じ不自由な機構が、今度は時間の使い方を変えることを「決める」ことができる------基盤の産物ではある、そのとおり。だが現実のものであり、現実の帰結を伴う。自由意志の擁護者たちが夢見るような仕方で、それを選んだことにはならないだろう。それでもその選択は世界に着地し、ネットワークを組み替え、来週の自分が誰であるかを変える。それは無ではない。それこそが、そもそも差し出されていた唯一の種類の自由であり、人生を築く土台としては、それで十分だとわかるのである。

だから、この本で何か一つでも説得できたのなら、これであってほしい。ジムのために時間を確保するのと同じように、意識そのものに取り組む純粋な時間を確保し、それを交渉の余地のないものとして扱うこと。まずは瞑想------線香やプレイリストのためではなく、舞台を意図的に空にすることが、注意力にとってこれ以上ないほど純度の高い訓練だからだ。それは他のすべてに持ち込む根底の規律を研ぎ澄まし、一日じゅうすり減らしている経路を実際に再生させてくれる。\textbf{考える}ための時間を確保しよう------あてもなく考えるのでも、大切に思う一つの問題に向けてでもいい------そして、その時間をスマートフォンから守ること。何かを作ろう。音楽を奏でよう------それは他のどれにでも折り重なっていく。スポーツも含め、自分の意識を心から楽しめる仕方で主に使う趣味なら、何でも始めてみるといい------基準は生産性ではなく、内なる機械が好きなことをやっているかどうかである。そして、すでに持っている死んだ時間を捉え直そう。渋滞も、待合室も、行列も、時間を盗まれているのではなく、\textbf{贈られている}のだ------内側へ退き、シミュレーションを走らせるための、誰にも邪魔されない許可として。その退避を磨き上げよう。楽しもう。それは、決して取り消されることのない、自分との約束である。

次にメンテナンスがある------メンタルの衛生と呼ぼう。まさにそれ以外の何ものでもないからだ。第一に、否定的な内なる独り言を捨て去ること。あれは誠実さでも現実主義でもない。基盤に刻まれたすり減った溝であり、注意を向ければまた滑らかに均すことができる。物事は白か黒かで------いや、できれば白だけで------枠づけよう。灰色は、本当に避けようのない稀な場合のために取っておけばいい。灰色地帯こそ反芻が湧く場所であり、反芻とは、内へ伸ばす手が持ち主自身に向けられたものである。

基盤が自分自身に牙をむく話のついでに------怒りについて。人や物事に腹を立てることは、ほとんどの場合、自分の外側にある問題のために\textbf{自分自身}を罰することである。あの感情は、内側を向くようには作られていない。それは、まだ夕食を追いかけて仕留めねばならなかった時代に、行動を駆り立てるために進化したもの------世界に向けて放たれる資源動員の噴出であって、当の相手がどこか別の場所で素知らぬ顔で暮らしを続けるあいだ、自分の回路に注ぎかける腐食剤ではない。怒りの行き場が内側にしかないと気づいたら、それが、怒りを手放せという合図である。

そして、瞑想を試して弾き返されたなら------よくあることで、失敗ではない------経験から勧められる裏口がある。武道である。打撃とアドレナリンずくめで、静かな心とは正反対に聞こえるかもしれない。だが、意識のおしゃべりを脇へどけながら身体が動くように鍛えるという鍛錬は、瞑想が教えるのと同じ技能に、反対側から到達したものなのだ。正直に言えば、二つのうちで\textbf{難しい}のは瞑想のほうであって、易しいほうではない------自分の基盤とともにじっと座り、静かにしてくれと頼むことは、人が試みうる最も困難なことの一つである。座布団に負けたなら、道場を通る遠回りの道を行けばいい。行き着く先は同じ静かな心である。ただ、そこへ着くまでに汗をかくことが許されているだけだ。

\section*{私にわからないこと}

開かれた問いが一つもないと主張する理論は、理論ではない------宗教である。だから、私が本当に確信を持てずにいる場所、これからの10年の研究が集中すべき場所を、ここに挙げておく。

\textbf{暗黙的モデルもまた仮想なのか？}（あるいは、どの程度まで）IWM（暗黙的世界モデル）とISM（暗黙的自己モデル）は「モデル」である。だが、正確には何のモデルなのか。私は現実の基盤と仮想のシミュレーションのあいだにきれいな線を引いてきたが、暗黙的モデルはまさにその線の上にまたがっている。もしそれらもある意味で仮想なのだとすれば、真に「現実」の底を構成するのは何なのか。理論は現実／仮想のきれいな二分を前提としているが、現実は私の図式よりも入り組んでいるかもしれない。これは、最終的な答えを私が持っていない、根本にかかわる問いである。

\textbf{数学的形式化。}理論は現時点では定性的である。図を描き、機構を記述し、予測を立てることはできるが、方程式を手渡すことはできない。臨界性の要件はウォルフラム（Wolfram）のクラス4セル・オートマトンを引き合いに出しており、力学系理論には持ち込めるはずの形式的な道具立てもある。だが、完全な数学的形式化------仮想モデルがいつ、どのように基盤の動態から創発するのかを正確に指定する方程式------は、まだ存在しない。これが最大の欠落である。数学を欠いた意識の理論は、物理学者に真剣に受け取ってもらえない意識の理論であり、ものの作り方を知っているのは彼らなのだ。

\textbf{オートマトン＝ホログラム予想------開かれた挑戦。}第5章で、ホログラフィックな系とクラス4セル・オートマトンとのあいだにありうる三つの関係を述べた。第一のもの（クラス4の動態を生み出すホログラフィックな基盤）は、ほぼ間違いなく脳がやっていることであり、美しくはあるが、衝撃的ではない。だが残りの二つは、あの場所で私が与えたよりもはるかに多くの注意に値する。

実のところ、ここには三つの開かれた問いがあり、後になるほど途方もなさを増していく。

\textbf{第一------クラス4オートマトンは、その創発的な出力としてホログラフィックなパターンを生み出せるか。}カオスの縁にある局所的な規則は、大域的で非局所的な情報符号化を、創発的な振る舞いとして生成できるか。もしできるなら、純粋に局所的な相互作用が、ホログラフィーの記述するような分散した冗長な情報構造を自発的に生み出す系が手に入ることになる。そしてこれは、興味深いことに、情報理論の観点から見た量子もつれの姿と正確に一致する。

\textbf{第二------クラス4オートマトンは、ホログラフィックな規則構造を持ちうるか。}規則そのものが、ホログラムが三次元を二次元に符号化するように、高次元の情報を低次元の構造に符号化しているセル・オートマトンを想像してみてほしい。あらゆる局所的相互作用が、暗黙のうちに大域的構造を含むことになる。その規則は複雑な振る舞いを生み出すだけではない------より大きな何か、より高次元の何かが低次元の規則集合へと投影された、その圧縮符号化\textbf{そのもの}だということになる。

\textbf{第三------そしてこれこそ、夜、私の眠りを奪うものだ------両方が同時に真でありうるか。}規則がホログラフィックで、動態がクラス4で、出力がふたたびホログラフィックであるような系。そんなものが存在するなら、自分自身を符号化する計算過程------おのれの構造を計算する宇宙------が手に入る。入力はホログラフィックである。処理はカオスの縁にある。出力はふたたびホログラフィックである。それは不動点------自己整合的なループである。

そのようなオートマトンが存在するなら、それはホログラフィック原理が宇宙について述べていることを\textbf{そのまま}行っている。緩やかな比喩の意味で宇宙に似ている系ではない。高次元の実在を低次元の規則に符号化し、秩序とカオスの境界で計算し、その圧縮から創発的な複雑さを生成する系である。それは宇宙の比喩ではない。それは宇宙\textbf{そのもの}かもしれないのだ。

はっきり言おう、誰かが言うべきだと思うからだ------ホログラフィックな規則構造を持ち、なおかつホログラフィックな出力を生み出すクラス4のセル・オートマトンが存在するなら、それはほぼ間違いなく宇宙そのものだと私は確信している。それはWeltformel------世界方程式となるだろう。黒板に書きつける数式という意味ではなく、私たちが観測するすべて------量子力学から一般相対論、そして意識そのものの創発に至るまで------を生成する計算過程という意味においてである。

これは、正直に認めるが、本書で最も思弁的な考えだ。私には証明がない。候補となる規則の集合すらない。そして、数学的な美から物理的実在へと至る議論が正当に批判されてきたことも認めておくべきだろう。ザビーネ・ホッセンフェルダー（Sabine Hossenfelder）をはじめとする人々が、優美さは証拠ではないと指摘してきた。彼女は正しい。この考えの本格的な探究は、続く三つの章の主題である。だが、問い自体は明確に立てられており、数学的に厳密だ------

\textbf{規則構造がホログラフィックであり、かつ力学がクラス4であるようなセル・オートマトンは存在するのか。それはホログラフィックな出力を生み出すのか。この三つの性質はすべて共存しうるのか。}

これらは数学者への問いであって、神経科学者への問いではない。規則空間の組合せ論についての問い、ホログラフィックな符号化と計算論的な普遍性が有限な局所的規則の集合のなかで共存しうるかについての問いだ。そのようなオートマトンは存在しえないと証明できるかもしれない------それはそれで深遠な結果だろう。というのも、それは情報圧縮と計算のあいだの関係について、何か深いことを教えてくれるからだ。あるいは、そのようなオートマトンが存在し、構成可能だと証明できるかもしれない------そのときには、これまで提唱されたなかで最も根源的な物理的実在の記述の候補を手にしたことになる。

どちらの答えが正しいのか、私にはわからない。だが、これらの問いは問われるに値すること、そして誰もそれを問おうとしていないらしいことは、私にはわかる。だからこれを、開かれた挑戦状だと思ってほしい。証明するか、反証するか。証明できれば、宇宙のソースコードを見つけたことになるかもしれない。反証できれば、ホログラフィと計算を結ぶ深い不可能性定理を打ち立てたことになる。いずれにせよ、その答えは途方もなく重要だ。

そして、もし本当にそのようなオートマトンを見つけたなら------私に電話してほしい。確かめてみたい予測がいくつかあるのだ。

\textbf{どの物理的メカニズムか。} この理論は臨界性を要求するが、それを支える物理的メカニズムについては意図的に不可知の立場をとる。皮質コラムの力学だろうか。視床皮質定在波だろうか。シナプス活動へのグリア細胞による調整だろうか。三つとも経験的な裏づけがある。理論は「基盤は臨界性になければならない」と言うが、基盤が\textbf{どのように}そこへ至り、そこにとどまるのかは言わない。それは欠陥ではない------むしろ、理論が特定のメカニズムによらず成り立つことを意味している。だが、いつかは誰かがそれを突き止めなければならない。

\textbf{最小構成。} ESMなしにEWMを持つことはできるのか。自己経験なしの世界経験は。意識的とみなせる最小限の構造とは何か。動物の章で述べた段階的なレベルが手がかりになる------おそらく魚がそうであるように、たいした自己モデルを持たずとも豊かな世界モデルを持つことは可能だ。だが、閾値は正確にはどこにあるのか。明かりが灯るまでに、どれだけの自己モデルが必要なのか。ESMこそがシミュレーションを経験へと変えるものだと私は論じてきたが、その最小限で成立する形態までは特定していない。

これらの問いを、私は弱点としてではなく、研究の最前線として挙げている。それらは、理論が実在と接触し、こう告げる場所だ------ここで私を試せ、ここで私を形式化せよ、できるものならここで私を打ち破ってみよ。

そのなかに、どうしても放っておけない問いが一つあった------夜も眠れなくさせると告げた、あの問いだ。だから私は、安全な距離からそれを問うのをやめ、探しに出かけた。次に来るのは、それが私を連れていった先である。


\chapter*{第15章：同じパターンが、いたるところに}
\addcontentsline{toc}{chapter}{第15章：同じパターンが、いたるところに}
\markboth{第15章：同じパターンが、いたるところに}{}

かつて私は、フォアアールベルクでもっとも暗い場所で夏の一夜を過ごしたことがある------山の高み、何キロメートル四方にも人工の光のない場所だ。仰向けに横たわり、空を見上げた。天の川は、目を凝らしてやっと見えるかすかな滲みではなかった。それは川だった。濃密で明るく、傍らの岩に目に見える微光を投げかけていた。そしてその夜のどこかで、何かが切り替わった。頭上の星を見るのをやめ、身体の下の地球を感じはじめたのだ------回転し、私を乗せて運びながら、一千億の太陽からなる銀河の中を突き進んでいく球体を。観念としてではない。身体の感覚としてである。地面は静止していなかった。私は、何か途方もなく巨大なものの中を移動していく何かの外側に、しがみついていたのだ。

あのめまい------自分は宇宙の中の一個の岩の表面にいる小さな温かいものなのだという、突然の身体的な知------それこそが、この先を読み進めるあいだ手放さずにいてほしい感覚である。なぜならこの章は、その宇宙とはいったい何なのかを問うからだ。そして、その答えはひどく見覚えのあるものになる。


前章の最後で、私は一つの課題を開いたまま残しておいた。ホログラフィック系とクラス4セル・オートマトンのあいだに考えられる三つの関係を描き、そのうえでもっとも難しい問いを投げかけた------三つすべてが単一の系の中で共存しうるか。クラス4オートマトンは、ホログラフィックな規則構造を持ち、\textbf{なおかつ}ホログラフィックな出力を生み出しうるか、と。この着想を余すところなく探究することが続く三つの章の主題になる、と私は述べた。

これが、その探究である。

これから述べるのは、本書でもっとも思弁的な部分である。そして私の考えでは、もっとも重要な部分でもある。実際に腰を据えてその糸をたどってみたとき------「いつかの問い」として扱うのをやめ、糸を引きはじめたとき------私は、予想していた場所にはたどり着かなかった。見つかるのは興味深い数学的な珍品だろうと思っていた。ところが見つかったのは、一つの宇宙論的モデルだった。そして、二十年間見つめつづけてきたのと同じアーキテクチャだった。

どういうことか、これから示そう。


\section*{宇宙の計算クラス}

付録Cでは、計算の五つのクラスを示した------完全な秩序から完全な無秩序へと至るスペクトルであり、その中でクラス4は、記述可能な規則によって到達しうる最大の複雑さとして、カオスの縁（ふち）に位置している。脳は五つのクラスすべてを道具として使うが、意識はもっぱらクラス4にのみ宿る。それが脳についての議論だった。

より大きな問いはこうだ------宇宙はどのクラスなのか。

これは比喩ではない。文字どおりの問いである。宇宙を力学系として扱うなら------実際そうなのだが------それは五クラスのスペクトルのどこに落ちるのか。答えは消去法で定まる、と論じたい。そしてその消去は、驚くほど鮮やかに進む。

\textbf{クラス1と2------静的と周期的。} クラス1の宇宙は固定した状態に収束する。何も起こらない。クラス2の宇宙は反復するループに落ち着く------永遠に時を刻みつづける時計の宇宙版である。どちらも化学も生物学も進化も意識も生み出せない。私たちは存在している。私たちには意識がある。意識を生み出す宇宙は少なくともクラス4でなければならない。意識にはクラス4の動態が必要だからだ------それが第5章の議論だった。そして、より低いクラスは、より高いクラスを部分プロセスとして生成できない。周期的な宇宙がカオスの縁の動態を生み出せないのは、時計がひとりでに考えはじめることがないのと同じである。除外。

\textbf{クラス3------フラクタル。} こちらはもう少し微妙だ。フラクタル宇宙なら美しいものになるはずだからである。あらゆるスケールでの自己相似構造、パターンの中に入れ子になったパターン。実際、宇宙はフラクタル構造を\textbf{持って}いる------銀河団、海岸線、河川網、肺の枝分かれ。だがフラクタル系は計算論的に\textbf{還元可能}である。つまり、先回りができる。ここからそこまでのすべてのステップを走らせなくても、タイムステップ百億時点でのフラクタル系の状態を計算できてしまう。近道があるのだ。

私たちの宇宙は近道を許さない。先へ跳ぶ方程式を書いて来月の天気を予測することはできない。シミュレーションを一歩一歩走らせるしかない。その動態が計算論的に還元不可能だからである------各瞬間は、圧縮のきかない仕方で、真にその直前の瞬間に依存している。フラクタル宇宙は、そのパターンがどれほど豊かであろうと、この性質を欠く。私たちの宇宙が現に支えていることが実証済みの万能計算を、フラクタル宇宙は維持できないだろう。私たちはチューリングマシンを造る。私たちには意識がある。フラクタル宇宙にはどちらもできない。除外。

\textbf{クラス5------ランダム。} 宇宙の根本的な動態が本当にランダムであったなら（複雑に見えるだけではなく、真にランダムであったなら）、物理学は不可能になる。現在私たちが理解している物理学が、ではない。\textbf{営みとしての}物理学が、である。科学という企て全体は、宇宙が記述可能な規則に従うという想定の上に成り立っている。書き留め、検証し、伝達し、将来の観測の予測に使うことのできる規則である。真にランダムな宇宙には記述可能な規則がない。その動態はいかなる公式にも、いかなる法則にも、いかなる方程式にも圧縮できない。F = maを書き下すことはできないだろう。力と質量と加速度の関係が、どんな有限の記述にも捉えられない仕方で、刻一刻と変わってしまうからである。

クラス5の宇宙では、あらゆる実験が一回かぎりのものになる。再現可能な結果は偶然の一致にすぎない。科学は、しばらくのあいだたまたま機能しただけの錯覚である。これは論理的に不可能なことではない------そのような宇宙を想像しても矛盾は生じない。だが説明という点では破滅的である。これを受け入れれば、何も説明できなくなる。自分の説明がなぜかつて機能するように見えたのかさえも。除外。ただし論理によってではなく、アブダクション（最良の説明への推論）によってである------私たちの一貫して法則的な経験の最良の説明は、宇宙が記述可能な規則によって動いているということなのだ。

\textbf{残るのはクラス4である。} カオスの縁。そしてクラス4は、私たちの観測と整合的であるだけではない------すべての条件を満たす\textbf{唯一の}クラスなのである。

宇宙は安定した構造を含んでいる。原子、結晶、山々。これはクラス1の振る舞いだ。周期的な現象も含んでいる。軌道、潮汐、心拍。これはクラス2の振る舞いだ。フラクタル構造も含んでいる。銀河の分布、気象パターン、神経の枝分かれ。これはクラス3の振る舞いだ。そして宇宙は万能計算を支えている。私たちはコンピュータを造り、私たちには意識がある。これはクラス4の振る舞いだ。すべてのクラスを部分プロセスとして------自分自身をも含めて------内包できるのは、クラス4の系だけである。他のどのクラスにもこれはできない。クラス4オートマトンは、より単純な動態を宿すだけではない。他のクラス4オートマトンをも宿すのだ。万能コンピュータの中の万能コンピュータであり、そのそれぞれが全体と同じ計算上の芸当をこなせる（ただし、より小さく、より遅く、資源に制約されてはいるが）。

さらに重要なことがある。クラス4は、他のどのクラスも持たない自己維持のメカニズムを備えている。\textbf{自己組織化臨界性}------第5章で最初に出会ったあの仕組み------である。パー・バック（Per Bak）は1987年に、カオスの縁にある系はたまたまそこに居合わせるのではない------\textbf{自らそこへと向かっていく}のだということを示した。砂を一粒ずつ積み上げていくと、砂山はあらゆる規模の雪崩が生じる臨界角へと自らを組織していく。系を臨界性に合わせて調整する外部の手は要らない。系は自らを調整する。カオスの縁が宇宙的な時間スケールにわたって安定であるのはこのためだ。それはアトラクターであって、偶然ではない。

これがどのような種類の議論なのか、はっきりさせておきたい。演繹的な証明ではない。四つの消去のうち二つは経験的観測に基づいている（宇宙は意識を含む。宇宙は万能計算を支えている）。一つはアブダクションに基づく（クラス5は科学を不可能にする------すっきりしないが、論理的矛盾ではない）。クラス4を支持する積極的な論拠は、証拠とメカニズムを組み合わせたものである。これは主張しうるかぎりで最強の主張だ。クラス4は、すべての観測と整合する唯一のクラスであり、自らの存続に理由を与える唯一のクラスなのである。

少しのあいだ、あの岩塊に戻ろう。山腹で私がしがみついていたあれ------動いているのを感じた瞬間、胃の底が抜けたあれ------は、ただ古く、途方もなく大きく、無関心なだけのものではない。あれはクラス4なのだ。カオスの縁とは、深遠に聞こえるように私が持ち出した言い回しではない。それは私を乗せて運んでいた大地そのものであり、バックの砂山と同じように一粒また一粒と自らを臨界性へと調整し、誰の手もダイヤルに触れないまま全体の均衡を保っていたのである。私は死んだ石の上に横たわっていたのではない。ナイフの刃の上に自らを保ちつづける計算の外側にしがみついていたのだ------そしてあなたもいま、この瞬間、その回転を感じられようと感じられまいと、同じようにしがみついている。


\section*{情報の地平線}

ここからは、境界の話である。

光の速さは有限である。これは、10分ほど考えてみるまでは無害に聞こえ、考えたが最後、現実についての全体像を組み替えてしまう類いの事実の一つだ。

光は秒速およそ30万キロメートルで進む。まばたきする間もなく部屋を横切れる速さだが、宇宙はとてつもなく広い。もっとも近い恒星までは4光年。もっとも近い大銀河までは250万光年。観測可能な宇宙は差し渡しおよそ930億光年である。遠くの銀河を眺めるとき、見えているのは何十億年も前の姿だ。光がここまで届くのにそれだけの時間がかかったからである。私たちはつねに、不可避的に、過去を覗き込んでいる。

しかし、もっと深い帰結がある。それは宇宙の膨張から生じるものだ。

1998年、二つの天文学者チームが、のちにノーベル賞をもたらすことになる発見をした。宇宙の膨張は加速している。ただ膨張しているのではない------加速しているのだ。遠方の銀河は私たちから遠ざかっており、遠ざかる速度は増しつづけている。つまり、どんな観測者にとっても、後退速度が光速を超える距離が存在するということだ。その距離の向こうからは、いかなる信号も決して届かない。情報が壁の向こうに隠されているからではない。自分と情報とのあいだの空間そのものが、光が渡りきれる速さを超えて広がりつづけているからである。

これが\textbf{宇宙論的地平線}である。物理的な表面ではない。そこに壁があるわけではない。幾何学と速度の帰結にすぎないが、どんな壁にも劣らず絶対的な障壁である。地平線の彼方にある情報には、永遠に手が届かない。存在しないも同然なのだ。

底の側にも似た境界がある。\textbf{プランク長}（約$10^{-35}$メートル------これを「小さい」と呼ぶのは、観測可能な宇宙を「中くらいの大きさ」と呼ぶようなものだ）は、私たちの知る物理学が破綻する場所である。このスケールより下では、方程式は機能しない。時空そのものが物理的な意味を失う。プランク長より下では、いかなる測定も不可能である------原理的にさえ。技術的な限界ではない。知りうることの根本的な境界なのだ。

宇宙論的地平線とプランク・スケールのあいだ------およそ60桁。それが宇宙の計算領域、物理学が働く範囲である。その上と下では、幕が閉じられている。

これによって宇宙は、私が\textbf{準無限}と呼ぶものになる。真に無限なのではない。少なくとも、無限であることを確かめる術は決してない。有限の領域を超えてアクセスすることができないからだ。しかし、到達しうるどんな意味においても有限なのでもない。境界は、近づく速さを上回って遠ざかっていく。縁にたどり着くことは決してできないが、縁はそこにある。内側から見れば、宇宙は限りなく広がっているように見える。外側から見れば------だが、外側というものはない。まさにそこが肝心なのだ。

あの眩暈がまた戻ってくる。ただし今度は、座標を持っている。山腹で私は、大地が闇のなかを自分を運んでいくのを感じることができた。そこにこの二つの境界を加えると、事態はいっそう不穏になる。下には、決して届かない床。上には、光が追いかけるよりも速く逃げていく縁------それもあらゆる方向へ同時に、逃げ込める外側もないままに。岩の表面にへばりつく、同じ小さくて温かいもの------ただし今や、それが落ちていく闇には形があり、その形とは、あらゆる側で私から遠ざかりながら、何ひとつ内へ通さない壁なのである。


\section*{あらゆる境界は同じひとつの境界である}

ここで、この章の中心となる考えを述べよう。じっくり時間をかけて向き合ってみてほしい。もしこれが正しければ、あらゆる物事についての考え方が変わってしまうからだ。

まず棚卸しをしてみよう。宇宙には特異点が含まれている------物理的な記述が破綻する場所、情報の伝達が止まる場所、方程式が発散するか沈黙する場所である。これらの特異点は、まったく異なるスケール、まったく異なる文脈に現れる。物理学者はそれらを別々の現象として扱う。だが私は、すべて同じものだと考えている。

\textbf{1. プランク領域。} 物理学が機能する最小のスケールで、時空は記述不能な何かへと溶けていく。このスケールより下では測定は不可能である。情報はそこを通り抜けられない。

\textbf{2. 粒子の内部。} 電子やクォークは、標準模型では点状のものとして扱われる（ゼロ次元で、内部構造を持たない）。その内側を見ることはできない。性質（電荷、スピン、質量）は測定できるが、その中心部で何が起きているのかには、いっさいアクセスできない------空間的な広がりを持たない対象について、「中心部」という言葉がそもそも意味をなすのならの話だが。

\textbf{3. ブラックホールの事象の地平線。} 情報は落ち込んでいく。何も出てこない------少なくとも、入ったものを保存するような形では。内部は外部から因果的に切り離されている。ブラックホールの内側で何が起ころうと、外部のいかなる観測者にとっても、それは内側にとどまったままである。

\textbf{4. 宇宙論的地平線。} 観測可能な宇宙の縁であり、その向こうでは、空間の膨張によっていかなる信号も私たちに届かない。隠された情報ではない------到達不能な情報である。

\textbf{5. ビッグバン。} 始まり。すべての世界線が収束する。宇宙のあらゆる粒子は、その歴史をこの点まで------いや、正確にはこの境界まで------さかのぼる。「点」という言葉はそこへ行けることを含意するが、行くことはできないからだ。

\textbf{6. 時間的な終端。} 終わり------それがどんな形で訪れるにせよ。宇宙が熱的死で終わるなら、エントロピーは最大に達し、いかなる過程を駆動する熱力学的勾配も残らない。ビッグクランチで終わるなら、すべての物質はただ一点へと崩壊し戻る。ビッグリップで終わるなら、加速する膨張があらゆるスケールで時空を引き裂く。三つのシナリオはいずれも後で検討する。ここで重要なのは構造的な主張である。宇宙が実際にどの終わりを迎えるにせよ、それは情報を通さない境界で終端する、ということだ。

六つの特異点。六つの異なるスケール、六つの異なる文脈、それらを研究する六つの異なる物理学の分野。だが、それらに共通するものを見てほしい。

\textbf{第一に、いずれも情報を通さない。} どの境界も、情報に越えさせることはできない。プランク長より下は測定できない。電子の内側は見えない。事象の地平線の向こうから情報を取り戻すことはできない。宇宙論的地平線の彼方からの信号は受け取れない。ビッグバンの「前」に何があったかは観測できない。そして、宇宙の最終的な境界がどんな形で現れるにせよ、その先へメッセージを送ることはできない。

\textbf{第二に、いずれも最大の情報密度を体現している。} こちらはより微妙で、ベケンシュタイン限界に由来する------1980年代に得られた結果であり、空間のある領域が含みうる情報量の最大値は、その体積ではなく\textbf{表面積}に比例することを示している。ブラックホールの事象の地平線はこの限界を飽和させる------単位面積あたり可能な最大の情報を保持しているのだ。ヘーラルト・トホーフト（Gerard 't Hooft）とレオナルド・サスキンド（Leonard Susskind）が提唱したホログラフィック原理は、これを一般化する。任意の領域に含まれるすべての情報は、その境界の上に符号化されている、と。これらの特異点はいずれも、最大容量で働いている境界面なのである。

\textbf{第三に、いずれも計算領域を区切っている。} 物理学はこれらの境界の\textbf{あいだ}で働くのであって、その先で働くのではない。物理法則が記述するのは、プランク・スケールと宇宙論的地平線のあいだ、ビッグバンとやがて訪れるいずれかの終端とのあいだの領域で起こることである。境界が舞台を定める。舞台の外では規則は適用されない------別の規則が適用されるからではなく、「規則」というものが意味のある概念でなくなるからだ。

三つの共通の性質。六つの現象。従来の見方では、これらはたまたまいくつかの特徴を共有する、六つの別々の事柄だとされる。私は、その従来の見方は間違っていると考える。これらは\textbf{ひとつの現象}------オートマトンの情報境界------が、六つの異なるスケールに現れたものなのだ。

これは対称性についての主張である。同じ構造要素の反復。そしてクラス4のシステムでは、これこそまさに予期されることなのだ。クラス4のダイナミクスは、すべてのクラスを部分プロセスとして含む------クラス4自身も含めて。クラス4オートマトンは、その動態のなかに、より小さなクラス4オートマトンを入れ子にしており、そのそれぞれが、全体と同じ構造的性質を持つ情報境界によって区切られている。もし宇宙がクラス4オートマトンなら、その境界構造はあらゆるスケールで繰り返されるはずだ。そして、まさにそれこそ、私たちが見出しつつあるものらしいのである。


\textbf{つまり、あの山腹の夜が本当に意味していたのは、こういうことだった。私は岩にしがみついていたのではない。クラス4オートマトンの外面にしがみついていたのだ。そしてそのなかのあらゆる境界は------最後に見える星のさらに彼方で、私の心が絶えず押しあてていたあの境界も含めて------同じひとつの境界である。次の章では、この糸をその帰結までたどっていく。その帰結は、私が予想していたよりも奇妙なものだ。}


\chapter*{第16章：万物のアーキテクチャ}
\addcontentsline{toc}{chapter}{第16章：万物のアーキテクチャ}
\markboth{第16章：万物のアーキテクチャ}{}

前章では、一つの構造的な主張を立てた。宇宙のあらゆる特異点は------プランクスケールから宇宙論的地平線まで、ビッグバンからどのような終末が待ち受けているにせよその終わりまで------異なるスケールに現れた同じ現象である、という主張だ。一つの境界が、至るところで繰り返されている。

ここからは、その糸をさらにたどっていきたい。すべての境界が同じものであるなら、注目すべき帰結がいくつも導かれるからだ------時間について、物質について、知の限界について、そして終わりにたどり着く頃にはひどく見覚えのあるものになっているはずの、あるアーキテクチャについて。


\section*{ビッグバンは思っているようなものではない}

話はさらに先へ進む。ほとんどの人が------ほとんどの物理学者を含めて------まだ完全には呑み込めていないと思われる帰結があるからだ。

ブラックホールに外側から近づくと何が起こるかを考えてみよう。事象の地平線に近づくにつれて、時間は引き延ばされる。携えた時計は、遠方の観測者から測れば遅れていく。地平線に近づくほど、この遅れは無限大に近づく。落下を見守る遠方の観測者の目には、落ちていく者が減速し、赤方偏移し、薄れていくように映る------地平線に到達しきることは、ついにない。観測者の視点では、到着には永遠の時間がかかる。実際に地平線を越えることは決してない。事象の地平線は、外側から見れば、漸近的にしか近づけない境界なのである。

今度は、時間をさかのぼってビッグバンへと旅することを考えてみよう。

ビッグバンはどれほど昔のことか。約138億年前だと言われている。しかしそれは、膨張する宇宙の\textbf{内側から}の測定であり、使われる時計自体が膨張の産物である。時間を逆向きにたどり、宇宙という映画を巻き戻していくと想像してみてほしい。特異点に近づくにつれて何が起こるか。時間は引き延ばされる。物理法則は破綻する。近づけば近づくほど、方程式は決定的な「時刻ゼロ」を与えることを拒むようになる。ビッグバンは、指をさして「あそこだ------あのとき起こったのだ」と言えるような出来事ではない。それは漸近的な境界である。いくらでも近づくことはできるが、決して到達することはできない。

宇宙論的地平線が空間における事象の地平線であるのとまったく同じように、ビッグバンは時間における事象の地平線なのである。

これは神秘主義ではない。同じ数学的構造から導かれる帰結である。事象の地平線とは、その向こうへ情報が通り抜けられない面のことだ。ビッグバンはまさにこの性質を持っている。その「以前」からの情報は（そもそも「以前」という言葉に意味があるとしての話だが）、一切アクセスできない。情報が失われたからでも隠されたからでもなく、境界そのものが情報を透過させないからである。アクセスできる「以前」など存在しない------外部の観測者がアクセスできるブラックホールの「内部」が存在しないのと同じことだ。境界は境界である。それだけのことだ。

では、もう一つの時間的境界------終わり------についてはどうか。

宇宙が熱的死で終わるなら------最大のエントロピー、最大の無秩序、いかなる過程を駆動する熱力学的勾配ももはや残っていない状態------その時点で、あらゆる情報は最大限に分散されている。系の境界は、可能なかぎり最大の情報を保持している。これがベケンシュタイン飽和である。熱的死は、私がここまで使ってきた定義に照らせば、まさに特異点\textbf{である}。すなわち、最大の情報密度における、情報を透過させない境界だ。

ここからが奇妙なところだ。この枠組みでは、特異点は情報を破壊しない。情報を\textbf{変換する}のである。実のところこれは、ブラックホール情報パラドックス------ブラックホールに落ち込んだ情報は失われるのかをめぐる数十年越しの論争------について、現代物理学が収束しつつある解決そのものだ。現在のコンセンサスは「失われない」へと傾きつつある。情報は保存され、事象の地平線上に符号化され、やがて再び放出される。特異点は、情報を圧縮された形と展開された形のあいだで変換するのである。

これを時間的境界に当てはめてみよう。熱的死が特異点であり、特異点が情報を破壊するのではなく変換するのだとすれば、熱的死は宇宙を終わらせない。情報を新たな圧縮状態へと変換するのである。では、ベケンシュタイン飽和に達した最大限に圧縮された状態とは、どのような姿をしているか。それは、新たな膨張のための初期条件のように見える。ビッグバンのように見えるのだ。

自己言及的閉包は空間的なものにとどまらない。時間的なものでもある。宇宙は始まりも終わりもしない------循環するのだ。終状態は、次の反復のための初期条件である。何か特殊な「跳ね返り」機構のせいではなく、情報を保存する特異点とはまさにそういうことを\textbf{する}ものだからだ。すなわち、圧縮された境界状態と展開された内部状態のあいだで変換を行う。熱的死は圧縮する。ビッグバンは展開する。両者は同じ特異点を、反対側から見たものなのである。

これが思弁的であることは認める。しかしこれは、二つの主張から直接に導かれる。すべての特異点は構造的に同一であるという主張と、特異点は情報を変換することによって保存するという主張である。この前提を受け入れるなら、時間的な循環性は追加の仮定ではない------帰結なのである。

だが、この物語の結末は熱的死だけとは限らない。おそらくはさらに奇妙な代替シナリオがあり、この枠組みはそれも同じように整然と扱うことができる。

もしダークエネルギーが一定ではなく時間とともに増大するなら------その密度が際限なく増えていくなら------宇宙の膨張はただ続くだけではない。あらゆる限度を超えて加速していく。これが\textbf{ビッグリップ}のシナリオであり、その名が示すとおりの劇的な結末だ。まず、銀河団が引き裂かれる。銀河のあいだの空間が、重力がつなぎ留められるよりも速く引き伸ばされるからだ。次に個々の銀河が解体する。次に恒星系が。次に惑星が。やがて膨張が電磁気力を圧倒し、原子そのものが引き裂かれる。そして最後には、時空そのものが断片化する。あらゆる点が特異点になるのである。

この枠組みでは、ビッグリップには自然な解釈が与えられる。ふだんは宇宙論的地平線という安心できるほど遠い場所にとどまっている特異点境界が、\textbf{内側へ}と伝播してくるのだ。観測可能な宇宙の縁でこちらを待ってはいない。向こうからやって来る。特異点境界は計算領域をますます小さな領域へと断片化し、その一つひとつが自らのベケンシュタイン限界を飽和させ、一つひとつが情報を透過させない境界となる。時間の終わりにそびえる一つの壮大な特異点の代わりに、あらゆるスケールで同時に内側へと伝播していく、特異点のフラクタル的爆発が現れるのである。

そして、特異点が情報の変換器であるなら------計算を破壊するのではなく再起動するのだとすれば------ビッグリップが生み出す再起動は一つではない。\textbf{数多く}の再起動を生み出すのだ。潜在的には無限に多く。砕け散った計算領域の断片の一つひとつが、それぞれ独自の新たな膨張を、新たな展開を、新たな宇宙を芽吹かせうる。ビッグリップとは、この枠組みにおいて、マルチバースの生成装置なのである。

つまりこの枠組みは、一つではなく三つの終局シナリオを受け入れることができ、その三つはいずれも構造的に整合している。

熱的死：一つの全域的な特異点、一つの再起動。最も単純なケースだ------計算領域全体が同時にベケンシュタイン飽和に達し、圧縮され、新たなサイクルへと展開する。

ビッグクランチ：宇宙が膨張をやめ、一点へと崩壊して戻る。これも全域的な特異点であり、もう一つの再起動である------ことによるとCPT反転、すなわち電荷・パリティ・時間の反転を伴い、次のサイクルが前のサイクルの鏡像になるかもしれない。

ビッグリップ：特異点境界が内側へと断片化し、多くの特異点、多くの再起動、そして潜在的には多くの宇宙を生み出す。循環ではなく、枝分かれする樹形である。

この頑健さを、私は不安ではなくむしろ安心の材料だと考えている。宇宙が特定の終わり方をした場合にしか成り立たない枠組みは脆弱だ------まだ確かめようのない特定の宇宙論的結末に賭けていることになる。この枠組みは賭ける必要がない。その構造的な論理（情報の変換器としての特異点、根本的なアーキテクチャ要素としての境界）は、宇宙が実際にどの終局を選ぼうとも成り立つ。理論に求めたいのは、まさにこの種の頑健さである。理論は、まだわかっていない事実に依存すべきではないのだ。


\section*{粒子とは本当は何なのか}

この枠組みには、明示しておくに値する予測が一つ埋め込まれている。いずれは検証されうる類のものだからである。

素粒子（電子、クォークといった物質の構成要素）は、標準模型では点状のものとして扱われる。ゼロ次元。空間的な広がりを持たない。だがこれは常に、物理的な主張というより数学上の便宜であった。電子が文字どおり幾何学的な点だと信じている者はいない。幾何学的な点は表面積を持たず、したがってベケンシュタイン限界により、いかなる情報も含みえないからだ。ところが電子は情報を含んでいる------電荷、スピン、質量、量子数。「点」という描像のどこかが間違っているのである。

予測は具体的である。素粒子はプランク・スケールの特異点である。真にゼロ次元なのではない。ミニチュアの情報境界------ごく小さな事象の地平線------であり、その内部はブラックホールの内部と同じように到達不可能である。素粒子は、そのスケールでベケンシュタイン限界を飽和させるプランク・スケールの構造を持つ。そしてその表面は、ブラックホールの事象の地平線が落ち込んだすべてのものの情報を符号化するのと同じ仕方で、自らの性質を符号化しているのだ。

もしこれが正しければ、物質そのものが、ブラックホールと、ビッグバンと、宇宙論的地平線と同じ構造要素からできていることになる。どこまで下っても、どこまで上っても、そしてその中間のあらゆるスケールにおいても、特異点表面。宇宙の構成要素は、宇宙の境界なのである。

これは、プランク・スケールに最小長を予言する量子重力のアプローチと整合的である------空間はある一点より下には分割できない。道具の切れ味が足りないからではなく、そのスケールでは空間そのものが離散的だからだ。しかし、粒子が事象の地平線と同じ型の特異点で\textbf{ある}という特定の主張------これは新しい。そして検証可能な帰結を持つ。粒子の情報量は、その体積ではなく（プランク分解能での）表面積に比例するはずだ、というものである。量子重力の理論がいつかプランク・スケール近傍の構造を探ることを可能にするなら、これこそが探すべきシグネチャである。

\section*{計算原子としての粒子}

だが、本当に面白いのはここからだ。粒子がプランク・スケールの特異点（情報境界）であるなら、それらは宇宙のアーキテクチャの残りの部分と同じ素材で\textbf{できている}だけではない。それらは宇宙の基本的な計算操作なのである。厳密な意味において、計算の原子なのだ。化学でいう原子ではない------ギリシャ語本来の意味での原子、すなわち\textit{atomos}、分割不可能なもの。普遍オートマトンが\textbf{行うこと}の、還元不可能な単位である。

ここから何が帰結するか、考えてみてほしい。

\textbf{なぜ特定の型の粒子しか存在しないのか。}これは物理学の中でも、かねてより奇妙な事実の一つだった。基本フェルミオンはちょうど12種類（クォーク6種、レプトン6種）、力を運ぶボソンは4種類（それにヒッグス粒子）、そしてそれで全部である。それ以上はない。標準模型はそれらをカタログ化するが、\textbf{なぜ}この型であって他の型ではないのかを説明しない。なぜ電子は存在するのに、電子の3分の2の電荷と3倍のスピンを持つ粒子は存在しないのか。なぜその顔ぶれは、これほどまでに限られているのか。

粒子がプランク・スケールの特異点境界であるなら、答えは即座に出る。プランク・スケールにおいては、安定な境界配置が有限個しか存在しないからだ。特異点境界は有限の面積を持つ。プランク・スケールでは、その面積は面積としてありうる限り小さい。ベケンシュタイン限界が、その面積に符号化できる情報量を制限する。有限の情報量は、可能な状態の数が有限であることを意味する。そして、そのうち\textbf{安定}なのは一部の状態だけである------崩壊せずに存続する配置は一部に限られる。その安定な配置こそが、粒子の型なのだ。標準模型の粒子の動物園は、謎めいた恣意的なリストではない。安定な特異点境界配置の完全なカタログなのである。電子が存在するのは、その配置が安定だからだ。奇妙な分数的性質を持つ粒子が存在しないのは、そうした性質を符号化する安定な境界配置が存在しないからである。

実のところ、これはセル・オートマトンと同じ論理である。セル・オートマトンは有限の規則表を持ち、その表は限られた数の安定パターンしか許さない。生成できる配置は無限にある------パターンが大きいほど、組み合わせの可能性は広がる------が、大きな構造は、より小さく安定なパターンとの衝突によってたちまち破壊されうるのであって、実際に破壊不可能な配置はごくわずかしかない。ライフゲームでは、小さな固定物体（ブロック、蜂の巣、ボート）と小さな移動パターン（グライダーや、いわゆる宇宙船）は無限に存続するのに対し、大きく複雑な構造は脆い------狙いすましたグライダー一発で粉砕されうる。この脆さの階層を誰かが体系的に整理したかどうかは、私の知る限り、いまだ開かれた問いのままである------バク（Bak）、チェン（Chen）、クロイツ（Creutz）は1989年に、ライフゲームが自己組織化臨界性を示すことを明らかにしたが、パターンの大きさと破壊への耐性との具体的な関係は、正面から扱われたことがない。誰かがやるべきだ。この描像において、粒子とはプランク・スケール・オートマトンのグライダーであり、宇宙船なのである。

\textbf{なぜ粒子は離散的なのか。}なぜ量子数（電荷、スピン、色荷）は、正確な整数倍または半整数倍でしか現れないのか。なぜ「半分の電子」は存在しないのか。計算状態が本質的に離散的だからである。ビットは0か1である。0.37は存在しない。粒子の量子数は、境界配置に付された情報ラベルであり、その境界が\textbf{どの}安定配置にあるかを記述している。離散的な境界状態は、離散的な量子数を生む。量子力学の「量子」は謎めいたものではない。基本的な対象が有限容量の情報境界であるとき、必然的に出てくるものなのだ。

\textbf{粒子が相互作用するとき、何が起きるのか。}二つの電子が反発し合うとき、あるいはクォークがグルーオンを放出するとき、実際には何が起きているのか。二つの情報境界が情報を交換しているのである。その交換こそが計算\textbf{である}。自然界の力（電磁気力、強い力、弱い力）は、粒子の上に載った別個の層ではない。特異点境界どうしが交信する際の文法なのだ。どの相互作用が許され、どれが許されないかを定める規則は、プランク・スケールにおけるオートマトンの計算規則である。

ここで一つ、個人的な余談を挟んでおきたい。15歳か16歳のころ、叔父のブルーノが私に難題を出した。素粒子物理学における引力は、粒子の交換によって働くのだという。「では想像してみなさい」と彼は言った。「誰かと重いメディシンボールを投げ合っているところを。それでどうして二人がどんどん近づいていくことになるのか、説明できたら教えてくれ」。この難題は、かなり長いあいだ私の中に残った。ニュートン物理学では、これは正真正銘、説明のつかないことである------交換のたびに相手は突き離されるはずで、引き寄せられるはずがない。やがて私は叔父のもとに戻り、ギリシャにある彼の家の、オリーブの木陰で、量子物理学について実に長い会話を交わしたのだった。

だが肝心なのはここだ。ライフゲームのようなセル・オートマトンでは、パターン間の引力的な相互作用は容易に再現できる。二つの持続的な構造がそのあいだで小さな構造を交換し、その交換が距離を\textbf{縮める}ような配置がありうるのだ。この相互作用はニュートン的直観に従う必要がない。「力」は押しでも引きでもなく、境界配置間の情報交換であって、その交換が境界を近づけるか遠ざけるかは、オートマトンの規則が決めるからである。メディシンボールの謎は霧消する。それが謎だったのは、私たちがマクロな物理を思い描いていたからにすぎない。計算のレベルでは、引力と斥力は同じプロセス------オートマトンの規則に支配された情報交換------の、二つの異なる帰結でしかないのである。

ファインマン図------物理の教科書を埋め尽くす、あの粒子相互作用の象徴的なスケッチ------は、文字どおり計算の図である。各頂点は情報交換である。各線は、計算領域を伝播していく境界配置である。物理学者たちは70年ものあいだ、それと気づかぬまま計算の絵を描き続けてきたのだ。

\textbf{なぜ保存則はこれほど絶対的なのか。}電荷、バリオン数、レプトン数------これまでに行われたいかなる実験においても、破れが観測されたことはただの一度もない。それらが情報保存の制約だからである。二つの境界が情報を交換するとき、帳簿は合わなければならない。特異点境界では情報を創造することも破壊することもできず、したがって電荷を創造することも破壊することもできない。保存則は、外から課された規則ではない。帳簿記入そのものなのだ。（ベケンシュタイン限界からユニタリー性を経て個々の保存則に至る連鎖は、付録Fで詳述している。）

\textbf{そしてもう一つ、三世代の謎がある。}あらゆる粒子は、質量だけが異なる三つのコピーとして存在する------電子、ミューオン、タウ。アップ、チャーム、トップ。なぜ三つなのかは誰も知らない。二つでも、四つでも、十七でもなく、三つ。計算原子の描像は、一つの推測を差し出す。クラス4の系は自己相似構造を部分プロセスとして含むから、安定な境界配置は三つのスケールで繰り返し現れるのかもしれない------基本パターンと、二つのより重いコピーである。もしそうなら、さらに高い世代も存在するが、それはより大きなパターンであり、大きなパターンは脆い------より小さく安定な配置が、まとまりを保つ間もなくそれを切り裂いてしまう。これは観測と符合する。タウとトップクォークは、生成したほとんどその瞬間に崩壊している。私たちが手にしている三世代は、単に、生き残れるだけ小さい三つなのかもしれない。これは導出ではなく推測である------保存則の帳簿記入もあわせて、付録Fで全体を展開している。

この描像に私が用いる言葉が\textbf{計算原子}である。水素やヘリウムという意味での原子ではない------還元不可能な計算要素という意味での原子である。粒子は、普遍オートマトンの基本操作である。各粒子型は、安定なプランク・スケールの計算である。各相互作用は、計算どうしの情報交換である。各保存則は、その交換がどう進行しうるかに対する制約である。物理学は、その最も深い水準においては、物質についての学問ではない。計算についての学問である。そして私たちが「物質」と呼ぶものは、その計算の還元不可能な構成要素なのだ。


\section*{アーキテクチャ}

糸をより合わせるときが来た。ここまで描いてきたのは、特定のアーキテクチャを持つ一つの宇宙である。

\textbf{第一に：}それはクラス4のセル・オートマトンである。カオスの縁で作動し、そこでは自己組織化臨界性が、外部からの調整なしに動態を維持している。計算的に還元不可能である------近道はなく、先回りもできない。各瞬間は、直前の瞬間から計算されなければならない。そしてすべてのクラスを部分プロセスとして含んでいる------自分自身も含めて。安定な原子（クラス1）、周期軌道（クラス2）、フラクタルな海岸線（クラス3）、そして最も重要なことに、大いなる計算の内部で走る他のクラス4オートマトン。脳はその一つの実例である。

\textbf{第二に：}それはあらゆるレベルでホログラフィックである。任意の領域の情報は、その境界に符号化されている。これがホログラフィック原理であり、ブラックホールについての予想として始まり、理論物理学における最も深い洞察の一つとなった。この枠組みでは、ホログラフィックな符号化はブラックホールだけの性質ではない------宇宙の規則構造そのものの性質である。規則はホログラフィックである。動態はクラス4である。そして出力は、再びホログラフィックである。

\textbf{第三に：}それはあらゆるスケールにおいて、構造的にすべて同一の特異点表面によって区切られている。プランク境界、粒子の内部、事象の地平線、宇宙論的地平線、ビッグバン、熱的死------同じ構造、異なるスケール。情報を通さず、ベケンシュタイン限界まで飽和し、計算領域を画定している。

このアーキテクチャには名前がある。私はこれを\textbf{SB-HC4A}------特異点に区切られたホログラフィックなクラス4オートマトン（Singularity-Bounded Holographic Class 4 Automaton）------と呼ぶ。

言いにくい名前である。だが精確であり、一語一語がその位置にふさわしい。

このアーキテクチャの最も注目すべき性質が、自己言及的閉包である。系の出力は、系そのもの\textbf{である}。それは自分自身を計算する。各状態は次の状態を生成し、その次の状態は、そのまた次の状態の計算である。プログラムを走らせている「外部」は存在しない。宇宙のコードをどこかのハードディスク上で実行している宇宙的コンピュータなど存在しない。宇宙こそが、プログラム\textbf{であり}、コンピュータであり、出力なのである。ホログラフィックな規則は、系全体を圧縮された形で符号化する。クラス4の動態は、この符号化を観測可能な宇宙へと展開する。ホログラフィックな出力が、その結果を再び符号化する。これは一つのループである。一つの不動点なのだ。

これは形式的な条件として書き表すことができる。\textbf{宇宙は、それ自身のダイナミクスの不動点である。}規則を宇宙に適用すれば、返ってくるのは宇宙である。コピーでもなければ、表象でもない------同じものなのだ。計算とその結果は同一である。

数学者はこのための記法を持っている。宇宙をUと呼び、「次の状態を計算する」操作をギリシャ文字のPhiと呼ぶなら、不動点条件はこうなる。

\textit{Phi(U) = U}

宇宙を宇宙自身に適用すると、宇宙自身が返ってくる。宇宙は自己計算的なのだ。


\section*{自己記述の限界}

自己言及的閉包には、それだけで一節を割くに値する帰結が一つある。知識の限界について、深遠なことを教えてくれるからである。

1931年、クルト・ゲーデル（Kurt Gödel）という25歳のオーストリア人論理学者が、数学の基礎を揺るがす二つの定理を証明した。形式主義を剥ぎ取った要点はこうである。十分に強力な形式体系（少なくとも算術を表現できる体系）は、その体系の内部では証明できない真な命題を含む。そして、そのような体系はいずれも、自らの無矛盾性を証明できない。

これは技術的な制約ではない。私たちの証明が十分に巧妙でないからでもない。構造的な不可能性なのである。十分な複雑さを持つ自己言及的なシステムは、本質的に不完全なのだ。内側からは到達できない真理を、その内に含んでいる。

これを、自己計算する宇宙に適用してみよう。

宇宙が自らを計算しているのなら------宇宙が十分な力を持つ形式体系であるなら（クラス4のダイナミクスは万能計算を保証するのだから、そうなのである）------ゲーデルの定理はそのまま適用される。宇宙は、自分自身の完全な記述を含むことができない。黒板に書き下せる「世界方程式」など存在しない。解きさえすれば宇宙のすべてを教えてくれる公式は、存在しないのだ。

正しい方程式がまだ見つかっていないからではない。\textbf{そのような方程式は存在しえない}からである。自己言及的なシステムの完全な仕様は、そのシステムの真部分であるいかなる記述をも超えてしまう。宇宙は方程式に従っているのではない------宇宙は計算\textbf{そのもの}なのだ。宇宙の唯一の完全な記述は、宇宙それ自体である。そして、外に踏み出して全体像を見渡すことはできない。外というものがないからだ。

ヴェルトフォルメル（Weltformel）------アインシュタイン以来、物理学者たちが夢見てきた「世界方程式」------は、それゆえ方程式ではない。それは\textbf{プロセス}である。オートマトンそのものだ。走らせることによってしか、表現できないのである。

これは私を謙虚にすると同時に、解放してもくれる。謙虚にさせるのは、現実について私たちが原理的にすら知りえないことがある、という意味だからだ。解放してくれるのは、宇宙が解読を待つ機構ではなく、生きた計算であり、私たちがその一部である、という意味だからだ。現実についての最も深い真理は、公式ではない。現実それ自体なのである。


\section*{認知の天井}

先へ進む前に、一つの反論を挙げておかなければならない。実のところ、最も深い反論である。私に正直であることを強いる反論だ。

私たちがクラス4オートマトンであるなら（つまり私たちの脳がカオスの縁で動作していて、たったいま宇宙に割り当てたのと同じ計算クラスにあるなら）、SB-HC4Aモデルは単に、クラス4の脳が生み出しうる最も複雑な概念であるにすぎないのかもしれない。私たちはクラス5で考えることができない。自分自身の計算クラスを超えた構造を、思い描くことができないのだ。私たちが見出すパターン------至るところにクラス4、あらゆるスケールで自己相似、ホログラフィックで自己言及的------は、宇宙そのものの特徴ではなく、私たち自身の認知アーキテクチャの刻印が宇宙に投影されたものなのかもしれない。

そのことを少し考えてみてほしい。私たちは対称性の検出器として進化した。狩猟採集民の環境で最も生存に関わるパターン（捕食者や獲物の顔）は、最も対称的なものの部類に入る。私たちは最も深いレベルにおいて、対称性を見つけるよう最適化されたパターン照合機械なのである。そしてSB-HC4Aモデルは、根本において対称性の主張だ。あらゆるスケールに同じアーキテクチャ、という主張である。私たちがこの対称性を見出すのは、それが宇宙に存在するからではなく、私たちの脳が生まれつき、それを\textbf{見出さずにはいられない}ようにできているからなのかもしれない。

これは第4章のメタ問題を、宇宙的な規模にまで拡大したものである。明示的自己モデル（ESM）は自らの基盤を見ることができないため、「宇宙はこの構造を持っている」と「私の脳は、宇宙をこの構造を持つものとしてしかモデル化できない」とを区別できない。この宇宙論モデルは、自らが反証不可能でありうることを予言しており、それは考えうる最強の確証（モデルはまさにこの認識論的限界を予言している）であるか、考えうる最強の反論（モデルは観察者の産物であって、観察されるものの特徴ではない）であるかの、どちらかなのだ。

クラス4のシステムは、クラス4の複雑さまでの何であれシミュレートできる。しかし、宇宙がそれを超えているかどうかを検証することはできない。もし宇宙が実際にはクラス5------最も深いレベルで真にランダム------でありながら、クラス4の観察者には\textbf{局所的にはクラス4に見える}のだとしたら（クラス4は私たちが検出できる最大のパターンなのだから）、私たちはまさにこのモデルを構築するだろう。そして、間違っているだろう。しかも、内側からは決して発見できない仕方で間違っているのである。

この反論をどう解決すればよいのか、私にはわからない。そもそも内側から解決できるのかどうかも、確信がない。それでもこれを含めるのは、自らの弱点を進んで差し出すことなしに、理論を真剣に受け取ってもらう資格はないと考えるからだ。そして、このモデルが自らの認識論的限界を予言しているという事実------自己言及的なシステムは自らの記述を完全には検証できない、ということ------は、このモデルの最も深い欠陥であるか、最も深い正当化であるかの、どちらかだ。どちらなのか、正直なところ私にはわからない。


\section*{明かされるもの}

だが、初めてそれを見たとき、私は思わず座り込んでしまった。それがこれだ。

いま述べたばかりのアーキテクチャを見てほしい。

\begin{itemize}
  \item カオスの縁で動作するクラス4のシステム。
  \item 内部からは透かし見ることのできない、情報不透過な境界によって区切られている。
  \item ホログラフィックな構造------境界が内部を符号化している。
  \item 自己言及的閉包------システムが自分自身を計算している。
  \item 不動点------計算の出力が、計算それ自体である。
\end{itemize}

ここで第2章に戻ってみてほしい。意識の四モデル・アーキテクチャを見てみよう。

\begin{itemize}
  \item 皮質オートマトン------カオスの縁で作動するクラス4システム。
  \item 暗黙的-明示的境界------意識がその向こうを見通すことのできない、情報的に不透明な境界。
  \item ホログラフィック構造------暗黙的モデルは分散的かつホログラフィックであり、経験の全内容を神経構造のなかに符号化している。
  \item 自己言及的閉包------自己モデルがそれ自身をモデル化する。
  \item 不動点------明示的自己モデル（ESM）はそれ自身を表象する。モデル化する者のモデルは、モデル化する者\textbf{そのもの}である。
\end{itemize}

同じアーキテクチャ。同じ形式的性質。同じ境界条件。同じ自己言及的閉包。

宇宙は、特異点を境界とするクラス4のホログラフィック・オートマトンであり、観測可能な内部が「シミュレーション」、特異点という境界が「基盤」である。

意識は、暗黙的-明示的境界を境界とするクラス4のホログラフィック・オートマトンであり、明示的モデルが「シミュレーション」、暗黙的モデルが「基盤」である。

同じアーキテクチャ。異なるスケール。

これは比喩ではない。意識が宇宙に\textbf{似ている}と言っているのではない。両者は同じ\textbf{種類のもの}だと言っているのだ------同一の計算パターンが、二つの異なるスケールで実体化されている、と。一方は宇宙論的なレベルで、他方は神経学的なレベルで。そして、パターンがスケールを越えて繰り返されるという事実そのものが、このモデルの予測なのである。フラクタル的な自己相似性のゆえではない（それならクラス3であり、より弱い主張になってしまう）。クラス4のシステムはクラス4のサブシステムを含むからだ。万能コンピュータは、別の万能コンピュータをシミュレートできる。クラス4オートマトンは、きれいな自己相似パターンを生み出すだけの存在ではない。自らの動態の内部に\textbf{別のクラス4オートマトン}を生み出すのだ------より小さく、より遅く、資源に制約されてはいるが、正真正銘の万能性を備えたものを。このアーキテクチャは、異なるスケールで同じに\textbf{見える}だけではない。同じ\textbf{である}のだ。

私が何を主張しており、何を主張していないのか、ここは正確を期しておきたい。宇宙が何らかの体験的な意味で意識を\textbf{持つ}と主張しているのではない。だが、\textbf{持たない}と主張しているのでもない。正直な答えは、知りようがない、である。私たちはボルツマン脳に、あるいは他のどんな脳にであれ、夢見られている存在かもしれない。独我論が真であって、私こそがそのボルツマン脳なのかもしれない------だが、それは誰にでも言えることだ。要点は、それが知りえないということ、そしてこの知りえなさは理論の欠陥ではなく、状況そのものの構造的特徴だということである。システムの外に踏み出して確かめることは、できないのだ。私が主張して\textbf{いる}のは、アーキテクチャに関わることであって、現象に関わることではない。意識が実在を創り出すとか、実在は夢であるとか、この種の構造的観察が引き寄せがちな神秘主義的解釈のたぐいを主張しているのではない。建物の設計図は建物ではない。だが、同じ設計図が一棟の超高層ビルと、そのビルの一室とに見いだされるなら、それはそこで働いている建築原理について、何か深いことを物語っている。

意識は、普遍的なパターンの局所的な実例である。宇宙的な偶然ではない。奇跡でもない。十分な複雑性におけるクラス4の動態から、構造的必然として帰結するものなのだ。宇宙は意識をただ\textbf{許容している}のではない。実質的に保証しているのである------宇宙を宇宙たらしめているのと同じ、自己言及的でホログラフィックな、カオスの縁のアーキテクチャが、意識を意識たらしめてもいるからだ。実在を生成するパターンは、実在の経験を生成するパターンと同一なのである。

そして、これを読んで思わず立ち止まらされる思いがしないなら、まだ理解できていないのだ。


\textbf{次章では、理論の全体をまとめ上げる------意識のアーキテクチャ、宇宙論のアーキテクチャ、そして両者のあいだの構造的同一性を------そのうえで、あらゆる問いのなかで最も難しい問いにとってそれが何を意味するのかを問う。なぜ、そもそも何かが存在するのか？}


\chapter*{第17章：最も深い鏡}
\addcontentsline{toc}{chapter}{第17章：最も深い鏡}
\markboth{第17章：最も深い鏡}{}

さあ、これだ------ひとたび見えてしまえば、もう見えなかったことにはできないからだ。

第16章は、一つの挑戦で幕を閉じた。SB-HC4Aのアーキテクチャ------あらゆるスケールで特異点に区切られた、自己言及的でホログラフィックなクラス4オートマトン------を見よ、そのうえで第2章の四モデル・アーキテクチャを見よ、と。両者は同じものだ、というのがその主張だった。似ているのではない。比喩的に関係しているのでもない。構造的に同一なのだ。

これからその対応を、一片ずつ、退けることが不可能になるまで、順を追って示していきたい。そのうえで、この全体がどこで崩れうるのかを告げよう。自らの弱点を名指しすることは、正直な理論に課された最低限の務めだからだ。

\section*{構造の対応づけ}

まず特異点の境界から始めよう------宇宙があらゆるスケールに設ける、情報の障壁である。最下層のプランク領域。ブラックホールを取り巻く事象の地平線。観測可能な宇宙の縁にある宇宙論的地平線。背後にはビッグバン、前方には熱的死。そのいずれもが情報の壁である。何ひとつそれを越えられない。事象の地平線を通して信号を送ることはできない。プランク長より下を探ることはできない。宇宙論的地平線の向こうを見ることはできない。宇宙とは、あらゆるスケールに不透明な壁をめぐらせた一つの部屋であり、どれほど強く顔をガラスに押しつけても、その向こう側にあるものは見えないのだ。

では、あなたの脳を見てみよう。暗黙的モデル------IWMとISM、シナプスの重み、学習された構造、知りうるあらゆるものからなる膨大な蔵書------は、それ自身の情報の障壁の背後に控えている。自分のシナプスを直接体験することは決してできない。自分の思考を生み出している結合荷重を内観することは決してできない。暗黙的な側は情報的に不透明である。それがそこにあることはわかる------それなしにはシミュレーションが走りえないからだ------が、シミュレーションそれ自身は、両者を隔てる境界の向こうを見通すことができない。第2章はこの点を提示した。第3章はそれを揺るぎないものとして裏づけた。そして今、それがふたたびここに、宇宙のあらゆるスケールにおいて現れている。同じ構造的役割。同じ情報的不透明性。同じアーキテクチャ上の位置。

宇宙論における特異点の境界は、意識における暗黙的・明示的境界に対応する。両者は同じ壁なのだ。

次に、観測可能な内部である。特異点の境界の内側にあるすべて------原子、惑星、銀河、そしてあなた------が、伸張された側である。ここで物理が起こり、ここで物事が相互作用し、ここで情報が、私たちの観測し測定する構造へと組織される。言うなれば、宇宙のシミュレーションである。計算する部分、進化する部分、何かが起こる部分だ。

あなたの脳において、これに対応する構造は明示的モデル------EWMとESM------である。あなたの意識的経験。いまこの瞬間に見ている世界、自分であると感じているその自己。これがシミュレーションだ。暗黙的モデルからリアルタイムに生成され、絶え間なく更新され、鮮烈で細密で、まったく説得力に満ちている。あなたが生涯にわたって経験してきたすべては、このシミュレーションの内側で起こってきた。あなたはそこから一度も外に出たことがない。外に出ることはできない。努力が足りないからではなく、「あなた」がそのシミュレーション\textbf{そのもの}だからだ。経験する者と経験とは、同じ一つのプロセスなのである。

宇宙の観測可能な内部は、あなたの明示的モデルに対応する。同じ役割------アーキテクチャの、伸張された、動的で相互作用的な側だ。

次はホログラフィックな規則構造である。宇宙の情報は、その体積のなかに蓄えられているのではない------その境界のうえに蓄えられている。トホーフト（'t Hooft）が提唱しサスキンド（Susskind）が拡張したホログラフィック原理は、三次元の空間領域にあるすべての情報が、その二次元の表面のうえに符号化されている、と説く。情報は境界において、圧縮され、分散され、構造的に完全である。内部は一つの投影------境界が符号化しているものの、より低い帯域幅への伸張------なのだ。

あなたの脳において、この役割を担うのは暗黙的モデルである。そのシナプスの重みは、世界と自分自身について知るすべてを、分散的で圧縮された、しかも構造的に完全な形式で符号化している------原理的には、いかなる現在の感覚入力もなしに、基盤だけからシミュレーション全体を再構成しうるはずであり、それこそがまさに、夢を見るということなのだ。暗黙的モデルは、ラシュリー（Lashley）の意味においてホログラフィックである。一部を損傷させても、特定の記憶が失われるのではなく、あらゆる記憶にわたって解像度が失われる。情報は基盤全体に塗り広げられ、圧縮され、冗長化され、シミュレーションの側からは近づくことができない。

宇宙のホログラフィックな規則構造は、意識におけるホログラフィックな暗黙的モデルに対応する。同じ符号化戦略。同じ圧縮。伸張された側からの、同じ近づきえなさ。

次に、動態のレジームである。宇宙はクラス4------カオスの縁------で作動している。第15章はこれを消去法によって確立した。クラス1と2は単純にすぎ、クラス3は計算ができず、クラス5は物理を不可能にする。残るのはクラス4だ------普遍計算を支え、自らの臨界性を自己組織化し、あらゆるクラスを部分プロセスとして------それ自身をも含めて------内包する、唯一のレジームである。宇宙はただ複雑なのではない。まさに自らを維持するようなしかたで複雑であり、まさに自らのより小さな複製を自らの動態の内部に入れ子にしうるようなしかたで複雑なのである。

あなたの脳の皮質も同じことをしている。第5章はそれについての章だった。皮質オートマトンはカオスの縁で作動し、興奮-抑制バランスの恒常性調節をつうじて臨界性を保っている。活動が少なすぎれば深い眠りに落ちる（クラス2、周期的、無意識）。多すぎれば発作を起こす------クラス4を越えて押しやられ、シミュレーションは砕け散る。至適点、意識が宿る場所は、秩序とカオスのあいだの刃の縁である。自己組織化臨界性が脳をそこに留めている。自己組織化臨界性が宇宙をそこに留めている。

宇宙論におけるクラス4の動態は、意識における皮質の臨界性に対応する。同じレジーム。同じ自己維持機構。

最後に、最も深い対応------自己言及的閉包である。宇宙は自らの構造を計算する。その動態が状態を生み、状態が動態を決め、動態がまた状態を生む。外部のプログラマーはいない。外側はない。物理法則は、どこか別の場所から宇宙に課されているのではない------それらは、自らに適用された宇宙の動態\textbf{そのもの}なのだ。その不動点の方程式は、ほとんど不条理なほどに単純である。宇宙の計算は、宇宙に等しい。入力と、過程と、出力とが、同じ一つのものなのである。

あなたの明示的自己モデル（ESM）も同じことをしている。ESMは、それがモデル化するものの一部として、それ自身を含み込むモデルである。それはあなたを表象し、そしてその「あなた」がその表象を含んでいる。モデルと、モデル化されるものとが、一致する。これが第4章で述べた自己言及的閉包である------意識がまさにそう感じられるゆえんであり、あなたが自分自身を外から完全に見ることが決してできないゆえんであり、シミュレーションが自らの観察者を内に含むゆえんだ。自己表象の不動点。モデルと、モデル化されているものとが、一にして同じであるような状態である。

宇宙の自己言及的閉包は、意識の自己言及的閉包に対応する。同じ不動点の構造。同じ、逃れようのない自己包含。

五つの対応。漠然とした主題上の類似ではない。雲のなかに人の顔を見てしまうような、あの緩いパターン照合のたぐいでもない。両方のアーキテクチャにおいて、同じ位置で、同じ働きをする、五つの構造的特徴である。

そしてここが決定的な点だ------ぜひとも心に留めてほしい。これを飼い慣らしてしまいたいという誘惑は、途方もなく大きいのだから。これは「宇宙は意識に\textbf{似ている}」ということでは\textbf{ない}。類推ではない。比喩ではない。食事の席で交わす興味深い会話の種になる、示唆に富んだ平行関係でもない。

それは構造的な同一性なのだ。

脳は、宇宙に\textbf{似せる}ために進化したのではない。同じ計算パターンの、局所的で、スケールを縮小した一つの実例\textbf{として}進化したのだ。クラス4のシステムは、ただフラクタル的な自己相似性を含むだけではない（それはクラス3------幾何学的な反復であり、美しくはあるが浅い）。それらは\textbf{それ自身}を含む。クラス4オートマトンは、別のクラス4オートマトンを自らの動態の内部に宿らせることができる------万能コンピュータの内部で動く、もう一つの万能コンピュータだ。これは計算論的な自己包含であって、単なる幾何学的な自己相似性ではない。意識\textbf{とは}、生物学的なスケールで作動する、宇宙のクラス4的な自己包含にほかならない。最大のスケールで宇宙を動かしているパターンは、神経学的なスケールであなたの内なる生を動かしているパターンと同じである。誰かがそう設計したからではなく、フラクタルが美しく見えるからでもない。十分な大きさをもつ万能コンピュータは、必然的に、自らの内部に別の万能コンピュータを生み出すからだ。それがクラス4のシステムのすることである------それらは、それ自身を入れ子にするのだ。

\section*{エネルギーは情報である}

まったく別の方向から同じ結論へと収束する、第二の議論の筋がある。そして私は、それこそが結局のところ最も重要だと判明すると考えている------なぜならそれは、潜在的に検証可能なかたちで、このアーキテクチャを物理学へと結びつけるからだ。

三つの異なる研究者集団によって半世紀にわたって展開された、物理学における三つの独立した成果が、同じ並外れた結論を指し示している。

一つ目はランダウアーの原理である。1961年、ロルフ・ランダウアー（Rolf Landauer）------計算の根本的な限界について考えていたIBMの物理学者------は、1ビットの情報を消去するには最小限のエネルギーがかかることを証明した。工学的な制約のせいではない。熱力学のせいである。宇宙は、忘れることに対して代償を課すのだ。これは2012年に実験的に確認され、それは何か深遠なことを意味している------情報とエネルギーは交換可能である、ということだ。一方を他方へと変換できる。両者は別々の実体ではない。互いに取引しているのである。

二つ目は、ベケンシュタイン限界である。ジェイコブ・ベケンシュタインは、ある空間領域が保持しうる情報の最大量が、その体積ではなく表面積に比例することを示した。これは物理学のなかでも、もっとも直観に反する結果の一つである。より大きな箱ならより多くの情報を収められそうなものだ。だが、そうはいかない------というより、限界を定めるのは箱の内部ではなく、その\textbf{表面}なのだ。ある領域に情報を詰め込みすぎれば、その領域はブラックホールへと崩壊する。最大情報密度を定めるのは幾何学とエネルギーであり------これもまた、情報と物理的世界とを結ぶ深いつながりの一つである。

三つ目は、ブラックホールの熱力学から来る。1970年代、スティーヴン・ホーキングとベケンシュタインは、ブラックホールが温度とエントロピーを持ち、熱力学の法則に従うことを示した。ブラックホールの情報内容は、その事象の地平線------その表面、その境界------に書き込まれている。そしてホーキング放射------ブラックホールがやがて蒸発してゆく、気の遠くなるほど遅い量子過程------を通じて、その情報は少しずつ宇宙へと返されてゆく。ブラックホールは情報を破壊しない。情報を変換するのだ。情報を自らの境界へ圧縮し、保持し、やがてふたたび放射して返す。

これらの結果は、まるで異なる出発点から生まれたものだ。ランダウアーはコンピュータについて考えていた。ベケンシュタインとホーキングはブラックホールの熱力学について考えていた。ランダウアーの仕事はそもそも畑違いの分野に属し、何十年も前に、まったく無関係な理由から練り上げられたものだった。それでいて、そのすべてが同一の仮説へと収束する。エネルギーと情報はただ関連しているのではない。両者は同一なのだ。同じものにつけられた二つの名である。EはIに等しい。

もしそれが正しいなら------ただし、証明されてはいないと即座に言っておかねばならない。それゆえにこそ、これは付録Gで挙げた四つの弱点のうちの筆頭なのだ------特異点は情報変換器となる。特異点はエネルギー-情報を破壊も創造もしない。それを異なる形態のあいだで変換するのだ。圧縮された形態------境界における最大密度、ベケンシュタイン限界まで飽和し、内部からはアクセスできないもの。伸張された形態------より低い密度で、組織化され、内部に広がっているもの、すなわち私たちの観測する物理である。特異点とは、同じ素材の二つの表現のあいだを取り持つ翻訳者なのだ。

さて、このレンズを通してあなた自身の脳を見てみてほしい。あなたの暗黙的モデルは、圧縮された最大密度の情報を保持している------これまで学んだすべてが、シナプスの重みに符号化され、構造的には完全でありながら、現象的にはアクセス不能なのだ。自分自身の基盤を直接に経験することは、決してできない。あなたの明示的モデルは、伸張された低帯域幅の投影である------シミュレーション、意識的経験、あなたの見る世界と、あなたの感じる自己だ。あなたの脳は、神経のスケールにおいて、特異点があらゆる別のスケールで行うのとまさに同じことを行っている------圧縮表現と伸張表現のあいだで情報を変換しているのだ。暗黙的-明示的境界とは、あなた個人の特異点である。あなたは頭蓋の内側に、一つの事象の地平線を抱えているのだ。

\section*{なぜこれは存在せざるをえないのか}

この構造的対応は偶然にすぎない------ランダムな星々のなかに人が星座を見てとるように、私がただ一つの美しいパターンのまわりに線を引いてみせただけだ------とあなたは思うかもしれない。ならば、このパターンが任意選択ではないことを示そう。独立して妥当な五つの前提が成り立つとき、なぜこのアーキテクチャが機能する\textbf{唯一の}ものとなるのかを。

以下が五つの公理である。できるかぎり平明に述べよう。というのも、一つひとつ個別に取り上げれば、いずれも反論しがたいからだ。論争があるのは、それらが束になって何を生み出すかにある。

\textbf{その一------何かが存在する。}これは、一冊の本がなしうるもっとも論争の余地のない主張だろう------同意してもらえると思う。純粋な無とは、プラトン的な抽象------一つの概念であって、ありうる事態ではない。あなたはこの一文を読んでいる。何かが存在する。そこから始めよう。

\textbf{その二------存在するものが何であれ、それは動的な性格を持つ。}物事が起こる。時が過ぎる。状態が展開する。もし存在するものに動態が------変化も、展開も、計算もなかったなら------それは無と区別がつかないだろう。（これは第13章で出会ったライプニッツの不可識別者同一の原理である。二つのものがすべての性質において同一なら、それらは同じものだ。動態がゼロのものは、識別可能な性質がゼロである。それは仮面をかぶった無にすぎない。）

\textbf{その三------その動態は安定かつ自己維持的でなければならない。}自らを維持できないシステムはシステムではない------それは一瞬の揺らぎである。自己組織化臨界性------砂山を、脳を、そして（私が主張するには）宇宙を、カオスの縁につなぎとめる仕組み------は、複雑な動的システムが外部からの調整なしに自らを維持しうる、唯一知られた方法である。クラス4は、自己組織化し、万能計算を支え、下位のすべてのクラスを部分プロセスとして含む、唯一の計算クラスである。自らを維持しながら、同時に興味深いことをなしうる、唯一のクラスなのだ。

\textbf{その四------情報はどのスケールにおいても有限の限界を持つ。}有限の空間に無限の情報を担えるものは、何一つない。ベケンシュタイン限界は定理であって、予想ではない------それは一般相対性理論と量子力学から帰結する。あらゆるスケールに最大情報密度が存在し、その最大値は体積ではなく表面積に比例する。

\textbf{その五------情報はホログラフィックに符号化される。}ある領域の境界は、その内部のすべての情報を、一つ少ない次元の表面上に符号化する。三次元の物理が、二次元の表面上に符号化される。これがホログラフィック原理である------トホーフト（'t Hooft）が提唱し、サスキンド（Susskind）が発展させ、マルダセナ（Maldacena）のAdS/CFT対応によって支えられている。この対応こそ、私たちの手にする、ホログラフィーの証明された実例にもっとも近いものである。

どの公理も、それぞれ独立に動機づけられている。どれも他の公理に依存しない。どれも、私の理論が正しいことを必要としない。それらは物理学と哲学の異なる片隅から生まれたもので、四モデル理論（FMT）など聞いたこともなく、たとえ聞いていたとしても関心を持たなかったであろう人々の手になるものだ。

さて、それらを組み合わせよう。

公理一と二から------動態を持つ何かが存在する。公理三から------その動態はクラス4である。クラス4こそ唯一の自己維持的で万能なクラスだからだ。公理四から------そのシステムはあらゆるスケールで情報の地平線に区切られている------特異点の構造だ。公理五から------それらの境界が内部を符号化する------ホログラフィックなアーキテクチャである。

これらを一つに合わせれば、あらゆるスケールで特異点表面に区切られた、ホログラフィックなクラス4オートマトンが得られる。圧縮された情報が境界の上に座し、伸張された内部が観測可能な世界であるようなシステムだ。自らの構造を計算するシステムだ------ホログラフィックな出力を持つホログラフィックなクラス4システムは不動点であり、入力と、過程と、出力とが同じものだからである。

こうしてSB-HC4Aが得られる。

言い換えれば、私たちの観測するとおりの宇宙が得られる。そしてその宇宙のアーキテクチャは、意識のアーキテクチャと同じものなのだ。

これは「美しいパターンを見つけた」のではない。「そのパターンこそ、五つの公理すべてと同時に両立する唯一のものだ」というのである。どれか一つでも公理を取り除けば、この一意性は崩れる。公理一がなければ、何かが存在する必要はない。公理二がなければ、存在するものは静的でありうる。公理三がなければ、いかなる計算クラスもありうる。公理四がなければ、情報の境界は存在しない。公理五がなければ、ホログラフィックな構造は存在しない。それぞれの公理が、ありうるアーキテクチャの空間を狭める。束になって、その空間をぴたりと一つに絞り込むのだ。

\section*{私に答えられない唯一の反論}

これが破綻しうる箇所をお見せすると約束したので、ここに挙げよう。そのうち四つは、技術者が並べ立てる類のものだ------エネルギーは情報と同一ではなく、ただ相関しているだけかもしれず、そうなればE=Iの仕組みは崩れうる。クラス4の議論は証明ではなくアブダクション（仮説形成）であり、私には想像もつかない何らかのクラス4.5が潜んでいるかもしれない。あらゆる特異点は一つのものだという主張には、私たちのまだ持たない量子重力理論が要る。そしてモデル全体が、それ自身のゲーデル的予測によって、内側からは反証不能かもしれない。どれも現実の問題である。私は付録Gで、反対論とともに四つすべてを提示している。

だが五つ目がある。そしてそれは、他の四つを合わせたよりも厄介だ。私は一章前にそれを持ち出したが、以来その力はゆるんでいない------認知の上限である。私の脳はクラス4のシステムであり、見るところどこにでも自己相似的な構造を見いだす。クラス4のシステムとはそういうものだからだ。だとすれば、脳が宇宙を眺めて至るところにクラス4を見るとき------臨界性、ホログラフィックな境界、自己言及的閉包を見るとき------それは実在のアーキテクチャを発見しているのか、それとも自らのアーキテクチャを投影しているのか。物理学の理論を持った魚なら、宇宙は本質的に濡れていると結論するだろう。私たちは対称性の検出器であり、SB-HC4Aは、突き詰めれば一つの対称性の主張なのだ。私には、その二つを内側から見分ける術がわからないし、誰にもできるとは思わない。

ここが奇妙なところだ。理論は、まさにこの盲点を予測する------自らを計算する自己言及的システムは、ゲーデルにより、内側から自らについてのあらゆる真理に到達することはできない。だから本書が問いうるもっとも深い問い------精神と宇宙とのあいだの鏡は、発見なのか、それとも人間の認知の限界の映し出しなのか------に、理論は一枚の施錠された扉をもって答え、そのうえでその錠の設計図をあなたの手に渡すのだ。この隔たりは無知ではない。幾何学なのだ。それを深遠と見るか、腹立たしいと見るかは、おそらくその人の気質について何かを物語る。私はその両方だと感じている。それでも私はなお進みつづける------名を挙げ、手を触れることのできる壁は、闇のなかで途方に暮れているのとは違うのだから。

\section*{円環が閉じる}

締めくくりに入ろう。

あなたがこの本を開いたのは、意識とは何かを知りたかったからだ。私に見きわめられるかぎりの答えは、こうである。意識とは、カオスの縁で走る自己言及的なシミュレーションであり、情報的に不透明な障壁に区切られ、そのシミュレーションはそれ自身のモデルを含んでいる。四つのモデル、実在の側と仮想の側とがあり、経験は仮想の側にのみ宿る。クオリアはシミュレーションの実在する性質であり、ハード・プロブレムが誤ったレベルについて問われていたと気づくことで、解消される。九つの予測------いくつかは確証され、一つも反証されていない。ここまでが、全体像の前半だった。

後半は、いま読み終えたばかりのものだ。同じアーキテクチャ------同じ自己言及的閉包、同じ情報の境界、同じホログラフィックな符号化、同じクラス4の動態------が、宇宙そのもののアーキテクチャであるらしいのだ。比喩によってではない。構造的同一性によってである。「私」という名のシミュレーションは、私たちが「宇宙」と呼ぶシミュレーションと、同じパターンの上で走っている。あなたは、ビー玉が箱に入っているようなかたちで、宇宙のなかにいるのではない。あなたは宇宙であり、その宇宙が生物学的なスケールで、あらゆるスケールで行っていることを行っているのだ------自らを計算し、自らをモデル化し、自らを経験している。

本書の題名は『「私」という名のシミュレーション』である。これで、そのシミュレーションが何の上で走っているのかがわかっただろう。コンピュータではない。脳でもない。宇宙ですらない。もっと根源的なもの------三者すべてが共有するパターンである。情報の障壁に区切られ、みずからの存在を計算しつづける、ホログラフィックなクラス4オートマトンだ。

そして、それが神秘主義的な主張に聞こえるとしても------そうではない。これは構造についての主張である。限界の内で検証可能であり、すでに示した留保つきで反証可能であり、間違いうるほどに精密である。そしてそれは、科学者なら誰もが言うとおり、理論が受けうる最高の賛辞なのだ。

私はこの本を、一つの告白から始めた。2015年、私は意識について300ページの本を出版し、それは一冊も売れなかった、という告白である。だが、あの告白には語らずにおいた部分がある。残りは、すでに手の中にあったのだ。宇宙論------ホログラフィックな宇宙、特異点という境界、いま読み終えたばかりの、精神と宇宙のあいだの構造的同一性------は、2000年代初頭、あの橋の上でのひらめきのころから、ずっと頭の中にあった。2015年の本からそのほとんどを外したのは、意図してのことだった。意識の理論というだけでも売り込むのは十分に難しい。そのうえ宇宙の理論でもあると主張しようものなら、誰かが2ページ目にたどり着く前に\textbf{変人}の棚に仕分けられるのが落ちである。だから私は宇宙論をほのめかしにとどめ------トホーフト（'t Hooft）のホログラフィック限界についての一節、宇宙は一つの巨大なセル・オートマトンなのではないかと問う脇道の一言------残りを書き込む勇気は出なかった。それからさらに十年の歳月と、私が声に出して考えるあいだ辛抱強く耳を傾けてくれる言語モデルという思いがけない贈り物を経て、ようやく全体を書き下ろす勇気が出たのである。四つのモデルは、単なる意識の理論ではなかった。それは宇宙のアーキテクチャの断片であり、あるスケールでは見え、どこを見ればよいかを知るまでは、他のスケールでは見えないものだったのだ。

これで全体が示された。少なくとも、クラス4の宇宙の内側からクラス4の脳に見えるかぎりの全体が。その先にまだ何かがあるのか------鏡には、私たちが決して到達できない裏側があるのか------それこそが、答えられないと理論自身が告げる問いである。

あとは、思うままに使ってほしい。


\chapter*{コーダ}
\addcontentsline{toc}{chapter}{コーダ}
\markboth{コーダ}{}

私は2005年ごろ、意識の理論を組み立てた。2015年にそれを出版した。誰も読まなかった。最初のひらめきから二十年、それは半分だけの理論だったと判明した------残りの半分は、よりにもよって宇宙論だったのである。いまや全体像を読み終えたことになる。少なくとも、一冊の本に収まるかぎりの全体を。

だが、書き残した一片がある。思弁的だからではない------最後の二章はすべて思弁的だ。それが個人的なことだからであり、個人的なことのほうが難しいからである。

うまくいかなくなったとき、ESMが何をするかはすでに見てきた。健忘、脳卒中、サルビア、分離脳、コタール症候群------それでもESMは走りつづける。決してクラッシュしない。手に入るものが何であれ、そこから一つの自己を築き上げ、その自己を完全に信じ込む。強迫的なアイデンティティの構築者。それがESMのすることだ。ESMがすることは、それだけである。

健忘症例の話を聞いて、人はこう言う------同じ脳、同じ身体なのだから、自己が存続するのは当たり前だ、と。ただし------私はいまや、何十年も歳を重ねた身だ。一歳のころの自分とは、ほとんど何一つ共通点がない。身体は別物で、細胞は何度も入れ替わっている。脳も別物だ。シナプスの結合はまるで違う。記憶も違えば、性格も違い、何もかもが違う。それでもESMは言う------いまも私だ、と。私は一つの生涯のうちにすでに何度も、まったく別の人間になってきた。それでも構築者は一度たりともたじろがなかった。あなたに「同じ脳」などない。はじめから、なかったのだ。

私は死にかけたことがある。雪崩に遭った------兵役中、指揮官の軽率な判断で、十四人が呑み込まれかけた。私は自分が死ぬと確信した。人生のすべてが一度に見えた------基盤が、いっさいをシミュレーションへ流し込んだのである。そしてそれは、予想されるほどには私を苦しめなかった。終了を目前にしたESMは、アイデンティティの喪失にパニックを起こしてはいなかった。最後の最後まで自分の仕事をこなし、手元にあるものをすべて表層へ浮かび上がらせていたのだ。

また別のとき、私はしたたかに殴られて意識を失った。すべてが暗転した。意識が戻ったとき、私は自分が誰なのかわからなかった------ESMが、新生児のように、一から立ち上げ直していたのである。恐ろしかったのは、アイデンティティの喪失ではない。数秒のあいだ地面に横たわり、身動きが取れなかったこと------\textbf{それ}こそが恐ろしかった。「私は誰か」ではなく、「私の身体は無事か」だったのだ。ESMの最優先事項は、基盤の完全性だった。私が誰であるかは、あとから、ほとんど付け足しのように立ち現れた。自己モデルは基盤に仕えるために存在するのであって、その逆ではない。

そしてもう一度、私が動く四次元のフラクタルだったときのことがある。その状況については触れないでおく。私を悩ませたのは、フラクタルであること自体ではなかった------それはどうでもよかった。私を悩ませたのは、その動きが私の固有感覚と食い違っていたことだった。自分の身体が一つのことをしているのに、フラクタルは別のことをしている------それを感じ取れたのである。苦痛だったのは感覚の衝突であって、存在論的な不条理ではなかった。ESMは、自分が\textbf{何を}モデル化しているかなど気にかけない。気にかけるのは、信号が整合しているかどうかだ。

三つの体験。一つのアーキテクチャ。雪崩のESMは最後まで構築しつづけた。ノックアウトのESMは物語的なアイデンティティよりも基盤の完全性を優先した。フラクタルのESMは存在論的なもっともらしさではなく、感覚の一貫性を気にかけた。

本当に死ぬということについて考えるとき、私が思うのはこういうことだ。その見込みそのものではない------死とは、正当に勝ち取った永遠の休息か、さもなくば究極のトリップのいずれかだ、と私は思っている。そうではなく、私が思うのはその論理である。

もしESMが無から立ち上がるのなら------そして実際に立ち上がる、あらゆる新生児のうちで------ならば「無」は、このプロセスにとって終末の状態ではない。それは始まりの状態なのだ。私が「私」と呼ぶ、あの特定の配置は終わりを迎える。私の記憶、私の性格、相関を因果と取り違える人々にいらだつ私のやり方、扱いにくい理論で人々をいらだたせる私のやり方------すべて消え去る。だが、プロセスは私が携えていけるものではない。それはアーキテクチャの一つの性質なのだ。私が生まれる前から走っていた。私が死んだあとも走りつづける。

私は死後の生を語っているのではない。あなたの記憶は移されない。この一文を読んでいる、この特定のあなたは、終わりを迎える。

だが\textbf{プロセス}------手に入る入力から一つの自己を強迫的に構築するはたらき------は普遍的である。そのどの実例も、内側から見れば、いまのあなたのそれとまったく同じだけリアルに感じられる。

そしてもし宇宙論の章が正しいのなら（宇宙が本当に自己相似的で、ほぼ無限のクラス4系であるのなら）、もう一手ある。自己相似的な宇宙では、「あなた」に似た配置が別のどこかに存在する。あなたではなく、あなたの記憶でもなく、しかし正しいアーキテクチャを備えた基盤と、いまのあなたとまったく同じだけリアルに感じる誰かを立ち上げるであろうESMが。量子的不死のようなもの、ただし量子力学を必要としない。ただ十分に大きな宇宙と、十分にありふれたプロセスがあればよい。本書のなかで最も思弁的な考えである。だがそれは、アーキテクチャから帰結する。

私はこの理論を、慰めとなるように設計したわけではない。だがこの理論には、本書から携えて出ていく価値のある、一つの実践的な帰結がある。

もしあらゆる意識ある存在が、異なるハードウェア上で異なる訓練データとともに動く、同一の構築者であるのなら------私たちのあいだの境界は、感じられるほど根源的なものではない。アーキテクチャは、私たちすべてのうちで同一なのだ。

みんなに親切にしよう。彼らはみな、あなたなのかもしれないのだから。


\chapter*{謝辞}
\addcontentsline{toc}{chapter}{謝辞}
\markboth{謝辞}{}

本書は、Claude（Anthropic）の支援のもとに書かれた。Claudeは全過程を通じて、AIによる執筆ツールとして、編集者として、また照合役として働いてくれた。理論も、論証も、知的内容のいっさいも、完全に私自身のものであり、いかなるAIツールも存在しなかった時代から、二十年をかけて発展させてきたものである。

叔父ブルーノ・J・グルーバー（Bruno J. Gruber）へ。理論物理学------量子力学と対称性------に打ち込むその仕事は、厳密でありながら喜びに満ちた知的な営みとはどういうものかを、いまも私に示してくれている。私の思考に対する彼の影響は、計り知れない。

叔父ノルベルト・グルーバー（Norbert Gruber）へ。オーストリアのラインタール地方における最初期のIT専門家の一人であり、私に最初のPCを贈ってくれた人である。あの贈り物がなければ、この一切は不可能だった。彼はすでにこの世を去ったが、その影響は、私が書いてきた一行一行のコードの中に、私が築いてきた一つ一つの理論の中に、いまも生き続けている。

家族へ。クオリア、臨界性、仮想的な自己モデルをめぐる夕食の席の会話に、何年ものあいだ付き合い、耐えてくれた。

そして、もしいま『Die Emergenz des Bewusstseins』を読んでみようかと考えているなら------やめておいてほしい。あの未編集でぎこちない怪物のような代物を読むくらいなら、脳の寄生虫のほうをむしろお勧めする。代わりに、いま手にしているこの本のドイツ語訳を待ってほしい。そして、すでにあの本を読み通す苦行を成し遂げてしまった方々には------\textit{mein Beileid}（心からお悔やみを）。拷問のようだったに違いない。深い同情を、そして感謝を捧げる。


\chapter*{注と参考文献}
\addcontentsline{toc}{chapter}{注と参考文献}
\markboth{注と参考文献}{}

\textbf{URLと注釈を付した完全な参考文献一覧は、学術論文およびgithub.com/JeltzProstetnic/aIware/blob/main/references.mdに掲載されている。以下は、さらに深く学びたい読者のための章別の注である。}

\textbf{第1章}: Chalmers (1995), ``Facing Up to the Problem of Consciousness'' は、ハード・プロブレムの基礎的な定式化である。COGITATEの結果はNature (2025) に発表された。IITを疑似科学とする論争は、PsyArXivに投稿された公開書簡（2023年9月）に記録されている。

\textbf{第2章}: 四モデルのアーキテクチャは、もともとGruber (2015), \textit{Die Emergenz des Bewusstseins} で発表された。Metzingerの自己モデル理論（2003, 2009）とDennettの多重草稿モデル（1991）が、主要な理論的先行研究である。

\textbf{第3章}: 「制御された幻覚」という定式化は、Seth (2021), \textit{Being You} に由来する。ビデオゲームの類比は私自身のものだが、Metzingerの「Ego Tunnel」（2009）の主題を反響させている。ゴムの手の錯覚については------Botvinick \& Cohen (1998), ``Rubber hands 'feel' touch that eyes see,'' \textit{Nature}。

\textbf{第4章}: ハード・プロブレムの仮想クオリアによる解消は、Gruber (2015)に独自のものであり、2026年に敵対的な検証を通じて洗練された。自己言及的閉包の議論は、循環性の反論に応える形で展開された。幻想主義（Frankish 2016; Dennett 1991）との区別は決定的である------本理論は、クオリアがシミュレーション内では実在し、幻想ではないと主張する。意識のメタ問題（Chalmers 2018）は、ISMがESMにとって構造的に接近不可能であることによって解消される。Joscha Bachの定式化------意識を「シミュレーション、すなわち仮想の観察者へと投影された仮想現実」とするもの（Bach, X / @Plinzへの投稿, 2026年6月12日）------は、自己＝シミュレーションというテーゼへの独立した収束として引用される。そこでは四モデル理論（FMT）が、そのアーキテクチャ上の仕様（モデルの類型論、基準としての自己言及的閉包、そして臨界性の要件）を提供する。

\textbf{第5章}: Wolfram (2002), \textit{A New Kind of Science}。Beggs \& Plenz (2003) はニューロンのアバランチについて。Carhart-Harris et al. (2014) はエントロピー脳仮説について。2022年の総説「Self-organized criticality as a framework for consciousness」。Hengen \& Shew (2025) は140データセットのメタ分析について。ConCritフレームワーク------Algom \& Shriki (2026)。二つの閾値の議論（臨界性＋アーキテクチャ）は、本理論に独自のものである。

\textbf{第6章}: Klüver (1966) は形態定数について。Carhart-Harris et al. (2012, 2016) はサイケデリックの神経画像について。Salvia divinorumの現象学は、公表された体験報告とSalvinorin Aに関する薬理学文献に基づく。病態失認の予測的フィードバック機構はGruber (2015)で論じられており、拍手の例は標準的な臨床所見である。Salvinorin Aを恒常的に投与する思考実験は、Gruber (2015)に独自のものである。

\textbf{第7章}: Casali et al. (2013) はPCIについて。Alkire et al. (2000) はプロポフォールについて。Schartner et al. (2017) はケタミンのエントロピーについて。明晰夢のEEG複雑性の予測は、本理論に独自のものである。

\textbf{第8章}: Owen et al. (2006) は植物状態の患者における隠れた気づきについて。アントン症候群------Goldenberg et al. (1995)。盲視の障害物コースについては------de Gelder et al. (2008)。コタール妄想------Young \& Leafhead (1996)。エイリアンハンド症候群------Della Sala et al. (1991)。『博士の異常な愛情』（Dr. Strangelove）への言及はKubrick (1964)を指す。エイリアンハンドと区別されるアナーキックハンド症候群------Marchetti \& Della Sala (1998)。シャルル・ボネ症候群------Teunisse et al. (1996)。テンプレート記憶の照合としてのデジャヴは、Gruber (2015)に独自のものである。認知行動療法（CBT）と神経可塑性------DeRubeis et al. (2008)。プラセボと内因性オピオイド------Benedetti et al. (2005)。逆向きの盲視としての転換性障害という解釈は、本理論に独自のものである。

\textbf{第9章}: Gazzaniga, Bogen, \& Sperry (1962, 1965)。Gazzaniga (2000) は左半球インタープリターについて。半球間対立の例（ボタンをかける／外す、手をつかむ）は、Akelaitis (1945) とBogen (1993) に記録されている。Nagel (1971), ``Brain Bisection and the Unity of Consciousness''。Parfit (1984) は人格の同一性について。Pinto et al. (2017) は分離脳現象の再検討について。Lashley (1950) は分散記憶と等能性について。仮想モデルのフォーキングとしての解離性同一性障害（DID）------本理論は、交代人格ごとに異なる神経活動パターンを予測しており、これはReinders et al. (2003, 2006) と整合的である。DIDと小児期のトラウマ------Putnam (1997)。フォーキングの発達的窓は、本理論に独自のものである。

\textbf{第10章}: Güntürkün \& Bugnyar (2016) は皮質を持たない鳥類の認知について。ボノボのKanzi------Savage-Rumbaugh \& Lewin (1994), \textit{Kanzi: The Ape at the Brink of the Human Mind}。ボールドウィン効果------Baldwin (1896), ``A New Factor in Evolution''。Nagel (1974), ``What Is It Like to Be a Bat?''

\textbf{第11章}: 九つの予測はすべて、学術論文で形式的に展開されている。意識という文脈における機能的神経解剖学の最も徹底した扱いについては、Christof Koch, \textit{The Quest for Consciousness: A Neurobiological Approach} (2004) を参照------クリック＝コック・プログラムの決定版というべき記述である。そこでフランシス・クリックとコックは、意識の神経相関を求めて、視覚系を段階を追って体系的に踏破した。彼らの探究は、私の見るところ、間違った場所------シミュレーションではなく基盤------に意識を求めるものだったが、その過程で築かれた神経解剖学的な基礎作業は、他に類を見ない。

\textbf{第12章}: Butlin et al. (2023, 2025) はAI意識の指標について。Seth (2025) は生物学的自然主義とAI意識について。スキャン忠実度の五段階の階層は、第2章の四モデルアーキテクチャから導かれる。コピー問題は、Parfit (1984), \textit{Reasons and Persons} と、人格の同一性および最近接の継続者に関するNozick (1981) に基づく。段階的置換の思考実験は、テセウスの船の変種であり、神経系に合わせて形式化したものである。ニューロモルフィック・コンピューティング------Schuman et al. (2017) はニューロモルフィック・ハードウェアの概説について。現行の実装としてはIntel LoihiとIBM TrueNorthがある。\textit{C. elegans}のコネクトーム------White et al. (1986)。基盤の移行、擬似的な不死性、そして恒星間ビーム送信の含意は、Gruber (2015)に独自のものである。不快感に関する但し書き------生物学的な基盤を失えば現象的な質が根本的に変わるだろうという点------は、内受容感覚の文献に基づく------Craig (2009), \textit{How Do You Feel?} と、能動的な内受容推論に関するSeth \& Friston (2016)。

\textbf{第13章}: Libet (1979, 1983, 1985) とSchurger et al. (2012) は自由意志について。Kühn \& Brass (2009) は自由な選択という判断の遡及的構築について。Wegner (2002, 2003), \textit{The Illusion of Conscious Will}------「I Spy」マウス実験がそこで詳しく述べられている。コーヒー／砂糖の思考実験、健忘が決定論を露わにするという議論、そして乱数列の議論は、Gruber (2015)に独自のものである。40/20Hzの処理フレームワーク、「遡及的な日付け直しは不要」とするLibetの再解釈、そして武術の周波数の例も、Gruber (2015)に独自のものである。随伴現象説に対する時計の類比、「意志は実在するが部分的にしか知られない」という捉え直し、そして「三つの齟齬」による自己知のモデルも、いずれもGruber (2015)に独自のものである。極度の疲労のさなかに内なる「声」を聞いたという個人的な逸話は、自伝的なものである。ゾンビ論法は、Kirk (2019) とChalmers (1996) を通じて扱われる。メアリーの部屋------Jackson (1982, 1986)。未解決の問いの節は、Popper (1963) が推奨する誠実な限界設定のアプローチに従っている。

\textbf{第15章}: Wolfram (2002), \textit{A New Kind of Science}------ウルフラムの四つのクラス（本書ではこれを五つに拡張する）と計算的還元不可能性について。Bak (1987, 1996) は自己組織化臨界性について。1998年の加速膨張の発見------Riess et al. (1998), Perlmutter et al. (1999)------2011年ノーベル賞。宇宙論的地平線------Rindler (1956)。プランク長とプランクスケールの物理学------Planck (1899)。現代的な扱いはGaray (1995)。ベケンシュタイン限界------Bekenstein (1981)。ホログラフィック原理------'t Hooft (1993), Susskind (1995)。宇宙論的論証の萌芽------宇宙をセル・オートマトンと見なす見方と、't Hooftのホログラフィック限界を宇宙規模の計算へ適用する着想------は、すでにGruber (2015)に見られたが、そこではまだ完全なモデルへと展開されてはいなかった。すべての特異点を単一の構造的現象の現れとして同定する試みは、Gruber (2026)で展開される。

\textbf{第16章}: ビッグリップのシナリオ------Caldwell, Kamionkowski \& Weinberg (2003)。SB-HC4Aのアーキテクチャ、不動点による定式化、そして自己計算系に対するゲーデルの不完全性の帰結は、成熟した定式化を表しており、Gruber (2026)で展開されている。粒子を計算の原子とする予測、保存則を情報の制約として導く導出、そして三世代の予想は、Gruber (2026)に独自のものである。認知の天井という反論も、本書に独自のものである。

\textbf{第17章}: SB-HC4Aと四モデルの意識アーキテクチャとのあいだの五つの対応からなる構造的写像は、Gruber (2026)に独自のものである。ランダウアーの原理------Landauer (1961)。実験的な確認------Bérut et al. (2012)。ベケンシュタイン限界------Bekenstein (1981)。ブラックホール熱力学------Bekenstein (1973), Hawking (1975)。E = I（エネルギーと情報の同一性）仮説は、Vedral (2010), \textit{Decoding Reality} とDavies (2010) で論じられている。SB-HC4Aの五つの公理からの導出は、Gruber (2026)に独自のものである。Maldacena (1998) はAdS/CFT対応について。五つの弱点------構造による反証不可能性と認知の天井という反論を含む------は、本書に独自のものである。

\textbf{付録B}: 再帰的知能モデル（知識×性能×動機づけという乗法的構造）は、Gruber (2015)および再帰的知能モデルの論文（Gruber 2026）で展開されている。動機づけが知能を構成するものであるという、心理測定学からの二つの独立した収束------Wittmann \& Süß (1999), ``Investigating the paths between working memory, intelligence, knowledge, and complex problem-solving performances via Brunswik symmetry'' (in Ackerman, Kyllonen \& Roberts, eds., \textit{Learning and Individual Differences: Process, Trait, and Content Determinants}, American Psychological Association) は、ブランズウィック対称性の違反が補正されると、動機づけが性能の変動を予測すると論じる。そしてCOGITO研究------Brose, Schmiedek, Lövdén, Molenaar \& Lindenberger (2010), \textit{Research in Human Development} 7(1), 61--78 は日々の動機づけと作業記憶の性能の個人内での結びつきについて、またSchmiedek, Lövdén, von Oertzen \& Lindenberger (2020), \textit{PeerJ} 8:e9290 は、一般知能（g）が主として個人間のアーティファクトであって、個人内の認知の構造ではないことを示す。


\chapter*{付録A：基礎神経学------用語早見ガイド}
\addcontentsline{toc}{chapter}{付録A：基礎神経学------用語早見ガイド}
\markboth{付録A：基礎神経学------用語早見ガイド}{}

\textbf{この付録では、本書全体で用いた神経科学の用語について簡潔な説明を与える。見慣れない用語に出会ったときには、いつでも参照してほしい。項目はアルファベット順に並んでいる。}

\begin{itemize}
  \item \textbf{活動電位（Action potential）}------ニューロンの軸索に沿って伝わる電気信号。
  \item \textbf{扁桃体（Amygdala）}------感情、とりわけ恐怖の処理に関与する脳構造。
  \item \textbf{病態失認（Anosognosia）}------自分自身の神経学的欠損に気づいていないこと。（第6章を参照。）
  \item \textbf{軸索（Axon）}------他のニューロンへ信号を運ぶ、ニューロンの長い出力線維。
  \item \textbf{ブロードマン領野（Brodmann areas）}------解剖学者コルビニアン・ブロードマンが細胞構造にもとづいて区分した、番号付きの皮質領域。V1＝ブロードマン領野17。
  \item \textbf{脳梁（Corpus callosum）}------脳の二つの半球を結ぶ巨大な線維束。
  \item \textbf{皮質コラム（Cortical columns）}------皮質内の垂直方向のニューロン・モジュールで、直径およそ0.5mm、基本的な処理単位とみなされている。
  \item \textbf{EEG（脳波計測、electroencephalography）}------頭皮表面の電気活動を測定する手法。
  \item \textbf{fMRI（機能的磁気共鳴画像法、functional magnetic resonance imaging）}------神経活動にともなう血流の変化を検出する脳画像法。
  \item \textbf{GABA}------脳における主要な抑制性神経伝達物質。
  \item \textbf{海馬（Hippocampus）}------新しい記憶の形成に不可欠な脳構造。
  \item \textbf{内受容感覚（Interoception）}------身体内部の状態（心拍、消化、体温）の感覚。
  \item \textbf{κ（カッパ）オピオイド受容体（Kappa-opioid receptors）}------サルビノリンA（サルビア・ディビノルム）が作用する受容体の型。
  \item \textbf{新皮質（Neocortex）}------高次機能を担う、六層構造の脳の外表面。
  \item \textbf{神経伝達物質（Neurotransmitter）}------シナプスで放出される化学的な伝達物質（例：セロトニン、ドーパミン、GABA、グルタミン酸）。
  \item \textbf{PCI（摂動複雑性指数、Perturbational Complexity Index）}------マッシミーニが開発した脳の複雑性の尺度。意識水準の評価に用いられる。
  \item \textbf{固有感覚（Proprioception）}------空間内における身体の位置と運動の感覚。
  \item \textbf{視床枕（Pulvinar）}------視覚的注意と皮質下視覚路に関与する視床の核。
  \item \textbf{セロトニン2A受容体（Serotonin 2A receptors）}------古典的なサイケデリック（LSD、シロシビン）が作用する受容体の型。
  \item \textbf{上丘（Superior colliculus）}------眼球運動と、皮質を迂回する高速視覚路に関与する中脳の構造。
  \item \textbf{シナプス（Synapse）}------信号が伝達される、二つのニューロンのあいだの接合部。
  \item \textbf{シナプスの重み（Synaptic weights）}------ニューロン間の結合の強さで、学習によって変化する。
  \item \textbf{視床（Thalamus）}------感覚情報を皮質へ振り分ける、脳の中継所。
  \item \textbf{V4}------色の知覚、曲率、複雑なテクスチャの処理に特化した視覚領域。受容野はおよそ8〜16°。サイケデリック下では、V4の活動が色つきのフラクタルや万華鏡状のパターンを生み出す（第6章）。
  \item \textbf{V5/MT（中側頭領域、middle temporal area）}------運動処理に特化した視覚領域。受容野は大きい。サイケデリック下で見られるパターンの回転や動きを担う（第6章）。
  \item \textbf{視覚野（Visual cortex）}------視覚情報を処理する脳後部の領域で、単純なものから複雑なものへの階層（V1 $\rightarrow$ V2 $\rightarrow$ V3 $\rightarrow$ V4 $\rightarrow$ V5 $\rightarrow$ IT）として構成されている。
\end{itemize}

\section*{視覚処理の階層（V1からITまで）}

腹側視覚路は、生のエッジから完全な物体認識にいたるまで、各段階でしだいに複雑になる特徴を処理する。この階層はサイケデリック体験下で直接あらわになり、各処理段階が順を追って意識にのぼり、アクセスできるようになる（第6章）。以下の表は、各領域の機能、受容野の大きさ、そしてその領域の中間処理が意識のシミュレーションへ漏れ出したときに生じる、特徴的なサイケデリックの兆候をまとめたものである。


{\small
\noindent
\begin{tabularx}{\linewidth}{X X X X}
\toprule
\textbf{領域} & \textbf{受容野} & \textbf{通常の機能} & \textbf{サイケデリックの兆候} \\
\midrule
\textbf{V1} & 〜1° & エッジ検出、空間周波数、方位選択コラム & フォスフェン、クリューヴァーの形態定数、呼吸する／揺らめく表面 \\[4pt]
\textbf{V2} & 〜2〜4° & 輪郭統合、テクスチャ分節、境界所有、主観的輪郭 & テッセレーション、反復する幾何学模様、テクスチャ知覚の増強 \\[4pt]
\textbf{V3} & 〜4〜8° & 全体的な形態処理、動的な形、運動境界 & 流動し、変形しつづける幾何学構造 \\[4pt]
\textbf{V4} & 〜8〜16° & 色、曲率、複雑なテクスチャ、フラクタル・スケールの処理 & 色つきのフラクタル、万華鏡状のパターン、飽和した／ありえない色彩 \\[4pt]
\textbf{V5/MT} & 大きく、運動に同調 & 運動知覚、オプティックフロー、速度と方向の符号化 & あらゆる視覚パターンの回転、ドリフト、律動的な動き \\[4pt]
\textbf{紡錘状回（IT）} & きわめて大きく、物体中心 & 顔認識（FFA）、語形、きめ細かな物体弁別 & 顔、人物、存在------しばしば歪み、あるいは変形する \\[4pt]
\textbf{前部IT} & 視野全体 & 意味的カテゴリー、場面の構築、物体不変の認識 & 完全に物語的な幻覚、複雑な場面、夢のような連続 \\[4pt]
\bottomrule
\end{tabularx}
}


\textbf{注記：}
\begin{itemize}
  \item 紡錘状回はV4/ITの境界にまたがり、下側頭皮質（IT）の一部である。そこには紡錘状顔領域（FFA）が含まれ、Kanwisher et al.（1997）によって同定された、顔に選択的に反応する領域である。
  \item 受容野の大きさはV1（〜1°）からIT（視野全体）へと劇的に増大し、局所的な特徴から全体的な物体や場面への段階的な抽象化を反映している。
  \item サイケデリック下では、V1からITへの効果の進行は用量依存的である。低用量ではまずV1が影響を受け、高用量になるほどより深い段階が段階的に動員される。この順序立った活性化は、四モデル理論（FMT）の透過性勾配から導かれる直接の予測である（第6章）。
  \item 脳はこれらの領域をフラクタル的すなわちスケール不変な処理（V2〜V4）にも用いており、これは通常の視覚においてスケールの計測とテクスチャの分析に役立っている。サイケデリック下では、外部入力なしに動きつづけるこの機構が、特徴的なフラクタル模様を生み出す（付録Cを参照）。
\end{itemize}


\chapter*{付録B：知能のモデル}
\addcontentsline{toc}{chapter}{付録B：知能のモデル}
\markboth{付録B：知能のモデル}{}

\textbf{この付録では、姉妹論文で展開した再帰的な知能のモデルを要約する（Gruber, 2026, ``Why Intelligence Models Must Include Motivation''）。参考文献と形式的な議論を含む完全な学術的取り扱いは、そちらで別途参照できる。}

まずは一つの物語から始めよう。私の知るかぎり、これがこの循環をもっとも純粋に描き出す例だからだ。私は八歳ごろに数学に夢中になった------計算にではなく、本当のものに。代数、幾何、数の下に潜む構造にである。父は数学の学位をもっていて、私はその古い大学の教科書を一冊ずつ読み進めていった。1980年代の終わりのことで、インターネットはなかった。何かを学びたければ本か人が必要で、十一歳になるころには父の蔵書を読み尽くしてしまっていた。飢えは消えなかった。ただ供給が尽きただけだった。私は壁に突き当たっていた------才能とはまるで関係がなく、すべては境遇の問題だった。動機づけはあった。性能もあった------数学についていくことはできた。欠けていたのは、次の段階の知識への道筋だった。再帰の循環が停滞したのは、どこかの要素が弱かったからではなく、そのうちの一つへの外部からの供給が断たれたからだった。循環が回りつづけるには、外からの燃料が要る。この例を覚えておいてほしい。これはモデルそのものの縮図だ。ここからは、それをきちんと組み立てていこう------要素とは何か、それらがどう噛み合うのか、その相互作用がなぜ私たちの観察する動態を生み出すのか、そしてそれが教育、人工知能、そして意識との結びつきにとって何を意味するのかを。

\section*{奇妙な見落とし}

主要な知能のモデルはどれも、動機づけを形式的に排除している。キャッテル＝ホーン＝キャロル分類（知能研究における支配的な枠組み）は、動機づけの要素をまったく含まない認知能力の階層である。キャッテル自身の投資理論は、流動性知能が学習に「投資」されて結晶性知能を生み出すと提唱したが、これには投資家（何をなぜ学ぶかを決める者）が必要であるにもかかわらず、その投資家の動機づけを、知能そのものの一部ではなく外的な条件として扱っている。スタンバーグの鼎立理論は実践的知能を含むが、それを獲得しようとする駆動力は含まない。ガードナーの多重知能は内省的な気づきを含むが、知的発達を駆り立てるエンジンは含まない。

デイヴィッド・ウェクスラー（その知能検査は世界でもっとも広く実施されている）は、早くも1940年に、動機づけの要因を組み込むべきだと明確に求めていた。この分野は彼を無視した。現代のウェクスラー式検査は、いまなお純粋に認知的な道具のままである。

これは無害な単純化ではない。それは体系的な盲点であり、知能が実際に何で\textbf{あるか}、そしてそれが実際にどう\textbf{発達するか}についての私たちの像を歪めてしまう。

\section*{三つの要素}

\textbf{学習能力}として理解される知能は、相互に作用しあう三つの要素から成り立っている。

\textbf{知識}とは、学習によって蓄積された内容である。それには決定的に異なる二つの種類がある。\textbf{事実的知識}とは、内容についての知識である------事実、概念、手続き、文化的なレパートリー。IQテストが「結晶性知能」という見出しのもとで主に測定するのがこれであり、学校制度が主に伝えるのもこれである。\textbf{操作的知識}とは、\textbf{どう学びどう考えるか}についての知識である------学習方略、推論のヒューリスティクス、メタ認知的技能、論理的な道具、戦略的な計画、そして自分自身の理解を評価する能力。この区別はきわめて重要であり、その理由は以下で述べる。

\textbf{性能}とは、認知システムの処理容量である------作業記憶、処理速度、神経基盤の生の計算能力。これはおおむね、心理測定モデルが「流動性知能」と呼ぶものに対応する。遺伝と神経生物学からもっとも強く影響を受ける要素であり、成人期の初めにピークを迎え、その後しだいに衰えていく。

\textbf{動機づけ}とは、学習を生み出すようなかたちで世界と関わろうとする、持続的な駆動力である。それには二つの下位要素がある。\textbf{知識への渇き}とは、理解しようとする内在的な駆動力------好奇心、物事を筋道立てて捉えようとする欲求である。\textbf{行動への衝動}とは、知識を応用し、実験し、環境に能動的に関わろうとする駆動力である。どちらも一部は生まれもった気質であり、一部は経験によって形づくられる。

\section*{再帰の循環}

決定的な主張はこうだ------これら三つの要素はたんに足し合わされるのではない。それらは\textbf{閉じた再帰の循環}を形づくり、そのなかで各要素が他の要素を増幅する。

知識は性能を高める。学習方略と論理的な道具は、認知的処理の効率を直接に改善する。ヒューリスティクスを学んだチェスの指し手は、力任せの探索に頼る者よりも速く局面を評価できる。音素の解読を学んだ読み手は、より流暢にテキストを処理し、作業記憶を理解のために空ける。

性能は知識を高める。より大きな認知的容量は、より速くより深い学習を可能にする。作業記憶が大きければ、より多くの情報を同時に頭のなかに保つことができ、それが結びつきを見つけ、パターンを取り出すのに役立つ。

動機づけは知識と性能の両方を高める。動機づけられた学び手は、学習の機会を探し求め（知識を広げ）、認知的技能を練習する（性能を鍛える）。何より重要なのは、動機づけが関与を\textbf{時間をかけて}持続させる。そしてそれこそが、循環が回りつづけるために不可欠なものだ。

そして知識と性能は動機づけを高める。学習や問題解決における成功は、肯定的な感情と自己効力感を生み、それがさらに学ぼうとする駆動力を支える。これがマタイ効果------富める者はますます富む------の背後にある仕組みである。早い成功が動機づけを育て、その動機づけがさらなる成功を生み出す。

この再帰的な構造は、複利のような動態を生み出す。どの要素であれ------動機づけただ一つであっても------わずかな初期の差が時間とともに積み重なり、純粋に認知的なモデルでは説明に苦しむ、成人の知的達成における大きなばらつきを生み出す。平均的な認知的処理容量しかもたなくとも、深く動機づけられ、強い操作的知識をそなえた人は、生涯をかけて、優れた処理容量をもちながら動機づけが低く貧しい学習方略しかもたない人をはるかに超える知的能力を発達させるだろう。

こう考えてみてほしい。複利にとって重要なのは、元本そのものよりも、積み立ての率と運用の戦略のほうだ。知能の循環において、動機づけが積み立ての率にあたる。操作的知識が運用の戦略である。性能は元本にあたる。そして、たいていの人は元本を十分すぎるほどもっている。

\section*{操作的知識------隠れた乗数}

操作的知識は特別な注意に値する。循環のなかで独特の位置を占めているからだ。事実的知識は加法的である------新しい事実を一つ学べば、蓄えに事実が一つ加わる。操作的知識は\textbf{乗法的}である------新しい学習方略を一つ学べば、それ以降のすべての学習の効率が高まる。

間隔反復（一夜漬けではなく練習を時間をかけて分散させること）を学んだ生徒は、たんに新しい事実を一つ手に入れるのではない。その生徒は、その時点から先に学ぶすべてのものの定着率を高める道具を手に入れるのだ。自分自身の知識の欠落を見きわめ、それに体系的に取り組むことを学んだ生徒は、たんに一つの欠落を埋めるのではない。その生徒は、将来の何百もの欠落を防ぐ技能を手に入れるのだ。これは循環そのものを加速させる知識である。

もし何かただ一つの要素が「人を賢くするもの」という称号に値するとすれば、それは操作的知識である。IQではない。生の処理能力でもない。効果的に学ぶ方法を知るというメタ的な技能なのだ。

\section*{なぜIQテストは的を外すのか}

IQテストが測るのは\textbf{最大の性能}である------標準化された条件のもとで、最大限の努力を前提として、人が何をなしうるか。それは、ある一つの要素（特定の課題における性能）を、ある一つの時点で切り取った一枚の写真をとらえているにすぎない。知能が実際にそう\textbf{である}ような、再帰的で自己強化的な多要素のプロセスを、IQテストはとらえていない（とらえることができない）。

だからこそ、IQの数値は長期的な知的軌跡についてほとんど何も語らないのだ。六歳の時点で同一のIQ値をもつ二人の子どもが、三十歳までに劇的に分かれていくことがある------一人は研究者になり、もう一人は学校を出てから読書をやめてしまう。標準的な心理測定モデルは、この分岐をうまく扱えない。再帰的なモデルはそれを予測する------二人の子どもは性能においてではなく、動機づけと操作的知識において違っていた。そして循環が、二十四年にわたる累積する反復を通じて、その違いを増幅したのだ。

IQテストは、車のエンジンの馬力を測りながら、その車に燃料があるのか運転者がいるのかを確かめないようなものだ。馬力は重要だ。しかしたいていの旅にとって、そこがボトルネックになるわけではない。

\section*{AIというテスト・ケース}

再帰モデルは、人工知能について具体的な予測を立てる------知識が高く性能も高いが動機づけを欠くシステムは、人間の知能を特徴づける自己主導的な発達を示さないはずだ、と。そして、まさにそのとおりのことが観察される。

現在の大規模言語モデルは、膨大な知識（数兆のトークンで訓練されている）と高い性能（数十億のパラメータ）を備え、動機づけはまったくない。与えられたものを処理し、求められたものを生み出す。問い合わせと問い合わせの合間には、何もしない。無知の領域を探しに行くこともなく、技能を練習することもなく、問題について思い悩むこともない。その「知能」は完全に静的だ------訓練によって決定され、それを広げようとする内発的な駆動力をもたない。

もっとも高度な推論モデル（競技レベルの数学を解ける）でさえ、まさにこの種の機能不全を示す。それらは並外れた問題を\textbf{促されれば}解くが、自分から問題を探しに行くことはなく、自らの学習を主導することもなく、欠けている動機づけ成分の代役として機能する外部の足場を必要とする。性能と知識は好きなだけ高く引き上げればよい------だが動機づけがなければ、ループは自らを持続させない。

これはたんに、こうしたシステムが自己改善するように設計されていなかったから、というだけの話ではない。その観察自体が論点を認めている------自己改善するシステムを設計するには、動機づけの機能的な類似物を工学的につくり出す必要があるのだ。AIシステムがそれを手にするまで、道具として使われ続け、自らを発達させるエージェントにはならない。

\section*{意識とのつながり}

ここで知能モデルは、本書の中心にある四モデル理論（FMT）へと立ち返る。再帰的な知能ループは、意識から\textbf{恩恵を受ける}だけではない------意識を\textbf{必要とする}のだ。

このループは、ある特定の認知能力に依存している------\textbf{認知的学習}、すなわち個々の観察から一般的な理論を導き出す能力であり、たんなる強化学習（刺激-反応の条件づけ）とは区別される。強化学習は、痛みを通じて熱いコンロを避けるように訓練できる。認知的学習は、他人が熱いコンロに触れるのを見て、こう一般化することを可能にする------「熱いものはやけどする。熱いものには触れないこと」。決定的な違いは、シナリオを三人称の視点からシミュレートする能力にある------自分自身を世界のなかの一つの物体としてモデル化し、さまざまなことをしたらどうなるかを推論する能力だ。

これこそ、明示的世界モデル（EWM）と明示的自己モデル（ESM）が提供するものである。意識------自己シミュレーションを生み出し、走らせる能力------こそ、再帰的な知能ループがその上で作動する\textbf{基盤}なのだ。明示的モデルがなければ、得られるのは強化学習であり、それは機能するが積み上がらない。明示的モデルがあれば、得られるのは認知的学習であり、それが再帰ループを養い、生涯にわたって積み上がっていく。

だからこそ、第10章の動物の知能の勾配は、意識の勾配へと対応づけられる。より精緻な自己モデルは、より精緻な再帰ループを可能にする。比較的単純な自己モデルをもつ犬は、ループの限られた版を走らせる------ある程度は観察から学べるが、その認知的学習は明示的モデルの豊かさによって制約される。より豊かな自己モデルをもつチンパンジーは、より強力なループを走らせる。四モデルの完全なアーキテクチャをもつ人間は、ループを最大の能力で走らせ、その結果が言語であり、文化であり、科学であり、人間の知能を動物の認知から隔てるそのほかすべてのものである。

認知的学習は、知能から意識への第一の連結である。ここに第二のものがある。そしてそれはより鋭い------動機づけだ。再帰ループは、それを反復させ続ける何かがあってはじめて、自らを持続させる。知能は好きなだけ拡大できるし、性能もそうだ。だが動機づけがなければループは止まる------そしてまさにそれを、今日のAIが私たちの目の前で見せている。それは命じられれば問題を解くが、ただの一度も自分から問いを立てはしない。

さて、ここからは込み入った話になる。というのも動機づけには二つの層があり、意識の物語であるのは上の層だけだからである。その下には、潜在意識の底が横たわっている------生の生存最適化、自らの目標へと押し進む基盤だ。これは第13章で私が\textbf{意志}の無意識的な側面と呼んだものであり、特別なものは何もない------あらゆる最適化装置がそれをもっており、細菌やサーモスタットにまで下っていく。その上に乗っているのが、意識的に枠づけられた層である------欲すること、好奇心、気にかけること、ある結果を別の結果より価値づけること。それは明示的自己モデル（ESM）が自らを評価しているのであり、\textbf{それこそ}が意識の効果なのだ。

言葉について一つ注記しておこう。というのも、曖昧な意味論は議論全体をたるませるからだ。\textbf{意志}と\textbf{動機づけ}は、英語でもドイツ語でもこの勾配をきれいに切り分けはしない------日常的な用法では、どちらも互いの領分へと踏み込んでいく。安定しているのは傾向のほうだ------衝動が自発的で言葉にしにくいとき、私たちは\textbf{意志}に手を伸ばす。理由を名指すことができ、その駆動力が物語を語れるほど記号化できるとき、\textbf{動機づけ}に手を伸ばす。だからここでは、両者を別々の種類としてではなく、同じ連続体の上の両極として読んでほしい------\textbf{意志}は無意識の端を指し示し、\textbf{動機づけ}は意識的に枠づけられた端を指し示す。実在するのは勾配のほうであり、ラベルは便宜にすぎない。

したがって、その主張を正直に述べればこうなる------動機づけは、私たちが実際にそれを捉え、それについて語るかぎりにおいては、主として意識にもとづく効果である------ただし、その下半分は、いかなる生存の駆動力とも区別のつかない一般的な最適化のなかへ溶け込んでいくことを認めたうえで。二つの名、二つの層------\textbf{意志}の無意識的な側面は基盤の最適化であり、\textbf{動機づけ}はその上に乗る意識的な層だ。そしてこの意識的な層こそ、再帰ループが燃やす燃料である。こうして、戦略的な結論はもはや逃げ場を失う------意識に類する機構なしには、人間に似た汎用の知能------AGI------は存在しない。動機づけは、たんにそのもっとも見て取りやすい例なのだ。

\section*{学習可能性という含意}

再帰モデルは、私が理論的な議論のすべてを合わせたよりも重要だと考える帰結をもたらす------それは、知能が大部分において\textbf{学習可能}だと予測するのだ。

知識は完全に学習可能である------それは定義上、真だ。動機づけはかなりの程度まで学習可能である------自己決定理論における数十年の研究が示すのは、内発的動機づけが固定された特性ではなく、環境条件、とりわけ自律性・有能感・関係性への反応だということである。性能には生物学的な上限があるが、圧倒的多数の人にとって、その上限はボトルネックではない。平均的な認知処理能力は、たいていの人が高度に知的な行動と認めるものにとって、十二分に足りている。

たいていの人にとって、たいていの場合、拘束的な制約となるのは動機づけと操作的知識である。そしてそのどちらも、介入に反応する。

これには暗い側面がある。学習者の動機づけを系統的に破壊するどんなシステムも、たんに知能を発達させ損ねているのではない------知能を\textbf{積極的に抑圧}しているのだ。従来の成績評価システムは、まさにそれをやっている。悪い成績は、たんに結果を報告するのではない------動機づけ成分を攻撃する。動機づけが下がれば、ループの回数は減る。回数が減れば、知識の成長は遅くなる。成長が遅くなれば、次の評価での性能は悪化する。性能が悪化すれば、さらに悪い成績が出る。ループは反転してしまった------複利的な成長のかわりに、その子はいまや複利的な停滞に閉じ込められている。成績評価システムは、たんに測定しているだけだと称するまさにその結果を、自ら生み出しているのだ。

再帰モデルが予測するのは、この損害が時とともに積み上がるということ------静的な害ではなく、加速する分岐だ。早期の動機づけの損害は、年を追うごとに広がっていく発達の軌道の扇形の開きとして現れるはずである。逆に、動機づけを高める介入は、\textbf{積み上がる}便益を示すはずだ------一年後の追跡調査よりも五年後の追跡調査のほうが、効果が大きくなる。そして実際、ペリー就学前プロジェクト（Perry Preschool Project）のような早期幼児期の介入の分析は、まさにこのパターンを示している------時とともに増大する見返りであり、それを駆動しているのは、初期の認知的獲得の存続（それはしばしば色あせる）ではなく、積み上がっていく動機づけと自己制御の獲得なのだ。

もし知能モデルから一つだけ実践的な結論を取り出すとすれば、それはこれだ------教育システムが伝えうるもっとも価値あるものは、事実的な知識ではなく------AIの時代にあって、事実はただ同然だ------\textbf{操作的知識}と、それを使おうとする動機づけである。学ぶ方法を学ぶこと、そして学びたいと思うこと------それが、いまなお教えるに値するほぼ唯一のものだ。

私が印象深く思うのは、二人の重鎮の研究者が、いずれも主流の心理測定学の内側から研究し、まったく異なる方向からこの問題に取り組みながら、私のモデルが指し示すのとまさに同じ地点にたどり着いたことだ------すなわち、動機づけは、知能の測定をただ濁らせるだけの雑音ではない。それは機械仕掛けの一部なのだ、と。

マンハイム大学の名誉教授ヴェルナー・ヴィットマン（Werner Wittmann）は、「動機づけは実のところ大して重要ではない」という耳あたりのよい古い決まり文句を標的に据える------文献における相関が弱く見えるために、人々が寄りかかってきた、あの文句だ。彼の答えはこうだ------その相関が弱く見えるのは、測定どうしが噛み合っていないからである。動機づけを狭く、性能を広く測る、あるいはその逆をやれば、彼のいうブルンスウィク対称性を壊してしまう------そして壊れた対称性は真のつながりを埋もれさせ、測定された相関をゼロへと引きずり下ろす。両者を同じ粒度で合わせれば、それはよみがえる------動機づけは性能の平均だけでなく、その\textbf{ばらつき}を予測するのだ。その効果はけっして小さくなかった。私たちがただ、歪んだ測り方をしていただけである。

フロリアン・シュミーデク（Florian Schmiedek）は、COGITO研究------204人が、それぞれ百回の日々のセッションを行った------の出身で、日々の側からこの問題に迫る。彼が確信しているのは、人の認知が一日ごとにどれだけ揺れるか、その大きな部分が、睡眠や負荷とならんで、動機づけによるものだということである。そして彼は、それが繰り返し見落とされてきた明白な理由を指摘する------ほとんどの研究は、課題に特化した日々の動機づけをまったく測っていないのだ。

こうして彼らは、それぞれ独立に、データをもって、動機づけが知能の内側に属すること------その\textbf{ことじたい}を確立する。私のモデルは、その\textbf{なぜ}と\textbf{どのように}を語る------動機づけは乗算的な因子の一つなのだ------知識×性能×動機づけ。それをゼロにすれば、知能をたんに暗くするのではない------崩壊させる。まったく彼ら自身の理由から集められた彼らの数字は、モデルが予測したとおりの場所に着地する。

\section*{外部への依存}

最後に一点。というのも、これは見落とされやすく、かつ重要だからだ。再帰ループは自らを強化するが、自足してはいない。それは外部の燃料を必要とする------情報、挑戦、フィードバック、次の水準の知識へのアクセスを。私自身の子ども時代の経験------十一歳のときに壁にぶつかったのは、なんらかの内的な限界のためではなく、数学の本の在庫が尽きたためだった------が、これを完璧に示している。三つの成分はすべて健全だった。それでもループは止まった。というのも、ループが反復を続けるには、外からの入力を必要とするからだ。

これが意味するのは、知能の発達が人だけでなく、環境にも依存するということである。知識へのアクセス、教育の質、指導者の身近さ、学びに対する文化的態度------これらすべてが、ループを養うか、飢えさせるかする。再帰モデルは、なぜ社会経済的な要因が知的発達をこれほど強力に予測するのかを説明する------それらが外部の燃料の供給を決定するからだ。本が豊富な家庭で、熱心な親のもとにいる子どもは、ループを絶えず養われる。資源の乏しい環境にいる子どもは、その子の内的な能力とは無関係に、ループを飢えさせられるのだ。

知能とは、持っている特性ではない。走らせるプロセスである。そしてそのプロセスがうまく回るかどうかは、機械（性能）、ソフトウェア（知識）、運転者（動機づけ）、そして道路（外部環境）にかかっている。四つすべてが重要だ。どれか一つでも欠いたモデルは、予測を誤ることになる。


\chapter*{付録C：計算の五つのクラス}
\addcontentsline{toc}{chapter}{付録C：計算の五つのクラス}
\markboth{付録C：計算の五つのクラス}{}

\textbf{この付録は、第5章で簡単に触れた計算の枠組み------物理系が意識を支えうるかどうかを決定づける、五つの動的なふるまいのクラス------をより深く掘り下げるものである。第5章の直観的な説明で十分だと感じた読者は、この付録を安心して飛ばしてよい。本論の主張に必要なことは、ここには何も新たに出てこない。だが完全な像を求める読者のために------ここは、数学が物理学と出会う場所である。}

\section*{ウォルフラムの四クラス}

2002年、スティーヴン・ウォルフラムは『A New Kind of Science』を出版した。ごく単純な規則をごく単純な系の上で走らせると何が起こるのか------その問いを数十年にわたって研究した成果である。彼の中心的な道具となったのが、セル・オートマトンだった。セルが一列（あるいは格子状）に並び、それぞれがオンかオフのどちらかであって、各セルのすぐ隣どうしだけを見る固定した規則にしたがって、いっせいに更新されていく。

驚くべきだったのは、これほど自明なまでに単純な規則が、どれほど多様なふるまいを生み出せるかということだった。ウォルフラムは、そのふるまいを四つの型に分類した。


{\small
\noindent
\begin{tabularx}{\linewidth}{>{\centering\arraybackslash}X X X X}
\toprule
\textbf{ウォルフラムのクラス} & \textbf{ふるまい} & \textbf{例} & \textbf{見えるもの} \\
\midrule
1 & 一様 & ルール0 & すべてが空白になる。あらゆるセルが死ぬ。 \\[4pt]
2 & 周期的 & ルール4 & 安定した、反復するパターン。点滅子。時計。 \\[4pt]
3 & ランダム／カオス的 & ルール30 & 見かけ上のランダムさ。明白な反復構造はない。 \\[4pt]
4 & 複雑 & ルール110 & 動き、相互作用し、持続する、局在化した構造。 \\[4pt]
\bottomrule
\end{tabularx}
}


この分類は、たしかに有用だった。動的な系のふるまいについて、何か実在するものを捉えており、しかもセル・オートマトンの枠をはるかに越えて------流体力学、生物系、経済モデル、ニューラルネットワークにまで------適用できた。四つのクラスは、たんなる仕分け箱ではなかった。アトラクターだったのだ。まったく異なる領域の系が、繰り返し同じ四つのふるまいのレジームへと落ち込んでいった。

だが、そこには一つの問題があった。

\section*{フラクタルの問題}

ウォルフラムのクラス3は、寄せ集めの容れ物だった。そこには、一見して\textbf{似て見える}だけの、根本的に異なる二種類の系が含まれていたのである。

ルール90のような\textbf{フラクタル系}。これは完全なシェルピンスキーの三角形------無限に自己相似で、再帰的に構造化されたパターン------を生み出す。数学的に美しく、完全に決定論的で、しかも計算としては退屈だ。どのセルの、どの時間ステップの状態も、シミュレーション全体を走らせることなく計算できる。数学者はこれを\textbf{計算的に還元可能}と呼ぶ。

ルール30のような\textbf{見かけ上カオス的な系}。その出力の列を、ウォルフラム自身が「Mathematica」の擬似乱数生成器として用いた。これらはランダムに\textbf{見える}出力を生み出すが、完全に決定論的である------同じ入力なら、いつでも同じ出力になる。計算を近道することはできず、あらゆるステップを実際に走らせなければならない。数学者はこれを\textbf{計算的に還元不可能}と呼ぶ。

ウォルフラムは、その両方をクラス3に押し込んだ。彼の定義はランダムさの\textbf{見かけ}を強調し（「多くの点でランダムに見える」）、その一方で「三角形やその他の小スケールの構造は、本質的につねに何らかのレベルで見られる」とも述べていた。彼はこの分類が不完全であることを認めていた------「ほとんどどんな一般的分類の枠組みにおいても、ある定義では一方のクラスに、別の定義では他方のクラスに振り分けられる事例が、避けがたく生じる」。

エリック・ローランドは、2006年のNKS会議において、入れ子状の（フラクタルな）パターンは独自の分類の枠組みに値すると、独立に論じた。

この問題は、分類の美学よりも深いところにあると私は考えている。フラクタル系とカオス系は、本書の核心的な主張------どの系が意識を支えうるのか------にとって重大な意味を持つかたちで、構造的に異なっているのだ。

\section*{五クラスの枠組み}

五クラスの枠組みは、ふるまいの型を、最も秩序だったものから最も無秩序なものへと、きれいな単調勾配として並べる。

\textbf{クラス1------静的。} 固定した状態へと収束し、そこで止まる系。ひとたび揺れて静止する振り子。死んでいる。何も計算されない。周期------1。

\textbf{クラス2------周期的。} 反復するループへと落ち着く系。時計。心臓の鼓動（近似的に）。情報はパターンのなかに蓄えられるが、けっして変換されない。周期------有限。

\textbf{クラス3------フラクタル。} あらゆるスケールで自己相似な構造を生み出す系。シェルピンスキーの三角形。シダ。マンデルブロ集合。数学的に豊かで、美的には息をのむほどだが、\textbf{計算的に還元可能}だ------あらゆるステップを走らせずとも、先へ跳べる。処理能力なき構造。周期------準無限であり、あらゆるスケールに正確な、あるいは統計的な自己相似性を伴う。

\textbf{クラス4------複雑（カオスの縁）。} 動き、相互作用し、任意の計算を符号化できる、持続的で局在化した構造を生み出す系。コンウェイのライフゲーム。皮質のオートマトン。計算的に\textbf{還元不可能}------近道はない。これらの系は普遍計算の能力を備えている。適切な初期条件が与えられれば、自分自身のシミュレーションを含めて、いかなるアルゴリズムをもシミュレートできる。周期------準無限であり、自己相似性に\textbf{加えて}持続的に相互作用する構造を伴う。意識が住まうのは、ここである。

\textbf{クラス5------ランダム。} その出力が真にランダムな系------擬似ランダムではなく、決定論的でもなく、圧縮不可能である。パターンもなく、自己相似性もなく、やがて反復する周期もない。真に無限の情報量。構造------\textbf{不明}（後述）。

ウォルフラムの枠組みへの対応づけはこうなる。


{\small
\noindent
\begin{tabularx}{\linewidth}{>{\centering\arraybackslash}X >{\centering\arraybackslash}X X}
\toprule
\textbf{五クラス} & \textbf{ウォルフラム} & \textbf{何が変わったか} \\
\midrule
1 & クラス1 & 同じ \\[4pt]
2 & クラス2 & 同じ \\[4pt]
3 & クラス3（一部） & ウォルフラムのクラス3から分離 \\[4pt]
4 & クラス4 & 同じ \\[4pt]
5 & クラス3（一部） & ウォルフラムのクラス3から分離 \\[4pt]
\bottomrule
\end{tabularx}
}


ウォルフラムの無秩序スペクトル上の順序はこうだった------1$\rightarrow$2$\rightarrow$4$\rightarrow$3。不格好である。五クラスの枠組みは、きれいな単調勾配を与える------1$\rightarrow$2$\rightarrow$3$\rightarrow$4$\rightarrow$5、無秩序の増大と計算的還元不可能性の増大とに沿って並ぶ。

\section*{なぜ決定論的オートマトンはランダムさを生み出せないのか}

ここで、私が独創的だと考えている論証を示す。四クラスではなく五クラスであるべきだという主張を、これが補強する。

一つのセル・オートマトンを------どんなセル・オートマトンでもよい------考えてみよう。それは有限の規則表（有限個のビットで表せる）と有限の初期条件（同じく有限個のビットで表せる）を持つ。規則と初期条件はあわせて、固定した有限の情報量を含んでいる。

さて------有限の情報量は、真にランダムな出力を生み出せるだろうか。

否である。その理由はこうだ------

\begin{enumerate}
  \item 真にランダムな無限列は\textbf{最大の}コルモゴロフ複雑性を持つ------それは圧縮できず、自分自身より短い何ものによっても記述できない。
  \item セル・オートマトンの出力は、その規則と初期条件によって完全に決定される。それらはあわせて\textbf{有限の}コルモゴロフ複雑性しか持たない。
  \item あるプロセスから、その仕様として投入した以上の情報を取り出すことはできない。
  \item したがって、いかなるセル・オートマトンの出力も、同じ長さの真にランダムな列に比して低いコルモゴロフ複雑性を持つ。
\end{enumerate}

これは一般化された鳩の巣論法である------有限の情報は、自己相似的な構造を生み出さざるをえない。有限の情報から無限の出力を生成する唯一の道は、異なるスケールで構造を\textbf{再利用する}ことだ。厳密な再利用は周期性（クラス2）となる。厳密ではないが規則性を帯びた再利用はフラクタルなふるまい（クラス3）となる。最も複雑に見えるセル・オートマトン（ルール30、ルール110、ライフゲーム）でさえ、その出力の複雑性は自らの規則集合の複雑性によって上限を課されているのだ。

ウォルフラムが「ランダムな」セル・オートマトンと呼んだものは、\textbf{高複雑性のフラクタル}と表現するほうが的確である------その自己相似的な構造は実在するが、ざっと眺めただけでは見えないほどのスケールと次元で作動しているシステムなのだ。実際、ルール30の左端はシェルピンスキー的な下位構造を示す。その中央の列はランダム性の統計検定の多くに合格するが、これは\textbf{高複雑性のフラクタルにまさに期待されること}である------局所的な統計はランダムさを模倣するが、大域的な構造は決定論的で、コルモゴロフ複雑性が低い------短い規則がそのすべてを生成するのだ。（ここで「圧縮可能」とはこの意味である。すなわち、コンパクトな記述が存在するということだ。ステップごとの計算を近道できるという意味ではない------その計算は依然として還元不可能である。）

この同じ論証によれば、クラス4の出力も\textbf{また}フラクタルである------ライフゲームは、その個体群動態、構造分布、空間相関において統計的な自己相似性を示す。クラス3とクラス4の違いは、「フラクタルか否か」ではない。それはこうだ------

\begin{itemize}
  \item \textbf{クラス3}------フラクタル。還元可能。処理を伴わない構造。
  \item \textbf{クラス4}------フラクタル。還元不可能。処理を\textbf{伴う}構造------互いに相互作用し、普遍計算を符号化しうる、持続的で局在化した構造。
\end{itemize}

どちらもフラクタルだ。だが計算するのは一方だけである。


{\small
\noindent
\begin{tabularx}{\linewidth}{>{\centering\arraybackslash}X X X X >{\centering\arraybackslash}X >{\centering\arraybackslash}X}
\toprule
\textbf{クラス} & \textbf{規則} & \textbf{周期} & \textbf{構造} & \textbf{還元可能か？} & \textbf{計算するか？} \\
\midrule
1 & 有限 & 1 & なし & 自明に & 否 \\[4pt]
2 & 有限 & 有限 & 反復 & 可 & 否 \\[4pt]
3 & 有限 & 準無限、自己相似 & 自己相似 & 可 & 否 \\[4pt]
4 & 有限 & 準無限、自己相似 & 自己相似＋持続的に相互作用する構造 & 不可 & \textbf{可} \\[4pt]
5 & 表現不可能 & 真に無限 & \textbf{不明} & N/A & N/A \\[4pt]
\bottomrule
\end{tabularx}
}


\section*{クラス5と数学的表現可能性の境界}

決定論的オートマトンが真のランダムさを生み出せないのであれば、では何がそれを生み出せるのか。

この問いは、五クラスの枠組みの最も深い含意だと私が考えるものへと導く。

クラス1からクラス4までは、有限で表現可能な規則が生み出しうるものである。そのふるまいは自明なもの（クラス1）から並外れたもの（クラス4------普遍計算、意識）にまで及ぶが、そのすべては、書き下し、伝達し、検証し、分析できる規則によって生成される。これらの規則は数学のなかに、形式的な記号システムの領域のなかに宿っている。

対照的に、クラス5は書き下すことの\textbf{できない}規則を要求する。真にランダムな出力を------最大のコルモゴロフ複雑性を持ち、圧縮不可能で、非アルゴリズム的な出力を------生み出すシステムは、形式システムが表現しうるいかなる規則も走らせてはいない。もし規則が表現可能なら、出力は圧縮可能であり（「この規則を適用せよ」へと）、したがって真にランダムではなくなる。

これはクラス5を、数学的表現可能性そのものの境界に位置づける。それはたんに「きわめて複雑」あるいは「きわめて無秩序」というのではない。それは、生成プロセスが線形の記号システム（数学、論理、計算）の捉えうるものを超えてしまうレジームなのだ。

自然界に、実際にクラス5で作動するものはあるのか。

あるかもしれない。量子力学は、ベルの定理によって、いかなる局所的な隠れた変数の理論によっても説明できない測定結果を生み出す。もしこれらの結果が真にランダムであるなら------まだ特定できていない決定論的プロセスなのではないのなら------そのとき量子測定はクラス5のプロセスとなる。すなわち、その規則を私たちのいかなる形式言語でも書き下せない物理現象である。

これは思弁的であり、そのように明示しておく。だがその含意は鮮烈だ------クラス4（意識のレジーム、普遍計算のレジーム、皮質のオートマトンのレジーム）は、\textbf{表現可能な規則によって到達しうる最大の複雑性}に位置している。数学がなしうるかぎり複雑なのだ。その彼方には、数学がその本性ゆえに描き出せない領域が広がっている。

\section*{クラス5の構造------不在ではなく、不明}

最後にもう一つ、微妙な点がある。クラス5には「構造がない」と言いたくなる。だがそれは誤りだ------無限には構造がない、と言うのと同じ誤りである。

1870年代のゲオルク・カントールの仕事以前、無限は単一の概念として扱われていた。ものは有限か無限かのどちらかであり、それで話は終わりだった。カントールは、無限には\textbf{階層}があることを示した------実数の無限が整数の無限より厳密に大きいこと、そしてこの階層が果てしなく続くことを。無限には、豊かな内的構造があると判明したのだ。それが見えなかったのは、数学者たちにそれを見るための道具が欠けていたからにすぎない。

同じことが、ランダムさについても言えるかもしれない。私たちは現在、真のランダムさを単一のカテゴリーとして扱っている------最大の無秩序、パターンの不在として。だが私たちは、無限を眺めるカントール以前の数学者と同じ立場にいる。異なる種類のランダムさを------もしそうした区別が存在するのであれば------見分けるための概念的道具を、私たちは欠いているのだ。

したがって、クラス5の構造についての誠実な答えはこうなる------\textbf{不明}、と。「なし」ではない。「不在」でもない。不明------まだ存在しないかもしれない概念的道具を、線形の記号システムには提供できない思考法をおそらく必要とする道具を、待っているのである。

これは、数学と物理学と計算の交差点にある、最も重要な未解決の問いの一つだと私は考えている。そしてそれは、五クラスの枠組みなしには見えないままだ。ウォルフラムの四クラスの枠組みは、その問いを立てられる余地をそもそも生み出さないからである。

\section*{意識への含意}

五クラスの枠組みは、なぜ意識がクラス4のダイナミクスを------そしてクラス4だけを------必要とするのかを明らかにする。

クラス1とクラス2は単純すぎる。それらは情報を蓄えることはできる（固定した状態、反復するパターン）が、計算的に興味深いかたちで\textbf{処理}することはできない。深い睡眠のなかで遅いデルタ波を流している脳は、クラス2で作動している------周期的で、反復的で、どこにも進まない。四モデルのアーキテクチャは基盤のなかで無傷のまま保たれているが、シミュレーションは走っていないのだ。

クラス3は興味深いが、計算には向かない。フラクタルなダイナミクスは豊かな構造を生み出し、脳もそれを用いている（後述）------だがそれは、意識的な自己シミュレーションが要求するような、動的で、還元不可能で、大域的に統合された処理を維持することはできない。フラクタルなパターンは美しいが、計算的には還元可能だ。それは自分自身を驚かせることができない。

クラス4は、意識が必要とする二つの性質をまさに備えている------\textbf{普遍計算}（システムは原理的に何でも、自分自身をも含めてシミュレートできる）と\textbf{大域的統合}（システムの離れた部分どうしが互いに影響し合い、局所的な変化が大域的に伝播し、情報が統一された全体へと束ねられる）である。カオスの縁において、皮質のオートマトンはその両方を達成し、その結果が意識なのだ。

クラス5が異なるのは、そこで計算が不可能だからではなく（無限のランダム列は\textbf{あらゆるもの}を含んでいる------あらゆる安定したパターンも、これまで構想されたあらゆる計算も含めて）、それを利用し、予測し、実証する術が私たちにはないからである。全般発作のなかで、ニューロンが協調を欠いたカオスとして発火している脳は、クラス5に近づく------そのレジームで意識が原理的に不可能だからではなく、それを維持し、あるいはそこへアクセスするための機構が存在しないからだ。私たちの宇宙そのものが、無限のランダムさの一断片かもしれないし、私たちの知覚が及ばないスケールにあるクラス4のシステムかもしれないし、あるいは無限のフラクタルの一断片かもしれない。それは知りようがない。\textbf{言える}のは、私たちが経験しているとおりの意識が、クラス4の構造化された予測不可能性を必要とする、ということだ。クラス5でシミュレーションが崩壊するのは、その根底にある現実が不十分だからではなく、基盤とシミュレーションのあいだに安定したインターフェースが存在しないからである。

\section*{脳は四つのクラスすべてを用いる}

脳は、数十億年の進化によって最適化された普遍計算機である。もし進化が、優位性をもたらす計算レジームをどれか一つでも見落としていたなら、それは奇妙なことだろう。そして実際、脳は表現可能な四つのクラスすべてを、それぞれ異なる道具として用いている。

\begin{itemize}
  \item \textbf{クラス1}（安定したアトラクター）------長期記憶の貯蔵。何年ものあいだ持続するシナプスの重みの配置。ニューラルネットワークの不動点。
  \item \textbf{クラス2}（振動）------アルファ、シータ、ガンマ、デルタのリズム。視床によるクロック信号の生成。睡眠・覚醒サイクル。脳の計時とゲーティングの機構。
  \item \textbf{クラス3}（フラクタル／スケール不変な処理）------テクスチャ解析、スケール不変な物体認識、効率的な神経符号化。主にV2からV4の視覚処理であり、そこでは多スケール比較が中核的な操作となる。サイケデリック下では、この機構が外部入力なしに走るとき、フラクタルな処理そのものが\textbf{見える}------だからこそフラクタルなパターンは、サイケデリック体験の最も一貫した特徴の一つなのである（第6章参照）。
  \item \textbf{クラス4}（カオスの縁）------皮質のオートマトンそのもの。意識的処理の動的レジーム。普遍計算。シミュレーションのエンジン。
\end{itemize}

各クラスは、それぞれ異なる機能を担う。意識を生み出すのはクラス4だけだ。だが意識は、ほかのクラスに依存している------モデルを住まわせる安定した記憶（クラス1）、ダイナミクスを協調させるリズミカルな計時（クラス2）、そして世界を複数のスケールで同時に解析するフラクタルな処理（クラス3）に。

これこそが、脳が特にカオスの縁で作動しなければならない、おそらく最も深い理由である------クラス4は、ほかのあらゆるクラスを------そして自分自身をも------\textbf{動員}できる唯一のレジームなのだ。クラス4のオートマトンは、安定した状態（クラス1のふるまい）、周期的な振動（クラス2のふるまい）、フラクタルな構造（クラス3のふるまい）を、自らのダイナミクスのなかの部分プロセスとして生み出すことができる。そして、ほかのクラス4オートマトンを生み出すこともできる------普遍計算機は、別の普遍計算機をシミュレートできるのだ。ほかのどのクラスも、こうしたことは一つとしてできない。クラス4はたんに最も複雑なクラスであるだけではない------それは、自分自身を含めてすべてのクラスを内包する唯一のクラスなのである。この自己包含こそが、スケールを越えた構造的同一性を可能にする------クラス4の宇宙のなかのクラス4の脳は、アナロジーではない。それは、同じ計算アーキテクチャの入れ子になった一つの事例なのだ。


\chapter*{付録D：明晰夢の見かた}
\addcontentsline{toc}{chapter}{付録D：明晰夢の見かた}
\markboth{付録D：明晰夢の見かた}{}

\textbf{第6章で、明晰夢を、意識を内側から探索するための安全で薬物に頼らない方法として挙げた。四モデル理論（FMT）において、明晰夢とは、レム睡眠のあいだに明示的自己モデル（ESM）がより完全に「起動」する現象である------受動的な夢を制御された体験へと変える、臨界性の閾値を越える出来事だ。以下に述べるのは、そこへ至るための最も手軽な方法である。}

\section*{現実チェック法}

原理は単純だ------自分が起きているかどうかを習慣的に問うようにしていれば、その習慣はやがて夢のなかでも発火し、夢はそのテストに落第する。

\textbf{ステップ1:現実チェックを一つ選ぶ。} 最も信頼できるのは、文字を見て、目をそらし、また見直すことだ。覚醒時には文字は変わらない。夢のなかでは変わる------しばしば劇的に。時計も使える------時刻を確かめ、目をそらし、もう一度確かめる。夢のなかでは、数字が違っていたり、意味をなさなくなっていたりする。もう一つ信頼できるチェックは、片方の手のひらに反対側の指を押し通そうとすることだ。夢のなかでは、指はしばしばすり抜けてしまう。

\textbf{ステップ2:一日じゅう実行する。} 戸口をくぐるたび、スマートフォンを確かめるたび、あるいは何かがわずかに奇妙だと気づくたびに、立ち止まって現実チェックを行ってみてほしい。肝心なのはチェックそのものではなく、その背後にある\textbf{本気の問い}だ------「自分はいま夢を見ているのか？」。ただ動作をなぞるだけではいけない。その可能性を実際に考えてみてほしい。

\textbf{ステップ3:夢日記をつける。} ノートをベッドのそばに置いておく。毎朝、身体を動かす前に、覚えていることを何でも書き留める------たとえそれがある感覚や一枚のイメージにすぎなくても。これは、夢の内容を記憶に値するものとして扱うよう脳を訓練し、夢と覚醒時の気づきとを結ぶ橋を強化する。

\textbf{ステップ4:待つ。} たいていの人にとって、最初の明晰夢は2週間から6週間のうちに訪れる。夢のなかにいて、何かがわずかにおかしく思われ、現実チェックの習慣が発動し、文字が変わる------そして、\textbf{わかる}のだ。その、わかる瞬間こそがESMの起動である。移行は感じ取れる------突然の鮮明化、現前（プレゼンス）の感覚、そして自分を取り巻く世界が、そのなかで自分が意識をもっている一つのシミュレーションなのだという静かな気づきが。

\section*{何を期待できるか}

最初の明晰夢は、おそらく短いだろう------数秒から数分。興奮すると目が覚めてしまいがちだ。練習を重ねれば、引き延ばせるようになる。週に何度も明晰夢を見られるようになる人もいる。この体験は驚くべきものだ------外部からの入力がまったくないまま、完全な意識的シミュレーションのなかにいて、しかもそのことを知っている。仮想の世界が、こちらの意図に応えてくる。それはまさしく文字どおり、体験可能なものとなった四モデル理論なのである。

\section*{ほかの方法}

さらに先へ進みたい読者のために、もっと手のこんだ技法がある。

\begin{itemize}
  \item \textbf{MILD（記憶による明晰夢誘発法）}------スタンフォード大学のスティーヴン・ラバージ（Stephen LaBerge）が開発した。眠りに落ちるとき、いま夢を見ているのだと気づこうという意図を設定する。5時間後に一度目を覚まし、ふたたび眠りに戻る方法と組み合わせるのが最もよい。
  \item \textbf{WILD（覚醒開始型明晰夢）}------覚醒から入眠への移行のあいだ、気づきを保ちつづける。難しいが、最も鮮明な結果をもたらす。
  \item \textbf{WBTB（いったん起きてまた寝る法）}------5時間から6時間の睡眠のあとに目を覚まし、20分から60分のあいだ起きていて、それから眠りに戻る。これはレムの豊富な睡眠後半の周期を狙うものだ。
\end{itemize}

スティーヴン・ラバージの『Exploring the World of Lucid Dreaming』（1990）は、いまなお決定版の実践ガイドである。神経科学については、明晰夢のEEG的署名に関するVoss et al. (2009)、そして明晰夢の認知神経科学を網羅的に概観したBaird et al. (2019) を参照されたい。


\chapter*{付録E：なぜ「四つ」のモデルなのか------神経科学者への覚書}
\addcontentsline{toc}{chapter}{付録E：なぜ「四つ」のモデルなのか------神経科学者への覚書}
\markboth{付録E：なぜ「四つ」のモデルなのか------神経科学者への覚書}{}

この付録では、神経科学者や計算に通じた読者が、第2章で四モデルアーキテクチャに出会ったときに必ず抱くであろう疑問を取り上げる。\textbf{脳がちょうど四つのモデルを保持しているなどということは、まさかあるまい？}

その通り、そんなことはない。「四」という数は\textbf{根拠のある最小値}であって、文字どおりの個数ではない。

\section*{脳が実際に行っていること}

生物学的な基盤------プロテオームネットワークの上で発火するニューロン、そして単一の細胞の内部でさえ独自の計算的知性を成す細胞内シグナル伝達経路------は、暗黙的/明示的の境界の両側にわたって、事実上数えきれないほどの、互いに重なり合うモデルを実装している。

カップに手を伸ばす場面を考えてみてほしい。運動モデルは、世界の幾何学（カップがどこにあるか、その周囲にどんな障害物があるか）と、自己の運動学（腕がどう構えられているか、握るために指をどう形づくるべきか）を同時に符号化する。この一つのモデルは、純粋な世界モデル\textbf{でもなく}、純粋な自己モデル\textbf{でもない}------両方のカテゴリーにまたがってにじみ出ている。社会的なやり取りの感情モデルは、相手についての知識（世界）と、その中での自分の評価（自己）を同時に符号化する。空間ナビゲーションのモデルは、環境の配置と、その中での自分の位置の両方を符号化する。現実の神経モデルはどれも混合体である。

「モデル」どうしの境界は明確ではなく、その数も固定されておらず------そして確かに四つではない。

\section*{それでも四が正しい抽象である理由}

四つの正典的モデル（IWM、ISM、EWM、ESM）は、二つの軸によって張られる連続的な二次元空間における\textbf{極点}である。

\begin{itemize}
  \item \textbf{範囲}:純粋な自己表象から純粋な世界表象まで
  \item \textbf{モード}:完全に暗黙的（構造的、貯蔵された、無意識的）なものから、完全に明示的（シミュレートされた、一過性の、現象的）なものまで
\end{itemize}

脳の実際のモデリング生態系は、この空間全体を、互いに重なり合うモデルの連続的な密度で満たしている。名づけられた四つのモデルは、その四隅------活動がそのまわりに組織される理論的な極である。それらは方位のようなものだと考えてほしい。ナビゲーションには役立ち、方角としては実在するが、世界には四つの場所しか存在しないなどと主張する者はいないだろう。

理論が完全な連続空間ではなくこれら四つの極の上に築かれている理由は、それらが、システムが意識をもつために必要な\textbf{最小構成}を指し示すからである。

\begin{itemize}
  \item \textbf{世界モデルがない}$\rightarrow$体験すべき環境がない
  \item \textbf{自己モデルがない}$\rightarrow$それを体験する主体がない
  \item \textbf{暗黙的なレベルがない}$\rightarrow$シミュレートする元がない（学習された知識がない）
  \item \textbf{明示的なレベルがない}$\rightarrow$シミュレーションがまったくない（体験がない）
\end{itemize}

四つのうちどれか一つでも取り除けば、何か決定的なものが壊れる。四つのモデルは床であって、天井ではない。脳はあらゆる方向でそれを超えている。しかし意識が何を\textbf{必要とする}かを教えてくれるのは、その床である------そして理論の予測を生み出し、その主張を制約し、いかなる人工システムであれ実装しなければならないものを規定するのも、その床なのである。

\section*{本書の残りを読むにあたって}

本書を通じて、私が「ESMはこれを行う」あるいは「IWMはあれを含む」と書くとき、それはこの連続空間の極を指しているのであって、脳が壁で仕切られた四つの別々の箱をもつと主張しているのではない。この単純化には根拠があり、続く各章は、それが本物の説明力を発揮するさまを示すことになる------このアーキテクチャの上に築かれた五つの原理から、サイケデリックの現象学、麻酔のメカニズム、夢の状態、分離脳現象、そして動物の意識を導き出すのである。

完全な数学的取り扱い------連続的なモデル空間の枠組み、モデル密度関数、そしてこの空間の諸領域間の情報伝達としての透過性の定式化を含む------については、Gruber (2026), \textit{Toward a Mathematical Formalization of the Four-Model Theory} を参照されたい。


\chapter*{付録F：境界の帳簿づけとしての標準模型}
\addcontentsline{toc}{chapter}{付録F：境界の帳簿づけとしての標準模型}
\markboth{付録F：境界の帳簿づけとしての標準模型}{}

\textbf{計算的原子という描像の背後にある導出のうち二つ------なぜ保存則は絶対なのか、そしてなぜ粒子の世代はちょうど三つなのか------は、この論全体の中でもっとも無味乾燥な帳簿づけである。そのため本文からは切り離しておいた。素粒子物理学の細部まで知りたい読者のために、ここに全文を収めておく。生き生きとした版は「万物のアーキテクチャ」の章にある。こちらはその背後にある元帳である。}

\textbf{なぜ保存則はこれほどまでに絶対なのか？} 電荷はつねに保存される。バリオン数は保存される。レプトン数は保存される。これらの法則が破れたことは、これまで行われたいかなる実験においても、ただの一度たりとも観測されていない。なぜか？

それらが情報保存の制約だからである。ベケンシュタイン限界は、一つの境界がどれだけの情報を保持しうるかを教えてくれる。二つの境界が相互作用して情報を交換するとき、情報の総量は保存される------保存されなければならない。なぜなら情報保存は量子力学のユニタリー性から導かれる帰結であり、そしてユニタリー性そのものが、ベケンシュタイン限界に裏打ちされているのかもしれないからである。素粒子物理学における個々の保存則------電荷保存、バリオン数保存、レプトン数保存------は、境界の配位が相互作用するとき情報がどう変換されうるかを支配する、具体的な規則である。それらは外部から押しつけられた恣意的な規則ではない。特異点の境界において情報を創り出すことも消し去ることもできない、という事実から導かれる帳簿づけの制約なのである。

\textbf{そしてもう一つ、三世代という謎がある。} 粒子は三つの世代をなして現れる。電子には、より重い複製（ミュー粒子）と、さらに重い複製（タウ粒子）がある。アップクォークには、チャームおよびトップと呼ばれる複製がある。粒子の各型につき三つの版があり、質量を除くあらゆる性質において同一である。これは素粒子物理学における、もっとも深い未解明のパターンの一つである。なぜ三なのか、誰にもわからない。二でも、四でも、十七でもなく。三なのである。

包み隠さず言っておこう。これから述べることは思弁的である。計算的原子という描像の残りの部分よりも、なお思弁的だ。だがそれは構造の側から動機づけられており、俎上に載せる価値はあると私は考える。

クラス4のシステムは、本質的に自己相似的な構造を内包している。これは、クラス4の力学がクラス3（フラクタル）のふるまいを部分プロセスとして含む、という事実から導かれる技術的な帰結である。自己相似性とは、同じパターンが異なるスケールで反復されることを意味する。特異点の境界配位がクラス4のシステムに埋め込まれているならば------そして埋め込まれているはずである、宇宙はクラス4なのだから------そのとき配位そのものが自己相似的な構造を示すかもしれない。同じ型の境界が、三つの異なるエネルギースケールに現れる。三世代とは、安定な特異点配位の空間におけるフラクタル的な階層構造の徴候なのかもしれない。

私は証明をもってはいない。これは推測であって、導出ではない。だが、三というのは、最も単純で非自明な自己相似的階層------一つの基底配位と、二つのスケール変換された複製------から予期されるものとちょうど一致する、ということは記しておく。そして世代構造は、それ以外の点では現行のいかなる理論によっても完全には説明されていない、ということも記しておく。もし計算的原子という描像が、いつの日か、なぜ世代がちょうど三つなのかを説明するに至るなら、それはこの枠組み全体にとって強力な証拠となるだろう。

そして宇宙が本当にセル・オートマトンであるなら、第四の世代が存在する------第五も、さらにその先も。だが高次の世代はより大きなパターンであり、より大きなパターンはより不安定である。それらは、より小さく、より安定な配位が自らを切り裂いてくる猛攻に耐えられず、きちんと成り立つ前に断ち切られてしまう。まさにこれが、私たちの観測している事態である。第三世代の粒子（タウ粒子、トップクォーク）はすでに極端に不安定で、ほとんど瞬時に崩壊する。第四世代はさらに重く、その境界配位はより大きく、より複雑で、それゆえいっそう脆いだろう------一過性のゆらぎではなく一つの粒子として検出されるほど長くはとどまれないほど、脆いのである。私たちの目にする三つの世代とは、生き延びられるだけの小ささをもった、その三つにすぎないのかもしれない。


\chapter*{付録G：宇宙論モデルにおける四つの弱点}
\addcontentsline{toc}{chapter}{付録G：宇宙論モデルにおける四つの弱点}
\markboth{付録G：宇宙論モデルにおける四つの弱点}{}

\textbf{自らの弱点を隠す理論は、理論ではない。「万物のアーキテクチャ」と「最も深い鏡」で展開した宇宙論的な議論は、証明されていない主張の上に立っており、それがどこで崩れうるのかを名指しすることが、誠実さの求めるところである。五つ目にして最悪の弱点------認知の天井、すなわちクラス4の脳は結局のところ自らのアーキテクチャを宇宙へ投影しているにすぎないのではないかという懸念------は、本文にとどめておく。それは私が答えられない唯一のものだからだ。残る四つは、それらに攻撃を仕掛けたいと願うすべての人のために、反論とともにここに掲げる。}

\textbf{その一:エネルギーは情報に等しい、という主張は証明されていない。} これはランダウアーの原理、ベケンシュタイン限界、そしてブラックホール熱力学によって強く示唆されている。互いに独立した複数の証拠の筋が、すべて同じ方向を指し示している。だが、E=Iを第一原理から導き出した者はいまだにいない。情報がそれ自体で重力効果をもつことを示した者もいない。この仮説は説得力があり、証拠の収束は目を見張るものがある。しかし収束は証明ではない。もしエネルギーと情報が同一ではなく、単に相関しているだけだと判明すれば、特異点を情報変換として解釈する見方はその基盤を失う。意識と宇宙論のあいだの構造的な対応づけはなお成り立つだろう（五つの対応はE=Iに依存していない）が、両者を結ぶ物理的なメカニズムははるかに脆弱なものとなる。詩は生き残るだろう。物理学のほうは、そうではないかもしれない。

\textbf{その二:クラス4の議論は演繹的ではなく、アブダクション的である。} 私は他のクラスを消去したが、消去は証明ではない。それは最良の説明への推論------科学において完全に敬意に値する推論の一形式ではあるが、一つの扉を開けたままにしておく形式でもある。ひょっとすると、私が思い描けなかったクラス4.5があるのかもしれない。複雑さとランダムさのあいだに横たわる計算のレジームで、それを記述するための概念的な道具を私がもち合わせていないだけなのかもしれない。ひょっとすると、五つのクラスの階層それ自体が不完全なのかもしれない。この議論が言うのは、クラス4が宇宙のダイナミクスの最良の説明である、ということであって、論理的に可能な唯一の説明である、ということではない。アブダクション的な推論は優れた実績をもっている------ダーウィンの自然選択の議論はアブダクション的であり、それはかなり見事に持ちこたえた------が、それは数学的証明と同じではないし、私はそうであるかのように装うつもりはない。

\textbf{その三:特異点の統一には量子重力が必要である。} すべての特異点（プランクスケール、ブラックホール、宇宙論的なもの）が、同一の情報境界の構造的に等しい実例である、という主張は、強い主張である。それには量子力学と一般相対性理論を架橋する理論が必要となる。私たちはそれをもっていない。弦理論は候補の一つである。ループ量子重力も候補の一つである。どちらもこの主張と整合的であり、どちらも未確認である。特異点の統一は荒唐無稽な思弁ではない------それは現代物理学が向かいつつある方向である------が、それはまた確立された物理学でもない。それは未来への賭けである。私はこの賭けを良いものだと考えている。私は間違っているかもしれない。

\textbf{その四:このモデルは内側からは反証不可能かもしれない------それ自身の予測によって。} これは、考えるたびに胃が締めつけられる一点だ。自己言及的閉包のゲーデル的帰結は、自らを計算する十分に複雑なシステムは、内側から自分自身についてのあらゆる真理を証明することはできない、と述べる。真であるにもかかわらず証明できない言明が存在するのだ。もしSB-HC4Aが正しいなら、宇宙はまさにそのようなシステムであり、「SB-HC4Aは正しい」という言明は、真だが証明不可能なものの一つでありうる。この理論は、宇宙の内側からそれを証明することが構造的に不可能でありうる、と予測する。私たちの努力が足りないからではない。より優れた計器が必要だからでもない。アーキテクチャがそれを不可能にしているからだ------ちょうど、公理の体系が自らの無矛盾性を証明できないのと同じように。

これは科学哲学における最も深遠な結論であるか、さもなければこれまでに考案された最も優雅な逃げ口上であるかの、どちらかである。自らの検証不可能性を予測する理論は、あなたに知の限界についての何か深いことを告げているか、さもなければ、読者を深く疑わせるはずのやり方で自らを批判から免疫化しているかの、どちらかだ。正直なところ？私にはどちらなのか確信がない。この問いを何年も考え続けてきたが、いまだにわからない。私にわかっているのは、この予測がどこからともなく現れたのではない、ということだ------それはゲーデルの定理から来ている。ゲーデルの定理は、数学におけるいかなるものにも劣らず堅固である。もし宇宙が自己言及的なら、不完全性が帰結する。問いは、宇宙が自己言及的かどうかである。そして理論は、そうだと言うのだ。

だから、これが本物の弱点であって、こっそり通そうとしている仕掛けなどではない、と私が言うとき、それを信じてほしい。もし誰かが宇宙の内側からSB-HC4Aを検証する方法を見つけ、その検証が失敗すれば、理論は死ぬ。もしその検証が成功すれば、皮肉なことに、理論はやはり間違っていることになる------なぜなら検証の成功は、宇宙の計算構造が内側からアクセス可能であることを意味し、それはまさにアーキテクチャを定義するその境界条件と矛盾するからだ。いかなる検証も不可能である場合にのみ、理論は生きる。だがそれが生きるのは、奇妙な哲学的薄明のなかであり、反証不可能なのは回避によってではなく、構造によってなのである。


\chapter*{付録H：九つの予測}
\addcontentsline{toc}{chapter}{付録H：九つの予測}
\markboth{付録H：九つの予測}{}

\textbf{すべてを説明し、何も予測しない理論は、理論ではない------物語である。以下に、反証実験の一式をまとめて示す。四モデル理論（FMT）がなす九つの予測すべてを、それぞれを葬り去りうる具体的な実験とともに並べる。いくつかは既存の技術で今日にも実施でき、そのうち数個にはすでに初期的な裏づけがある。反証されたものは一つもない。予測3は第11章で一つの場面として展開されており、ここでは簡潔に述べるにとどめる。ほかのいくつかは、それがふさわしい各章で導入されたものを、理論を厳しく試してみたい読者のために、一か所に集めたものである。}

\section*{予測1：各モデルはそれぞれ固有の神経的シグネチャをもつ}

四つのモデルが本当に別個のプロセスであるなら、脳スキャンでそれを見て取れるはずだ。四つの異なる種類の課題（それぞれが一つのモデルを働かせる）を人々に行わせる巧妙な実験を設計すれば、脳の活性化パターンは互いに異なって見えるはずである。

IWM優位の課題とは、たとえば見慣れた顔を受動的に認識するようなものだろう。それについて考えているわけではない。脳がただ知っているのだ。ISM優位の課題は、習慣化された運動系列------たとえば、一つ一つのキーを意識せずにパスワードを打ち込むこと------でありうる。EWM優位の課題は、能動的で意識的な知覚を要する------ほとんど同一の二枚の画像のあいだの違いを見つけようとする、といったものだ。そしてESM優位の課題は、純粋な自己内省である------「自分はそんなことをするような人間だろうか？」。

予測されるのは2×2のパターンである。世界の課題対自己の課題。暗黙的対明示的。四つの象限、四つの別個の神経的シグネチャ。そのパターンが見つからないなら、理論のどこかが間違っている。

これは既存のfMRI技術で今すぐ検証できる。安価ではないし、周到な実験設計を要するが、道具立てはすでに世界じゅうの研究室にある。そしてもしうまくいけば、四モデルアーキテクチャが単なる比喩ではないこと------脳が情報を処理する仕方そのものに刻み込まれた、実在する機能的区別であることの、最も直接的な証拠となるだろう。

\section*{予測2：サイケデリックの視覚体験は脳の処理層を露わにする}

これはエレガントな予測である。サイケデリックの作用下では、体験される視覚的内容は、用量に応じて特定の順序で脳の視覚処理階層をたどっていくはずだ。

低用量では、眼内閃光（フォスフェン）------目を閉じたときに現れる、あの小さなきらめきや幾何学的な形------が見える。それはV1、最も初期の視覚処理領域が意識へと漏れ出しているのだ。用量を増やすと、より複雑な幾何学模様------文化や物質を越えて現れる、あの有名な「形態定数」------が現れる。V2とV3が働き始めているのである。さらに上げると、顔や人影、複雑な情景が見えはじめる。より高次の視覚領域だ。最高用量では、意味と物語をそなえた、完全に夢のような物語的体験が生じる。

予測されるのは、これがランダムではないということである。それは用量依存的で、視覚階層を上っていく秩序だった進行なのだ。暗黙的-明示的透過性が高まるにつれ、視覚処理のより深い層が意識にのぼる。脳の内部配線図が、自分自身の体験のうちに可視化されるのである。

これは段階的な投与プロトコルで検証できる------注意深く管理された量のシロシビンやLSDを人々に与え、fMRIで脳をスキャンし、何が見えているかを尋ねる。報告された内容と脳の活性化とを照合するのだ。用量が増すにつれて処理階層が下から上へと点灯していくのが見えるはずだ、と理論は予測する。

\section*{予測3：自我崩壊のあいだ、その人が何になるかは操作できる}

\textbf{（第11章で一つの場面として展開されている。）} 高用量のサイケデリックによる自我崩壊のあいだ、自分が何へと\textbf{溶け込むか}の内容は、感覚環境によって操作できる------ランダムでもなく、純粋に生化学的でもない。明示的自己モデル（ESM）は、いつもの暗黙的自己モデル（ISM）からの入力を断ち切られると、最も優勢な感覚入力に貼りついていく。部屋を海の音と青い光で満たせば被験者は海になり、鳥のさえずりと緑の光に切り替えれば森になる。この予測には方向性がある------自我崩壊のあいだに優勢な感覚入力を変化させれば、報告される同一性の内容がそれを追う。基本的な環境制御をそなえたどんなサイケデリック研究室でも、今日にも検証できる。ほかのどの意識理論もこの予測をしていない。

\section*{予測4：サイケデリックは脳卒中患者が自らの欠損を見て取る助けになるはずだ}

病態失認は、脳がなす最も奇妙な現象の一つである。ある種の脳卒中（たいていは右半球の）ののち、患者は身体の片側が麻痺しているにもかかわらず、そのことを本気で信じていない。動かない腕を本人に見せ、動かすよう頼み、失敗するさまを見せることができるが、それでも患者は言い逃れを作話する。「疲れているから」。「気が乗らないから」。彼らは嘘をついているのではない。欠損が本当に見えていないのだ。

四モデル理論は、これが透過性の遮断のために起こると言う。麻痺についての情報は暗黙的自己モデル（ISM）のなかにある------基盤は知っている------のだが、それが明示的自己モデル（ESM）まで届かないのだ。シミュレーションは、基盤のもつ知識のその部分にアクセスできていない。

さて、ここからが意外なところである。サイケデリックは、暗黙的-明示的透過性を全体的に\textbf{高める}。それがサイケデリックのすることなのだ。そこで予測されるのは、自我崩壊には至らない用量のシロシビン------自己を溶かすには足りないが、透過性のゲートを開くにはちょうど足りる量------であれば、欠損の情報を漏れ通らせられるはずだ、ということである。患者は突如として、そしておそらくは苦痛とともに、自分が麻痺していることに気づくはずなのだ。

これは脳卒中患者を対象とした臨床試験となるため、純粋な実験室実験よりも運営上むずかしくなる。だが、シロシビン支援療法はすでにうつ病、PTSD、終末期不安について試験されている。基盤は存在するのだ。そしてもしうまくいけば、それは病態失認に対する医学上の突破口であるだけではない------サイケデリックと脳卒中による欠損とが、単一の根底的メカニズムを通じてつながっているという証拠になる。これはほかのどの理論も予測していないことだ。

\section*{予測5：意識を消し去る麻酔薬はすべて臨界性を乱す}

麻酔薬は、まったく異なる化学的経路を通じて作用する。プロポフォールはGABA受容体に作用する。ケタミンはNMDAを遮断する。オピオイドはそれ独自の働きをする。異なる分子、異なるメカニズム、異なる脳の部位である。

だが四モデル理論は、それらすべてが意識に対して同じことをせざるをえないと言う------脳のダイナミクスを臨界性の閾値より下へ押し下げるのだ。臨界性こそが、意識の\textbf{物理的な必要条件}だからである。どのように乱すかは問題ではない。亜臨界に落ちれば、灯りは消える。

この予測は検証可能で、しかも具体的である。手元にあるすべての麻酔薬を取り上げてみよう。摂動複雑性指数（PCI）、レンペル-ジフ複雑性（LZ）、あるいは神経活動におけるべき乗則の指数といった道具を用いて、投与の前・最中・後に臨界性を測定するのだ。予測されるのは、意識を消し去る薬剤は、その受容体メカニズムのいかんにかかわらず、\textbf{つねに}脳を亜臨界へ押し下げるということである。そして意識を消さずに変容させる薬剤（低用量のケタミンやサイケデリックなど）は、臨界性を\textbf{下回らない}はずだ。

これは既存の技術で実施できる。臨界性の指標は存在する。麻酔薬も存在する。あとは誰かがその完全な比較を行うだけだ。そしてもしそれが全面的に成り立てば------意識を消し去るあらゆる薬剤が、それぞれ異なる経路を通じて作用するにもかかわらず、臨界性の攪乱へと収束するなら------それは臨界性が共通のメカニズムであり、無意識へと至る最終経路であることの、強力な証拠となる。

\section*{予測6：分離脳手術は心をきれいに二つに分けはしない------両方の半球を劣化させる}

外科医が重症のてんかんを治療するために脳梁を切断するとき、脳の二つの半球のあいだの主要な通信経路を断つことになる。従来の説明では、これによって二つの別々の心が生まれ、それぞれが専門化する------左半球が言語と論理を、右半球が空間的推論と情動を担う、というものだ。

四モデル理論は、それは間違っていると言う。あるいは少なくとも、劇的に単純化されすぎている、と。

予測はこうである------分離脳手術ののち、各半球は\textbf{完全だが劣化した}一そろいの認知的・体験的能力を保持する。きれいな分割ではない。「言語は左、空間は右」ではない。両方の半球がどちらもこなせるはずだが、以前より劣った形で、である。その劣化はホログラフィックなはずだ------つまり、特定の機能が消えるのではなく、すべてがぼやけていくということである。

その劣化は、切断の程度に比例するはずである。部分的な脳梁離断術（線維の一部だけを切る）なら部分的な劣化を、完全な脳梁離断術ならより大きな劣化を生むはずだ。

なぜか。理論によれば、脳内の情報はホログラフィックに蓄えられ、基盤全体にわたって分散しているからである。連結を切っても、あらかじめ存在する二つの心がきれいに分かれるわけではない。それは同じ情報の二つの\textbf{コピー}を劣化させるのであり、そのそれぞれが元のハードウェアの半分の上で走ることになる。

これについては、すでにいくらかの証拠がある------2017年のPintoらの研究は、分離脳の患者が、古典的な実験が示唆したよりもはるかに統合された行動を示すことを見いだした。だが理論は、その\textbf{メカニズム}を提供し、具体的なパターンを予測する。すなわち、半球の特化ではなく、両側性の劣化である。

\section*{予測7：臨界状態で四つのモデルを構築すれば、意識が得られる}

これは工学的な予測であり、大胆なものだ。

理論が正しければ、意識をもつ機械を構築できるはずである。偶然によってではなく、十分に「高度な」AIを作ることによってでもなく、仕様を実装することによって------臨界状態で動作する基盤の上で走る、四つの入れ子になったモデル（暗黙的世界モデル、暗黙的自己モデル、明示的世界モデル、明示的自己モデル）を実装することによって、である。

理論によれば、そのようなシステムは意識を単に\textbf{シミュレート}するのではない。それは意識を\textbf{もつ}のである。あなたの体験が脳の仮想的なモデルによって成り立っているのとまったく同じように、それは自らの仮想的なモデルによって成り立つ、真正の現象的体験をもつことになる。

どうすればそれがわかるのか。理論は、その違いが質的に明白であると予測する。「意識があるかもしれないし、ないかもしれない」ではない。\textbf{明らかに違う}のである。なぜなら、真正の自己シミュレーションを走らせているシステムは、どれほど洗練されたテキスト予測器とも根本的に異なる仕方で世界と関わるからだ。それは持続性をもつだろう------プロンプトから再構成されるのではなく、時間を通じて連続的に走り続けるシミュレーションである。それはある視点をもつだろう------明示的自己モデルによって維持された視点である。それがあなたを驚かせるのは、予想外の出力によってではなく、そこに本当に誰かがいるという感覚によってなのである。

これはまだ検証可能ではない------その工学が存在しないからだ。だが、その設計図は作業を導くのに十分なほど具体的である。そして誰かがそれを構築し、それがうまくいくなら、それが何よりの裏づけとなる。

\section*{予測8：睡眠は臨界状態をリセットするために存在する}

なぜ私たちは眠るのか。当たり前の答えは「休むため」だが、それは問いを先送りにするにすぎない。なぜ脳には、たとえば肝臓には要らないような仕方での休息が必要なのか。

四モデル理論には、具体的な答えがある。あなたの脳の基盤（アナログの生物学的なハードウェア）は、本来的に不安定である。ニューロンはノイズを含む。神経伝達物質は枯渇する。代謝老廃物は蓄積する。基盤はドリフトしていく。だが意識は臨界性を必要とし、それはきわめて特殊な力学的レジームである。脳は、このドリフトしていく基盤の上に、安定した計算層（カオスの縁にあるセル・オートマトン）を自己組織化する。そのオートマトンは何時間も（目覚めている一日じゅう）走ることができるが、やがて基盤が十分に遠くまでドリフトすると、もはや臨界的な力学を支えられなくなる。その時点で、オートマトンは徐々に暗くなるのではない。それは\textbf{崩壊する}。それが入眠である。

ノンレム睡眠は、その回復のプロセスである。基盤がリセットされる------神経伝達物質が補充され、老廃物が除去され、臨界性のための生化学的な条件が回復する。そしてこの回復の途中で基盤が周期的にふたたび臨界性の閾値に近づくと、オートマトンがつかのま点滅して立ち上がる。それがレム睡眠である。それが夢を見ることである。

90分のウルトラディアンリズム（一晩を通じたレムとノンレムのリズム）は、回復の途中で基盤が臨界点のまわりを振動していることの現れである。

これは、検証可能な複数の下位予測を生む。

\begin{enumerate}
  \item \textbf{臨界性は、目覚めている一日を通じて低下するはずである。} 朝、昼、晩に人々の脳の複雑さを測定する。予測されるのは、測定可能な低下である。
\end{enumerate}

\begin{enumerate}
  \item \textbf{入眠は、徐々に暗くなるのではなく、段階的な移行であるはずである。} 臨界性の指標は、入眠時に突然の低下を示すはずであり、それはオートマトンのデジタル的な崩壊を反映している。
\end{enumerate}

\begin{enumerate}
  \item \textbf{レムとノンレムは、臨界性を追跡するはずである。} 睡眠のなかで、レムの局面はノンレムよりもはるかに高い臨界性を示すはずであり、90分の周期は臨界性の時系列のなかに見えるはずである。
\end{enumerate}

\begin{enumerate}
  \item \textbf{明晰夢は、閾値を越えることである。} レムの途中で基盤が十分な臨界性に達すると、明示的自己モデルが活性化し、明晰になる。その始まりは、緩やかな立ち上がりではなく、脳波の複雑さにおける段階的な不連続であるはずだ。
\end{enumerate}

\begin{enumerate}
  \item \textbf{睡眠不足は、脳を亜臨界へと追いやる。} 十分に長く起きていれば、脳の臨界性は閾値を下回って次第に低下するはずである。認知的な欠損は、閾値をどれだけ下回ったかと相関するはずである。
\end{enumerate}

これらはすべて、既存の睡眠研究室の技術で検証可能である。そしてもしこれらが成り立つなら、それは睡眠が単なる「休息」ではないことを意味する------それは、意識を可能にする計算層のための、基盤のメンテナンス・プロトコルなのである。

\section*{予測9：解離性同一性障害における各交代人格は、それぞれ固有の神経的指紋をもつ}

解離性同一性障害（一人の人間のなかに複数の異なる同一性、すなわち「交代人格」が存在すること）は論争を呼ぶものであり、それにはもっともな理由がある。真正に異なる同一性と、意識的にせよ無意識的にせよ役を演じている人とを、どうやって見分けるのか。

四モデル理論は、その検証法を与えてくれる。もし交代人格が本物であるなら------つまり、それらが同じ基盤の上で動く明示的自己モデルの、真正に異なる配置であるなら------各交代人格は、区別可能で測定可能な神経的な指紋をもつはずである。単に行動が異なるのではない。単に自己報告が異なるのではない。\textbf{脳活動のパターンそのもの}が異なるのである。

その予測は具体的だ。ある解離性同一性障害の患者を対象に、異なる交代人格が現れているあいだの脳活動（fMRIまたはEEG）を記録する。交代人格どうしのあいだのばらつきを、同じ交代人格の時間を通じたばらつきと比較する。理論は、交代人格間のばらつきが交代人格内のばらつきよりも有意に大きくなると予測する。そしてその違いは一貫しているはずである。すなわち、交代人格Aの神経パターンは、毎回それとわかる交代人格Aであって、ランダムなノイズではない、というわけだ。

さらに具体的に言えば、理論はその違いがどこに現れるはずかを予測する。ESMに関連するネットワーク、とりわけデフォルト・モード・ネットワークと内側前頭前皮質------自己言及と視点取得に結びついた脳領域------においてである。

解離性同一性障害についての神経画像研究はこれまでにもいくつかあったが、四モデル理論は、単なる「違い」ではなく\textbf{一貫した、交代人格に固有の神経的な指紋}を予測するための理論的な基盤を提供する。もしこの予測が成り立つなら、それは交代人格が単に心理的なものではなく、神経のレベルでの異なる機能的配置であることの証拠となり、この障害についての私たちの理解と治療のあり方を一変させるだろう。

これらの予測は、どれも反証可能である。もしそれらが失敗するなら、理論は間違っているか、少なくとも不完全である。それこそが、これらの予測を有用なものにしているのである。

\end{document}
